\documentclass[11pt, a4paper]{ctexart}
\usepackage{../common/ncd-style}

\hypersetup{
  pdftitle={新概念丹麦语 4 · Holdning og Debat},
  pdfauthor={Jaehaerys},
  pdfsubject={Ny Koncept Dansk 4}
}

\begin{document}

% =====================================================================
% Book Cover Page
% =====================================================================
\begin{titlepage}
  \centering
  \vspace*{2.5cm}
  {\color{NCDRed}\rule{\linewidth}{2.5pt}}\\[1.2cm]
  {\fontsize{34}{40}\selectfont\bfseries\color{NCDNavy} 新概念丹麦语 4}\\[0.5cm]
  {\huge\itshape\color{NCDTeal} Ny Koncept Dansk 4 · Holdning og Debat}\\[0.8cm]
  {\large\bfseries\color{NCDGray} 思辨篇（全 20 课完整版） \quad|\quad 对标 NCE 3 \& DU3 Modul 4 (B1+ $\rightarrow$ B2-)}\\[1.2cm]
  {\color{NCDRed}\rule{\linewidth}{2.5pt}}\\[2cm]

  \begin{tcolorbox}[
    colback=NCDBgBlue,
    colframe=NCDNavy,
    width=0.90\linewidth,
    arc=4pt,
    boxrule=0.8pt,
    left=14pt, right=14pt, top=12pt, bottom=12pt
  ]
    \textbf{\color{NCDNavy}本册学习目标：}实现从“生活叙事与制度描述”到“观点建构、修辞微调与社会辩论”的 B1+ $\rightarrow$ B2- 跃升（每课 300--450 词）。每课采用“双篇递进思辨语篇（口语微调 / 报刊专栏 / 政策辩论）+ DU3 Modul 4 实战任务”双轨编排，重点攻克 \textbf{灵魂情态小品词（\textit{jo, da, nok, vist, lige, vel, bare, egentlig}）}、\textbf{名词化（\textit{-ning/-else}）与前置长分词定语}、\textbf{非真实条件虚拟句}、\textbf{严谨因果与限制连词（\textit{som følge heraf, medmindre, såfremt}）} 与 \textbf{\textit{Debatindlæg} 辩论写作框架}。
  \end{tcolorbox}

  \vfill
  {\large\color{NCDGray} 编写体例：导读 $\rightarrow$ 思辨双语篇与语音 $\rightarrow$ 屈折词汇表 $\rightarrow$ 语法与修辞精讲 $\rightarrow$ 句型操练 $\rightarrow$ Modul 4 实战与详解}
\end{titlepage}

% Table of Contents
\tableofcontents
\clearpage

% =====================================================================
% Unit 1
% =====================================================================
\ncdunit{1}{Nuancer og Stilskift}{灵魂语气与语体转换}{Lektion 1--5：情态小品词 I（jo/da/nok/vist）· 情态小品词 II（lige/vel/bare/egentlig/ellers）· 口语转书面语 · 礼貌让步反驳 · 25 个高频习语与比喻}

\ncdlesson{1}{Hvad mener danskerne egentlig, når de siger "det er fint"?}{丹麦人说“挺好”时到底是什么意思？}
{DU3 Modul 4（B1+ $\rightarrow$ B2-）｜对标 NCE 3 \& Modultest 4 \textit{Læsning Opgave 1--2 \& Skrivning Opgave 1 (Email)}}
{\textbf{灵魂情态小品词 I}：\textit{jo}（你我皆知/共识）· \textit{da}（当然/略带反驳/惊叹）· \textit{nok}（大概/恐怕/安心预测）· \textit{vist}（听说/好像）· 情态小品词在主句与从句中的\textbf{中场语序定位}}
{欢迎来到《新概念丹麦语 4 · 思辨篇》！从 B1 迈向 B1+/B2-，决定地道程度的分水岭往往不在高深长词，而在看似不起眼的“情态小品词（\textit{modalpartikler}）”。丹麦人说话从不直接把话说死，常把真正立场藏在 \textit{jo, da, nok, vist} 里：少一个词语气生硬，用错一个词意味全反。吃透本课，你不仅能听懂丹麦同事的“弦外之音”，写 Modul 4 私人长邮件时也能瞬间摆脱僵硬的教科书腔！}

% ---------------------------------------------------------------------
\begin{foerduleaeser}
带着下面两个问题阅读 Tekst A 和 Tekst B：
\begin{enumerate}
  \item \textit{Hvorfor misforstår Emma først Henriks kommentar ”det er vist meget fint” i Tekst A, og hvordan forklarer Sofie betydningen af ”vist” og ”jo”?}（在 Tekst A 中，艾玛最初为什么误解了亨里克那句“det er vist meget fint”？索菲是如何解释 \textit{vist} 和 \textit{jo} 的含义的？）
  \item \textit{Hvorfor er Kasper bekymret før fredagsmødet i Tekst B, og hvordan bruger Freja ”jo”, ”da” og ”nok” til at berolige ham?}（在 Tekst B 中，卡斯珀为什么在周五会议前感到担忧？弗蕾娅如何用 \textit{jo}、\textit{da} 和 \textit{nok} 来宽慰他？）
\end{enumerate}
\end{foerduleaeser}

% ---------------------------------------------------------------------
\ncdsubsec{二、Tekster（口语微调 / 报刊与思辨双语篇）}

\subsection*{Tekst A：\textit{Den skjulte betydning af ”det er fint”}（职场文化随笔 · 约 235 词）}

\begin{tipbox}{语篇说明}
本篇聚焦丹麦职场中极易引发跨文化误解的“含蓄表达”。阅读时请紧盯两种色彩标记：\textbf{\color{NCDNavy}深蓝色加粗的 \textit{jo / da}}（诉诸共识或温和反驳/鼓励）与\textbf{\color{NCDRed}红色加粗的 \textit{nok / vist}}（推测、宽慰或委婉存疑）。
\end{tipbox}

Da Emma flyttede til Aarhus og begyndte i en dansk virksomhed, troede hun, at hun hurtigt forstod sine kolleger. Efter sin første store præsentation spurgte hun sin kollega Henrik, hvad han syntes om hendes udkast. Henrik smilte roligt og svarede: ”Det er \dabold{vist} meget fint, men vi kan \da{jo} lige kigge på tidsplanen i morgen.”

Emma blev lettet og troede, at Henrik var helt overbevist. Ved kaffemaskinen forklarede hendes veninde Sofie dog, at danske småord ofte gemmer på vigtige nuancer. ”Når Henrik siger \dabold{vist}, betyder det faktisk, at han ikke er helt sikker, eller at han tvivler en smule,” sagde Sofie. ”Og med \da{jo} minder han dig høfligt om noget, som I begge to godt ved --- nemlig at tidsplanen stadig mangler arbejde.”

Efterhånden opdagede Emma, hvor stor forskel disse fire småord gør i hverdagen. Bruger en dansker \da{da}, ligger der ofte en mild indvending, forbavselse eller opmuntring i tonen, som i: ”Du kommer \da{da} til mødet, ikke?” eller ”Det var \da{da} ærgerligt!” Siger kollegaen derimod \dabold{nok}, udtrykker det enten en beroligende formodning eller en forsigtig indrømmelse: ”Det går \dabold{nok} godt,” eller ”Du har \dabold{nok} ret, men vi bør vente lidt.”

Samtidig må man passe på ordstillingen. I hovedsætninger står småordene efter det første verbum, mens de i ledsætninger sniger sig ind foran verbet, ligesom \textit{ikke}. Når man først lærer disse nuancer at kende, bliver danskernes sprog pludselig langt lettere at afkode.

\begin{transbox}
\small
\textbf{\color{NCDTeal}[Tekst A 中文精译]}\\
当艾玛搬到奥胡斯并在一家丹麦企业开始工作时，她原本以为自己很快就能听懂同事们的话。在完成第一次大型汇报后，她问同事亨里克对自己的方案初稿怎么看。亨里克平静地微笑着回答：“这大概挺不错的吧，不过我们明儿当然可以顺便再看一下时间表。”

艾玛松了一口气，以为亨里克已经完全被说服了。然而在咖啡机旁，她的好友索菲却解释说，丹麦语里的小品词常常藏着关键的语气微差。“当亨里克说 \textit{vist} 的时候，其实意味着他还不太确定，或者稍微有点疑虑，”索菲说道，“而用 \textit{jo}，他是在礼貌地提醒你一件你们俩心知肚明的事——那就是时间表还需要再打磨。”

渐渐地，艾玛发现这四个小词在日常生活中能带来天壤之别。当丹麦人使用 \textit{da} 时，语气里往往带有一丝温和的反驳、惊讶或鼓励，比如：“你肯定会来开会的吧？”或者“那可真是太遗憾了！”相反，如果同事说 \textit{nok}，通常表达一种令人安心的推测或谨慎的让步：“一切大概会顺利的”，或者“你说得恐怕没错，但我们还是应该再等等”。

与此同时，人还必须留心语序。在主句中，这些小词位于第一个变位动词之后；而在从句中，它们会像 \textit{ikke} 一样溜到变位动词前面。一旦掌握了这些细微差别，丹麦人说话背后的真实意图一下子就变得清晰可读了。
\end{transbox}

\subsection*{Tekst B：\textit{Intern mail: Nervøs før fredagsmødet}（对标 DU3 Modul 4 私人与职场回信 · 约 130 词）}

\begin{emailbox}{来信｜Fra: Kasper \quad Til: Freja \quad Emne: Lidt nervøs før præsentationen}
\textbf{Hej Freja}\\
Jeg skal præsentere vores nye budget på fredag, og jeg er ærlig talt lidt bekymret. Sidst jeg sendte et forslag til direktøren, skrev han bare, at det \dabold{vist} så fint ud, og så hørte jeg ikke mere. Tror du, at han var skeptisk?\\
\textbf{Venlig hilsen, Kasper}
\end{emailbox}

\begin{emailbox}{回复｜Fra: Freja \quad Til: Kasper \quad Emne: Sv: Lidt nervøs før præsentationen}
\textbf{Hej Kasper}\\
Tag det helt roligt! Det betyder \da{jo} ikke automatisk, at han var utilfreds --- han havde \dabold{vist} bare utrolig travlt den uge. Du ved \da{jo}, at vores chef sætter pris på konkrete tal, og dine analyser er \da{da} altid grundige og klare! Hvis vi øver os en halv time i morgen, skal det \dabold{nok} gå helt fint på fredag. Vi tales ved!\\
\textbf{Bedste hilsner, Freja}
\end{emailbox}

\textbf{Tekst B 阅读核对（Rigtigt / Forkert / Kort svar，答案见第六部分）：}
\begin{enumerate}
  \item \textit{Direktøren gav Kasper en lang og meget begejstret tilbagemelding sidste gang.} \underline{\hspace{2.2cm}}
  \item \textit{Hvordan forklarer Freja, at direktøren ikke skrev mere efter det sidste forslag?} \\\underline{\hspace{5.2cm}}
  \item \textit{Hvad foreslår Freja at gøre i morgen, og hvilken sikkerhed udtrykker hun med ”skal det nok gå”?} \underline{\hspace{5.2cm}}
\end{enumerate}

\begin{tipbox}{本课发音与语调要点}
\begin{itemize}
  \item \textbf{轻读 vs 重读决定词性}：情态小品词在句中一律\textbf{轻读且不带喉塞音（\textit{stød}）}：\textit{jo} \ipa{jo}、\textit{da} \ipa{da}、\textit{nok} \ipa{nʌɡ}、\textit{vist} \ipa{vesd}。一旦重读，词义立刻改变：\textit{\textbf{Jo!}} \ipa{ˈjoːˀ}（用于反驳否定问句的“不，恰恰相反/是的”）、\textit{\textbf{da}} \ipa{ˈdaː}（“那时”）、\textit{\textbf{nok}} \ipa{ˈnʌɡ}（“足够”，如 \textit{tid nok}）。
  \item \textbf{极易忽略的静音字母}：\textit{egentlig} \ipa{ˈeːˀənli}（\textbf{g 和 t 均不发音}）；\textit{indvending} \ipa{ˈenˌvɛneŋ} 与 \textit{indrømmelse} \ipa{ˈenˌʁømˀəlsə}（前缀 \textit{ind-} 里的 \textbf{d 不发音}）；\textit{ærgerligt} \ipa{ˈɛɐ̯ˀɐli}（\textbf{g 和 t 均不发音}！）。
  \item \textbf{核心词重音与软 d}：\textit{formodning} \ipa{fʌˈmoːˀðneŋ}（重音在第二音节，含软 d）；\textit{overbevist} \ipa{ˈʌʊ̯ɐbeˌviːˀsd}；\textit{tidsplan} \ipa{ˈtˢiðsˌpʰlæːˀn}。
\end{itemize}
\end{tipbox}

\begin{transbox}
\small
\textbf{\color{NCDTeal}[Tekst B 中文对照]}\\
\textbf{卡斯珀的邮件}：你好，弗蕾娅。我周五要汇报我们的新预算，老实说有点担心。上次我把一份提案发给总经理时，他只回了一句说看上去大概还行，之后就没下文了。你觉得他当时是持怀疑态度吗？诚挚的问候，卡斯珀。

\textbf{弗蕾娅的回复}：你好，卡斯珀。完全放轻松！这当然并不自动意味着他不满意——他那周大概只是忙得不可开交罢了。你本来也知道我们老板最看重扎实的数据，而且你的分析本来向来就非常细致清晰嘛！如果我们明天抽出半小时一起排练一下，周五肯定没问题的。回头聊！谨致问候，弗蕾娅。
\end{transbox}

% ---------------------------------------------------------------------
\ncdsubsec{三、Gloser（核心生词表）}

\begin{tipbox}{词表重点}
\textit{jo, da, nok, vist} 作为情态小品词时均无词形变化，但其语义高度依赖上下文语境与重音位置。名词 \textit{en nuance} 与 \textit{en indrømmelse}（让步）是 Modul 4 讨论语气、写作审题时的必备核心词；动词 \textit{at tvivle} 固定搭配介词 \textit{på}（怀疑某事）。
\end{tipbox}

\begin{glosertable}
小品词 & jo & \textit{(不变，情态小品词)} & 你我皆知的，毕竟，本来嘛 \\
小品词 & da & \textit{(不变，情态小品词)} & 当然，难道不是吗，可真是 \\
小品词 & nok & \textit{(不变，情态小品词)} & 大概，恐怕，放心会……的 \\
小品词 & vist & \textit{(不变，情态小品词)} & 听说，好像，大概（含不确定） \\
\midrule
名词 & en nuance & \textit{(nuancen, nuancer, nuancerne)} & 细微差别，语气微差 \\
名词 & et udkast & \textit{(udkastet, udkast, udkastene)} & 草稿，初稿（单复同形） \\
名词 & en tidsplan & \textit{(tidsplanen, tidsplaner, -erne)} & 时间表，进度计划 \\
名词 & en indvending & \textit{(indvendingen, indvendinger, -erne)} & 异议，反对意见 \\
名词 & en formodning & \textit{(formodningen, formodninger, -erne)} & 推测，假设，猜想 \\
名词 & en indrømmelse & \textit{(indrømmelsen, indrømmelser, -erne)} & 让步，承认 \\
名词 & en bekymring & \textit{(bekymringen, bekymringer, -erne)} & 担忧，顾虑 \\
\midrule
动词 [har] & at mene & \textit{(mener, mente, har ment)} & 认为，意思是指，意欲表达 \\
动词 [har] & at tvivle (på) & \textit{(tvivler, tvivlede, har tvivlet)} & 怀疑，拿不准 \\
动词 [har] & at berolige & \textit{(beroliger, beroligede, har beroliget)} & 安抚，使安心 \\
动词 [har] & at afkode & \textit{(afkoder, afkodede, har afkodet)} & 解码，领会弦外之音 \\
\midrule
形容词 & overbevist & \textit{(overbevist, overbeviste)} & 信服的，确信的 \\
形容词 & skeptisk & \textit{(skeptisk, skeptiske)} & 持怀疑态度的 \\
\end{glosertable}

% ---------------------------------------------------------------------
\ncdsubsec{四、Grammatik \& Stilistik（核心语法、语气与修辞精讲）}

\begin{grammatikbox}{语法点 1：四大灵魂情态小品词（\textit{jo, da, nok, vist}）语用潜台词全景对照}
这四个单音节词在英文里没有直接一一对应的词，因此常被初学者忽略，却决定着句子听起来像“冷冰冰的机器人”还是“懂分寸的丹麦人”：

\vspace{0.3em}
\renewcommand{\arraystretch}{1.25}
{\small
\begin{tabularx}{\linewidth}{@{} >{\bfseries\raggedright\arraybackslash}p{1.6cm} >{\raggedright\arraybackslash}p{3.6cm} >{\raggedright\arraybackslash}p{4.3cm} >{\raggedright\arraybackslash}X @{}}
\toprule
\textbf{小品词} & \textbf{核心语用功能} & \textbf{潜台词与中文神韵} & \textbf{典型例句与语气解析} \\
\midrule
\da{jo} & 诉诸双方共识\newline （提醒已知事实） & “你我都知道的呀 / 毕竟……嘛”（暗含：这事没必要再争了） & \textit{Vi har \da{jo} aftalt mødet i går.}（我们昨天毕竟约好开会了呀。）\newline \textit{Det tager \da{jo} tid at lære dansk.} \\
\midrule
\da{da} & 温和反驳、催促、安慰或惊叹 & “当然啦 / 可不是嘛 / 怎么不……呢？”（带情绪温度与互动感） & \textit{Det var \da{da} en god idé!}（这主意可真不错！）\newline \textit{Du er \da{da} ikke sur?}（你该不会生气了吧？） \\
\midrule
\dabold{nok} & 说话人自身较有把握的推测、宽慰或让步 & “大概会…… / 放心吧会…… / 恐怕确实如此，但……” & \textit{Han skal \dabold{nok} nå toget.}（他肯定赶得上火车，放心。）\newline \textit{Det er \dabold{nok} dyrt, men kvaliteten er høj.} \\
\midrule
\dabold{vist} & 转述听说、模糊记忆或委婉存疑 & “听说…… / 我好像记得…… / 恐怕有点……（留余地）” & \textit{Han kommer \dabold{vist} først i morgen.}（他好像明天才来，我不太确定。）\newline \textit{Mødet var \dabold{vist} lidt langt.} \\
\bottomrule
\end{tabularx}
}
\vspace{0.3em}

\textbf{高频对比辨析}：\ding{172} \dabold{nok} vs \dabold{vist}：说 \textit{Det går \dabold{nok} godt} 是安慰对方“肯定没问题”；说 \textit{Det gik \dabold{vist} meget godt} 则是“听说还行吧（但说话人自己心里可能没底或有所保留）”。\ding{173} \da{jo} vs \da{da}：\textit{Du ved \da{jo}, at...} 强调“你本来就知道”；\textit{Du ved \da{da}, at...} 往往在反问或反驳对方忘了这件事（“你明明知道的呀！”）。
\end{grammatikbox}

\begin{grammatikbox}{语法点 2：语序铁律——情态小品词在句子“中场（Midterfeltet）”的精确落点}
情态小品词与否定词 \textit{ikke}、时间副词 \textit{altid / ofte} 同属于丹麦语拓扑句法中的\textbf{中场副词（\textit{sætningsadverbialer}）}，位置错一步整句就变味：
\begin{enumerate}
  \item \textbf{主句（Hovedsætning）中的 V2 定位}：\ding{172} 若主语在句首，小品词紧跟在\textbf{第一个变位动词}之后：\textit{Hun kommer \da{nok} i dag.}（不可以说 \ding{55}\ \textit{Hun nok kommer}）。\ding{173} 若时间/地点或从句占领句首触发 V2 倒装，小品词放在\textbf{动词 + 主语}之后：\textit{I morgen kommer hun \da{nok}.}
  \item \textbf{遇上轻代词宾语与 \textit{ikke} 谁先谁后？}：\ding{174} 没有重读的短代词宾语（\textit{det, ham, hende, mig, sig}）必须穿过中场跑到小品词\textbf{前面}：\textit{Jeg sagde \dabold{det jo} til dig!} / \textit{Han husker \dabold{det vist} ikke.} \ding{175} 情态小品词通常排在否定词 \dabold{ikke} \textbf{前面}：\textit{Det er \da{jo ikke} så svært.} / \textit{Hun kommer \dabold{nok ikke} til festen.}
  \item \textbf{从句（Ledsætning）语序铁律}：在 \textit{at, fordi, når, hvis, som} 等引导的从句中，小品词必须紧贴主语之后、跑到\textbf{变位动词之前}：
    \begin{itemize}
      \item \ding{51}\ \textit{Han blev hjemme, fordi han \dabold{jo} havde travlt.}（\ding{55}\ \textit{...fordi han havde jo travlt.}）
      \item \ding{51}\ \textit{Hun fortalte, at direktøren \dabold{vist ikke} var helt tilfreds.}
    \end{itemize}
\end{enumerate}
\end{grammatikbox}

\begin{grammatikbox}{地道交际表达与文化贴士（Pragmatik på arbejdspladsen）}
\begin{itemize}
  \item \da{Det skal nok gå!} —— 这事肯定没问题的！（丹麦人安慰朋友、同事的万能金句）。
  \item \da{Det var da ærgerligt / dejligt!} —— 那可真遗憾 / 太棒了！（加 \textit{da} 让情绪共鸣自然真实）。
  \item \textbf{文化提醒}：在丹麦职场听评价时，如果对方用 \textit{vist} 或 \textit{nok} 搭配平淡形容词（如 \textit{Det ser vist okay ud}），多半意味着还有修改空间，千万别误以为对方已经“百分百满意”！
\end{itemize}
\end{grammatikbox}

% ---------------------------------------------------------------------
\ncdsubsec{五、Mønsterøvelser（句型与修辞操练）}

\textbf{练习 A：语序定位手术——将括号中的情态小品词插入句中唯一准确位置（注意 V2 倒装、轻代词与从句语序）}
\begin{enumerate}
  \item \textit{Jeg sagde det til dig i går.}（插入 \textit{jo}）$\rightarrow$
  \item \textit{I morgen når hun ikke at sende rapporten.}（插入 \textit{nok}）$\rightarrow$
  \item \textit{Hun blev overrasket, fordi hun ikke kendte reglerne.}（在从句中插入 \textit{jo}）$\rightarrow$
  \item \textit{De var meget glade for forslaget.}（插入 \textit{da} 表温和感叹/反驳）$\rightarrow$
  \item \textit{Direktøren rejser til Berlin på tirsdag.}（插入 \textit{vist} 表听说/不确定）$\rightarrow$
  \item \textit{Han tror, at mødet begynder klokken ti.}（在从句中插入 \textit{vist}）$\rightarrow$
  \item \textit{Du har glemt din paraply igen.}（插入 \textit{da} 温和提醒）$\rightarrow$
\end{enumerate}

\textbf{练习 B：语用潜台词选填（从 \textit{jo, da, nok, vist} 中选择最符合语境的词填空）}
\begin{enumerate}
  \item Du behøver ikke at være nervøs for eksamen --- det skal \underline{\hspace{2.2cm}} gå rigtig godt!
  \item Hvorfor spørger du mig om adressen? Du har \underline{\hspace{2.2cm}} været hos mig tre gange før!
  \item Jeg kan desværre ikke huske det præcise tal, men det kostede \underline{\hspace{2.2cm}} omkring 500 kroner.
  \item Nå, har du allerede mistet din cykel? Det var \underline{\hspace{2.2cm}} utrolig ærgerligt!
  \item Han lovede at ringe klokken to, men han har \underline{\hspace{2.2cm}} glemt det, fordi han sidder i møde.
  \item Kommer du ikke med til frokost? --- \underline{\hspace{2.2cm}}, selvfølgelig gør jeg det!（注意回答否定问句时的重读词！）
\end{enumerate}

\textbf{练习 C：中译丹实战（活用情态小品词传递潜台词与准确语序）}
\begin{enumerate}
  \item 放心吧，明天他大概不会迟到。（用 \textit{nok}）
  \item 我毕竟已经把邮件发给你了呀。（用 \textit{jo}）
  \item 听说我们的老板下周要去奥胡斯。（用 \textit{vist}）
  \item 你当然可以问同事这个问题！（用 \textit{da}）
  \item 她没来上班，因为她好像生病了。（在从句中用 \textit{vist}）
  \item 他大概是对的，但我们还需要再核对一下数据。（用 \textit{nok}）
\end{enumerate}

% ---------------------------------------------------------------------
\ncdsubsec{六、DU3 Modul 4 实战微任务与参考答案}

\begin{emailbox}{DU3 Modul 4 写作任务（对标 Skrivning Opgave 1 · 私人回复邮件 $\ge$ 125 ord）}
\textbf{Situation}: Din ven Mikkel er lige startet i et nyt job i København. Han har skrevet til dig, fordi hans danske chef sagde: ”Det er vist meget fint”, da Mikkel viste sit første projekt. Nu er Mikkel bange for, at chefen slet ikke kan lide hans arbejde.\\
\textbf{Opgave}: Skriv en beroligende og hjælpsom mail til Mikkel på \textbf{mindst 125 ord} (gerne 140--160 ord). Forklar ham, hvad danske småord som \textit{vist}, \textit{jo}, \textit{nok} og \textit{da} kan betyde i sådan en situation, giv ham et konkret råd til næste møde med chefen, og brug selv alle fire modalpartikler naturligt i din mail!
\end{emailbox}

\begin{answerbox}
\textbf{Før du læser 参考回答}
\begin{enumerate}
  \item \textit{Emma tror først, at Henrik er helt overbevist om hendes udkast, men Sofie forklarer, at ”vist” viser en lille tvivl eller usikkerhed, mens ”jo” høfligt minder om noget, de begge to ved --- at tidsplanen stadig mangler arbejde.}
  \item \textit{Kasper er bekymret, fordi direktøren sidst kun skrev, at hans forslag ”vist så fint ud”. Freja bruger ”jo” til at minde om direktørens travlhed og stil, ”da” til at rose Kaspers analyser og ”nok” til at forsikre ham om, at det går godt.}
\end{enumerate}

\textbf{Tekst B 阅读核对答案}
\begin{enumerate}
  \item \da{Forkert}（总经理只回复了一句 \textit{at det vist så fint ud}，之后没有再多说）
  \item \da{Hun mener, at han vist bare havde utrolig travlt den uge.}
  \item \da{Hun foreslår, at de øver sig en halv time i morgen, og viser tillid til, at det går godt.}
\end{enumerate}

\textbf{练习 A 答案}
\begin{enumerate}
  \item \textit{Jeg sagde det \da{jo} til dig i går.}（轻代词 \textit{det} 必须放在情态小品词 \textit{jo} 前面）
  \item \textit{I morgen når hun \dabold{nok} ikke at sende rapporten.}（V2 倒装后顺序：动词 + 主语 + \textit{nok} + \textit{ikke}）
  \item \textit{Hun blev overrasket, fordi hun \da{jo} ikke kendte reglerne.}（从句中 \textit{jo} 置于主语后、\textit{ikke} 与变位动词前）
  \item \textit{De var \da{da} meget glade for forslaget.}（主句中位于变位动词 \textit{var} 之后）
  \item \textit{Direktøren rejser \dabold{vist} til Berlin på tirsdag.}（位于第一变位动词 \textit{rejser} 之后）
  \item \textit{Han tror, at mødet \dabold{vist} begynder klokken ti.}（\textit{at} 引导从句，\textit{vist} 放在变位动词 \textit{begynder} 前面）
  \item \textit{Du har \da{da} glemt din paraply igen.}（位于助动词 \textit{har} 之后、过去分词 \textit{glemt} 之前）
\end{enumerate}

\textbf{练习 B 答案}
\begin{enumerate}
  \item \dabold{nok}（宽慰、有把握的安慰：“肯定没问题的”）
  \item \da{jo}（提醒双方早已知道的事实：“你毕竟来过三次了呀”）
  \item \dabold{vist}（记忆不确切或模糊估算：“大概五百克朗左右吧”）
  \item \da{da}（表达感叹与共情：“那可真是太遗憾了”）
  \item \dabold{nok} / \dabold{vist}（合理推测他恐怕是忘了，用 \textit{nok} 或表示揣测的 \textit{vist} 均可）
  \item \da{Jo}（注意此空用于反驳否定问句 \textit{Kommer du ikke...?}，须重读用 \textit{Jo}）
\end{enumerate}

\textbf{练习 C 答案}
\begin{enumerate}
  \item \textit{Bare rolig, i morgen kommer han nok ikke for sent.}
  \item \textit{Jeg har jo allerede sendt mailen til dig.}
  \item \textit{Vores chef skal vist til Aarhus i næste uge.}
  \item \textit{Du kan da sagtens stille din kollega det spørgsmål!}
  \item \textit{Hun kom ikke på arbejde, fordi hun vist er syg.}
  \item \textit{Han har nok ret, men vi skal lige tjekke tallene igen.}
\end{enumerate}

\textbf{DU3 Modul 4 高分范文（148 词，不含署名）}\\
\textit{Hej Mikkel}

\textit{Tak for din mail! Jeg kan godt forstå, at du blev lidt usikker efter mødet, men du behøver slet ikke at gå i panik. Danskere bruger nemlig ofte småord, som i starten kan virke lidt forvirrende.}

\textit{Når din chef siger ”det er \dabold{vist} meget fint”, betyder det \da{jo} ikke, at han synes, dit arbejde er dårligt. Ofte bruger danskere bare \dabold{vist}, når de gerne vil kigge lidt nærmere på detaljerne, før de siger endeligt ja. Du er \da{da} lige startet i firmaet, så ingen forventer, at alt er perfekt fra dag ét!}

\textit{Mit bedste råd er, at du spørger ham venligt næste gang: ”Har du eventuelt et par forslag til, hvad jeg kan forbedre i udkastet?” Danske chefer sætter \da{jo} stor pris på åbenhed, og hvis du tager det roligt og spørger nysgerrigt, skal det \dabold{nok} gå rigtig godt mellem jer.}

\textit{Jeg krydser fingre for dig på det nye arbejde!}

\textit{Bedste hilsner\\
\mbox{[Dit navn]}}

\small\textbf{名师点评}：148 词（不含署名）远超 Modul 4 Opgave 1 的 125 词底线。全信将 \textit{vist, jo, da, nok} 自然嵌入主句、否定句与从句，既有情绪共情（\textit{Du er da lige startet...}），又有实操建议与暖心收尾（\textit{skal det nok gå rigtig godt}），语域地道流畅，完全符合私人回复邮件的 B1+/B2- 评分标准。
\end{answerbox}
\ncdlesson{2}{Skal vi lige tage en kop kaffe?}{我们要不喝杯咖啡？}
{DU3 Modul 4（B1+ $\rightarrow$ B2-）｜对标 NCE 3 \& Modultest 4 \textit{Læsning Opgave 2 (Chat) \& Mundtlig Opgave 2 (Paropgave)}}
{\textbf{灵魂情态小品词 II}：\textit{lige}（稍微/顺便/降压）· \textit{vel}（反意求同“吧？”）· \textit{bare}（只是/尽管）· \textit{egentlig}（究竟/其实）· \textit{ellers}（要不然/除此之外）的语气软化与礼貌协商功能}
{在丹麦职场与日常对话中，直译中文的祈使句或英语的 \textit{Can you help me?} 常显得过于生硬，甚至带着命令感。学会用 \textbf{\textit{lige, vel, bare, egentlig, ellers}} 给请求“卸下压迫感”，是融入丹麦平权文化与通过 DU3 Modul 4 口语双人互动的关键。本课带你听懂咖啡机旁的真实协商艺术！}

% ---------------------------------------------------------------------
\begin{foerduleaeser}
带着下面两个问题阅读 Tekst A 和 Tekst B：
\begin{enumerate}
  \item \textit{Hvad vil Kasper have Emma til at hjælpe med i Tekst A, og hvordan giver han hende plads til at sige nej?}（在 Tekst A 中，卡斯珀想请艾玛帮什么忙？他是如何给对方留出拒绝余地的？）
  \item \textit{Hvorfor tøver Sofie med at tage med til fredagens borgermøde i Tekst B, og hvilket alternativ foreslår Thomas med ”ellers”?}（在 Tekst B 中，索菲为什么对参加周五的市民会议有些犹豫？托马斯用 \textit{ellers} 提出了什么替代方案？）
\end{enumerate}
\end{foerduleaeser}

% ---------------------------------------------------------------------
\ncdsubsec{二、Tekster（口语微调 / 报刊与思辨双语篇）}

\subsection*{Tekst A：\textit{Lige et øjeblik over kaffen}（职场口语对话 · 约 205 词）}

\begin{tipbox}{语篇说明}
本篇还原丹麦办公室咖啡机旁的真实对话。阅读时请特别关注\textbf{\color{NCDNavy}深蓝色加粗的 \textit{lige / bare / egentlig}}（用于放低姿态、发出邀请或温和切入）与\textbf{\color{NCDRed}红色加粗的 \textit{vel / ellers}}（用于寻求认同与提供退路）。
\end{tipbox}

\begin{tcolorbox}[enhanced, breakable, colback=NCDBgBlue, colframe=NCDNavy, coltitle=white, fonttitle=\bfseries, title={办公室咖啡间协商｜Kasper \& Emma}, arc=2pt, boxrule=0.8pt, left=10pt, right=10pt, top=7pt, bottom=7pt]
\small
\textbf{Kasper}: Hej Emma! Har du \da{egentlig} travlt lige nu, eller skal vi \da{lige} tage en kop kaffe?\\[3pt]
\textbf{Emma}: Hej Kasper! Jeg har \da{bare} et par mails tilbage, men en kaffepause lyder dejlig. Hvad tænker du på?\\[3pt]
\textbf{Kasper}: Jeg skal fremlægge vores nye spareplan for bestyrelsen i morgen. Kunne du \da{lige} kigge mine fem slides igennem bagefter? Det tager \dabold{vel} kun ti minutter?\\[3pt]
\textbf{Emma}: Jo, selvfølgelig vil jeg det! Men hvordan går det \dabold{ellers} med tallene? Du virkede lidt presset i går.\\[3pt]
\textbf{Kasper}: Ja, det var lidt stramt, men nu ser det fint ud. Du må \da{bare} sige til, hvis du hellere vil gøre det efter frokost. \dabold{Ellers} kan jeg også sende udkastet til dig på mail, så du \da{bare} læser det, når det passer dig.\\[3pt]
\textbf{Emma}: Nej, lad os \da{bare} tage det klokken tolv ved mit skrivebord. Vi skal \da{lige} sikre, at konklusionen står skarpt, ikke?\\[3pt]
\textbf{Kasper}: Jo, præcis! Det er \dabold{vel} heller ikke nødvendigt at tage alle bilagene med i selve præsentationen?\\[3pt]
\textbf{Emma}: Nej, det er det bestemt ikke. Hent du kaffen, så rydder jeg \da{lige} mit bord imens!
\end{tcolorbox}

\begin{transbox}
\small
\textbf{\color{NCDTeal}[Tekst A 中文精译]}\\
\textbf{卡斯珀}：嗨，艾玛！你现在其实忙吗，还是我们要不要顺便喝杯咖啡？\\
\textbf{艾玛}：嗨，卡斯珀！我只剩几封邮件要处理，不过去喝杯咖啡休息一下听起来很棒。你在想什么事呢？\\
\textbf{卡斯珀}：我明天得向理事会汇报我们的新节约方案。待会儿你能稍微帮我看一下我那五页幻灯片吗？大概也就花十分钟吧？\\
\textbf{艾玛}：好呀，我当然愿意！不过除此之外，数据那块推进得怎么样？你昨天看起来压力有点大。\\
\textbf{卡斯珀}：是啊，当时有点赶，不过现在看上去不错了。如果你更愿意午饭后再看，尽管告诉我。要不然我也可以把草稿邮件发给你，等你方便的时候再看就行。\\
\textbf{艾玛}：不用，我们十二点就在我办公桌那儿直接看吧。我们稍微确认一下结论写得够犀利，对吧？\\
\textbf{卡斯珀}：对，正是！反正也没必要把所有附件都放进正式演示里吧？\\
\textbf{艾玛}：是的，完全没必要。你去倒咖啡吧，我顺便把桌子收拾一下！
\end{transbox}

\subsection*{Tekst B：\textit{Chat om fredagens borgermøde}（对标 DU3 Modul 4 社交短信协商 · 约 140 词）}

\begin{emailbox}{微信 / 短信对话｜Sofie \& Thomas}
\textbf{Thomas}: Hej Sofie! Skal vi ikke tage sammen hen til borgermødet om nye cykelstier på fredag?\\[3pt]
\textbf{Sofie}: Hej Thomas! Jeg vil \da{egentlig} gerne med, men mødet slutter \dabold{vel} ret sent, gør det ikke? Jeg skal nemlig tidligt op lørdag morgen.\\[3pt]
\textbf{Thomas}: Det varer vist til omkring klokken ni, men du kan \da{bare} gå i pausen efter første time, hvis du bliver træt. \dabold{Ellers} kan vi \da{lige} mødes til en hurtig kop kaffe klokken fem først, så vi når at drøfte vores spørgsmål inden mødet.\\[3pt]
\textbf{Sofie}: Åh, det er en super idé! Jeg tjekker \da{lige} min kalender, når jeg kommer hjem, og så skriver jeg til dig i aften. Vi kan \da{vel} mødes ved biblioteket?\\[3pt]
\textbf{Thomas}: Ja da! Sig \da{bare} til, hvis klokken fem bliver for tidligt for dig.
\end{emailbox}

\textbf{Tekst B 阅读核对（Rigtigt / Forkert / Kort svar，答案见第六部分）：}
\begin{enumerate}
  \item \textit{Sofie vil slet ikke med til borgermødet om cykelstier, fordi emnet ikke interesserer hende.} \underline{\hspace{2.2cm}}
  \item \textit{Hvad foreslår Thomas, at Sofie kan gøre, hvis mødet bliver for langt for hende?} \\\underline{\hspace{5.2cm}}
  \item \textit{Hvornår lover Sofie at vende tilbage med endeligt svar til Thomas?} \underline{\hspace{4.5cm}}
\end{enumerate}

\begin{tipbox}{本课发音与语调要点}
\begin{itemize}
  \item \textbf{五大软化小品词的口语极简吞音}：\textit{lige} 读 \ipa{ˈliːə}（\textbf{g 完全不发音}，快读时常缩成短促的 \ipa{li}）；\textit{egentlig} 读 \ipa{ˈeːˀənli}（\textbf{g 与 t 双静音}！）；\textit{bare} 读 \ipa{ˈbɑːɐ}（常弱化为 \ipa{bɑ}）；\textit{vel} 在句中轻读 \ipa{vɛl}，但在句末疑问标签 \textit{..., vel?} 中音调上扬；\textit{ellers} \ipa{ˈɛlˀɐs}（带清晰喉塞音）。
  \item \textbf{对话连读与弱读节奏}：\textit{Skal vi lige} 常连读为 \ipa{sga vi ˈliːə}（\textit{skal} 的 l 常脱落）；\textit{Det tager vel} \ipa{de ˈtˢɑːˀ vɛl}（\textit{tager} 的 g 不发音！）。
  \item \textbf{核心词重音与软 d}：\textit{et borgermøde} \ipa{ˈbɒʊ̯ˀɐˌmøːðə}（\textit{borger} 中的 \textit{g} 滑音化为 \ipa{ʊ̯}，\textit{møde} 为软 d）；\textit{behagelig} \ipa{beˈhæːˀəli}；\textit{et udkast} \ipa{ˈuðˌkʰasd}。
\end{itemize}
\end{tipbox}

\begin{transbox}
\small
\textbf{\color{NCDTeal}[Tekst B 中文对照]}\\
\textbf{托马斯}：嗨，索菲！这周五我们要不要一起去参加关于新建自行车道的市民会议？\\
\textbf{索菲}：嗨，托马斯！我其实很想去，不过会议大概会结束得挺晚的，对吧？我周六一大早还要早起呢。\\
\textbf{托马斯}：听说会持续到九点左右，但如果你觉得累，完全可以在第一小时结束后的休息时间先走呀。要不然我们也可以五点先稍微碰头喝杯快咖啡，这样赶在开会前就能讨论好我们的问题。\\
\textbf{索菲}：啊，这主意太棒了！我回家后稍微查一下日程表，今晚给你回信。我们大概约在图书馆门口碰面吧？\\
\textbf{托马斯}：那当然！如果五点对你来说太早，尽管告诉我一声就行。
\end{transbox}

% ---------------------------------------------------------------------
\ncdsubsec{三、Gloser（核心生词表）}

\begin{tipbox}{词表重点}
本课五个核心小品词 \textit{lige, vel, bare, egentlig, ellers} 在句中作语气词时都没有词形变化；但务必注意 \textit{at gennemgå}（浏览、审核）和 \textit{at fremlægge}（汇报、呈交）是不规则复合动词，过去时和完成时分别随 \textit{gå} 和 \textit{lægge} 变化！
\end{tipbox}

\begin{glosertable}
小品词 & lige & \textit{(不变，情态小品词)} & 稍微，顺便，一下子（软化请求） \\
小品词 & vel & \textit{(不变，情态小品词)} & 大概……吧，难道不是吗（求同） \\
小品词 & bare & \textit{(不变，情态小品词)} & 只是，尽管，放心……好了 \\
小品词 & egentlig & \textit{(不变，情态小品词)} & 其实，究竟，说起来 \\
副词 & ellers & \textit{(不变，副词/小品词)} & 要不然，否则；除此之外 \\
\midrule
名词 & et øjeblik & \textit{(øjeblikket, øjeblikke, -ne)} & 片刻，一会儿 \\
名词 & en kaffepause & \textit{(kaffepausen, -pauser, -erne)} & 喝咖啡小憩时间 \\
名词 & en spareplan & \textit{(spareplanen, -planer, -erne)} & 节支方案，紧缩计划 \\
名词 & en bestyrelse & \textit{(bestyrelsen, bestyrelser, -ne)} & 理事会，董事会 \\
名词 & et borgermøde & \textit{(borgermødet, -møder, -erne)} & 市民议事会，居民大会 \\
名词 & et bilag & \textit{(bilaget, bilag, bilagene)} & 附件，附录（单复同形） \\
名词 & en konklusion & \textit{(konklusionen, -er, -erne)} & 结论 \\
\midrule
动词 [har] & at fremlægge & \textit{(fremlægger, fremlagde, har fremlagt)} & 汇报，展示，提交（方案） \\
动词 [har] & at gennemgå & \textit{(gennemgår, gennemgik, har gennemgået)} & 过一遍，仔细审核 \\
动词 [har] & at drøfte & \textit{(drøfter, drøftede, har drøftet)} & 探讨，商议（偏书面） \\
动词 [har] & at sikre & \textit{(sikrer, sikrede, har sikret)} & 确保，保障 \\
\midrule
形容词 & krævende & \textit{(不变，现在分词形容词)} & 要求苛刻的，费心力的 \\
\end{glosertable}

% ---------------------------------------------------------------------
\ncdsubsec{四、Grammatik \& Stilistik（核心语法、语气与修辞精讲）}

\begin{grammatikbox}{语法点 1：五大语气软化小品词（\textit{lige, vel, bare, egentlig, ellers}）语义与语用矩阵}
如果说 Lektion 1 的 \textit{jo, da, nok, vist} 负责交代“信息真伪与双方共识”，那么本课的五个小品词就是丹麦社交润滑剂——专门用来\textbf{降低指令侵略性、缓和冲突与制造亲切感}：

\vspace{0.3em}
\renewcommand{\arraystretch}{1.25}
{\small
\begin{tabularx}{\linewidth}{@{} >{\bfseries\raggedright\arraybackslash}p{1.65cm} >{\raggedright\arraybackslash}p{3.4cm} >{\raggedright\arraybackslash}p{4.3cm} >{\raggedright\arraybackslash}X @{}}
\toprule
\textbf{小品词} & \textbf{核心降压功能} & \textbf{语用潜台词} & \textbf{地道高频句式} \\
\midrule
\da{lige} & 弱化请求分量\newline （“稍微 / 顺便”） & “就花你一小会儿，不是大事” & \textit{Kan du \da{lige} åbne vinduet?}\newline \textit{Jeg tjekker \da{lige} min kalender.} \\
\midrule
\dabold{vel} & 礼貌寻求确认\newline （“大概……吧？”） & “我猜是这样，对不对？”（暗含期待对方点头） & \textit{Det koster \dabold{vel} ikke meget?}\newline \textit{Du kommer med, \dabold{vel}?}（句末标签） \\
\midrule
\da{bare} & 授权、安慰或轻描淡写（“尽管 / 只是”） & “没关系，随便这么做就好 / 别多想” & \textit{Tag \da{bare} den sidste bolle!}\newline \textit{Du må \da{bare} ringe, hvis der er noget.} \\
\midrule
\da{egentlig} & 缓冲突兀提问或委婉转折（“究竟 / 其实”） & “顺便好奇问一句 / 说实在的……” & \textit{Hvad synes du \da{egentlig} om planen?}\newline \textit{Jeg vil \da{egentlig} gerne med, men...} \\
\midrule
\dabold{ellers} & 留出体面退路或关切寒暄（“要不然 / 另外”） & “如果不方便也没事，我们还有 B 方案” & \textit{\dabold{Ellers} kan vi mødes i morgen.}\newline \textit{Hvordan går det \dabold{ellers} på jobbet?} \\
\bottomrule
\end{tabularx}
}
\vspace{0.3em}

\textbf{避坑对比}：\ding{172} 询问对方意愿时，说 \textit{Hvad mener du?} 听起来像审讯；加上 \textit{egentlig}（\textit{Hvad mener du \da{egentlig}?}）立刻变成真诚好奇的探讨。\ding{173} \textit{ellers} 放句首表示“否则/要不然”时，\textbf{触发 V2 倒装}：\textit{\dabold{Ellers} \textbf{kan} vi tage det i morgen.}（绝不能写成 \ding{55}\ \textit{Ellers vi kan...}）。
\end{grammatikbox}

\begin{grammatikbox}{语法点 2：高情商请求的“三步软化法”与句末 \textit{, vel?} 的回答规则}
\begin{enumerate}
  \item \textbf{从生硬祈使句升级到丹麦礼貌体}：
    \begin{itemize}
      \item 生硬直译：\ding{55}\ \textit{Hjælp mig med opgaven!}（带有“命令”感，只适合极其紧急情况）
      \item 柔和邀请：\ding{51}\ \textit{Kunne du \da{lige} hjælpe mig med opgaven, hvis du har tid?}（过去式情态动词 \textit{Kunne} + \textit{lige} 双重软化）
      \item 体贴留退路：\ding{51}\ \textit{Du kan \da{bare} kigge på den i morgen --- \dabold{ellers} finder vi en anden løsning.}
    \end{itemize}
  \item \textbf{句末确认标签 \textit{, vel?} 与 \textit{, ikke?} 的细微区别及回答法}：
    \ding{172} \textit{Det er dyrt, \textbf{ikke}?}（说话人几乎确定，只等对方附和）；\textit{Det er ikke for dyrt, \dabold{vel}?}（说话人希望得到“不贵”的安抚，带有小心翼翼试探感）。\newline
    \ding{173} \textbf{回答否定式 \textit{vel?} 疑问句注意 \textit{Jo} vs \textit{Nej}}：若对方问 \textit{Det tager ikke lang tid, vel?}（不会花很久吧？），如果你同意“确实不花很久”，请回答 \da{Nej, slet ikke}；如果你想反驳“不，其实要花挺久”，则用重读的 \dabold{Jo, det tager faktisk en time!}
\end{enumerate}
\end{grammatikbox}

% ---------------------------------------------------------------------
\ncdsubsec{五、Mønsterøvelser（句型与修辞操练）}

\textbf{练习 A：高情商软化改写——按括号提示用指定小品词软化下列生硬句子（注意语序与 V2 倒装）}
\begin{enumerate}
  \item \textit{Send udkastet til mig i dag.}（改用 \textit{Kan du lige...} 软化为礼貌请求）$\rightarrow$
  \item \textit{Hvad arbejder du med nu?}（在动词和主语后插入 \textit{egentlig}，并改得更柔和）$\rightarrow$
  \item \textit{Hvis du ikke har tid, finder vi en ny dag.}（以 \textit{Ellers} 开头改写后半句，注意 V2 倒装）$\rightarrow$
  \item \textit{Du må bruge min computer.}（在祈使句或情态句中插入 \textit{bare} 表热情许可）$\rightarrow$
  \item \textit{Du har ikke glemt mødet i morgen?}（用句中 \textit{vel} 改写为委婉确认问句）$\rightarrow$
  \item \textit{Hvordan går det med dit nye arbejde?}（插入 \textit{ellers} 表示“除此之外最近如何”）$\rightarrow$
\end{enumerate}

\textbf{练习 B：语境选词填空（从 \textit{lige, vel, bare, egentlig, ellers} 中选填，注意第 6 题双空！）}
\begin{enumerate}
  \item Undskyld, kan jeg \underline{\hspace{2.2cm}} låne din kuglepen et øjeblik?
  \item Du skal ikke bekymre dig om regningen --- jeg skal \underline{\hspace{2.2cm}} betale kaffen i dag!（表“尽管让我来就行”）
  \item Det bliver \underline{\hspace{2.2cm}} ikke regnvejr i eftermiddag, gør det?
  \item Jeg troede, at han boede i Odense, men han kommer \underline{\hspace{2.2cm}} fra Esbjerg.
  \item Lad os mødes klokken tre, og \underline{\hspace{2.2cm}} kan vi altid ringes ved i aften.
  \item Vil du \underline{\hspace{2.2cm}} holde døren et sekund? Du behøver \underline{\hspace{2.2cm}} ikke at vente på mig bagefter.
\end{enumerate}

\textbf{练习 C：中译丹实战（自然融入五大语气软化词）}
\begin{enumerate}
  \item 我们要不要顺便去咖啡馆喝杯咖啡？（用 \textit{lige}）
  \item 说起来，你周末究竟打算做什么？（用 \textit{egentlig}）
  \item 如果你有问题，尽管给我发消息就行。（用 \textit{bare}）
  \item 这次考试大概不会太难吧？（用 \textit{vel}）
  \item 我们可以今天见面，要不然就周四再聊。（用 \textit{ellers}）
  \item 顺便问一下，除此之外你的项目进展得怎么样？（用 \textit{ellers}）
\end{enumerate}

% ---------------------------------------------------------------------
\ncdsubsec{六、DU3 Modul 4 实战微任务与参考答案}

\begin{emailbox}{DU3 Modul 4 写作任务（对标 Skrivning Opgave 1 · 私人协商邮件 $\ge$ 125 ord）}
\textbf{Situation}: Du skal holde et vigtigt oplæg på dit studie eller arbejde næste tirsdag, men du er usikker på, om din tekst er klar og naturlig. Din danske kollega eller studieven, Sofie, er altid meget hjælpsom, men hun har også meget travlt for tiden.\\
\textbf{Opgave}: Skriv en høflig og uformel mail til Sofie på \textbf{mindst 125 ord} (gerne 140--160 ord). Spørg forsigtigt, om hun har tid til at læse dit udkast igennem eller mødes over en kop kaffe, giv hende en nem mulighed for at sige nej, og brug alle fem ord (\textit{lige, vel, bare, egentlig, ellers}) naturligt i mailen!
\end{emailbox}

\begin{answerbox}
\textbf{Før du læser 参考回答}
\begin{enumerate}
  \item \textit{Kasper vil gerne have Emma til at kigge hans fem slides igennem inden bestyrelsesmødet. Han giver hende plads til at sige nej ved at sige ”Du må bare sige til, hvis du hellere vil gøre det efter frokost” og foreslå at sende udkastet på mail i stedet.}
  \item \textit{Sofie tøver, fordi mødet slutter ret sent, og hun skal tidligt op lørdag morgen. Med ”ellers” foreslår Thomas, at de kan mødes til en hurtig kop kaffe klokken fem først for at drøfte deres spørgsmål inden mødet.}
\end{enumerate}

\textbf{Tekst B 阅读核对答案}
\begin{enumerate}
  \item \da{Forkert}（索菲表示她其实很想去：\textit{Jeg vil egentlig gerne med}，只是担心周五结束太晚影响周六早起）
  \item \da{Han foreslår, at hun bare kan gå i pausen efter første time, hvis hun bliver træt.}
  \item \da{I aften, når hun har tjekket sin kalender derhjemme.}
\end{enumerate}

\textbf{练习 A 答案}
\begin{enumerate}
  \item \textit{Kan du \da{lige} sende udkastet til mig i dag?}（用 \textit{Kan du lige...} 彻底消除命令感）
  \item \textit{Hvad arbejder du \da{egentlig} med nu?}（在主语 \textit{du} 之后插入 \textit{egentlig}）
  \item \textit{\dabold{Ellers} finder vi bare en ny dag.}（\textit{Ellers} 放句首后动词 \textit{finder} 紧跟占第 2 位）
  \item \textit{Du må \da{bare} bruge min computer.}（或 \textit{Brug bare min computer!}，\textit{bare} 传达热情大方）
  \item \textit{Du har \dabold{vel} ikke glemt mødet i morgen?}（主句中 \textit{vel} 放在 \textit{ikke} 前面，句末语调上扬）
  \item \textit{Hvordan går det \dabold{ellers} med dit nye arbejde?}（位于主语 \textit{det} 之后表顺便关切）
\end{enumerate}

\textbf{练习 B 答案}
\begin{enumerate}
  \item \da{lige}（“稍微借用一下”，降低打扰感）
  \item \da{bare}（“放宽心，由我来付吧”）
  \item \dabold{vel}（反意疑问句寻求安慰确认：“下午大概不会下雨吧？”）
  \item \da{egentlig}（点出事实背后的真实情况：“其实来自埃斯比约”）
  \item \dabold{ellers}（提供替代方案：“要不然我们也可以晚上通电话”）
  \item \da{lige}；\da{bare}（第一空用 \textit{lige} 请求顺手扶一下门；第二空用 \textit{bare} 宽慰对方“你完全不用等我”）
\end{enumerate}

\textbf{练习 C 答案}
\begin{enumerate}
  \item \textit{Skal vi lige tage en kop kaffe på caféen?}
  \item \textit{Hvad skal du egentlig lave i weekenden?}
  \item \textit{Hvis du har spørgsmål, kan du bare sende mig en besked.}
  \item \textit{Denne eksamen bliver vel ikke alt for svær?}
  \item \textit{Vi kan mødes i dag; ellers snakker vi bare på torsdag.}
  \item \textit{Hvordan går det ellers med dit projekt?}
\end{enumerate}

\textbf{DU3 Modul 4 高分范文（146 词，不含署名）}\\
\textit{Hej Sofie}

\textit{Hvordan går det \dabold{ellers} med dig og dit nye projekt? Jeg håber, at du ikke har alt for meget om ørerne i denne uge.}

\textit{Jeg skriver \da{egentlig}, fordi jeg skal holde et vigtigt oplæg næste tirsdag, og jeg er stadig lidt usikker på sproget i min konklusion. Jeg ved godt, at du har travlt, men kunne du \da{lige} kigge mit korte udkast igennem? Det tager \dabold{vel} kun ti minutter, og vi kunne måske tage en hurtig kop kaffe i kantinen efter frokost?}

\textit{Hvis din kalender er helt fyldt i morgen, må du \da{bare} sige nej --- det forstår jeg selvfølgelig fuldt ud! \dabold{Ellers} kan jeg også sende teksten på mail, så du \da{bare} læser den, hvis du får et roligt øjeblik før weekenden.}

\textit{Sig \da{bare} til, hvad der passer dig bedst!}

\textit{De bedste hilsner\\
\mbox{[Dit navn]}}

\small\textbf{名师点评}：146 词（不含署名）流畅自然地达到了 Modul 4 私人协商邮件的 B1+/B2- 水准。开篇用 \textit{ellers} 寒暄切入，中段用 \textit{egentlig} 温和引出请求、\textit{lige} 与 \textit{vel} 降低对方负担，最后用 \textit{bare} 和句首倒装的 \textit{Ellers kan jeg...} 留出充分退路，既展现了对情态小品词的纯熟驾驭，又极其贴合丹麦职场平等体贴的沟通礼仪。
\end{answerbox}
\ncdlesson{3}{Fra talesprog til skriftsprog: En nyhedsartikel bliver til}{从口语采访到书面报道}
{DU3 Modul 4（B1+ $\rightarrow$ B2-）｜对标 NCE 3 \& Modultest 4 \textit{Læsning Opgave 3 \& Skrivning Opgave 2 (Debatindlæg)}}
{\textbf{语体转换（Stilskift）三大核心手术}：消除口语左移外置（\textit{Venstreekstraposition}）与填充词（\textit{altså, ligesom, sådan}）· \textbf{动词短语名词化（Nominalisering）} · \textbf{松散并列句升格为定语从句与书面衔接词}}
{平时跟丹麦人聊天时，“这个文化中心嘛，它其实超棒的”这样带复指代词和口头禅的松散句式非常自然；但写 DU3 Modul 4 的观点辩论帖（\textit{Debatindlæg}）或阅读报刊评论时，口语腔会让文章显得松垮低龄。本课带你亲历一场真实的媒体转化：看看记者如何把一段街头采访口语，浓缩成结构紧凑、客观有力的 B2- 新闻报道！}

% ---------------------------------------------------------------------
\begin{foerduleaeser}
带着下面两个问题阅读 Tekst A 和 Tekst B：
\begin{enumerate}
  \item \textit{Hvilke typiske talesprogstræk bruger Lukas i radiointerviewet i Tekst A, når han beskriver det nye kulturhus og prisen på projektet?}（在 Tekst A 的街头采访中，卢卡斯描述新文化中心和项目花费时使用了哪些典型的口语特征？）
  \item \textit{Hvordan har journalisten omformuleret Lukas' udtalelser til skriftsprog i avisartiklen i Tekst B?}（在 Tekst B 的正式新闻报道中，记者是如何将卢卡斯的口头表达改写成紧凑书面语的？）
\end{enumerate}
\end{foerduleaeser}

% ---------------------------------------------------------------------
\ncdsubsec{二、Tekster（口语微调 / 报刊与思辨双语篇）}

\subsection*{Tekst A：\textit{Direkte gadeinterview: ”Det nye kulturhus, det er altså ret vildt!”}（街头采访实录 · 约 185 词）}

\begin{tipbox}{语篇说明}
本篇为地方电台记者玛雅（Maja）在街头采访居民卢卡斯（Lukas）的原声实录。阅读时请用“语体显微镜”锁定\textbf{\color{NCDNavy}深蓝色标记的口头填充词与随性用词}（\textit{altså, ligesom, mega, snakker}）以及\textbf{\color{NCDRed}红色标记的口语左移外置结构}（先抛名词再加多余代词 \textit{det / de}）。
\end{tipbox}

\begin{tcolorbox}[enhanced, breakable, colback=NCDBgBlue, colframe=NCDNavy, coltitle=white, fonttitle=\bfseries, title={街头采访实录｜Maja (Journalist) \& Lukas (Beboer)}, arc=2pt, boxrule=0.8pt, left=10pt, right=10pt, top=7pt, bottom=7pt]
\small
\textbf{Maja}: Hej Lukas! Hvad synes du egentlig om kvarterets nye kulturhus, som åbnede i sidste måned?\\[3pt]
\textbf{Lukas}: Jamen, \dabold{det nye kulturhus, det} er \da{altså} blevet et helt vildt godt sted! Førhen var der jo ingenting her i området, og \dabold{de unge, de} kedede sig ofte om aftenen. Nu kommer folk bare ned og \da{snakker} sammen, fordi de \da{ligesom} gerne vil lære deres naboer at kende.\\[3pt]
\textbf{Maja}: Men projektet har jo været ret dyrt for kommunen. Hvad tænker du om det?\\[3pt]
\textbf{Lukas}: Ja, når man skal renovere sådan en gammel fabriksbygning, så koster det selvfølgelig \da{mega mange} penge. Men min kone og jeg, \dabold{vi} bruger stedet tre gange om ugen til både kor og lektiecafé. Så \da{sådan set} synes jeg helt klart, at pengene er godt givet ud, fordi folk faktisk mødes på tværs af alder.
\end{tcolorbox}

\begin{transbox}
\small
\textbf{\color{NCDTeal}[Tekst A 中文精译]}\\
\textbf{玛雅}：嗨，卢卡斯！对于街区上个月新开张的文化中心，你其实觉得怎么样？\\
\textbf{卢卡斯}：哎呀，那个新文化中心嘛，它可真是变成了一个超级棒的地方！以前咱们这一片儿啥也没有，那些年轻人嘛，他们晚上经常无聊得很。现在大家就是过来聚着聊聊天，因为人们差不多都挺想认识自己邻居的。\\
\textbf{玛雅}：不过这个项目对市政厅来说开销也挺高的。你怎么看这一点？\\
\textbf{卢卡斯}：嗯，当你非得去翻修那样一座老厂房时，那当然得花巨多钱啦。不过我和我太太嘛，我们每周去那儿三次，参加合唱团和作业辅导角。所以从某种意义上来说，我觉得这钱花得完全值，因为大家真的能跨年龄段聚到一起。
\end{transbox}

\subsection*{Tekst B：\textit{Nyt kulturhus styrker det lokale fællesskab}（本地报刊书面报道 · 约 180 词）}

\begin{emailbox}{报刊新闻改写版｜Aarhus Lokalavis \quad Kategori: By og Kultur}
Et nyåbnet kulturhus i den gamle fabriksbygning har på kort tid udviklet sig til et \da{centralt samlingspunkt} for kvarterets beboere. Hvor området tidligere manglede fritidstilbud til især unge, skaber huset nu faste rammer for aktiviteter, \dabold{som} fremmer dialog og socialt samvær på tværs af generationer.

\dabold{Renoveringen af} den historiske bygning har \da{medført betydelige udgifter} for kommunen, hvilket før åbningen skabte politisk debat. \dabold{Ifølge flere lokale beboere} har investeringen imidlertid allerede vist sin værdi. Mange borgere \da{påpeger, at} de faste korøvelser og lektiecaféer trækker folk ud af hjemmene flere gange om ugen.

Byrådet vil nu \da{drøfte}, hvordan lignende projekter kan støtte sammenhængskraften i andre bydele, uden at de kommunale budgetter overskrides. Den samlede \da{vurdering} blandt de adspurgte er dog tydelig: Kulturen skaber ny energi i gadebilledet.
\end{emailbox}

\textbf{Tekst B 阅读与语体对比核对（Rigtigt / Forkert / Kort svar，答案见第六部分）：}
\begin{enumerate}
  \item \textit{Ifølge artiklen er alle politikere helt enige om, at renoveringen var billig.} \underline{\hspace{2.2cm}}
  \item \textit{Hvilket navneord (substantiv) bruger journalisten i Tekst B i stedet for Lukas' sætning ”når man skal renovere sådan en gammel fabriksbygning”?} \underline{\hspace{4.8cm}}
  \item \textit{Hvad skal byrådet nu drøfte på baggrund af erfaringerne med kulturhuset?} \\\underline{\hspace{5.2cm}}
\end{enumerate}

\begin{tipbox}{本课发音与语调要点}
\begin{itemize}
  \item \textbf{口语填充词的快速弱读}：\textit{altså} 在口语中常迅速带过读成 \ipa{ˈalˀsʌ}（甚至略去 \textit{t}）；\textit{ligesom} 读 \ipa{ˈliːsʌm}（\textbf{g 不发音}）；\textit{sådan} 常吞音缩读为 \ipa{ˈsʌn}（\textit{d} 完全消失）。
  \item \textbf{书面高阶词的软 d 与静音辅音}：\textit{betydelig} \ipa{beˈtˢyðˀəli}（中间软 d 带喉塞音，词尾 \textbf{g 不发音}！）；\textit{en udgift} \ipa{ˈuðˌɡefd}（前缀 \textit{ud-} 为软 d）；\textit{ifølge} \ipa{iˈføljə}（\textit{g} 化为滑音 \ipa{j}）。
  \item \textbf{报刊核心名词重音}：\textit{en renovering} \ipa{ʁenoˈveːˀɐŋ}（重音在 \textit{-ve-}）；\textit{et samlingspunkt} \ipa{ˈsɑmleŋsˌpʰɔŋˀd}；\textit{en vurdering} \ipa{vuɐ̯ˈdeːˀɐŋ}；\textit{at drøfte} \ipa{ˈdʁøfdə}。
\end{itemize}
\end{tipbox}

\begin{transbox}
\small
\textbf{\color{NCDTeal}[Tekst B 中文对照]}\\
在老厂房内新近开放的一座文化中心，在短时间内已经发展成为街区居民的核心聚会场所。此前该区域尤其缺乏面向年轻人的业余活动空间，而如今该中心为跨越代际的对话与社交往来提供了固定平台。

对这座历史建筑的翻修为市政厅带来了可观的开支，这在开放前曾引发政治争论。然而根据多位当地居民的反馈，这项投资已经展现出它的价值。许多市民指出，每周定期开展的合唱排练与作业辅导角每周数次把人们吸引出家门。

市议会接下来将商讨如何让类似项目在不超出市政预算的前提下，支持其他街区的社区凝聚力。不过受访者的整体评价非常明确：文化生活正为街景注入新的活力。
\end{transbox}

% ---------------------------------------------------------------------
\ncdsubsec{三、Gloser（核心生词表）}

\begin{tipbox}{词表重点}
书面写作的核心秘诀在于\textbf{“动词名词化”与“准确引述动词”}：把口语里的 \textit{at renovere} 变成 \textit{en renovering}（注意以 \textit{-ing} 结尾的名词一律为通性 \textit{en}！）；用 \textit{at påpege}（指出）和 \textit{at drøfte}（商讨）取代千篇一律的 \textit{at sige / snakke om}。
\end{tipbox}

\begin{glosertable}
名词 & et talesprog & \textit{(talesproget, talesprog, -ene)} & 口语（单复同形） \\
名词 & et skriftsprog & \textit{(skriftsproget, skriftsprog, -ene)} & 书面语（单复同形） \\
名词 & en renovering & \textit{(renoveringen, renoveringer, -erne)} & 翻修，修缮 \\
名词 & en udgift & \textit{(udgiften, udgifter, udgifterne)} & 开支，费用，支出 \\
名词 & et samlingspunkt & \textit{(samlingspunktet, -punkter, -erne)} & 聚集点，核心场所 \\
名词 & en beboer & \textit{(beboeren, beboere, beboerne)} & 居民，住户 \\
名词 & en dialog & \textit{(dialogen, dialoger, dialogerne)} & 对话，交流 \\
名词 & en vurdering & \textit{(vurderingen, vurderinger, -erne)} & 评估，评价 \\
名词 & en sammenhængskraft & \textit{(sammenhængskraften, kun ental)} & 社会凝聚力 \\
\midrule
动词 [har] & at påpege & \textit{(påpeger, påpegede, har påpeget)} & 指出，点明 \\
动词 [har] & at drøfte & \textit{(drøfter, drøftede, har drøftet)} & 探讨，商议（正式） \\
动词 [har] & at medføre & \textit{(medfører, medførte, har medført)} & 带来，导致（后果） \\
动词 [har] & at fremme & \textit{(fremmer, fremmede, har fremmet)} & 促进，推动 \\
动词 [har] & at omformulere & \textit{(omformulerer, -ede, har -et)} & 改写，重新措辞 \\
动词 [har] & at overskride & \textit{(overskrider, overskred, har overskredet)} & 超出，逾越（预算/界限） \\
\midrule
形容词 & betydelig & \textit{(betydeligt, betydelige)} & 相当大的，显著的，可观的 \\
介词 & ifølge & \textit{(不变，介词)} & 根据，依据（后接消息来源） \\
\end{glosertable}

% ---------------------------------------------------------------------
\ncdsubsec{四、Grammatik \& Stilistik（核心语法、语气与修辞精讲）}

\begin{grammatikbox}{语法点 1：从口语到书面语的“四大降噪手术”全景对照}
在 DU3 Modul 4 阅卷标准中，考官极度看重\textbf{语域适配度（\textit{register}）}。请熟记下面四步改写手术，彻底清除作文里的口语“噪音”：

\vspace{0.3em}
\renewcommand{\arraystretch}{1.26}
{\small
\begin{tabularx}{\linewidth}{@{} >{\bfseries\raggedright\arraybackslash}p{2.6cm} >{\raggedright\arraybackslash}p{4.3cm} >{\raggedright\arraybackslash}p{4.4cm} >{\raggedright\arraybackslash}X @{}}
\toprule
\textbf{手术步骤} & \textbf{口语原生态（应避免写入文章）} & \textbf{书面升级策略} & \textbf{B2- 报刊级成句} \\
\midrule
\ding{172} 砍掉左移外置\newline （复指代词） & \textit{\dabold{Det nye kulturhus, det} er blevet populært.} & 删除多余代词 \textit{det / den / de / vi}，让主语直接接动词 & \textit{\da{Det nye kulturhus} er blevet et populært mødested.} \\
\midrule
\ding{173} 剔除填充词与情绪口语词 & \textit{Det koster \dabold{altså mega mange} penge, \dabold{ligesom}.} & 删掉 \textit{altså, ligesom, jo, sådan set}；用精确书面形容词替换 \textit{mega / ret vildt} & \textit{Projektet medfører \da{betydelige} udgifter.} \\
\midrule
\ding{174} 动词短语名词化（浓缩信息） & \textit{\dabold{Når man renoverer} bygningen, tager det tid.} & 将 \textit{når man + V} 浓缩为 \textbf{\textit{-ing / -else} 名词短语}作主语 & \textit{\da{Renoveringen af} bygningen kræver tid.} \\
\midrule
\ding{175} 串联从句与客观引述 & \textit{Folk \dabold{snakker} og siger, \dabold{og så} mødes de...} & 用关系从句（\textit{som / der}）及 \textit{Ifølge / påpege / imidlertid} 重组 & \textit{\da{Ifølge} borgerne styrker huset, \da{som} samler alle, dialogen.} \\
\bottomrule
\end{tabularx}
}
\end{grammatikbox}

\begin{grammatikbox}{语法点 2：书面客观化四招——告别全篇“\textit{Jeg synes}”与“\textit{Folk siger}”}
报刊文章和 \textit{Debatindlæg} 之所以显得说服力强，在于作者懂得把个人感受包装成有依据的公共讨论：
\begin{enumerate}
  \item \textbf{用 \da{Ifølge + [信源]} 开头引出论据（触发 V2 倒装！）}：\textit{\da{Ifølge lokale beboere} \dabold{skaber} kulturhuset ny energi.}（注意：\textit{Ifølge} 后面\textbf{绝不加 \textit{til} 或 \textit{af}}，直接跟人或机构名称！）。
  \item \textbf{四大高分引述动词取代 \textit{sige}}：\textit{påpeger}（指出）、\textit{fremhæver}（强调突出）、\textit{vurderer}（评估认为）、\textit{argumenterer for, at...}（论证主张……）。
  \item \textbf{动词名词化的介词黄金搭档 \textit{af}}：将及物动词变成名词后，原宾语必须由介词 \dabold{af} 引出：
    \begin{itemize}
      \item \textit{at renovere bygningen} $\rightarrow$ \textit{renoveringen \dabold{af} bygningen}（对建筑的翻修）
      \item \textit{at lukke biblioteket} $\rightarrow$ \textit{lukningen \dabold{af} biblioteket}（对图书馆的关闭）
    \end{itemize}
  \item \textbf{严防“中/英式左移外置”回潮}：写作复查时，用笔划掉所有类似 \textit{”Min nabo, han mener...”} 或 \textit{”Politikerne, de glemmer...”} 中的多余代词 \textit{han / de / det}！
\end{enumerate}
\end{grammatikbox}

% ---------------------------------------------------------------------
\ncdsubsec{五、Mønsterøvelser（句型与修辞操练）}

\textbf{练习 A：语体升级手术——将下列口语采访句改写为紧凑的书面报道句（消除多余代词与口头禅，并按提示浓缩）}
\begin{enumerate}
  \item \textit{Det nye bibliotek, det ligger altså lige ved stationen.}（消除左移外置与 \textit{altså}）$\rightarrow$
  \item \textit{De unge i byen, de mangler jo ligesom et sted om aftenen.}（消除左移外置与填充词）$\rightarrow$
  \item \textit{Når man renoverer den gamle skole, så koster det mega mange penge.}（将从句名词化为 \textit{Renoveringen af...}，用 \textit{medføre betydelige udgifter}）$\rightarrow$
  \item \textit{Mange naboer siger, at støjen fra gaden er et ret vildt problem.}（用 \textit{Ifølge...} 开头触发倒装，将 \textit{ret vildt} 升级为 \textit{betydeligt}）$\rightarrow$
  \item \textit{Folk kommer hen og snakker sammen, fordi de gerne vil hjælpe hinanden.}（改用书面动词 \textit{drøfter} 或名词 \textit{dialog}，并用关系从句连接）$\rightarrow$
  \item \textit{Projektet var dyrt, men det har jo sådan set givet et godt resultat.}（删去口语词并用 \textit{imidlertid} 替代或升级转折）$\rightarrow$
\end{enumerate}

\textbf{练习 B：书面语精准选词填空（从括号内选出最契合报刊书面语体的词汇）}
\begin{enumerate}
  \item Projektet har desværre medført \underline{\hspace{2.2cm}} (\textit{betydelige / mega store}) udgifter for kommunen.
  \item \underline{\hspace{2.2cm}} (\textit{Ifølge / Efter hvad}) flere lokale beboere er det nye torv blevet et stort aktiv for bydelen.
  \item Flere eksperter \underline{\hspace{2.2cm}} (\textit{påpeger / snakker om}), at det er vigtigt at investere i børns fritidsliv.
  \item \underline{\hspace{2.2cm}} (\textit{Renoveringen / Når man fixer}) af den gamle bygning begynder allerede næste måned.
  \item Byrådet skal på næste møde \underline{\hspace{2.2cm}} (\textit{drøfte / snakke lidt om}) fremtidens busruter.
  \item Den øgede trafik gennem midtbyen kan \underline{\hspace{2.2cm}} (\textit{medføre / gøre sådan at}) længere ventetid for cyklisterne.
\end{enumerate}

\textbf{练习 C：中译丹实战（练习紧凑书面报道与客观评述句式）}
\begin{enumerate}
  \item 这座新文化中心已成为居民们的核心聚集点。
  \item 根据多位专家的说法，这项方案将带来可观的开支。
  \item 对这座老建筑的翻修促进了各代际之间的对话。
  \item 市议会下周将商讨地方图书馆的未来。
  \item 多位市民指出，该街区缺乏面向年轻人的休闲活动项目。
  \item 然而，大多数受访者认为这笔投资是值得的。
\end{enumerate}

% ---------------------------------------------------------------------
\ncdsubsec{六、DU3 Modul 4 实战微任务与参考答案}

\begin{emailbox}{DU3 Modul 4 写作任务（对标 Skrivning Opgave 2 · Debatindlæg $\ge$ 100 ord）}
\textbf{Situation}: Kommunen overvejer at lukke det lokale bibliotek om aftenen for at spare penge. På et borgermøde sagde en nabo: \textit{”Det lokale bibliotek, det er altså super vigtigt! Hvis de lukker det om aftenen, så kan folk, der arbejder, jo ligesom slet ikke nå at låne bøger eller mødes der.”}\\
\textbf{Opgave}: Skriv et sagligt og overbevisende \textbf{debatindlæg} til lokalavisen på \textbf{mindst 100 ord} (gerne 120--145 ord). Omformuler naboens mundtlige pointe til klart skriftsprog uden fyldord og venstreekstraposition, og argumenter for, hvorfor aftenåbning bør bevares (brug gerne mindst én nominalisering med \textit{-ing}, ét relativled med \textit{som / der} og ét objektivt udtryk som \textit{Ifølge / påpege / medføre}).
\end{emailbox}

\begin{answerbox}
\textbf{Før du læser 参考回答}
\begin{enumerate}
  \item \textit{Lukas bruger venstreekstraposition med ekstra pronominer (”det nye kulturhus, det...”, ”de unge, de...”, ”min kone og jeg, vi...”), fyldord som ”altså”, ”ligesom” og ”sådan set” samt mundtlige udtryk som ”ret vildt” og ”mega mange penge”.}
  \item \textit{Journalisten fjerner de ekstra pronominer og fyldord, erstatter hverdagssnak med præcise ord som ”samlingspunkt”, ”betydelige udgifter” og ”påpeger”, og samler løse sætninger til nominaliseringer som ”Renoveringen af den historiske bygning” og relativsætninger.}
\end{enumerate}

\textbf{Tekst B 阅读与语体对比核对答案}
\begin{enumerate}
  \item \da{Forkert}（原文指出翻修开支在开放前引发了政治争论：\textit{skabte politisk debat}）
  \item \da{Renoveringen}（以动词名词化形式 \textit{Renoveringen af den historiske bygning} 浓缩了整句口语从句）
  \item \da{Hvordan lignende projekter kan støtte sammenhængskraften i andre bydele uden at overskride budgetterne.}
\end{enumerate}

\textbf{练习 A 答案（参考）}
\begin{enumerate}
  \item \textit{Det nye bibliotek ligger tæt på stationen.}（删去多余复指代词 \textit{det} 与口头禅 \textit{altså}）
  \item \textit{Byens unge mangler et samlingssted om aftenen.}（删去复指代词 \textit{de} 以及 \textit{jo, ligesom}）
  \item \textit{Renoveringen af den gamle skole medfører betydelige udgifter.}（将时间从句浓缩为名词短语作主语）
  \item \textit{Ifølge mange naboer udgør støjen fra gaden et betydeligt problem.}（\textit{Ifølge...} 放句首触发动词倒装，\textit{et betydeligt problem} 取代口语词）
  \item \textit{Borgerne mødes for at styrke dialogen med andre, som ønsker at hjælpe hinanden.}（改用书面名词与关系从句衔接）
  \item \textit{Projektet var dyrt, men har imidlertid givet et positivt resultat.}（删去口语填充词并用书面转折副词 \textit{imidlertid}）
\end{enumerate}

\textbf{练习 B 答案}
\begin{enumerate}
  \item \da{betydelige}（客观书面形容词，“可观的”，替换口语词 \textit{mega store}）
  \item \da{Ifølge}（书面引述介词，不加多余词直接接来源）
  \item \da{påpeger}（准确引述观点“指出”，替代泛化口语词）
  \item \da{Renoveringen}（书面名词短语 \textit{Renoveringen af...}，替换冗长从句）
  \item \da{drøfte}（正式商讨决策，替代 \textit{snakke lidt om}）
  \item \da{medføre}（正式书面因果动词“导致、带来”）
\end{enumerate}

\textbf{练习 C 答案}
\begin{enumerate}
  \item \textit{Det nye kulturhus er blevet et centralt samlingspunkt for beboerne.}
  \item \textit{Ifølge flere eksperter vil dette forslag medføre betydelige udgifter.}
  \item \textit{Renoveringen af den gamle bygning har fremmet dialogen mellem generationerne.}
  \item \textit{Byrådet skal drøfte det lokale biblioteks fremtid i næste uge.}
  \item \textit{Flere borgere påpeger, at kvarteret mangler fritidstilbud til de unge.}
  \item \textit{Imidlertid vurderer de fleste adspurgte, at investeringen er pengene værd.}
\end{enumerate}

\textbf{DU3 Modul 4 高分范文（128 词）}\\
\textit{\textbf{Bevar bibliotekets aftenåbning for hele lokalsamfundet}}

\textit{Kommunen overvejer i øjeblikket at afskaffe aftenåbningen på vores lokale bibliotek for at reducere udgifterne. Besparelsen kan på kort sigt virke fornuftig, men lukningen om aftenen vil \da{medføre} alvorlige konsekvenser for mange borgere.}

\textit{Et bibliotek er i dag langt mere end blot et sted, hvor man låner bøger; det fungerer som et \da{centralt samlingspunkt}, \dabold{der} styrker både integration og sammenhængskraft. \dabold{Ifølge} mange fuldtidsarbejdende beboere er aftentimerne den eneste reelle mulighed for at benytte lektiehjælp, læsesale og lokale arrangementer. En begrænsning af åbningstiden vil derfor ramme både børnefamilier og unge, \dabold{som} har brug for et roligt studiemiljø.}

\textit{Byrådet bør i stedet \da{drøfte} alternative løsninger, eksempelvis selvbetjent adgang efter klokken atten, så biblioteket fortsat kan skabe værdi i lokalområdet.}

\small\textbf{名师点评}：128 词（含标题）精准展示 B2- 级别 \textit{Debatindlæg} 的高分骨架。全篇彻底摒弃口语填充词与左移外置结构，运用 \textit{lukningen om aftenen} 与 \textit{En begrænsning af åbningstiden} 两处高密度名词化短语，并结合 \textit{Ifølge...}（触发 V2 倒装）、\textit{medføre}、\textit{drøfte} 及限制性/非限制性关系从句，结尾更提出具体建设性替代方案（\textit{selvbetjent adgang}），说服力极强！
\end{answerbox}
\ncdlesson{4}{Kunsten at være uenig på en høflig måde}{礼貌表达不同意见的艺术}
{DU3 Modul 4（B1+ $\rightarrow$ B2-）｜对标 NCE 3 \& Modultest 4 \textit{Skrivning Opgave 2 (Debatindlæg) \& Mundtlig Opgave 2 (Paropgave)}}
{\textbf{让步与反驳高阶句式}（\textit{Ganske vist..., men ikke desto mindre... / Det kan godt være, at..., men... / Selvom..., mener jeg dog...}）· 让步连接词与副词的\textbf{主从句语序与 V2 倒装} · 丹麦式“先认可、再温和反驳”的共识协商修辞}
{在丹麦职场讨论与 DU3 Modul 4 口试、写作中，直接甩出一句 \textit{”Det er jeg helt uenig i!”} 或 \textit{”Du tager fejl!”} 往往显得生硬且具攻击性。丹麦社会高度重视\textbf{共识文化（\textit{konsensuskultur}）}与平实客气的讨论氛围：高分表达讲究“先退一步承认对方合理之处，再稳健有力地提出保留或反驳”。吃透本课的 \textbf{\textit{Ganske vist..., men ikke desto mindre...}} 与 \textbf{\textit{Det kan godt være, at..., men...}}，你的不同意见就能既有礼貌又有锋芒！}

% ---------------------------------------------------------------------
\begin{foerduleaeser}
带着下面两个问题阅读 Tekst A 和 Tekst B：
\begin{enumerate}
  \item \textit{Hvorfor oplevede Marco det først som forvirrende, når hans danske kolleger var uenige på møder i Tekst A, og hvordan formulerer de i stedet saglig kritik?}（在 Tekst A 中，为什么马尔科最初觉得丹麦同事在会上表达不同意见的方式令人困惑？他们实际是如何提出客观批评的？）
  \item \textit{Hvad er Henrik og Sofie henholdsvis enige og uenige om i debatten om hjemmearbejdsdage i Tekst B?}（在 Tekst B 关于居家办公日的讨论中，亨里克和索菲各自赞同什么、又在哪些方面存在分歧？）
\end{enumerate}
\end{foerduleaeser}

% ---------------------------------------------------------------------
\ncdsubsec{二、Tekster（口语微调 / 报刊与思辨双语篇）}

\subsection*{Tekst A：\textit{Når et ”nej” pakkes ind i høflighed}（职场文化随笔 · 约 220 词）}

\begin{tipbox}{语篇说明}
本篇以意大利籍工程师马尔科（Marco）初入丹麦职场的经历切入，剖析丹麦人如何在保留体面与合作氛围的同时表达反对意见。阅读时请重点标出\textbf{\color{NCDNavy}深蓝色加粗的让步引导语}（先认可对方部分道理）与\textbf{\color{NCDRed}深红色加粗的反驳转折接口}（注意 \textit{Ganske vist} 与 \textit{ikke desto mindre} 触发的 V2 倒装！）。
\end{tipbox}

Da Marco begyndte som ingeniør i en virksomhed i Odense, troede han i lang tid, at alle på afdelingen var helt enige om næsten alt. På projektmøderne hørte han nemlig sjældent nogen sige direkte, at et forslag var dårligt. Først efter nogle måneder opdagede han, at hans danske kolleger ofte udtrykte skarp uenighed gennem faste høflige mønstre.

Når en kollega for eksempel foreslog en billig løsning på et teknisk problem, svarede Emma roligt: ”\da{Det kan godt være, at} den model er billigere her og nu, \dabold{men} på lang sigt kræver den langt mere vedligeholdelse.” Dermed anerkendte hun først kollegaens pointe, før hun venligt pegede på svagheden.

Senere lærte Marco, at man i skriftlige vurderinger og formelle debatter ofte bruger endnu stærkere strukturer til at skabe balance. En teamleder forklarede ham princippet med et klassisk eksempel: ”\da{Ganske vist} sparer vi penge i år, \dabold{men ikke desto mindre} risikerer vi at miste vigtige kunder næste år.” Her viser \textit{ganske vist}, at man tager modpartens argument alvorligt, mens \textit{ikke desto mindre} fastholder den afgørende indvending.

I dansk arbejdskultur handler det nemlig om at adskille person og sag. \da{Selvom} man kan være dybt uenig i en konklusion, \dabold{bør} tonen \dabold{dog} forblive respektfuld og saglig. Når man mestrer denne balance mellem indrømmelse og modargument, bliver man taget langt mere alvorligt end den, der blot råber højest.

\begin{transbox}
\small
\textbf{\color{NCDTeal}[Tekst A 中文精译]}\\
当马尔科开始在欧登塞的一家企业担任工程师时，有很长一段时间他都以为部门里所有人在几乎所有事情上都完全一致。因为在项目会议上，他很少听到有人直接说某个方案很糟糕。直到几个月后他才发现，丹麦同事其实常常通过固定的礼貌句式来表达鲜明的反对意见。

例如，当一位同事针对某个技术难题提出了一个廉价方案时，埃玛平静地回应道：“那个模型眼下确实可能更便宜，但从长远来看，它需要多得多的维护成本。”她借此先认可了同事的着眼点，然后才温和地指出其中的薄弱之处。

后来马尔科了解到，在书面评估和正式辩论中，人们还常常使用更有力度的结构来营造平衡感。一位团队负责人用一个经典例句向他解释了这一原则：“我们今年诚然省下了资金，然而尽管如此，明年我们却面临着失去重要客户的风险。”在这里，\textit{ganske vist} 表明自己认真对待了对方的论据，而 \textit{ikke desto mindre} 则坚定地守住了关键的反驳立场。

因为在丹麦职场文化中，关键在于把“人”和“事”分开。尽管大家在某项结论上可能存在深刻分歧，语气却仍然应当保持尊重且就事论事。一旦掌握了这种在让步与反论之间游刃有余的平衡艺术，你的意见就会比一味嗓门最大的人受到重视得多。
\end{transbox}

\subsection*{Tekst B：\textit{Debat på intranettet: Skal alle tilbage på kontoret?}（企业内网正反讨论 · 约 125 词）}

\begin{emailbox}{职场论坛讨论｜Henrik (Projektleder) \& Sofie (Analytiker)}
\textbf{Henrik:} Kære kolleger. Jeg synes, vi bør reducere antallet af hjemmearbejdsdage fra to til én om ugen. \da{Det kan godt være, at} mange arbejder koncentreret hjemme ved skrivebordet, \dabold{men} vi mister alt for meget af den spontane sparring ved kaffemaskinen. \da{Ganske vist} fungerer onlinemøder fint til faste opgaver, \dabold{men ikke desto mindre} svækkes vores fællesskab i hverdagen.

\textbf{Sofie:} Tak for dit oplæg, Henrik. Jeg forstår godt din bekymring for holdånden, og jeg er enig med dig i, at fysiske møder er værdifulde. \da{Selvom} uformelle snakke på gangen er vigtige, \dabold{mener} jeg \dabold{dog}, at to hjemmedage stadig bør bevares. \da{Ganske vist} er vi færre på kontoret om fredagen, \dabold{men til gengæld} sparer mange pendlere både tid og stress, hvilket gør os mere effektive resten af ugen.
\end{emailbox}

\textbf{Tekst B 阅读核对（Rigtigt / Forkert / Kort svar，答案见第六部分）：}
\begin{enumerate}
  \item \textit{Henrik mener slet ikke, at onlinemøder fungerer til nogen arbejdsopgaver.} \underline{\hspace{2.2cm}}
  \item \textit{Hvad er Henrik bekymret for, at virksomheden mister med to hjemmearbejdsdage?} \underline{\hspace{2.2cm}}
  \item \textit{Hvad peger Sofie på som den konkrete fordel ved, at medarbejdere arbejder hjemme to dage om ugen?} \underline{\hspace{2.2cm}}
\end{enumerate}

\begin{tipbox}{本课发音与语调要点}
\begin{itemize}
  \item \textbf{高阶让步连词与副词读音}：\textit{ganske vist} \ipa{ˈɡansɡə ˈvesd}（\textit{ganske} 词尾轻读 \ipa{ɡə}）；\textit{ikke desto mindre} \ipa{ˈeɡə ˈdɛsdo ˈmenˀdʁə}（重音落在 \textit{ikke}、\textit{des-} 与 \textit{min-} 上，带喉塞音）；\textit{alligevel} \ipa{ɑˈliːəˌvɛlˀ}（中间的 \textit{g} 完全不发音！）；\textit{selvom} \ipa{ˈsɛlˀˌʌm}（注意词尾 \textit{v} 在口语中常轻读或接近脱落）。
  \item \textbf{立场与争论高频词的软 \textit{d} 与静音字母}：\textit{uenig} \ipa{uˈeːˀni}（前缀 \textit{u-} 重读或次重读，词尾 \textit{g} 不发音！）；\textit{vedligeholdelse} \ipa{veðˈliːəˌhʌlˀəlsə}（第一个 \textit{d} 为软 \textit{d}，\textit{g} 与第二个 \textit{d} 均不发音）；\textit{at anerkende} \ipa{ˈanɐˌkʰɛnˀə}（\textit{-nd-} 中的 \textit{d} 不发音！）；\textit{saglig} \ipa{ˈsɑʊ̯li}（\textit{g} 滑音化，词尾 \textit{g} 不发音）。
  \item \textbf{礼貌语调提示}：说 \textit{Det kan godt være, at...} 时，重音轻轻放在 \textit{godt} 上，尾音微微下沉留顿，再平稳接出 \textit{men...}，听起来诚恳而不带挑衅感。
\end{itemize}
\end{tipbox}

\begin{transbox}
\small
\textbf{\color{NCDTeal}[Tekst B 中文对照]}\\
\textbf{亨里克（项目经理）}：各位同事好。我认为我们应当把每周居家办公的天数从两天减到一天。诚然，许多人在家里的书桌前能非常专注地工作，但我们在茶水间丧失了太多自发的灵感碰撞。虽然线上会议对于处理固定任务确实运作良好，然而尽管如此，我们日常的团队凝聚力却正在被削弱。

\textbf{索菲（分析师）}：谢谢你的提议，亨里克。我完全理解你对团队精神的担忧，我也同意你所说的面对面会议很有价值。尽管走廊里的非正式交谈很重要，但我依然认为应当保留两天居家办公。诚然我们周五在办公室的人少了一些，但作为回报，许多通勤员工既省下了时间又减轻了压力，这使得大家在一周其余时间里效率更高。
\end{transbox}

% ---------------------------------------------------------------------
\ncdsubsec{三、Gloser（核心生词表）}

\begin{tipbox}{词表重点}
表达“同意/不同意”时，介词搭配非常固定：和某人意见一致/不合用 \textit{med}（\textit{uenig med Henrik}），在具体观点上同意/不同意用 \textit{i}（\textit{uenig i forslaget / uenig i, at...}）。此外，名词化词 \textit{en vedligeholdelse}（维护）和 \textit{en anerkendelse}（认可）均为通性 \textit{en} 名词，动词 \textit{at modsige}（反驳）为不规则动词（随 \textit{sige} 变化为 \textit{modsagde, har modsagt}）。
\end{tipbox}

\begin{glosertable}
副词短语 & ganske vist & \textit{(不变，让步副词短语)} & 诚然，固然，确实如此 \\
副词短语 & ikke desto mindre & \textit{(不变，转折副词短语)} & 尽管如此，然而依然 \\
连词 & selvom & \textit{(不变，让步从属连词)} & 尽管，虽然（后接从句语序） \\
副词 & alligevel & \textit{(不变，让步副词)} & 尽管如此，还是，仍然 \\
副词 & dog & \textit{(不变，句中转折副词)} & 不过，然而（多置于动词后） \\
\midrule
名词 & en uenighed & \textit{(uenigheden, uenigheder, -erne)} & 分歧，意见不合 \\
名词 & en konsensus & \textit{(konsensussen, ingen flertal)} & 共识，一致意见 \\
名词 & en pointe & \textit{(pointen, pointer, pointerne)} & 要点，核心论点 \\
名词 & et modargument & \textit{(modargumentet, -argumenter, -erne)} & 反论据，反驳理由 \\
名词 & en sparring & \textit{(sparringen, ingen flertal)} & 业务切磋，思想碰撞 \\
名词 & en vedligeholdelse & \textit{(vedligeholdelsen, ingen flertal)} & 维护，保养 \\
名词 & en høflighed & \textit{(høfligheden, ingen flertal)} & 礼貌，客气 \\
\midrule
动词 [har] & at anerkende & \textit{(anerkender, anerkendte, har anerkendt)} & 认可，承认（对方合理之处） \\
动词 [har] & at indrømme & \textit{(indrømmer, indrømmede, har indrømmet)} & 承认，作出让步 \\
动词 [har] & at modsige & \textit{(modsiger, modsagde, har modsagt)} & 反驳，反口，相矛盾 \\
动词 [har] & at fastholde & \textit{(fastholder, fastholdt, har fastholdt)} & 坚持（观点、原则） \\
\midrule
形容词 & uenig & \textit{(uenigt, uenige)} & 不同意的（与某人 \textit{med} / 对某事 \textit{i}） \\
形容词 & saglig & \textit{(sagligt, saglige)} & 客观务实的，就事论事的 \\
\end{glosertable}

% ---------------------------------------------------------------------
\ncdsubsec{四、Grammatik \& Stilistik（核心语法、语气与修辞精讲）}

\begin{grammatikbox}{语法点 1：让步与反驳四大高阶句式与语序全景对照}
在 DU3 Modul 4 写作（尤其是 \textit{Debatindlæg}）和口语双人辩论中，“先让步再反驳”是最具说服力的修辞骨架。请严格区分以下四种句式的\textbf{正式度}与\textbf{V2 倒装位置}：

\vspace{0.3em}
\renewcommand{\arraystretch}{1.24}
{\small
\begin{tabularx}{\linewidth}{@{} >{\bfseries\raggedright\arraybackslash}p{3.1cm} >{\raggedright\arraybackslash}p{4.3cm} >{\raggedright\arraybackslash}X @{}}
\toprule
\textbf{让步反驳句式} & \textbf{语域与语序铁律} & \textbf{标准政论与协商例句} \\
\midrule
\da{Ganske vist...,}\newline \da{men ikke desto mindre...} &
\textbf{书面正式高阶}（报刊/辩论首选）\newline
\textit{Ganske vist} 放句首必 \textbf{V2 倒装}；\textit{men} 后紧接 \textit{ikke desto mindre} 也必 \textbf{V2 倒装}！ &
\textit{\da{Ganske vist} \dabold{sparer vi} penge nu, \da{men ikke desto mindre} \dabold{risikerer vi} store fejl senere.}\newline
（也可分两句：\textit{... \da{Ikke desto mindre} \dabold{mangler} vi tal.}） \\
\midrule
\da{Det kan godt være, at..., men...} &
\textbf{口语与半正式通用}（委婉留余地）\newline
\textit{at} 后守\textbf{从句语序}（\textit{ikke} 前移）；\textit{men} 后接正常主句 &
\textit{\da{Det kan godt være, at} planen \dabold{ikke er} perfekt, \da{men} den løser problemet hurtigt.} \\
\midrule
\da{Selvom...,}\newline \da{... dog / alligevel...} &
\textbf{标准主从让步句}\newline
\textit{Selvom} 内部不倒装；前置时后半主句必须 \textbf{V2 倒装}，且\textbf{绝不加 \textit{men}}！ &
\textit{\da{Selvom} møder tager tid, \dabold{mener} jeg \da{dog}, at dialogen er vigtig.}\newline
（\textit{dog} 放在主句动词或主语后） \\
\midrule
\da{Jeg forstår godt..., men...}\newline \da{Jeg er kun delvist enig i, at...} &
\textbf{高情商开篇表态句}\newline
专用于回应对方来信或辩论贴的第一句，立刻奠定客观礼貌基调 &
\textit{\da{Jeg er kun delvist enig i}, at alle bør møde på kontoret hver dag.} \\
\bottomrule
\end{tabularx}
}
\vspace{0.3em}

\textbf{语序灵活微调提示}：\ding{172} \textit{ganske vist} 和 \textit{ikke desto mindre} 也可以从句首移到主句动词（或主语）后面，以让语气更紧凑自然：\textit{Forslaget er \da{ganske vist} dyrt, men det har \da{ikke desto mindre} stor værdi.} \ding{173} 想在反驳句中突出“作为补偿/相反另一方面”，还可用 \da{men til gengæld...} 或 \da{men på den anden side...}（放句首或 \textit{men} 后均须 V2 倒装）。
\end{grammatikbox}

\begin{grammatikbox}{语法点 2：中英对比排雷——让步反驳最易踩中的四大扣分陷阱}
\begin{enumerate}
  \item \textbf{严禁把“虽然……但是……”直译成 \textit{Selvom..., men...}！}：这是中文母语者最顽固的母语迁移错误。在丹麦语中，\textit{Selvom}（从属连词）和 \textit{men}（并列连词）\textbf{不能同时出现在同一主从句里}！\newline
    \ding{55}\ \textit{Selvom det er dyrt, men vi køber det alligevel.}\newline
    \ding{51}\ \textit{Selvom det er dyrt, \dabold{køber} vi det \da{alligevel} / \da{dog}.}（后半主句直接动词倒装！）
  \item \textbf{\textit{Ganske vist} 与 \textit{Ikke desto mindre} 放句首绝不打英语逗号！}：英语习惯写 \textit{Admittedly, ... / Nevertheless, ...}；丹麦语绝不加逗号，且动词必须紧随其后占据第 2 位：\ding{51}\ \textit{\da{Ganske vist} \dabold{er} mødet langt, men \da{ikke desto mindre} \dabold{bør} vi deltage.}（绝对不能写成 \ding{55}\ \textit{Ganske vist, mødet er langt...}）。
  \item \textbf{记住 \textit{Det kan godt være, at...} 从句内的否定词位置}：在 \textit{at} 引导的从句里，\textit{ikke / aldrig} 必须跑到变位动词前面：\ding{51}\ \textit{Det kan godt være, at han \dabold{ikke har} tid i dag, men...}（不能写成 \ding{55}\ \textit{...at han har ikke tid}）。
  \item \textbf{分清 \textit{enig / uenig med}（对人）与 \textit{enig / uenig i}（对观点/从句）}：\newline
    \ding{51}\ \textit{Jeg er uenig \dabold{med} dig \dabold{i}, at forbuddet hjælper.}（我和你在“禁令有效”这件事上意见不合——先 \textit{med} 人，再 \textit{i} 事！）。
\end{enumerate}
\end{grammatikbox}

% ---------------------------------------------------------------------
\ncdsubsec{五、Mønsterøvelser（句型与修辞操练）}

\textbf{练习 A：高阶让步反驳精准填空（每空填入唯一准确的 1 个丹麦语词）}
\begin{enumerate}
  \item \underline{\hspace{2.2cm}} vist er hjemmearbejde meget fleksibelt, men vi mister noget af den spontane sparring.
  \item Forslaget lyder spændende, men ikke \underline{\hspace{2.2cm}} mindre må vi først vurdere de samlede omkostninger.
  \item Jeg er desværre uenig \underline{\hspace{2.2cm}} min chef i, at alle møder nødvendigvis skal være fysiske.
  \item Vi er kun delvist enige \underline{\hspace{2.2cm}}, at den nye arbejdsform styrker trivslen for alle medarbejdere.
  \item Det kan godt \underline{\hspace{2.2cm}}, at løsningen tager lidt tid at indføre, men den holder til gengæld længe.
  \item \underline{\hspace{2.2cm}} mange foretrækker at arbejde hjemme om fredagen, mener ledelsen dog, at kontoret er vigtigt.
  \item Argumentet har en god pointe. Ikke desto mindre \underline{\hspace{2.2cm}} planen stadig flere alvorlige svagheder. (rumme，现在时)
\end{enumerate}

\textbf{练习 B：让步反驳句型升级与语病修正（按括号要求改写，或修正中/英式常见错误）}
\begin{enumerate}
  \item 用 \textit{Ganske vist..., men ikke desto mindre...} 合并两句：\textit{Kaffemaskinen er dyr. Den skaber et bedre arbejdsmiljø.} $\rightarrow$
  \item 消除“虽然……但是”叠加病句：\textit{Selvom projektet kræver mange timer, men jeg synes, vi bør gennemføre det.} $\rightarrow$
  \item 用 \textit{Det kan godt være, at..., men...} 礼貌反驳下方生硬句子：\textit{Din idé er dårlig, for den koster alt for meget.}（先承认“这主意有创意”）$\rightarrow$
  \item 修正句首逗号与倒装错误：\textit{Ganske vist, planen har sine fordele. Ikke desto mindre, vi bør vente med beslutningen.} $\rightarrow$
  \item 修正从句语序与介词搭配：\textit{Jeg er uenig på dig om, at man kan ikke arbejde effektivt hjemmefra.} $\rightarrow$
  \item 用 \textit{Jeg er kun delvist enig i, at...} 温和改写：\textit{Jeg er helt uenig med, at onlinemøder altid er spild af tid.} $\rightarrow$
\end{enumerate}

\textbf{练习 C：中译丹实战（活用礼貌让步反驳句式与精准语序）}
\begin{enumerate}
  \item 我部分同意你的观点，但我们依然需要更多数据。
  \item 这项计划诚然很昂贵，然而尽管如此它却极具长远价值。
  \item 也许他说得没错，不过我认为我们还是应该等待。
  \item 尽管大家意见不合，团队还是保持着礼貌而务实的语气。
  \item 在数字化方案这个问题上，我与同事的看法并不一致。
  \item 尽管线上会议省时，它们却无法完全替代面对面的交流。
\end{enumerate}

% ---------------------------------------------------------------------
\ncdsubsec{六、DU3 Modul 4 实战微任务与参考答案}

\begin{emailbox}{DU3 Modul 4 Skrivning Opgave 2 实战微任务（Debatindlæg · minimum 100 ord）}
\textbf{Situation}: I lokalavisen har en læser skrevet et skarpt debatindlæg om moderne kontorer:\\
\textit{”Åbne storrumskontorer bør afskaffes med det samme! Man kan slet ikke koncentrere sig på grund af støj fra kolleger og telefoner. Alle bør i stedet have deres eget lukkede kontor.”}\\[3pt]
\textbf{Opgave}: Skriv et høfligt og sagligt \textbf{debatindlæg} på \textbf{mindst 100 ord} (gerne 130--155 ord), hvor du forholder dig til indlægget. Anerkend først læserens pointe, forklar derefter, hvorfor et totalt forbud mod storrumskontorer ikke nødvendigvis er den bedste løsning, og brug mindst \textbf{to} indrømmelsesstrukturer fra denne lektion (\textit{Ganske vist..., men ikke desto mindre... / Det kan godt være, at..., men... / Selvom...}).
\end{emailbox}

\begin{answerbox}
\textbf{Før du læser 参考回答}
\begin{enumerate}
  \item \textit{Marco blev først forvirret, fordi han sjældent hørte nogen sige direkte på møderne, at et forslag var dårligt. Hans danske kolleger formulerer i stedet saglig kritik ved først at anerkende modpartens pointe med udtryk som ”Det kan godt være, at..., men...” eller ”Ganske vist..., men ikke desto mindre...”.}
  \item \textit{Henrik og Sofie er enige om, at fysiske møder er værdifulde. Men Henrik vil skære ned til én hjemmearbejdsdag om ugen, fordi han frygter, at man mister spontan sparring og fællesskab, mens Sofie mener, at to hjemmedage bør bevares, da pendlere sparer tid og stress og dermed arbejder mere effektivt.}
\end{enumerate}

\textbf{Tekst B 阅读核对答案}
\begin{enumerate}
  \item \da{Forkert}（亨里克明确承认线上会议用于固定任务效果很好：\textit{Ganske vist fungerer onlinemøder fint til faste opgaver}）
  \item \da{Den spontane sparring ved kaffemaskinen og fællesskabet i hverdagen.}（他担心失去茶水间自发的业务交流以及日常团队凝聚力）
  \item \da{Pendlere sparer både tid og stress, hvilket gør medarbejderne mere effektive resten af ugen.}（通勤员工节省时间和减轻压力，从而提高一周其余时间的效率）
\end{enumerate}

\textbf{练习 A 答案}
\begin{enumerate}
  \item \da{Ganske}（固定搭配：\textit{Ganske vist...}“诚然……”）
  \item \da{desto}（固定转折短语：\textit{ikke desto mindre}“然而尽管如此”）
  \item \da{med}（表示“与某人意见不合”必须用介词 \textit{med}：\textit{uenig med min chef}）
  \item \da{i}（在具体观点或从句前表示赞同/不赞同，固定用介词 \textit{i}）
  \item \da{være}（委婉让步开篇句式：\textit{Det kan godt være, at...}）
  \item \da{Selvom}（引导让步从句，注意后半主句不再加 \textit{men}，而用 \textit{dog}）
  \item \da{rummer}（\textit{Ikke desto mindre} 放句首触发 V2 倒装，紧接现在时动词 \textit{rummer}）
\end{enumerate}

\textbf{练习 B 答案（参考）}
\begin{enumerate}
  \item \textit{\da{Ganske vist} er kaffemaskinen dyr, \dabold{men ikke desto mindre} skaber den et bedre arbejdsmiljø.}（注意两处 V2 倒装：\textit{er kaffemaskinen} 与 \textit{skaber den}）
  \item \textit{\da{Selvom} projektet kræver mange timer, synes jeg \dabold{dog}, at vi bør gennemføre det.}（删掉多余的 \textit{men}，用倒装及句中 \textit{dog} 或 \textit{alligevel} 呼应）
  \item \textit{\da{Det kan godt være, at} din idé er kreativ, \dabold{men} den koster desværre alt for meget.}（先肯定对方创意，再委婉提出费用问题）
  \item \textit{\da{Ganske vist har} planen sine fordele. \dabold{Ikke desto mindre bør} vi vente med beslutningen.}（去掉英语式逗号，并保证句首副词后立刻接变位动词）
  \item \textit{Jeg er uenig \dabold{med} dig \dabold{i}, at man \dabold{ikke kan} arbejde effektivt hjemmefra.}（修正介词 \textit{med / i}，并将从句内的 \textit{ikke} 移到 \textit{kan} 前面）
  \item \textit{\da{Jeg er kun delvist enig i, at} onlinemøder altid er spild af tid.}（用温和的部分同意替代全盘否定，并把 \textit{med} 改为 \textit{i}）
\end{enumerate}

\textbf{练习 C 答案}
\begin{enumerate}
  \item \textit{Jeg er delvist enig med dig, men vi har stadig brug for flere tal.}
  \item \textit{Ganske vist er denne plan meget dyr, men ikke desto mindre har den stor langsigtet værdi.}
  \item \textit{Det kan godt være, at han har ret, men jeg synes stadig, vi bør vente.}
  \item \textit{Selvom alle er uenige, bevarer teamet dog en høflig og saglig tone.}
  \item \textit{Jeg er uenig med min kollega i dette spørgsmål om den digitale løsning.}
  \item \textit{Selvom onlinemøder sparer tid, kan de ikke helt erstatte den fysiske samtale.}
\end{enumerate}

\textbf{DU3 Modul 4 高分范文（Debatindlæg · 135 词）}\\
\textit{\textbf{Storrumskontorer kan fungere med klare rammer}}

\textit{Jeg forstår godt læserens frustration over støj på arbejdspladsen, og jeg er delvist enig i, at konstant uro kan ødelægge koncentrationen. \da{Det kan godt være, at} lukkede enmandskontorer giver mere ro ved skrivebordet, \dabold{men} et totalt forbud mod åbne kontorer er efter min mening ikke den rigtige vej.}

\textit{\da{Ganske vist} kan telefonopkald forstyrre kollegerne, \dabold{men ikke desto mindre} spiller åbne kontorer en vigtig rolle for samarbejde og hurtig sparring i hverdagen. Når man sidder sammen, bliver det langt lettere at stille et kort spørgsmål og hjælpe nye medarbejdere ind i opgaverne. \da{Selvom} støjen er en reel udfordring, \dabold{bør} vi \dabold{dog} vælge bedre indretning frem for at bygge vægge mellem alle.}

\textit{Med faste stillerum, tydelige samtaleregler og fleksible hjemmearbejdsdage kan vi nemlig bevare både roen og fællesskabet på moderne arbejdspladser.}

\small\textbf{名师点评}：正文含标题共 135 词，完全超越 Modul 4 Opgave 2 的 100 词基准。文章精准示范了丹麦式“礼貌反驳三步走”：首段先共情对手痛点并用 \textit{Det kan godt være, at..., men...} 温和转向；中段连用 \textit{Ganske vist..., men ikke desto mindre...}（双倒装满分无误）与 \textit{Selvom..., bør vi dog...}（不加 \textit{men}）展开反论；末段以建设性折中方案收束，逻辑严密、语调得体，达到典型 B2- 优异水准。
\end{answerbox}
\ncdlesson{5}{Faste udtryk og billedsprog i hverdagen}{日常习语与比喻表达}
{DU3 Modul 4（B1+ $\rightarrow$ B2-）｜对标 NCE 3 \& Modultest 4 \textit{Læsning Opgave 1--3 \& Skrivning Opgave 1 (Email)}}
{\textbf{25 个高频丹麦语成语与习语}（如 \textit{have mange bolde i luften, slå to fluer med ét smæk, ramme hovedet på sømmet, et tveægget sværd}）· 习语在叙事与议论中的\textbf{冻结定指规则}、动词变位与语序定位}
{欢迎来到 Unit 1 的收官课！迈入 B1+ $\rightarrow$ B2- 后，衡量语言地道程度的重要标尺不再只是“句子没语法错误”，而是能否自然调用母语者脱口而出的\textbf{固定习语与比喻表达（\textit{faste udtryk og billedsprog}）}。无论是报刊专栏里的 \textit{et tveægget sværd}（双刃剑），还是私人邮件里吐槽自己 \textit{har mange bolde i luften}（分身乏术），本课精选 25 个最高频的丹麦语习语，让你的叙事与辩论瞬间生动鲜活！}

% ---------------------------------------------------------------------
\begin{foerduleaeser}
带着下面两个问题阅读 Tekst A 和 Tekst B：
\begin{enumerate}
  \item \textit{Hvordan reagerer Jonas og Freja i Tekst A, når deres emballageprojekt bliver mere kompliceret end forventet, og hvad mener Freja med ”ikke at kaste barnet ud med badevandet”?}（在 Tekst A 中，当包装项目比预期更复杂时，乔纳斯和芙蕾雅是如何应对的？芙蕾雅说的“不要把婴儿和洗澡水一起倒掉”是什么意思？）
  \item \textit{Hvilke gode råd giver Emma sin ven Mikkel i Tekst B, når han føler, at han har alt for travlt op til eksamen?}（在 Tekst B 中，当米克尔感到临近考试忙得焦头烂额时，艾玛给了他哪些具体建议？）
\end{enumerate}
\end{foerduleaeser}

% ---------------------------------------------------------------------
\ncdsubsec{二、Tekster（口语微调 / 报刊与思辨双语篇）}

\subsection*{Tekst A：\textit{Når arbejdet kræver is i maven}（职场叙事与思辨随笔 · 约 230 词）}

\begin{tipbox}{语篇说明}
本篇讲述乔纳斯（Jonas）与芙蕾雅（Freja）在项目推进中密集使用丹麦语日常习语的真实场景。阅读时请特别留意\textbf{\color{NCDNavy}深蓝色加粗的行动与节奏习语}（如一举两得、卷起袖子）以及\textbf{\color{NCDRed}深红色加粗的洞察与劝慰比喻}（如沉住气、杞人忧天、因噎废食）。
\end{tipbox}

Mandag morgen havde projektteamet virkelig \da{meget om ørerne}. Jonas skulle både forberede en præsentation for ledelsen og finde en ny leverandør af bæredygtig emballage, så han havde mildest talt \da{mange bolde i luften}. Heldigvis valgte hans kollega Freja at \da{smøge ærmerne op} og hjælpe til fra morgenstunden.

”Hvis vi besøger trykkeriet i Aarhus i eftermiddag og samtidig taler med designeren på vejen, kan vi \da{slå to fluer med ét smæk},” foreslog Freja. Jonas nikkede og sagde: ”God idé! Men vi må \da{holde tungen lige i munden}, når vi gennemgår kontrakten, for vi vil jo ikke \dabold{købe katten i sækken}.”

Ved middagstid opstod der dog et problem med prisen. En yngre praktikant blev nervøs og begyndte straks at \dabold{male fanden på væggen}. Han frygtede, at hele planen ville falde sammen. Freja bad ham dog om at \dabold{slå koldt vand i blodet}. ”Det her er bare \dabold{toppen af isbjerget}, hvis vi ikke tænker os om,” sagde hun roligt. ”Prisen er \dabold{et tveægget sværd}: Den nye emballage koster mere nu, men til gengæld styrker den vores image. Vi skal ikke \dabold{kaste barnet ud med badevandet} og droppe hele projektet på grund af én svær forhandling.”

Jonas smilede lettet. Freja havde \da{ramt hovedet på sømmet}. Ved at \da{stikke en finger i jorden} og ringe direkte til leverandøren fandt de et godt kompromis. Efter en travl uge lykkedes det endelig holdet at \da{komme i mål} med lanceringen.

\begin{transbox}
\small
\textbf{\color{NCDTeal}[Tekst A 中文精译]}\\
周一早晨，项目团队真是忙得不可开交。乔纳斯既要为管理层准备汇报，又要为环保包装寻找新供应商，所以说他同时在抛接好几个球（多线并行）一点也不为过。幸好，他的同事芙蕾雅从一早起就决定卷起袖子加油干，伸出援手。

“如果我们今天下午去拜访奥胡斯的印刷厂，顺便在路上和设计师谈谈，就能一举两得，”芙蕾雅提议道。乔纳斯点点头说：“好主意！不过我们审阅合同时必须全神贯注、小心谨慎，因为我们可不想稀里糊涂地盲目签约上当。”

然而到了中午，价格方面还是出了状况。一名年轻的实习生紧张起来，立刻开始往最坏处想（杞人忧天），担心整个计划会彻底泡汤。但芙蕾雅让他先冷静下来别急躁。“如果我们不好好盘算，这恐怕还只是冰山一角，”她平静地说，“价格本身是把双刃剑：新包装眼下成本更高，但反过来能提升我们的品牌形象。我们不应该因为一次艰难的谈判就因噎废食，把整个项目全都放弃掉。”

乔纳斯松了一口气，微笑了起来。芙蕾雅真是一针见血说中了要害。通过先摸清实际情况、再直接致电供应商，他们最终找到了一个很好的折中方案。经过忙碌的一周，团队终于顺利完成了产品发布的目标。
\end{transbox}

\subsection*{Tekst B：\textit{Mails mellem venner: Alt for travlt før eksamen}（私人邮件对答 · 约 140 词）}

\begin{emailbox}{来信｜Fra: Mikkel \quad Til: Emma \quad Emne: Lidt presset i denne uge}
\textbf{Hej Emma}\\
Jeg skriver lige hurtigt, selvom jeg har \da{mange bolde i luften} lige nu. Jeg er startet på nyt studiejob og skal også til modulprøve på fredag. I går prøvede jeg at tale med min vejleder om en ekstra opgave, men det var som at \dabold{tale for døve øren}. Jeg er bange for, at jeg ikke når det hele!\\
\textbf{Hilsen, Mikkel}
\end{emailbox}

\begin{emailbox}{回复｜Fra: Emma \quad Til: Mikkel \quad Emne: Sv: Lidt presset i denne uge}
\textbf{Hej Mikkel}\\
Bare rolig --- du må \dabold{have is i maven}! Når presset stiger, hjælper det ikke at \da{gå som katten om den varme grød}. Skriv en kort mail til din vejleder igen, og \da{gå lige til sagen}. Hvis I kan \da{mødes på halvvejen} om afleveringen, får du ro til prøven. Du skal nok \da{tage tyren ved hornene} i morgen, og når fredagen er overstået, kan du nyde weekenden og \dabold{have det som blommen i et æg}!\\
\textbf{Kærlig hilsen, Emma}
\end{emailbox}

\textbf{Tekst B 阅读核对（Rigtigt / Forkert / Kort svar，答案见第六部分）：}
\begin{enumerate}
  \item \textit{Mikkel føler, at hans vejleder lyttede meget opmærksomt til hans bekymring i går.} \underline{\hspace{2.2cm}}
  \item \textit{Hvordan anbefaler Emma, at Mikkel formulerer sin nye mail til vejlederen?} \underline{\hspace{2.2cm}}
  \item \textit{Hvad kan Mikkel se frem til i weekenden efter fredagens prøve ifølge Emma?} \underline{\hspace{2.2cm}}
\end{enumerate}

\begin{tipbox}{本课发音与语调要点}
\begin{itemize}
  \item \textbf{习语核心名词的特殊音变与静音字母}：
    \begin{itemize}
      \item \textit{ét smæk} \ipa{ˈed ˈsmɛɡ}（注意 \textit{ét} 重读表“仅仅一次”，\textit{smæk} 短促有力）
      \item \textit{et tveægget sværd} \ipa{ˈtˢveːˌɛɡəð ˈsvɛɐ̯ˀ}（\textit{sværd} 中的 \textbf{\textit{d} 完全不发音}，带长元音与喉塞音！）
      \item \textit{døve øren} \ipa{ˈdøːvə ˈøːɐn}（古旧复数 \textit{øren} 读作 \ipa{ˈøːɐn}；注意 \textit{døve} 是浊擦音 \ipa{v}）
      \item \textit{den varme grød} \ipa{dɛn ˈvɑːmə ˈɡʁøːˀð}（\textit{grød} 末尾是清晰软 \textit{d} 加喉塞音 \ipa{ð}）
    \end{itemize}
  \item \textbf{喉塞音与软化连读}：\textit{tungen lige i munden} \ipa{ˈtˢoŋən ˈliːə i ˈmɔnˀən}（\textit{-nd-} 中的 \textit{d} 不发音！）；\textit{på sømmet} \ipa{pʰʌ ˈsømˀəð}；\textit{i sækken} \ipa{i ˈsɛɡən}；\textit{på væggen} \ipa{pʰʌ ˈvɛɡən}；\textit{bølgelængde} \ipa{ˈbøljəˌlɛŋdə}（\textit{g} 滑音化为 \ipa{j}，第二个 \textit{d} 不发音）。
\end{itemize}
\end{tipbox}

\begin{transbox}
\small
\textbf{\color{NCDTeal}[Tekst B 中文对照]}\\
\textbf{米克尔的来信}：嗨，艾玛。虽然我现在正同时忙着好多件事，还是赶紧给你写封信。我刚入职了新的学生兼职，周五还要参加模块考试。昨天我试图和导师沟通一份额外作业的事，但感觉完全是在对牛弹琴。我真担心自己没法全赶完！问候，米克尔。

\textbf{艾玛的回复}：嗨，米克尔。放轻松——你一定要沉住气！压力上来的时候，绕弯子避重就轻是没有用的。再给导师写一封简短邮件，直接开门见山说明重点。如果你们能在交作业的时间上各让一步，你就能安心准备考试了。明天拿出勇气迎难而上吧，等周五一结束，你就可以好好享受周末，过得像鸡蛋黄在蛋壳里一样舒心惬意了！亲爱的，艾玛。
\end{transbox}

% ---------------------------------------------------------------------
\ncdsubsec{三、Gloser（核心生词表）}

\begin{tipbox}{词表重点}
记忆习语时务必连同其中的\textbf{动词不规则变化}和\textbf{名词定指词尾}一起背诵。例如 \textit{at slå}（击打）过去式与完成时为 \textit{slog, har slået}；\textit{at stikke}（伸出、刺入）为 \textit{stak, har stukket}；\textit{at smøge}（卷起）为 \textit{smøgede, har smøget}。注意 \textit{et udtryk}（表达）单复数同形，复数定指加双写 \textit{k}（\textit{udtrykkene}）。
\end{tipbox}

\begin{glosertable}
名词 & et billedsprog & \textit{(billedsproget, ingen flertal)} & 比喻性语言，形象修辞 \\
名词 & et udtryk & \textit{(udtrykket, udtryk, udtrykkene)} & 表达，习语（单复同形） \\
名词 & en emballage & \textit{(emballagen, emballager, -erne)} & 包装，外包装 \\
名词 & et projekt & \textit{(projektet, projekter, projekterne)} & 项目，课题 \\
名词 & et pres & \textit{(presset, ingen flertal)} & 压力，紧迫感 \\
名词 & et sværd & \textit{(sværdet, sværd, sværdene)} & 剑（单复同形） \\
名词 & et isbjerg & \textit{(isbjerget, isbjerge, isbjergene)} & 冰山 \\
名词 & en bølgelængde & \textit{(bølgelængden, -længder, -erne)} & 波长；合拍，默契 \\
名词 & en deadline & \textit{(deadlinen, deadlines, -ene)} & 截止期限 \\
名词 & et smæk & \textit{(smækket, smæk, smækkene)} & 一巴掌，一击 \\
\midrule
动词 [har] & at smøge op & \textit{(smøger, smøgede, har smøget op)} & 卷起（袖子） \\
动词 [har] & at ramme & \textit{(rammer, ramte, har ramt)} & 击中，打中，说中 \\
动词 [har] & at male & \textit{(maler, malede, har malet)} & 绘画，涂抹 \\
动词 [har] & at stikke & \textit{(stikker, stak, har stukket)} & 伸出，刺入，插入 \\
动词 [har] & at kaste & \textit{(kaster, kastede, har kastet)} & 扔，抛，倒掉 \\
\midrule
形容词 & tveægget & \textit{(tveægget, tveæggede)} & 双刃的，利弊兼具的 \\
形容词 & overvældende & \textit{(不变，现在分词形容词)} & 压倒性的，令人难以承受的 \\
\end{glosertable}

% ---------------------------------------------------------------------
\ncdsubsec{四、Grammatik \& Stilistik（25 个高频习语全景图谱与活用铁律）}

\begin{grammatikbox}{核心习语库 I：职场效率与沟通交锋（第 1--13 个高频习语）}
以下精选丹麦日常对话与报刊讨论中使用率极高的 25 个固定习语。本框先呈现\textbf{职场节奏}与\textbf{沟通立场}两大高频场景（1--13 号）：

\vspace{0.3em}
\renewcommand{\arraystretch}{1.2}
{\small
\begin{tabularx}{\linewidth}{@{} >{\bfseries\raggedright\arraybackslash}p{4.3cm} >{\raggedright\arraybackslash}p{3.9cm} >{\raggedright\arraybackslash}X @{}}
\toprule
\textbf{丹麦语固定习语} & \textbf{字面意象与地道中文} & \textbf{高频语境简例} \\
\midrule
\da{1. have mange bolde i luften} & 同时抛接多球（多线并行） & \textit{Jeg har mange bolde i luften i år.} \\
\da{2. slå to fluer med ét smæk} & 一拍打两蝇（一举两得） & \textit{Vi slår to fluer med ét smæk i dag.} \\
\da{3. have meget om ørerne} & 耳边事务缠身（焦头烂额） & \textit{Teamet havde meget om ørerne.} \\
\da{4. holde tungen lige i munden} & 舌头摆正（屏气凝神防出错） & \textit{Hold tungen lige i munden ved kontrakten!} \\
\da{5. smøge ærmerne op} & 卷起袖子（甩开膀子加油干） & \textit{Nu må vi smøge ærmerne op.} \\
\da{6. komme i mål (med noget)} & 冲过终点（顺利达成目标） & \textit{Vi kom endelig i mål med planen.} \\
\midrule
\dabold{7. gå lige til sagen} & 直奔主题（开门见山） & \textit{Lad os gå lige til sagen på mødet.} \\
\dabold{8. gå som katten om den varme grød} & 像猫绕热粥（拐弯抹角） & \textit{Han gik som katten om den varme grød.} \\
\dabold{9. ramme hovedet på sømmet} & 锤中钉头（一针见血） & \textit{Du ramte virkelig hovedet på sømmet!} \\
\dabold{10. tale for døve øren} & 对聋耳说话（对牛弹琴） & \textit{Vores råd var som at tale for døve øren.} \\
\dabold{11. mødes på halvvejen} & 在半路相遇（各让一步妥协） & \textit{Begge parter mødtes på halvvejen.} \\
\dabold{12. være på bølgelængde med} & 与……同频（心有灵犀/合拍） & \textit{Jeg er helt på bølgelængde med Sofie.} \\
\dabold{13. se gennem fingre med} & 透过指缝看（睁一只眼闭一只眼） & \textit{Man kan ikke se gennem fingre med snyd.} \\
\bottomrule
\end{tabularx}
}
\end{grammatikbox}

\begin{grammatikbox}{核心习语库 II：辩论洞察与心态哲理（第 14--25 个高频习语）}
接下来 14--25 号习语尤其适合在 DU3 Modul 4 \textit{Debatindlæg}、口试评论利弊及日常书信安慰朋友时画龙点睛：

\vspace{0.3em}
\renewcommand{\arraystretch}{1.2}
{\small
\begin{tabularx}{\linewidth}{@{} >{\bfseries\raggedright\arraybackslash}p{4.3cm} >{\raggedright\arraybackslash}p{3.9cm} >{\raggedright\arraybackslash}X @{}}
\toprule
\textbf{丹麦语固定习语} & \textbf{字面意象与地道中文} & \textbf{高频语境简例} \\
\midrule
\da{14. et tveægget sværd} & 双刃剑（利弊并存） & \textit{Ny teknologi er ofte et tveægget sværd.} \\
\da{15. toppen af isbjerget} & 冰山一角（问题远不止于此） & \textit{Fejlene er kun toppen af isbjerget.} \\
\da{16. købe katten i sækken} & 买袋里的猫（盲目买单上当） & \textit{Læs vilkårene, så du ikke køber katten i sækken.} \\
\da{17. male fanden på væggen} & 把魔鬼画上墙（杞人忧天/危言耸听） & \textit{Vi behøver ikke at male fanden på væggen.} \\
\da{18. kaste barnet ud med badevandet} & 连婴带水倒掉（因噎废食） & \textit{Kast ikke barnet ud med badevandet!} \\
\da{19. gå over åen efter vand} & 过河打水（舍近求远/多此一举） & \textit{At flyve til et kort møde er at gå over åen efter vand.} \\
\midrule
\dabold{20. have is i maven} & 肚里有冰（沉得住气/镇定） & \textit{Det kræver is i maven før eksamen.} \\
\dabold{21. slå koldt vand i blodet} & 往血里兑冷水（冷静别冲动） & \textit{Slå koldt vand i blodet og vent lidt!} \\
\dabold{22. stikke en finger i jorden} & 把手指探进土里（审时度势摸底） & \textit{Lad os lige stikke en finger i jorden først.} \\
\dabold{23. tage tyren ved hornene} & 揪住公牛角（果断迎难而上） & \textit{I morgen tager jeg tyren ved hornene.} \\
\dabold{24. få blod på tanden} & 牙上沾血（尝到甜头越干越起劲） & \textit{Efter første succes fik hun blod på tanden.} \\
\dabold{25. have det som blommen i et æg} & 如蛋中蛋黄（养尊处优/舒心惬意） & \textit{Gæsterne havde det som blommen i et æg.} \\
\bottomrule
\end{tabularx}
}
\end{grammatikbox}

\begin{grammatikbox}{语法点：习语活用的三大语法与语体铁律（Stilistik \& Grammatik）}
\begin{enumerate}
  \item \textbf{定指与单复数的“冻结规则（\textit{Frosset form}）”}：习语里的名词定指后缀与冠词往往从历史上冻结下来，绝对不能按自己的英文或中文习惯随便改写：
    \begin{itemize}
      \item \ding{51}\ \textit{slå to fluer med \dabold{ét} smæk}（数字 1 保留重音 \textit{ét}，不能写成 \ding{55}\ \textit{to fluer med en smæk}）。
      \item \ding{51}\ \textit{tale for døve \dabold{øren}}（固定保留古旧复数 \textit{øren}，绝不能改成现代复数 \ding{55}\ \textit{døve ører}！）。
      \item \ding{51}\ \textit{i luft\dabold{en}} / \textit{i mund\dabold{en}} / \textit{på sømm\dabold{et}} / \textit{i blod\dabold{et}}（这些定指词尾均不可省去或改成不定指）。
    \end{itemize}
  \item \textbf{动词照样随主语和时态变位，否定词 \textit{ikke} 按正常主从句占位}：习语中的动词要正常变位，中场副词（\textit{ikke, jo, egentlig}）也照常插入动词与名词之间：\newline
    \ding{51}\ \textit{Han \dabold{slår jo ikke} koldt vand i blodet.} / \textit{Fordi vi \dabold{ikke ville købe} katten i sækken.}
  \item \textbf{控制密度与语域匹配（避免“中式成语直译”与堆砌）}：千万不要生造中文或英文直译（如 \ding{55}\ \textit{det er et stykke kage}）。在 DU3 Modul 4 一篇 140 词的邮件或辩论短文中，恰当嵌入 \textbf{2 至 3 个}丹麦习语最显功力；连篇累牍堆砌反而会显得做作。
\end{enumerate}
\end{grammatikbox}

% ---------------------------------------------------------------------
\ncdsubsec{五、Mønsterøvelser（句型与修辞操练）}

\textbf{练习 A：高频习语精准补全（每空填入唯一准确的 1 个丹麦语核心词）}
\begin{enumerate}
  \item Jeg har simpelthen for mange \underline{\hspace{2.2cm}} i luften i denne måned med både arbejde, kursus og flytning.
  \item Hvis jeg henter pakken på vejen til vores møde, slår jeg to \underline{\hspace{2.2cm}} med ét smæk.
  \item Nu må vi holde tungen lige i \underline{\hspace{2.2cm}}, så vi ikke overser fejl i den nye kontrakt.
  \item Din kommentar ramte virkelig hovedet på \underline{\hspace{2.2cm}} --- det er præcis der, problemet ligger!
  \item Vi behøver ikke at male \underline{\hspace{2.2cm}} på væggen endnu; lad os først se tallene for næste kvartal.
  \item Selvom tidsplanen er stram, prøver teamet at slå koldt \underline{\hspace{2.2cm}} i blodet og bevare overblikket.
  \item Det føles desværre som at tale for døve \underline{\hspace{2.2cm}}, når vi foreslår ændringer til direktøren.
  \item Parterne var uenige i starten, men de valgte heldigvis at mødes på \underline{\hspace{2.2cm}} til sidst.
\end{enumerate}

\textbf{练习 B：习语语境改写与搭配修正（用括号内指定的习语升级平淡句，或修正冻结形态错误）}
\begin{enumerate}
  \item 用 \textit{smøge ærmerne op} 改写：\textit{Nu må vi begynde at arbejde rigtig hårdt for at blive færdige til tiden.} $\rightarrow$
  \item 用 \textit{gå som katten om den varme grød} 改写：\textit{Han talte hele tiden udenom og nævnte slet ikke det svære emne.} $\rightarrow$
  \item 修正习语里的名词词尾错误：\textit{Jeg følte, at jeg talte for døve ører, da ingen lyttede til mit forslag.} $\rightarrow$
  \item 修正冠词与定指冻结错误：\textit{Vi vil ikke købe en kat i sæk, så vi tjekker bilen grundigt først.} $\rightarrow$
  \item 用 \textit{et tveægget sværd} 改写：\textit{Kunstig intelligens har både store fordele og alvorlige risici for skolerne.} $\rightarrow$
  \item 用 \textit{tage tyren ved hornene} 改写：\textit{I morgen tør jeg endelig håndtere det svære problem direkte.} $\rightarrow$
\end{enumerate}

\textbf{练习 C：中译丹实战（灵活运用本课习语，注意动词变位与准确语序）}
\begin{enumerate}
  \item 我们今天下午可以在奥胡斯一举两得。
  \item 老师要求我们在审阅合同时保持沉着冷静。
  \item 你绝对不应该因噎废食而放弃整个计划。
  \item 与新同事合拍让我的日常工作轻松了许多。
  \item 在尝到第一次成功的甜头后，她变得更加干劲十足。
  \item 这个新规则其实只暴露了冰山一角。
\end{enumerate}

% ---------------------------------------------------------------------
\ncdsubsec{六、DU3 Modul 4 实战微任务与参考答案}

\begin{emailbox}{DU3 Modul 4 Skrivning Opgave 1 实战微任务（私人长邮件 · minimum 125 ord）}
\textbf{Situation}: Din ven Kasper har skrevet til dig, fordi han har meget travlt lige nu. Han har lige fået et nyt job, går til dansk om aftenen og er også i gang med at renovere sin lejlighed. Nu er han bange for, at han ikke kan nå det hele.\\
\textbf{Opgave}: Skriv en varm og opmuntrende mail til Kasper på \textbf{mindst 125 ord} (gerne 140--160 ord). Svar på hans bekymring, giv ham et par konkrete råd til at skabe overblik, og brug naturligt mindst \textbf{tre} faste udtryk eller billedsprog fra denne lektion!
\end{emailbox}

\begin{answerbox}
\textbf{Før du læser 参考回答}
\begin{enumerate}
  \item \textit{Jonas har meget om ørerne, men Freja smøger ærmerne op og foreslår, at de kan slå to fluer med ét smæk ved både at besøge trykkeriet og tale med designeren. Da en praktikant bliver nervøs, beder hun ham slå koldt vand i blodet og forklarer, at ”ikke at kaste barnet ud med badevandet” betyder, at man ikke skal opgive hele det gode projekt på grund af én svær forhandling.}
  \item \textit{Emma råder Mikkel til at have is i maven, ikke gå som katten om den varme grød, men skrive en kort mail igen og gå lige til sagen, så han og vejlederen kan mødes på halvvejen.}
\end{enumerate}

\textbf{Tekst B 阅读核对答案}
\begin{enumerate}
  \item \da{Forkert}（米克尔形容自己的沟通感觉就像对牛弹琴：\textit{det var som at tale for døve øren}）
  \item \da{Kort og lige til sagen.}（写一封简短邮件，直接开门见山说明核心问题）
  \item \da{At nyde weekenden og have det som blommen i et æg.}（好好享受周末，过得舒心惬意）
\end{enumerate}

\textbf{练习 A 答案}
\begin{enumerate}
  \item \da{bolde}（固定习语：\textit{have mange bolde i luften} 同时忙很多事）
  \item \da{fluer}（固定习语：\textit{slå to fluer med ét smæk} 一举两得）
  \item \da{munden}（固定习语：\textit{holde tungen lige i munden} 聚精会神防出错）
  \item \da{sømmet}（固定习语：\textit{ramme hovedet på sømmet} 一针见血）
  \item \da{fanden}（固定习语：\textit{male fanden på væggen} 杞人忧天）
  \item \da{vand}（固定习语：\textit{slå koldt vand i blodet} 保持冷静）
  \item \da{øren}（注意此习语固定用古复数形式 \textit{øren}，不能写成 \textit{ører}）
  \item \da{halvvejen}（固定习语：\textit{mødes på halvvejen} 各让一步折中）
\end{enumerate}

\textbf{练习 B 答案（参考）}
\begin{enumerate}
  \item \textit{Nu må vi \da{smøge ærmerne op} for at blive færdige til tiden.}（生动替代泛泛的“努力工作”）
  \item \textit{Han \da{gik som katten om den varme grød} og nævnte slet ikke det svære emne.}（注意动词 \textit{gå} 随过去时变为 \textit{gik}）
  \item \textit{Jeg følte, at jeg \dabold{talte for døve øren}, da ingen lyttede til mit forslag.}（必须保留冻结的古复数 \textit{øren}）
  \item \textit{Vi vil ikke \dabold{købe katten i sækken}, så vi tjekker bilen grundigt først.}（两个名词都必须保留定指形式 \textit{katten} 和 \textit{sækken}）
  \item \textit{Kunstig intelligens er \da{et tveægget sværd} for skolerne.}（用比喻浓缩“既有利又有弊”）
  \item \textit{I morgen \da{tager} jeg endelig \da{tyren ved hornene}.}（注意句首时间状语触发主句 V2 倒装 \textit{tager jeg}）
\end{enumerate}

\textbf{练习 C 答案}
\begin{enumerate}
  \item \textit{Vi kan slå to fluer med ét smæk i Aarhus i eftermiddag.}
  \item \textit{Læreren bad os om at have is i maven, når vi gennemgår kontrakten.}
  \item \textit{Du bør absolut ikke kaste barnet ud med badevandet og opgive hele planen.}
  \item \textit{At være på bølgelængde med nye kolleger gør min hverdag meget lettere.}
  \item \textit{Efter den første succes fik hun blod på tanden i sit arbejde.}
  \item \textit{Den nye regel viser faktisk kun toppen af isbjerget.}
\end{enumerate}

\textbf{DU3 Modul 4 高分范文（148 词，不含署名）}\\
\textit{Hej Kasper}

\textit{Tak for din besked! Jeg kan godt forstå, at du har \da{meget om ørerne} for tiden. Når man starter på et nyt arbejde og samtidig renoverer lejlighed og går på sprogskole, har man jo virkelig \da{mange bolde i luften} på én gang.}

\textit{Mit bedste råd er, at du først prøver at \dabold{slå koldt vand i blodet}. Du behøver ikke at \dabold{male fanden på væggen} eller få alt i lejligheden færdigt i denne uge. Prøv i stedet at lave en kort plan over de vigtigste opgaver. Måske kan du også \da{slå to fluer med ét smæk} ved at lytte til danske podcasts, mens du maler væggene?}

\textit{Hvis presset bliver for stort, så \da{tag tyren ved hornene} og tal åbent med din chef om arbejdstiden --- så skal I nok \da{mødes på halvvejen}. Når du først kommer gennem denne travle måned, får du helt sikkert ro på igen!}

\textit{Rigtig god arbejdslyst!}

\textit{Bedste hilsner\\
\mbox{[Dit navn]}}

\small\textbf{名师点评}：正文 148 词（不含署名）圆满达成 Modul 4 Opgave 1 的 125 词以上要求。全信以真挚体贴的口吻回应朋友的焦虑，自然融入了 \textit{have meget om ørerne, mange bolde i luften, slå koldt vand i blodet, male fanden på væggen, slå to fluer med ét smæk, tag tyren ved hornene, mødes på halvvejen} 七个习语与单元核心小品词 \textit{jo}，形象贴切而不堆砌，展现出极佳的 B1+/B2- 交际成熟度。
\end{answerbox}

% =====================================================================
% Unit 2
% =====================================================================
\ncdunit{2}{Velfærdsstatens Dilemmaer}{福利国家的困境与辩论}{Lektion 6--10：政策预测与被动预期 · 正反方对立框架 · 非真实条件与虚拟语气 · 严谨因果链条 · 社会议题动宾搭配}

\ncdlesson{6}{Hvem skal passe de ældre i fremtiden?}{未来谁来照料老人？}
{DU3 Modul 4（B1+ $\rightarrow$ B2-）｜对标 NCE 3 \& Modultest 4 \textit{Velfærdsstaten \& Fremtidsudsigter}}
{\textbf{政策预测与被动预期结构}（\textit{forventes at, skønnes at, vurderes at, står over for en alvorlig udfordring}）· 主语提升句型 vs 形式主语 \textit{Det forventes, at...} · 客观趋势预测排除主观推测陷阱}
{欢迎进入 Unit 2《福利国家的困境与辩论》！丹麦引以为傲的福利国家正面临严峻的人口结构考验。从本课起，我们将用地道的高阶报刊句型探讨公共政策：表达“预计将会……”“据评估……”时，丹麦语很少用啰嗦的 \textit{Folk tror, at...}，而是高频使用 \textbf{\textit{forventes at} / \textit{skønnes at}} 这类“被动预期提升结构”。掌握它们，是写出 Modul 4 政论高分句的第一步！}

% ---------------------------------------------------------------------
\begin{foerduleaeser}
带着下面两个问题阅读 Tekst A 和 Tekst B：
\begin{enumerate}
  \item \textit{Hvilken alvorlig udfordring står de danske kommuner over for frem mod 2035 ifølge Tekst A, og hvilke løsninger peger eksperter på?}（根据 Tekst A，到 2035 年丹麦各市镇面临什么严峻挑战？专家们指出了哪些解决思路？）
  \item \textit{Hvordan ser Sofie og Henrik forskelligt på brugen af velfærdsteknologi og frivillige i ældreplejen i Tekst B?}（在 Tekst B 中，索菲和亨里克对养老照护中使用福利科技和志愿者分别持怎样的不同看法？）
\end{enumerate}
\end{foerduleaeser}

% ---------------------------------------------------------------------
\ncdsubsec{二、Tekster（福利国家专题与正反辩论双语篇）}

\subsection*{Tekst A：\textit{Ældrebyrden eller fremtidens omsorg?}（政策分析短文 · 约 245 词）}

\begin{tipbox}{语篇说明}
本篇聚焦丹麦老龄化与养老照护危机。阅读时请特别关注\textbf{\color{NCDNavy}深蓝色加粗的被动预期结构}（\textit{forventes at, skønnes at, vurderes at}）与\textbf{\color{NCDRed}深红色加粗的挑战句型及从句语序}。
\end{tipbox}

Danmark er kendt for en stærk velfærdsstat, hvor det offentlige traditionelt har taget hovedansvaret for plejen af ældre borgere. Men i disse år ændrer befolkningens sammensætning sig hurtigt. Vi lever længere end før, mens færre unge kommer ind på arbejdsmarkedet. Som følge af denne demografiske udvikling \dabold{står de danske kommuner over for en alvorlig udfordring}.

Ifølge nye prognoser \da{forventes} antallet af borgere over 80 år \da{at stige} med mere end fyrre procent frem mod 2035. Samtidig \da{skønnes} manglen på social- og sundhedsassistenter \da{at vokse} markant i hele landet. Det \da{vurderes, at} der allerede om ti år vil mangle tusindvis af hænder på plejehjem og i hjemmeplejen, hvis ikke flere vælger et omsorgsfag.

Debatten handler derfor ikke længere kun om penge, men også om arbejdskraft og kvalitet. Nogle kommuner forsøger at rekruttere udenlandsk personale, mens andre satser på velfærdsteknologi som digitale sensorer og automatiske hjælpemidler. Sådanne løsninger \da{forventes at kunne} aflaste personalet i de travle morgentimer. Alligevel \da{menes} teknologien \dabold{ikke at kunne} erstatte den menneskelige kontakt, som mange ældre savner allermest.

Et stigende antal eksperter peger på, at fremtidens ældrepleje kræver et tættere samarbejde mellem det offentlige, pårørende og frivillige. Selvom ingen ønsker at opgive den universelle velfærd, \da{forventes} familiemedlemmer fremover \da{at spille} en større rolle i hverdagen. Spørgsmålet er blot, hvor grænsen går mellem professionel pleje og familiens ansvar.

\begin{transbox}
\small
\textbf{\color{NCDTeal}[Tekst A 中文精译]}\\
丹麦以强大的福利国家体制闻名，传统上由公共部门承担照料老年公民的主要责任。但近年来，人口结构正在迅速变化。人们比以往更加长寿，而进入劳动力市场的年轻人却越来越少。受这一人口变化趋势的影响，丹麦各市镇正面临严峻挑战。

根据最新预测，到 2035 年，80 岁以上公民的人数预计将增长超过百分之四十。与此同时，全国范围内护理与健康助理员的短缺估计将显著扩大。据评估，如果届时没有更多人选择照护行业，仅仅十年后养老院和居家护理领域就将缺少数千名工作人员。

因此，这场讨论已不再仅仅关乎资金，也关乎劳动力与服务质量。一些市镇尝试招募外籍护理人员，而另一些市镇则押注于智能传感器和自动辅助设备等福利科技。此类方案预计将能够在繁忙的清晨时段分担员工的工作压力。即便如此，人们认为技术依然无法取代许多老人最渴望的人际交流。

越来越多的专家指出，未来的养老照护需要公共机构、亲属与志愿者之间更紧密的协作。尽管没有人愿意放弃普惠福利制度，但预计未来家庭成员将在老人日常生活中扮演更重要的角色。问题仅仅在于，专业照护与家庭责任之间的界线究竟划在哪里。
\end{transbox}

\subsection*{Tekst B：\textit{Debat i lokalavisen: Teknologi og varme hænder}（正反读者来信 · 约 145 词）}

\begin{emailbox}{地方报辩论专栏｜Fra: Sofie (Pårørende) \& Henrik (Plejehjemsleder)}
\textbf{Sofie:} Min mor bor på et lokalt plejehjem, og jeg er bekymret for den nye spareplan. Kommunen vil installere flere skærme og sensorer om natten, fordi det \da{skønnes at spare} arbejdstimer. Det kan godt være, at teknologien hjælper lidt, men en skærm kan jo ikke holde en ensom ældre i hånden. Hvis vi skærer ned på personalet, \dabold{forventes} kvaliteten \dabold{at falde} yderligere.

\textbf{Henrik:} Som leder af et plejehjem forstår jeg godt Sofies bekymring. Vi \dabold{står} dog \dabold{over for} en konkret virkelighed, hvor det allerede er svært at besætte alle vagter. Digitale hjælpemidler \da{vurderes at kunne} frigøre tid, så de ansatte får mere ro til selve omsorgen. Samtidig \da{forventes} frivillige besøgsvenner \da{at give} ekstra nærvær i hverdagen, uden at de overtager de faglige opgaver.
\end{emailbox}

\textbf{Tekst B 阅读核对（Rigtigt / Forkert / Kort svar，答案见第六部分）：}
\begin{enumerate}
  \item \textit{Sofie mener, at natlige sensorer og skærme fuldt ud kan erstatte menneskelig omsorg.} \underline{\hspace{2.2cm}}
  \item \textit{Hvilket konkret problem oplever Henrik allerede på sit plejehjem?} \underline{\hspace{4.5cm}}
  \item \textit{Hvad forventes frivillige besøgsvenner at bidrage med ifølge Henrik?} \underline{\hspace{4.5cm}}
\end{enumerate}

\begin{tipbox}{本课发音与语调要点（Udtale \& Reduktioner）}
\begin{itemize}
  \item \textbf{被动预期动词的“词尾弱化与喉塞音”}：\textit{forventes} \ipa{fʌˈvɛnˀdəs}（重音在第二音节，带喉塞音 \textit{stød}，词尾 \textit{-es} 轻读 \ipa{əs}）；\textit{skønnes} \ipa{ˈsgønəs}（短元音无 \textit{stød}）；\textit{vurderes} \ipa{vuˈdeːˀɐs}（\textit{-eres} 融合为 \ipa{eːˀɐs}）；\textit{menes} \ipa{ˈmeːnəs}。
  \item \textbf{养老与人口核心词静音字母}：\textit{et plejehjem} \ipa{ˈpʰlɑi̯əˌjɛmˀ}（\textit{hjem} 词首 \textit{h} 完全不发音！）；\textit{en udfordring} \ipa{ˈuðˌfɒːˀdʁeŋ}（前缀 \textit{ud-} 软 \textit{d} \ipa{ð}）；\textit{en pårørende} \ipa{ˈpʰʌˌʁøːˀɐnə}；\textit{arbejdskraft} \ipa{ˈɑːbɑi̯ˀdsˌkʰʁɑfd}（\textit{arbejds-} 中 \textit{d} 极轻或脱落）。
  \item \textbf{短语连读节奏}：\textit{står over for} \ipa{sd̥ɒːˀ ˈʌʊ̯ˀɐ fʌ}（\textit{over} 中 \textit{v} 弱滑音化为 \ipa{ʊ̯}，\textit{for} 极轻读）；\textit{at aflaste} \ipa{ɑ ˈɑʊ̯ˌlasdə}（前缀 \textit{af-} 读作 \ipa{ɑʊ̯}）。
\end{itemize}
\end{tipbox}

\begin{transbox}
\small
\textbf{\color{NCDTeal}[Tekst B 中文对照]}\\
\textbf{索菲（老人亲属）}：我母亲住在本地一家养老院，我对新的节支计划感到担忧。市政厅打算在夜间安装更多屏幕和传感器，因为据估算这能节省工时。也许技术确实能帮上一点忙，但一块屏幕总不能握住孤独老人的手啊。如果我们削减护理人员，预计照护质量会进一步下降。

\textbf{亨里克（养老院院长）}：作为养老院院长，我完全理解索菲的担忧。然而，我们正面临一个现实局面：现在就已经很难排满所有班次了。经评估，数字辅助设备能够腾出时间，让员工更从容地专注于照护本身。同时，预计志愿探访者也能在日常生活中提供额外的陪伴关怀，而无需接管专业护理任务。
\end{transbox}

% ---------------------------------------------------------------------
\ncdsubsec{三、Gloser（核心生词表）}

\begin{tipbox}{词表重点}
注意 \textit{en pårørende}（亲属/陪护家属）是由现在分词转化而来的名词，单复数不定式同形（\textit{en pårørende / flere pårørende}）。动词 \textit{at skønne} 在报刊政策语境中专指“基于数据估算、推算”，不要和“好看的” \textit{skøn} 混淆；\textit{at stå over for} 表示“面临（挑战、抉择）”，固定接介词 \textit{for}。
\end{tipbox}

\begin{glosertable}
名词 & en ældrepleje & \textit{(ældreplejen, kun ental)} & 养老照护，老年护理 \\
名词 & et plejehjem & \textit{(plejehjemmet, plejehjem, plejehjemmene)} & 养老院 \\
名词 & en pårørende & \textit{(den pårørende, pårørende, de pårørende)} & 亲属，近亲属 \\
名词 & en udfordring & \textit{(udfordringen, udfordringer, udfordringerne)} & 挑战，难题 \\
名词 & en arbejdskraft & \textit{(arbejdskraften, kun ental)} & 劳动力，人力 \\
名词 & en mangel (på) & \textit{(manglen, mangler, manglerne)} & 缺乏，短缺（接 \textit{på}） \\
名词 & en udvikling & \textit{(udviklingen, udviklinger, udviklingerne)} & 发展，演变趋势 \\
名词 & en prognose & \textit{(prognosen, prognoser, prognoserne)} & 预测，预报 \\
名词 & en velfærdsteknologi & \textit{(velfærdsteknologien, -er, -erne)} & 福利科技 \\
名词 & et nærvær & \textit{(nærværet, kun ental)} & 陪伴，亲近感 \\
\midrule
动词 [har] & at forvente & \textit{(forventer, forventede, har forventet)} & 预期，预计 \\
动词 [har] & at skønne & \textit{(skønner, skønnede, har skønnet)} & 估算，估计，判断 \\
动词 [har] & at vurdere & \textit{(vurderer, vurderede, har vurderet)} & 评估，评价 \\
动词 [har] & at stå over for & \textit{(står, stod, har stået over for)} & 面临，面对（挑战/局势） \\
动词 [har] & at rekruttere & \textit{(rekrutterer, rekrutterede, har rekrutteret)} & 招募，招聘 \\
动词 [har] & at aflaste & \textit{(aflaster, aflastede, har aflastet)} & 减轻……的负担，分担 \\
\midrule
形容词 & demografisk / alvorlig & \textit{demografisk (-, -e) / alvorlig (-t, -e)} & 人口结构的 / 严峻的，严重的 \\
\end{glosertable}

% ---------------------------------------------------------------------
\ncdsubsec{四、Grammatik \& Stilistik（核心语法、政策句式与修辞精讲）}

\begin{grammatikbox}{语法点 1：政策预测与被动预期结构——“主语提升” vs “形式主语 \textit{Det}”}
在 Modul 4 阅读理解、图表评论与社会辩论写作中，需要客观转述机构预测、普遍预期或专家评估。丹麦语主要有以下两种等价但语体极为精炼的句式：

\vspace{0.3em}
\renewcommand{\arraystretch}{1.24}
{\small
\begin{tabularx}{\linewidth}{@{} >{\bfseries\raggedright\arraybackslash}p{3.0cm} >{\raggedright\arraybackslash}p{4.3cm} >{\raggedright\arraybackslash}X @{}}
\toprule
\textbf{结构类型} & \textbf{句法公式} & \textbf{新闻政论标准例句} \\
\midrule
\da{1. 主语提升结构}\newline （\textit{Subject-raising}） &
\textbf{主题名词 + \dabold{forventes / skønnes / vurderes / menes} + \textit{at} + 动词原形} &
\textit{Antallet af ældre \dabold{forventes at stige} frem mod 2035.}\newline
\textit{Teknologien \dabold{skønnes at spare} tid om natten.} \\
\midrule
\da{2. 形式主语 \textit{Det}}\newline （\textit{Formelt subjekt}） &
\textbf{\dabold{Det forventes / Det skønnes / Det vurderes}, \textit{at} + 完整从句} &
\textit{\dabold{Det vurderes, at} der vil mangle personale i fremtiden.}\newline
\textit{\dabold{Det forventes, at} kommunerne ansætter flere.} \\
\bottomrule
\end{tabularx}
}
\vspace{0.35em}

\textbf{进阶变形三大细节}：
\begin{itemize}
  \item \ding{172} \textbf{否定词 \textit{ikke} 的位置}：在主语提升结构中，否定词 \textit{ikke} 既可放在预期动词之后，也可放在 \textit{at} 之前，呈现“紧邻 \textit{at}”的书面节奏：\textit{Teknologien \da{menes ikke at kunne} erstatte mennesker.}（注意不要放到 \textit{at} 与动词原形之间去写成 *\textit{at ikke kunne}！虽然极少数口语偶见，B1+/B2- 正式政论应把 \textit{ikke} 置于 \textit{at} 之前）。
  \item \ding{173} \textbf{情态与被动不定式连用}：在预期结构中可以嵌套情态动词原形或被动不定式：\textit{Løsningerne \da{forventes at kunne} aflaste personalet.} / \textit{Flere plejehjem \da{forventes at blive bygget}.}
  \item \ding{174} \textbf{表示“面临挑战”的规范框架}：\dabold{at stå over for} 接具体问题或挑战，常见高分搭配为：\textit{Samfundet står over for en alvorlig / voksende / dobbelt udfordring}（面临严峻/日益加剧/双重挑战）或 \textit{stå over for et svært valg}（面临艰难抉择）。
\end{itemize}
\end{grammatikbox}

\begin{grammatikbox}{语法点 2：中英对比排雷——为什么不能照搬英语 \textit{is expected to}？}
从英语思维切换写丹麦语政论文时，下面四个常见失误在 Modul 4 考场最易失分：
\begin{enumerate}
  \item \textbf{严禁用状态被动 \textit{er forventet at} 代替 \textit{forventes at}！}：英语写 \textit{The number is expected to rise}，很多同学直译成 \ding{55}\ \textit{Antallet er forventet at stige}。但在丹麦语中，表达普遍政策预期与客观趋势趋势预测\textbf{必须用 \textit{-s} 被动}：\ding{51}\ \textit{Antallet \dabold{forventes} at stige.}（也绝对不能漏掉被动 \textit{-s} 写成主动 \ding{55}\ \textit{Antallet forventer at stige}——数量自己可不会“期待”！）。
  \item \textbf{“面临挑战”用 \textit{stå over for}，别漏掉 \textit{over}！}：英语是 \textit{face a challenge}，中文是“面对挑战”，容易漏掉介词或误写成 \ding{55}\ \textit{stå for en udfordring}（那变成“代表/负责某事”了！）。正确固定三词组是：\ding{51}\ \textit{stå \dabold{over for} en alvorlig udfordring}。
  \item \textbf{\textit{there is expected to be...} 怎么翻译？}：如果主语不确定（如“预计将有更多人退休”），不能像英语那样写 \textit{*Der forventes at være flere...}，最自然地道的丹麦语是切回\textbf{形式主语 \textit{Det} + 存在句 \textit{der}}：\ding{51}\ \textit{\dabold{Det forventes, at der} bliver mangel på personale.}
  \item \textbf{句首状语触发 V2 倒装}：如 \textit{Ifølge nye prognoser \dabold{forventes} antallet at stige.} / \textit{Som følge af denne udvikling \dabold{står} kommunerne over for en udfordring.}（介词短语占第一位，动词 \textit{forventes} / \textit{står} 立刻紧随其后！）。
\end{enumerate}
\end{grammatikbox}

% ---------------------------------------------------------------------
\ncdsubsec{五、Mønsterøvelser（句型与辩论操练）}

\textbf{练习 A：精准选填（用 \textit{forventes, skønnes, vurderes, står, over, at, det} 等词补全政策句）}
\begin{enumerate}
  \item Antallet af borgere over 80 år \underline{\hspace{2.2cm}} at stige kraftigt i de kommende ti år. (forvente)
  \item De danske kommuner \underline{\hspace{2.2cm}} over for en alvorlig udfordring i ældreplejen. (stå)
  \item \underline{\hspace{2.2cm}} vurderes, at der snart vil mangle flere tusinde sosu-assistenter.
  \item Den nye teknologi \underline{\hspace{2.2cm}} at kunne spare personalet for meget tid om natten. (skønne)
  \item Et nyt plejehjem forventes \underline{\hspace{2.2cm}} blive bygget i udkanten af byen næste år.
  \item En skærm menes ikke at \underline{\hspace{2.2cm}} erstatte den menneskelige kontakt i hverdagen. (kunne)
  \item På grund af færre unge på arbejdsmarkedet står samfundet \underline{\hspace{2.2cm}} for et svært valg.
\end{enumerate}

\textbf{练习 B：句型升级与诊错（将普通句改写为高阶被动预期句，或修正语病）}
\begin{enumerate}
  \item 主语提升改写：\textit{Eksperterne forventer, at behovet for hjemmepleje vokser.} $\rightarrow$ \textit{Behovet for hjemmepleje...}
  \item 转为形式主语 \textit{Det} 句：\textit{Nye robotter vurderes at hjælpe personalet i fremtiden.} $\rightarrow$ \textit{Det...}
  \item 句首状语倒装合并：\textit{Prognosen viser nye tal. Udgifterne til sundhed skønnes at stige.}（用 \textit{Ifølge prognosen...} 合并）$\rightarrow$
  \item 英语直译纠错：\textit{Antallet af ældre er forventet at stige, og kommunerne står for en stor udfordring.} $\rightarrow$
  \item 存在句预测改写（改写为 \textit{Det forventes, at der...}）：\textit{Vi forventer, at der kommer flere frivillige til plejehjemmene.} $\rightarrow$
\end{enumerate}

\textbf{练习 C：中译丹（活用政策预测与挑战表达）}
\begin{enumerate}
  \item 到 2035 年，丹麦 80 岁以上老人的数量预计将大幅增加。
  \item 由于人手短缺，各市镇正面临一项严峻的挑战。
  \item 据估算，福利科技将在未来能够减轻员工的负担。
  \item 人们评估，将来人们将需要更多志愿者来协助养老照护。
  \item 据专家称，家庭成员预计将在未来扮演更重要的角色。
\end{enumerate}

% ---------------------------------------------------------------------
\ncdsubsec{六、DU3 Modul 4 实战微任务与参考答案}

\begin{emailbox}{DU3 Modul 4 Skrivning Opgave 1 实战微任务（私人长邮件 · minimum 125 ord）}
\textbf{Situation}: Du har fået en e-mail fra din danske ven Kasper, hvis bedstemor lige er flyttet på plejehjem. Han overvejer selv at uddanne sig til social- og sundhedsassistent, men er lidt usikker.\\
\textbf{Opgave}: Læs Kaspers e-mail og skriv et svar på \textbf{mindst 125 ord}. Du skal svare på alle hans spørgsmål og gerne bruge mindst to forventningsstrukturer fra lektionen (\textit{forventes at, skønnes at, stå over for en udfordring}).

\vspace{0.3em}
\textit{Hej!\\
Jeg håber, du har det godt. Min bedstemor flyttede på plejehjem i sidste uge, og jeg kan mærke, hvor travlt personalet har. Hvordan ser du på ældreplejen i Danmark? Tror du, at teknologi kan løse manglen på personale, eller bør familien hjælpe mere til? Og hvad synes du om min idé om at blive sosu-assistent --- ville det være et godt job i fremtiden?\\
Mange hilsner, Kasper}
\end{emailbox}

\begin{answerbox}
\textbf{Før du læser 参考回答}
\begin{enumerate}
  \item \textit{Kommunerne står over for en alvorlig udfordring, fordi antallet af borgere over 80 år forventes at stige kraftigt, mens manglen på social- og sundhedsassistenter skønnes at vokse. Eksperter peger på rekruttering af udenlandsk personale, velfærdsteknologi og et tættere samarbejde med pårørende og frivillige.}
  \item \textit{Sofie er bekymret for, at teknologi ikke kan erstatte menneskelig kontakt, og mener, at kvaliteten forventes at falde ved besparelser. Henrik fremhæver derimod, at digitale hjælpemidler kan frigøre tid til omsorg, og at frivillige besøgsvenner kan give ekstra nærvær.}
\end{enumerate}

\textbf{Tekst B 阅读核对答案}
\begin{enumerate}
  \item \da{Forkert}（索菲明确表示屏幕无法握住老人的手，技术只能略微帮忙，不能取代人际接触）
  \item \da{Det er allerede svært at besætte alle vagter.}（他所在养老院已经很难排满所有班次）
  \item \da{Ekstra nærvær i hverdagen.}（提供日常生活中的额外陪伴，同时不接管专业护理职责）
\end{enumerate}

\textbf{练习 A 答案}
\begin{enumerate}
  \item \da{forventes}（主语提升被动预期结构：\textit{forventes at stige}）
  \item \da{står}（固定搭配：\textit{stå over for en alvorlig udfordring}）
  \item \da{Det}（形式主语引导完整从句：\textit{Det vurderes, at...}）
  \item \da{skønnes}（客观估算：\textit{skønnes at kunne spare}）
  \item \da{at}（主语提升不定式标记：\textit{forventes at blive bygget}）
  \item \da{kunne}（否定词 \textit{ikke} 放在 \textit{at} 之前，后接情态动词原形）
  \item \da{over}（固定三词搭配：\textit{stå over for}）
\end{enumerate}

\textbf{练习 B 答案（句型改写参考答案）}
\begin{enumerate}
  \item \textit{Behovet for hjemmepleje \dabold{forventes at vokse}.}（去掉主观施动者，用主语提升 \textit{-s} 被动表达客观预测）
  \item \textit{\dabold{Det vurderes, at} nye robotter kan hjælpe personalet i fremtiden.}（也可写为：\textit{Det vurderes, at nye robotter hjælper personalet i fremtiden.}）
  \item \textit{Ifølge prognosen \dabold{skønnes} udgifterne til sundhed at stige.}（句首 \textit{Ifølge prognosen} 占第 1 位，触发 \textit{skønnes} V2 倒装）
  \item \textit{Antallet af ældre \dabold{forventes} at stige, og kommunerne \dabold{står over for} en stor udfordring.}（不用 \textit{er forventet at}，且补齐 \textit{over}）
  \item \textit{\dabold{Det forventes, at der} kommer flere frivillige til plejehjemmene.}（用形式主语 \textit{Det} 接 \textit{der} 存在句）
\end{enumerate}

\textbf{练习 C 答案}
\begin{enumerate}
  \item \textit{Frem mod 2035 forventes antallet af ældre over 80 år i Danmark at stige kraftigt.}
  \item \textit{På grund af mangel på personale står kommunerne over for en alvorlig udfordring.}
  \item \textit{Velfærdsteknologi skønnes i fremtiden at kunne aflaste de ansatte i ældreplejen.}
  \item \textit{Det vurderes, at man fremover får brug for flere frivillige i ældreplejen.}
  \item \textit{Ifølge eksperter forventes familiemedlemmer at spille en større rolle i fremtiden.}
\end{enumerate}

\textbf{DU3 Modul 4 高分范文（154 词，不含署名）}\\
\textit{Hej Kasper}

\textit{Tak for din mail! Jeg kan godt forstå, at det gjorde stort indtryk på dig at se, hvor travlt personalet har på din bedstemors plejehjem.}

\textit{Efter min mening \da{står Danmark over for en alvorlig udfordring} på ældreområdet. Antallet af ældre borgere \da{forventes at stige} hurtigt de næste ti år, og derfor bliver der brug for både nye løsninger og flere hænder. Teknologi kan jo være en god hjælp ved praktiske opgaver, men jeg mener ikke, at skærme kan erstatte menneskeligt nærvær. Samtidig synes jeg, at familien gerne må hjælpe med hyggelige besøg, mens det professionelle ansvar stadig bør ligge hos det offentlige.}

\textit{Derfor synes jeg faktisk, at din idé om at blive sosu-assistent er rigtig spændende! Der \da{skønnes at være} stor efterspørgsel efter dygtige omsorgsmedarbejdere i mange år fremover, så du vil næsten altid være sikker på at finde et meningsfuldt job.}

\textit{Skriv endelig, hvad du beslutter dig for!}

\textit{Bedste hilsner\\
\mbox{[Dit navn]}}

\small\textbf{名师点评}：正文 154 词（不含署名），远超 Modul 4 Opgave 1 要求的 125 词底线。结构紧凑自然，完整回应了 Kasper 关于“现状看法”“科技 vs 家庭”“职业选择”三大问题；巧妙嵌入 Unit 1 情态小品词 \textit{jo} 与本课核心骨架 \textit{står over for en alvorlig udfordring}、\textit{forventes at stige} 以及 \textit{skønnes at være}，兼具朋友对话的真诚温度与社会议题的思辨深度。
\end{answerbox}
\ncdlesson{7}{SU: En investering i fremtiden eller en dyr gavebod?}{助学金 SU：未来投资还是昂贵福利？}
{DU3 Modul 4（B1+ $\rightarrow$ B2-）｜对标 NCE 3 \& Modultest 4 \textit{Skrivning Opgave 2 (Debatindlæg)}}
{\textbf{正反方观点对立框架}（\textit{Tilhængere fremhæver, at...} $\leftrightarrow$ \textit{Modstandere indvender derimod, at...}）· 辩论立场动词与副词 \textit{derimod / imidlertid} 的主从句语序 · 双视角平衡论证}
{丹麦大学生每月无需偿还的助学金 SU（\textit{Statens Uddannelsesstøtte}）是北欧福利制度最鲜明的象征之一，但“硕士阶段是否应改为贷款”几乎每年都会登上丹麦报刊头条。在 DU3 Modul 4 的博客辩论写作（\textit{Opgave 2}）和双人口试中，高分关键在于能用清晰、对仗、客观的丹麦语\textbf{把支持者与反对者的核心论点针锋相对地摆上台面}。}

% ---------------------------------------------------------------------
\begin{foerduleaeser}
带着下面两个问题阅读 Tekst A 和 Tekst B：
\begin{enumerate}
  \item \textit{Hvilke argumenter bruger henholdsvis tilhængere og modstandere i debatten om kandidat-SU ifølge Tekst A?}（根据 Tekst A，支持者和反对者在关于硕士阶段 SU 的争论中分别提出了哪些论据？）
  \item \textit{Hvorfor mener Maja i Tekst B, at SU er en god investering, og hvad foreslår Mikkel i stedet?}（在 Tekst B 中，玛雅为什么认为 SU 是一项良好的投资？米克尔又提出了什么替代方案？）
\end{enumerate}
\end{foerduleaeser}

% ---------------------------------------------------------------------
\ncdsubsec{二、Tekster（福利国家专题与正反辩论双语篇）}

\subsection*{Tekst A：\textit{Fri uddannelse til alle --- eller en uholdbar udgift?}（专题辩论综述 · 约 250 词）}

\begin{tipbox}{语篇说明}
本篇呈现丹麦社会关于“削减硕士生 SU”的经典正反交锋。阅读时请特别留意\textbf{\color{NCDNavy}深蓝色标记的观点引出动词}（\textit{fremhæver, indvender, argumenterer for, peger på}）与\textbf{\color{NCDRed}深红色标记的转折对比位置}（\textit{derimod, hvorimod, imidlertid}）。
\end{tipbox}

I Danmark får næsten alle studerende på videregående uddannelser hver måned udbetalt SU af staten uden at skulle betale pengene tilbage. Ordningen har i årtier været en hjørnesten i den danske velfærdsmodel. Men i de senere år har politikere og økonomer diskuteret livligt, om støtten til kandidatstuderende bør ændres til et studielån.

\da{Tilhængere af den nuværende model fremhæver, at} gratis uddannelse og SU skaber høj social mobilitet. Når unge fra alle samfundslag har råd til at studere uden at arbejde tredive timer om ugen ved siden af, kan de koncentrere sig om bøgerne. \da{Fortalere argumenterer desuden for, at} veluddannede borgere senere betaler høje skatter tilbage til fællesskabet. Derfor \da{betragter de} ikke SU som en udgift, men som en langsigtet investering i hele samfundets fremtid.

\da{Modstandere indvender derimod, at} kandidat-SU i praksis fungerer som en dyr gavebod til fremtidens højtlønnede. \da{Kritikere peger på, at} akademikere som regel tjener væsentligt mere end faglærte gennem hele arbejdslivet, \dabold{hvorimod} de ofte først kommer ud på arbejdsmarkedet sent i tyverne. Et lån på de sidste to studieår ville \dabold{derimod} opmuntre flere studerende til at færdiggøre studiet hurtigere og frigøre milliarder af kroner til folkeskolen eller erhvervsuddannelserne.

Debatten viser således et klassisk dilemma i velfærdsstaten: Skal fællesskabet sikre helt lige muligheder for alle gennem brede tilskud, eller bør støtten målrettes dem, der har mindst? Hvor den ene side ser tryghed og retfærdighed, ser den anden side et skævt system, der kræver reformer.

\begin{transbox}
\small
\textbf{\color{NCDTeal}[Tekst A 中文精译]}\\
在丹麦，几乎所有高等教育阶段的学生每月都能从国家领取助学金 SU，且无需偿还。几十年来，这一制度一直是丹麦福利模式的基石。但近年来，政治家和经济学家们热烈讨论，硕士研究生阶段的助学金是否应当改为助学贷款。

现行模式的支持者强调，免费教育和 SU 创造了高度的社会流动性。当来自各个社会阶层的年轻人都能负担得起学业，而不必每周兼职三十小时时，他们就能专心读书。此外，倡导者还论证说，受过良好教育的公民日后会向社会回馈高额税收。因此，他们并不将 SU 视为一笔开销，而是视为对整个社会未来的长期投资。

然而反对者反驳称，硕士阶段 SU 在实践中就像一场昂贵的“免费大派送”，补贴的恰恰是未来的高收入人群。批评者指出，大学毕业生在整个职业生涯中的收入通常显著高于技工群体，而他们往往要到将近三十岁才进入职场。相比之下，最后两年采用贷款方式反而会激励更多学生更快完成学业，并为公立中小学或职业教育释放出数十亿克朗的资金。

因此，这场争论折射出福利国家的一个经典两难困境：社会是应当通过普惠补贴确保所有人完全平等的机会，还是应当将资助精准投向最缺少资源的人？一方看到的是安全感与公平，另一方看到的却是一个亟待改革的失衡体制。
\end{transbox}

\subsection*{Tekst B：\textit{To læserbreve i universitetsavisen}（观点交锋对读 · 约 140 词）}

\begin{emailbox}{校园报刊正反来信｜Emne: Skal kandidat-SU afskaffes?}
\textbf{Maja:} Jeg kommer fra en familie uden akademisk baggrund, og uden SU havde jeg aldrig turdet starte på universitetet. \da{Jeg vil gerne fremhæve, at} støtten giver ro til fordybelse, så også unge uden rige forældre kan få en lang uddannelse. Hvis staten fjerner stipendiet på kandidaten, risikerer vi, at færre tør drømme stort af frygt for gæld.

\textbf{Mikkel:} Jeg forstår godt Majas pointe, men vi glemmer let prioriteringen. \da{Kritikere af systemet indvender} med rette, at en håndværker via sin skat betaler for en kommende jurists eller læges sidste studieår. Et rentefrit studielån på kandidaten er \dabold{derimod} både rimeligt og solidarisk, fordi de penge i stedet kan styrke daginstitutioner og erhvervsskoler.
\end{emailbox}

\textbf{Tekst B 阅读核对（Hvem siger hvad? / Kort svar，答案见第六部分）：}
\begin{enumerate}
  \item \textit{Hvem argumenterer for, at et bortfald af SU kan afskrække unge fra ikke-akademiske hjem på grund af gæld?} \underline{\hspace{2.2cm}}
  \item \textit{Hvilket økonomisk argument bruger Mikkel imod den nuværende kandidat-SU?} \\\underline{\hspace{5.0cm}}
  \item \textit{Hvad foreslår Mikkel at bruge de sparede penge til i stedet?} \underline{\hspace{4.5cm}}
\end{enumerate}

\begin{tipbox}{本课发音与语调要点（Udtale \& Reduktioner）}
\begin{itemize}
  \item \textbf{辩论阵营与立场核心词}：\textit{en tilhænger} \ipa{ˈtˢelˌhɛŋˀɐ}（重音在第一音节，\textit{-hæng-} 带 \textit{stød}）；\textit{en modstander} \ipa{ˈmoðˌsdɑnˀɐ}（\textit{mod} 软 \textit{d}，\textit{-nd-} 中 \textit{d} 不发音！）；\textit{en kritiker} \ipa{ˈkʰʁitˢikʰɐ}（注意第一音节短重音）。
  \item \textbf{态度动词与转折词}：\textit{fremhæve} \ipa{ˈfʁɑmˌhɛːˀvə}（\textit{frem-} 元音偏 \ipa{ɑ}！）；\textit{indvende} \ipa{ˈenˌvɛnˀə}（两处 \textit{d} 均不发音！）；\textit{derimod} \ipa{ˈdɛɐ̯iˌmoːˀð}；\textit{hvorimod} \ipa{ˈvɒːiˌmoːˀð}（词首 \textit{h} 不发音）。
  \item \textbf{教育福利高频词}：\textit{SU} \ipa{ˈɛsˈuːˀ}（重音在字母 U）；\textit{social mobilitet} \ipa{soˈɕæːˀl mobiliˈtˢeːˀd}；\textit{en gavebod} \ipa{ˈɡæːvəˌboːˀð}；\textit{et studielån} \ipa{ˈsd̥uːˀdiəˌlɔːˀn}；\textit{en gæld} \ipa{ˈɡɛlˀ}（\textit{d} 不发音！）。
\end{itemize}
\end{tipbox}

\begin{transbox}
\small
\textbf{\color{NCDTeal}[Tekst B 中文对照]}\\
\textbf{玛雅（支持保留 SU）}：我来自一个没有大学学历背景的家庭，如果没有 SU，我当初根本不敢报考大学。我想强调的是，这笔助学金能让人安心钻研学业，让父母并不富裕的年轻人也有机会读长学制学位。如果国家取消硕士阶段的助学金，我们就会面临这样的风险：越来越多人因为害怕背债而不敢胸怀远大梦想。

\textbf{米克尔（主张硕士改贷款）}：我完全理解玛雅的观点，但我们很容易忽视财政优先次序。制度批评者有理有据地反驳说，一名蓝领技工正在用自己交的税为未来的律师或医生支付最后两年的学业开销。相比之下，硕士阶段实行免息助学贷款既合情合理又体现团结互助，因为省下的资金可以转而用于加强托幼机构和职业学校。
\end{transbox}

% ---------------------------------------------------------------------
\ncdsubsec{三、Gloser（核心生词表）}

\begin{tipbox}{词表重点}
辩论双方称呼有清晰对比：\textit{en tilhænger / en fortaler}（支持者/倡导者，常接介词 \textit{af}）对立于 \textit{en modstander / en kritiker}（反对者/批评者，常接介词 \textit{af}）。动词 \textit{at argumentere} 后面必须接介词 \textit{for} 或 \textit{imod} 再加 \textit{, at} 从句；\textit{at betragte} 固定搭配 \textit{som}（把……视为……）。注意 \textit{en gæld}（债务）不可数且词尾 \textit{d} 不发音！
\end{tipbox}

\begin{glosertable}
名词 & en uddannelsesstøtte & \textit{(uddannelsesstøtten, kun ental)} & 教育助学金（简称 SU） \\
名词 & en gavebod & \textit{(gaveboden, gaveboder, gaveboderne)} & 免费大派送，过慷慨的福利 \\
名词 & en investering & \textit{(investeringen, investeringer, -erne)} & 投资 \\
名词 & en tilhænger (af) & \textit{(tilhængeren, tilhængere, -ne)} & 支持者，拥护者 \\
名词 & en modstander (af) & \textit{(modstanderen, modstandere, -ne)} & 反对者，对手 \\
名词 & en kritiker (af) & \textit{(kritikeren, kritikere, -ne)} & 批评者，评论家 \\
名词 & en mobilitet & \textit{(mobiliteten, kun ental)} & 流动性（\textit{social mobilitet}） \\
名词 & et studielån & \textit{(studielånet, studielån, studielånene)} & 助学贷款 \\
名词 & en gæld & \textit{(gælden, kun ental)} & 债务，欠款 \\
名词 & en kandidatuddannelse & \textit{(kandidatuddannelsen, -er, -erne)} & 硕士学位课程 \\
\midrule
动词 [har] & at fremhæve & \textit{(fremhæver, fremhævede, har fremhævet)} & 强调，突出指出 \\
动词 [har] & at indvende & \textit{(indvender, indvendte, har indvendt)} & 反驳，提出异议 \\
动词 [har] & at argumentere (for/imod) & \textit{(argumenterer, -ede, har argumenteret)} & 论证（支持/反对） \\
动词 [har] & at betragte (som) & \textit{(betragter, betragtede, har betragtet)} & 把……视作，看待 \\
动词 [har] & at finansiere & \textit{(finansierer, finansierede, har finansieret)} & 资助，为……提供资金 \\
动词 [har] & at afskaffe & \textit{(afskaffer, afskaffede, har afskaffet)} & 废除，取消（制度） \\
\midrule
形容词 & langsigtet / solidarisk & \textit{langsigtet (-, langsigtede) / solidarisk (-, -e)} & 长期的 / 团结互助的 \\
\end{glosertable}

% ---------------------------------------------------------------------
\ncdsubsec{四、Grammatik \& Stilistik（核心语法、政策句式与修辞精讲）}

\begin{grammatikbox}{语法点 1：正反方观点对立框架——三步把辩论写得针锋相对}
在 DU3 Modul 4 \textit{Opgave 2 (Debatindlæg)} 和口试双人辩论中，千万不能只会用 \textit{Nogle mener..., andre mener...}。下面这套“对立修辞矩阵”能立刻让你的论述具备报刊评论的锋芒与平衡感：

\vspace{0.3em}
\renewcommand{\arraystretch}{1.24}
{\small
\begin{tabularx}{\linewidth}{@{} >{\bfseries\raggedright\arraybackslash}p{2.5cm} >{\raggedright\arraybackslash}p{5.5cm} >{\raggedright\arraybackslash}X @{}}
\toprule
\textbf{论证环节} & \textbf{核心对仗框架} & \textbf{词汇替换库与注意要点} \\
\midrule
\da{1. 立正方}\newline （提出支持论点） &
\textit{\da{Tilhængere fremhæver, at}...}\newline
\textit{\da{Fortalere argumenterer for, at}...} &
后接从句（内部不倒装，\textit{ikke} 在动词前）；可用 \textit{betragte A som B} 表积极定性。 \\
\midrule
\da{2. 切反方}\newline （针锋相对反驳） &
\textit{\da{Modstandere indvender derimod, at}...}\newline
\textit{\da{Kritikere peger imidlertid på, at}...} &
注意 \dabold{derimod / imidlertid}（然而/反过来）紧跟主句动词之后！若置于句首，必须引发倒装。 \\
\midrule
\da{3. 视角并置}\newline （一句浓缩双立场） &
\textit{\da{Hvor} den ene side ser X, \dabold{betragter} den anden side det som Y.}\newline
\textit{... \da{hvorimod} de ofte kommer sent ud.} &
\textit{Hvor} 从句置于句首时，后半主句必须 \textbf{V2 倒装}（\textit{betragter den anden side}）；\textit{hvorimod} 从句内部不倒装！ \\
\bottomrule
\end{tabularx}
}
\end{grammatikbox}

\begin{grammatikbox}{语法点 2：转折位置与介词搭配排雷（避免高频扣分陷阱）}
写正反辩论时，最容易在以下四处语序与介词搭配上踩雷：
\begin{enumerate}
  \item \textbf{\textit{argumentere} 必须带介词 \textit{for / imod} 才能接从句！}：英语可以直接写 \textit{argue that...}，但丹麦语\textbf{绝不能写} \ding{55}\ \textit{Fortalere argumenterer, at SU er god.} 必须明确表态是支持还是反对：\ding{51}\ \textit{Fortalere argumenterer \dabold{for, at} SU er god.} / \textit{De argumenterer \dabold{imod, at} man fjerner støtten.}
  \item \textbf{\textit{derimod} 的两个黄金位置（别弄混倒装！）}：
    \begin{itemize}
      \item 放在主句动词后面（最推荐、最流畅）：\ding{51}\ \textit{Modstandere \dabold{indvender derimod}, at ordningen er for dyr.}
      \item 放在句首强调对立（必须 V2 倒装！）：\ding{51}\ \textit{\dabold{Derimod indvender} modstandere, at ordningen er for dyr.}（绝不能漏掉倒装写成 \ding{55}\ \textit{Derimod modstandere indvender...}！）。
    \end{itemize}
  \item \textbf{\textit{indvende} 与 \textit{afvise} 的语义区分}：\textit{at indvende, at...} 是“提出反驳意见/指出相反理据”，语气理性客观；\textit{at afvise, at...} 则是“断然否认或拒绝接受某种说法/方案”，语气更强硬。
  \item \textbf{\textit{Hvor den ene side..., ... den anden side...}（对比性 \textit{Hvor} 从句）}：这里的 \textit{hvor} 不是指地点“在哪里”，而是书面政论文中极高级的对比引导词“就……而言 / 在一方看来……而另一方……”，记得前半部分是从句，后半主句动词一定要紧接在逗号后面（V2 倒装）！
\end{enumerate}
\end{grammatikbox}

% ---------------------------------------------------------------------
\ncdsubsec{五、Mønsterøvelser（句型与辩论操练）}

\textbf{练习 A：精准选填（用 \textit{tilhængere, modstandere, fremhæver, indvender, derimod, for, som, hvorimod} 等词补全正反对比句）}
\begin{enumerate}
  \item \underline{\hspace{2.2cm}} af gratis SU mener, at alle unge bør have lige chancer for at studere.
  \item Mange fortalere \underline{\hspace{2.2cm}}, at uddannelse er en langsigtet investering for hele samfundet. (现在时)
  \item Modstandere \underline{\hspace{2.2cm}} derimod, at kandidat-SU er en alt for dyr udgift for staten. (现在时)
  \item Økonomer argumenterer \underline{\hspace{2.2cm}}, at et studielån kan få flere studerende hurtigere gennem studiet. (支持……)
  \item Nogle studerende vil gerne bevare stipendiet. Andre mener \underline{\hspace{2.2cm}}, at et lån er mere retfærdigt.
  \item Fortalere betragter SU \underline{\hspace{2.2cm}} en vigtig sikkerhed for unge fra ikke-akademiske hjem.
  \item Hvor den ene side ser tryghed, \underline{\hspace{2.2cm}} den anden side et uretfærdigt system til de rigeste. (ser)
\end{enumerate}

\textbf{练习 B：句型升级与诊错（把平淡表述改写为高阶正反辩论句，或修正语病）}
\begin{enumerate}
  \item 合并成对立句：\textit{Tilhængere fremhæver, at SU giver social mobilitet. Modstandere siger, at ordningen koster for meget.}（第二句用 \textit{indvender derimod} 改写）$\rightarrow$
  \item 英语直译纠错：\textit{Mange politikere argumenterer, at staten bør ændre reglerne for SU.} $\rightarrow$
  \item \textit{derimod} 句首倒装改写：\textit{Et lån er derimod en bedre løsning for samfundet.}（把 \textit{Derimod} 移至句首）$\rightarrow$
  \item 高阶 \textit{Hvor} 对比改写：\textit{Den ene side ser SU som en investering, og den anden side betragter den som en gavebod.}（用 \textit{Hvor...} 改写，注意主句倒装）$\rightarrow$
  \item 从句语序纠错：\textit{Kritikere peger på, at man bør ikke give SU til studerende, hvorimod nogle mener derimod, at støtten er nødvendig.}（修正两处语序与用词错误）$\rightarrow$
\end{enumerate}

\textbf{练习 C：中译丹（活用正反方观点对立框架）}
\begin{enumerate}
  \item 支持者强调，免费助学金使所有年轻人都能安心学习。
  \item 然而反对者反驳说，这项制度对纳税人来说花费太高。
  \item 批评者主张废除硕士助学金，转而提供免息助学贷款。
  \item 一些人把 SU 视为一种投资，而另一些人却将其视为昂贵的福利派送。
  \item 相比之下，支持者指出社会流动性在丹麦起着重要作用。
\end{enumerate}

% ---------------------------------------------------------------------
\ncdsubsec{六、DU3 Modul 4 实战微任务与参考答案}

\begin{emailbox}{DU3 Modul 4 Skrivning Opgave 2 实战微任务（Debatindlæg · minimum 100 ord）}
\textbf{Situation}: Peter Nørgaard har skrevet et debatindlæg på nettet om SU til kandidatstuderende.\\
\textbf{Opgave}: Skriv et debatindlæg på \textbf{mindst 100 ord}, hvor du diskuterer Peters synspunkt. Er du enig? Hvorfor? / Hvorfor ikke? Brug gerne både positive og kritiske modargumenter med lektionens rammer.

\vspace{0.3em}
\textit{\textbf{Gør kandidat-SU til et lån!}\\
Jeg synes, det er helt urimeligt, at almindelige skatteydere skal betale SU til unge på kandidatuddannelsen. De fleste med en lang uddannelse får jo gode lønninger bagefter. Hvis de selv måtte låne pengene til de sidste to år, ville de også blive hurtigere færdige med studiet. Derfor bør vi fjerne kandidat-SU nu!}
\end{emailbox}

\begin{answerbox}
\textbf{Før du læser 参考回答}
\begin{enumerate}
  \item \textit{Tilhængere fremhæver, at SU skaber høj social mobilitet og giver studerende ro til at fokusere på studiet, samt at de senere betaler høje skatter tilbage. Modstandere indvender derimod, at kandidat-SU er en dyr gavebod til fremtidens højtlønnede, og at et lån ville få flere til at studere hurtigere.}
  \item \textit{Maja mener, at SU er en god investering, fordi den giver unge fra ikke-akademiske familier mod og ro til at tage en lang uddannelse uden frygt for gæld. Mikkel foreslår i stedet et rentefrit studielån på kandidaten, så pengene kan bruges på daginstitutioner og erhvervsskoler.}
\end{enumerate}

\textbf{Tekst B 阅读核对答案}
\begin{enumerate}
  \item \da{Maja}（她指出取消硕士助学金会令没有学术背景的年轻人因害怕债务而不敢读长学制大学）
  \item \da{En håndværker betaler via sin skat for kommende juristers eller lægers sidste studieår.}（普通技工纳税补贴未来高收入律师或医生的最后几年学业）
  \item \da{Daginstitutioner og erhvervsskoler.}（用于加强托幼机构和职业学校）
\end{enumerate}

\textbf{练习 A 答案}
\begin{enumerate}
  \item \da{Tilhængere}（名词复数“支持者”，与后文动词复数意义及 \textit{af} 搭配）
  \item \da{fremhæver}（观点引出动词：\textit{fremhæver, at...}）
  \item \da{indvender}（反驳动词：\textit{indvender derimod, at...}）
  \item \da{for}（固定搭配：\textit{argumentere for, at...} 表示“论证支持……”）
  \item \da{derimod} / \da{imidlertid}（放在主句动词之后表达对立转折）
  \item \da{som}（固定搭配：\textit{betragte A som B}）
  \item \da{ser}（\textit{Hvor} 从句前置触发主句 V2 倒装：\textit{ser den anden side...}）
\end{enumerate}

\textbf{练习 B 答案（句型改写参考答案）}
\begin{enumerate}
  \item \textit{Tilhængere fremhæver, at SU giver social mobilitet, mens modstandere derimod indvender, at ordningen koster for meget.}（也可分两句写为：\textit{Tilhængere fremhæver, at SU giver social mobilitet. Modstandere indvender derimod, at ordningen koster for meget.}）
  \item \textit{Mange politikere argumenterer \dabold{for, at} staten bør ændre reglerne for SU.}（必须补上介词 \textit{for} 或 \textit{imod}）
  \item \textit{\dabold{Derimod er} et lån en bedre løsning for samfundet.}（\textit{Derimod} 在句首必须引发 V2 倒装）
  \item \textit{Hvor den ene side ser SU som en investering, \dabold{betragter} den anden side den som en gavebod.}（\textit{Hvor} 从句前置，主句动词 \textit{betragter} 紧跟逗号后）
  \item \textit{Kritikere peger på, at man \dabold{ikke bør} give SU til studerende, mens nogle \dabold{derimod mener}, at støtten er nødvendig.}（修正两点：\ding{172} 从句内 \textit{ikke} 放在情态动词 \textit{bør} 前；\ding{173} \textit{hvorimod} 与 \textit{derimod} 语义重复且从句不倒装，改为 \textit{mens... derimod mener}）
\end{enumerate}

\textbf{练习 C 答案}
\begin{enumerate}
  \item \textit{Tilhængere fremhæver, at gratis SU gør det muligt for alle unge at studere i ro.}
  \item \textit{Modstandere indvender derimod, at denne ordning koster skatteyderne for mange penge.}
  \item \textit{Kritikere argumenterer for at afskaffe kandidat-SU og i stedet tilbyde et rentefrit studielån.}
  \item \textit{Nogle betragter SU som en investering, mens andre derimod ser den som en dyr gavebod.}
  \item \textit{Fortalere peger derimod på, at social mobilitet spiller en vigtig rolle i Danmark.}
\end{enumerate}

\textbf{DU3 Modul 4 高分范文（144 词）}\\
\textit{\textbf{Kandidat-SU er stadig en klog investering}}

\textit{Jeg er kun delvist enig med Peter Nørgaard i hans indlæg om kandidat-SU. \da{Modstandere indvender derimod} ofte, at akademikere senere får høje lønninger, og at skatteyderne derfor ikke bør betale for de sidste to studieår. Det er da også et vigtigt argument, når velfærdsstaten skal prioritere sine penge.}

\textit{Ikke desto mindre \da{vil jeg fremhæve, at} kandidat-SU spiller en afgørende rolle for den sociale mobilitet i Danmark. Unge uden velstående forældre tør måske ikke tage en lang uddannelse, hvis de ved, at de står med en stor gæld bagefter. Derudover \da{argumenterer fortalere med rette for, at} højtuddannede betaler langt mere i skat gennem hele arbejdslivet og dermed betaler støtten tilbage flere gange.}

\textit{Derfor \da{betragter jeg} fortsat kandidat-SU \da{som} en fornuftig investering. Man kan jo i stedet stille krav om hurtigere studieaktivitet frem for at fjerne trygheden helt.}

\small\textbf{名师点评}：144 词完整回应辩论贴要求（远超 100 词下限）。文章先以退为进、复现了 Unit 1 的 \textit{Ikke desto mindre} 和小品词 \textit{da / jo}，接着熟练调度本课正反框架（\textit{Modstandere indvender derimod... / vil jeg fremhæve, at... / argumenterer for, at... / betragter som}），做到了有理有据、温和而有力地反驳对手观点，是典型的 B2- 级别优秀 \textit{Debatindlæg}。
\end{answerbox}
\ncdlesson{8}{4-dages arbejdsuge: Produktivitet eller ønsketænkning?}{四天工作制：生产力还是空想？}
{DU3 Modul 4（B1+ $\rightarrow$ B2-）｜对标 NCE 3 \& Modultest 4 \textit{Arbejdsliv \& Hypotetisk Argumentation}}
{\textbf{非真实条件句与虚拟语气}（过去完成时 \textit{havde/var} + 分词 $\rightarrow$ \textit{ville have/være} + 分词）· 无 \textit{hvis} 虚拟条件倒装（\textit{Havde man...}）· 情态虚拟（\textit{kunne have / burde have}）}
{近年来，“每周四天工作制”在丹麦市政厅和企业界掀起热潮：如果每周少上一天班，究竟能提高专注力与生产力，还是纯粹一厢情愿的空想？在复盘试点项目或反驳假设情景时，丹麦语依靠\textbf{“过去完成时 + \textit{ville have} + 分词”}构建严谨的非真实条件推演。本课帮你彻底摆脱英语 \textit{would have} 错用陷阱，精准说出与过去或现实相反的高分虚拟句！}

% ---------------------------------------------------------------------
\begin{foerduleaeser}
带着下面两个问题阅读 Tekst A 和 Tekst B：
\begin{enumerate}
  \item \textit{Hvorfor var forsøget med 4-dages arbejdsuge i en jysk kommune ikke blot en gave, og hvad ville der være sket uden strammere mødekultur ifølge Tekst A?}（根据 Tekst A，为什么日德兰某市镇的四天工作制试点并不只是“白给福利”？如果不收紧会议文化，原本会发生什么？）
  \item \textit{Hvordan adskiller Frejas kontorarbejde sig fra Lukas' hverdag på hospitalet i Tekst B?}（在 Tekst B 中，芙蕾雅的办公室工作与卢卡斯的医院日常有何不同？）
\end{enumerate}
\end{foerduleaeser}

% ---------------------------------------------------------------------
\ncdsubsec{二、Tekster（福利国家专题与正反辩论双语篇）}

\subsection*{Tekst A：\textit{Mindre tid på kontoret --- mere værdi?}（政策试点剖析 · 约 250 词）}

\begin{tipbox}{语篇说明}
本篇剖析丹麦地方政府与企业界关于“四天工作制”试点的真实经验与争论。阅读时请特别留意\textbf{\color{NCDNavy}深蓝色标记的非真实条件句}（\textit{Hvis ... havde/var..., ville ... have/være...}）与\textbf{\color{NCDRed}深红色标记的无 \textit{hvis} 倒装虚拟及情态推演}。
\end{tipbox}

Kan man arbejde fire dage om ugen til fuld løn uden at miste produktivitet? Spørgsmålet deler stadig arbejdsmarkedet i to lejre. I flere danske kommuner og it-virksomheder har man allerede gennemført et år lange forsøg med 4-dages arbejdsuge. Erfaringerne viser, at resultatet afhænger helt af, hvordan arbejdet bliver organiseret.

I en jysk kommune indførte man en ny struktur med kortere møder og faste koncentrationstimer uden afbrydelser. Efter forsøget konkluderede ledelsen, at \da{hvis afdelingen ikke havde ændret} sin mødekultur først, \da{ville projektet hurtigt være mislykkedes}. Medarbejderne fik desuden bedre trivsel, og sygefraværet faldt. \dabold{Havde kommunen fortsat} med den gamle rytme, \dabold{ville de ansatte næppe have nået} alle opgaver på fire dage.

Samtidig advarer økonomer imod at tro, at modellen virker i alle brancher. På et kontor kan man ofte skære overflødige møder væk, men på et plejehjem eller i en børnehave kan man ikke arbejde hurtigere. \da{Hvis hospitalerne havde indført} fire dages arbejdsuge uden at ansætte flere sygeplejersker, \da{ville ventelisterne blot være blevet} endnu længere.

Debatten handler derfor om realisme frem for ønsketænkning. Nogle eksperter mener, at man i stedet \dabold{burde have testet} mere fleksible vagtplaner fra begyndelsen. Uanset om man er for eller imod, kræver færre arbejdsdage ny ledelse --- ellers risikerer man blot, at folk bliver lige så udbrændte på fire dage som på fem.

\begin{transbox}
\small
\textbf{\color{NCDTeal}[Tekst A 中文精译]}\\
人能否每周只工作四天却拿全薪，同时又不损失生产力？这个问题至今仍将劳动力市场分为两大阵营。在丹麦多个市镇和 IT 企业中，人们已经开展了为期一年的四天工作制试点。经验表明，最终结果完全取决于工作是如何组织的。

在日德兰半岛的一个市镇，管理层推行了新架构，缩短会议时间并设立不受打扰的固定专注时段。试点结束后管理层总结道，如果该部门当初没有先改变自己的会议文化，项目很快就会失败。此外，员工的身心幸福感也有所提升，病假缺勤率随之下降。如果市镇当时继续沿用旧的工作节奏，员工们恐怕几乎不可能在四天内完成所有任务。

与此同时，经济学家警告人们不要以为这一模式在所有行业都行得通。在办公室里，人们往往可以砍掉多余会议，但在养老院或幼儿园里却无法“加快速度干活”。如果医院在没有多招护士的情况下就推行四天工作制，候诊名单原本只会变得更长。

因此，这场争论关乎现实可行性而非一厢情愿的空想。一些专家认为，人们当初本该从一开始就测试更灵活的排班方案。无论支持还是反对，减少工作天数都需要全新的管理方式——否则人们只会面临另一种风险：员工在四天里照样和五天一样精疲力竭。
\end{transbox}

\subsection*{Tekst B：\textit{To brancher, to virkeligheder}（职场人物对谈 · 约 140 词）}

\begin{emailbox}{广播辩论连线｜Freja (HR-konsulent) \& Lukas (Sygeplejerske)}
\textbf{Freja:} Hos os indførte vi 4-dages arbejdsuge sidste forår, og vores effektivitet er faktisk steget. \da{Hvis vi havde ventet} på den perfekte plan, \da{ville vi aldrig have opdaget}, hvor meget tid vi før spildte på lange statusmøder. Jeg tror også, at færre ville gå ned med stress, hvis flere firmaer turde prøve det.

\textbf{Lukas:} Det lyder jo fantastisk for jer på et kontor, Freja. Men på min afdeling har vi patienter døgnet rundt. \dabold{Var der blevet} indført 4-dages uge hos os i fjor uden ekstra hænder, \dabold{kunne vi ikke have dækket} nattevagterne. En model for alle ville simpelthen være urealistisk i sundhedsvæsenet.
\end{emailbox}

\textbf{Tekst B 阅读核对（Rigtigt / Forkert / Kort svar，答案见第六部分）：}
\begin{enumerate}
  \item \textit{Effektiviteten i Frejas virksomhed faldt, efter at de indførte 4-dages arbejdsuge.} \underline{\hspace{2.2cm}}
  \item \textit{Hvad opdagede Frejas firma først, efter at de prøvede 4-dages uge?} \underline{\hspace{4.5cm}}
  \item \textit{Hvad ville der være sket på Lukas' afdeling sidste år, hvis man havde indført 4-dages uge uden ekstra personale?} \underline{\hspace{4.5cm}}
\end{enumerate}

\begin{tipbox}{本课发音与语调要点（Udtale \& Reduktioner）}
\begin{itemize}
  \item \textbf{虚拟语气核心助动词链的连读弱化}：\textit{ville have} 在自然语流中极度轻读，\textit{have} 绝对不能重读成 \ipa{hæːvə}，而是弱化为 \ipa{hɑ}（或几乎融入上一词）：\textit{ville have nået} \ipa{ˈvilə hɑ ˈnɔːˀəð}；\textit{kunne have dækket} \ipa{ˈkʰunə hɑ ˈdɛɡəð}；\textit{havde} 常弱读为 \ipa{ˈhæðə} 或极短的 \ipa{hæð}。
  \item \textbf{职场与效率高频词}：\textit{en arbejdsuge} \ipa{ˈɑːbɑi̯ˀdsˌuːə}；\textit{produktivitet} \ipa{pʰʁoduɡtˢiviˈtˢeːˀd}；\textit{effektivitet} \ipa{efɛɡtˢiviˈtˢeːˀd}（重音都在末音节，带软化的长元音加 \textit{stød}）。
  \item \textbf{静音辅音与软 \textit{d}}：\textit{ønsketænkning} \ipa{ˈønsɡəˌtˢɛŋɡneŋ}；\textit{sygefravær} \ipa{ˈsyːəˌfʁɑːˌvɛɐ̯ˀ}（\textit{syge} 中的 \textit{g} 不发音！）；\textit{udbrændt} \ipa{ˈuðˌbʁɛnˀd}（\textit{ud-} 为软 \textit{d}，\textit{-nd-} 中的 \textit{d} 不发音！）；\textit{døgnet rundt} \ipa{ˈdʌi̯ˀnəð ˈʁɔnˀd}（\textit{øgn} 发 \ipa{ʌi̯ˀ}）。
\end{itemize}
\end{tipbox}

\begin{transbox}
\small
\textbf{\color{NCDTeal}[Tekst B 中文对照]}\\
\textbf{芙蕾雅（人力资源顾问）}：我们公司去年春天推行了四天工作制，工作效率其实反而提高了。如果我们当时一直等一份完美方案出来才动手，我们就永远不会发现以前在漫长汇报会上浪费了多少时间。我也相信，如果更多企业敢于尝试，因压力崩溃病倒的人就会更少。

\textbf{卢卡斯（医院护士）}：芙蕾雅，这对你们坐办公室的人听起来当然太棒了。但在我们科室，一天二十四小时都有病人需要照料。如果去年在没有增加人手的情况下就在我们这里推行四天工作制，我们原本根本排不开夜班。在医疗系统里，一刀切的方案简直不切实际。
\end{transbox}

% ---------------------------------------------------------------------
\ncdsubsec{三、Gloser（核心生词表）}

\begin{tipbox}{词表重点}
\textit{at mislykkes}（失败）与第 3 册收官课学的 \textit{at lykkes}（成功）互为反义异态动词：两者的完成时和过去完成时都\textbf{必须用 \textit{er / var}} 作助动词（\textit{det er/var mislykkedes}），因此在非真实条件句中说“原本会失败”必须写 \dabold{ville være mislykkedes}，切勿错用 \textit{ville have}！
\end{tipbox}

\begin{glosertable}
名词 & en arbejdsuge & \textit{(arbejdsugen, arbejdsuger, -ne)} & 工作周 \\
名词 & en produktivitet & \textit{(produktiviteten, kun ental)} & 生产力，生产率 \\
名词 & en ønsketænkning & \textit{(ønsketænkningen, kun ental)} & 一厢情愿的空想 \\
名词 & et sygefravær & \textit{(sygefraværet, kun ental)} & 病假缺勤 \\
名词 & et forsøg & \textit{(forsøget, forsøg, forsøgene)} & 试点，实验，尝试 \\
名词 & en mødekultur & \textit{(mødekulturen, -kulturer, -erne)} & 会议文化 \\
名词 & en afbrydelse & \textit{(afbrydelsen, afbrydelser, -erne)} & 打断，干扰 \\
名词 & en branche & \textit{(branchen, brancher, brancherne)} & 行业，领域 \\
名词 & en trivsel & \textit{(trivselen, kun ental)} & 身心愉悦感，工作幸福感 \\
名词 & en fleksibilitet & \textit{(fleksibiliteten, kun ental)} & 灵活性 \\
\midrule
动词 [har] & at indføre & \textit{(indfører, indførte, har indført)} & 引进，推行（制度） \\
动词 [har] & at effektivisere & \textit{(effektiviserer, -ede, har effektiviseret)} & 提高……的效率，精简 \\
异态动词 \dabold{[er]} & at mislykkes & \textit{(mislykkes, mislykkedes, \dabold{er mislykkedes})} & 失败，落空 \\
动词 [har] & at skære ned (på) & \textit{(skærer, skar, har skåret ned på)} & 削减，缩减 \\
动词 [har] & at spilde & \textit{(spilder, spildte, har spildt)} & 浪费（时间/金钱） \\
\midrule
形容词 & realistisk & \textit{(realistisk, realistiske)} & 现实可行的 \\
形容词 & udbrændt & \textit{(udbrændt, udbrændte)} & 精疲力竭的，职业倦怠的 \\
\end{glosertable}

% ---------------------------------------------------------------------
\ncdsubsec{四、Grammatik \& Stilistik（核心语法、政策句式与修辞精讲）}

\begin{grammatikbox}{语法点 1：非真实条件句与虚拟语气（\textit{Kontrafaktisk konjunktiv}）}
在辩论社会政策（如四天工作制试点、财政改革）时，常常需要推演“\textbf{当初如果那样做了，原本会发生什么}”或“\textbf{假如现在普遍这样，将会怎样}”。丹麦语不使用特殊动词词尾，而是靠\textbf{时态后移 + \textit{ville}} 精准区分两类虚拟：

\vspace{0.3em}
\renewcommand{\arraystretch}{1.24}
{\small
\begin{tabularx}{\linewidth}{@{} >{\bfseries\raggedright\arraybackslash}p{2.9cm} >{\raggedright\arraybackslash}p{4.5cm} >{\raggedright\arraybackslash}X @{}}
\toprule
\textbf{虚拟类型} & \textbf{标准句法公式} & \textbf{典型辩论例句} \\
\midrule
\dabold{1. 过去非真实条件}\newline （复盘与事实相反的过去） &
\textbf{条件句：}\textit{Hvis} + \dabold{havde / var + 分词}\newline
\textbf{结果句：}\dabold{ville + have / være + 分词} &
\textit{Hvis vi \dabold{havde ventet}, \dabold{ville} vi aldrig \dabold{have opdaget} det.}\newline
\textit{Hvis projektet ikke \dabold{var startet}, \dabold{ville} stressen \dabold{være steget}.} \\
\midrule
\da{2. 现在/将来假设}\newline （设想眼下未实现的方案） &
\textbf{条件句：}\textit{Hvis} + \da{过去时}\newline
\textbf{结果句：}\da{ville + 动词原形} &
\textit{Hvis alle firmaer \da{turde} prøve det, \da{ville} færre gå ned med stress.}\newline
\textit{Hvis jeg \da{var} leder, \da{ville} jeg teste modellen først.} \\
\bottomrule
\end{tabularx}
}
\vspace{0.35em}

\textbf{进阶修辞三大武器}：
\begin{itemize}
  \item \ding{172} \textbf{省略 \textit{hvis} 的虚拟倒装（最高分报刊句型）}：直接把 \textit{Havde / Var} 放到从句最开头，省略 \textit{hvis}，后面结果句同样紧跟 \textit{ville/kunne} 形成双倒装：\textit{\dabold{Havde kommunen} fortsat med den gamle rytme, \dabold{ville} de ansatte næppe \dabold{have nået} målet.}
  \item \ding{173} \textbf{情态虚拟（“本可以 / 本应该……”）}：把结果句的 \textit{ville} 换成 \textit{kunne}（本能够）或 \textit{burde}（本应该），直接构成：\dabold{kunne have + 分词} / \dabold{burde have + 分词}（例：\textit{Man \dabold{burde have testet} vagtplanerne først.} / \textit{Vi \dabold{kunne ikke have dækket} vagterne.}）。
  \item \ding{174} \textbf{助动词 \textit{have} vs \textit{være} 铁律不变}：变化/移动/结果状态动词（如 \textit{blive, være, stige, falde, mislykkes, flytte}）在虚拟条件句和结果句中照样用 \dabold{var} 和 \dabold{ville være}，不要一律写成 \textit{have}！
\end{itemize}
\end{grammatikbox}

\begin{grammatikbox}{语法点 2：中英对比排雷——杜绝四大虚拟语气高频错误}
\begin{enumerate}
  \item \textbf{严禁在 \textit{hvis} 条件句里面塞入 \textit{ville}！}：受部分口语化美式英语或错误类推影响，绝对不能写成 \ding{55}\ \textit{Hvis firmaet \dabold{ville have indført} ordningen før...}！记住黄金铁律：\textbf{\textit{ville} 永远只能出现在结果主句里，\textit{hvis} 从句内部只能用过去完成时（\textit{havde/var + 分词}）或一般过去时}：\ding{51}\ \textit{Hvis firmaet \dabold{havde indført} ordningen før, \dabold{ville} trivslen \dabold{have været} bedre.}
  \item \textbf{从句前置时的 V2 倒装与否定词 \textit{ikke} 位置}：
    \begin{itemize}
      \item 在 \textit{hvis} 从句里，\textit{ikke / aldrig} 依然放在 \textit{havde/var} 前面：\ding{51}\ \textit{Hvis afdelingen \dabold{ikke havde ændret} mødekulturen...}
      \item 而在省略 \textit{hvis} 的倒装条件句里，\textit{ikke} 放在主语后面：\ding{51}\ \textit{Havde afdelingen \dabold{ikke} ændret mødekulturen...}
      \item 在后面的结果主句中，\textit{ikke / næppe / aldrig} 放在 \textit{ville} 后面、\textit{have/være} 前面：\ding{51}\ \textit{..., \dabold{ville} vi \dabold{aldrig have opdaget} det.}
    \end{itemize}
  \item \textbf{别漏掉 \textit{have}！}：受英语缩略或中文影响，不能把过去虚拟结果句写成 \ding{55}\ \textit{Vi ville opdaget det}（漏了 \textit{have}）！只要表达“当初本会……”，\textit{ville} 和过去分词之间必须有 \dabold{have}（或 \dabold{være}）。
  \item \textbf{状态动词 \textit{blive} 要配 \textit{være} 还是 \textit{have}？}：注意被动或状态变化：\textit{Var der blevet indført...}（如果当时被推行了）以及 \textit{ville ventelisterne være blevet længere}（名单原本会变得更长），绝对不能用 \textit{havde blevet}！
\end{enumerate}
\end{grammatikbox}

% ---------------------------------------------------------------------
\ncdsubsec{五、Mønsterøvelser（句型与辩论操练）}

\textbf{练习 A：精准填空（使用正确的非真实虚拟结构：\textit{havde, var, ville, kunne, have, være} 等）}
\begin{enumerate}
  \item Hvis firmaet \underline{\hspace{2.2cm}} indført 4-dages uge sidste år, ville sygefraværet være faldet.
  \item Uden faste koncentrationstimer ville projektet desværre \underline{\hspace{2.2cm}} mislykkedes.
  \item Hvis hospitalet ikke \underline{\hspace{2.2cm}} ansat flere sygeplejersker, ville vagterne have manglet folk.
  \item \underline{\hspace{2.2cm}} kommunen spurgt de ansatte tidligere, ville forsøget have været nemmere.
  \item Hvis alle brancher arbejdede fire dage nu, \underline{\hspace{2.2cm}} nogle tjenester måske blive dyrere.
  \item Ledelsen burde \underline{\hspace{2.2cm}} testet den nye mødekultur før sommerferien.
  \item Hvis ventelisterne \underline{\hspace{1.8cm}} blevet længere sidste efterår, ville politikerne have stoppet forsøget.
\end{enumerate}

\textbf{练习 B：句型升级与纠错（将真实经历转为过去非真实假设，或改用无 \textit{hvis} 倒装及纠错）}
\begin{enumerate}
  \item 将真实过去事实改为过去虚拟句：\textit{Vi ændrede ikke mødekulturen. Derfor nåede vi ikke opgaverne.}（用 \textit{Hvis vi havde...} 改写）$\rightarrow$
  \item 省略 \textit{hvis} 改写为高阶倒装：\textit{Hvis kommunen havde planlagt vagterne bedre, ville personalet have været gladere.} $\rightarrow$
  \item 英语污染纠错：\textit{Hvis jeg ville have vidst det i går, havde jeg ikke taget overarbejde.} $\rightarrow$
  \item 助动词与语序纠错：\textit{Hvis forsøget ikke havde blevet en succes, ville ledelsen aldrig have fortsat.}（注意 \textit{blive} 搭配的助动词！）$\rightarrow$
  \item 用情态虚拟 \textit{kunne have} 改写：\textit{Uden ekstra personale var det umuligt for os at dække alle nattevagter i fjor.} $\rightarrow$ \textit{Uden ekstra personale...}
\end{enumerate}

\textbf{练习 C：中译丹（活用非真实条件句与虚拟语气）}
\begin{enumerate}
  \item 如果我们早点缩短会议，员工原本就会有更多时间专注于工作。
  \item 倘若没有额外的护士，医院原本就不可能推行四天工作制。
  \item 如果我当时是团队负责人，我早就先测试这个试点项目了。
  \item 如果企业没有改变会议文化，这个计划去年就已经失败了。
  \item 当初要是有更多人支持，这个想法本来会变得更现实。
\end{enumerate}

% ---------------------------------------------------------------------
\ncdsubsec{六、DU3 Modul 4 实战微任务与参考答案}

\begin{emailbox}{DU3 Modul 4 Skrivning Opgave 2 实战微任务（Debatindlæg · minimum 100 ord）}
\textbf{Situation}: Henrik Dalgaard har skrevet et debatindlæg i en lokalavis om 4-dages arbejdsuge.\\
\textbf{Opgave}: Skriv et debatindlæg på \textbf{mindst 100 ord}, hvor du diskuterer Henriks synspunkt. Er du enig? Hvorfor? / Hvorfor ikke? Inddrag mindst én hypotetisk eller kontrafaktisk sætning (\textit{Hvis..., ville...}).

\vspace{0.3em}
\textit{\textbf{4-dages arbejdsuge er urealistisk ønsketænkning!}\\
Jeg synes, det lyder dejligt at arbejde fire dage til samme løn, men det er kun for folk på et kontor. På sygehuse, i daginstitutioner og i butikker skal der jo være personale hele ugen. Hvis alle fik fri om fredagen, ville samfundet hurtigt bryde sammen. Derfor bør vi glemme drømmen om en 4-dages arbejdsuge og i stedet holde fast i 37 timer.}
\end{emailbox}

\begin{answerbox}
\textbf{Før du læser 参考回答}
\begin{enumerate}
  \item \textit{Fordi kommunen samtidig indførte kortere møder og faste koncentrationstimer uden afbrydelser. Hvis afdelingen ikke havde ændret sin mødekultur først, ville projektet hurtigt være mislykkedes.}
  \item \textit{Freja arbejder på et kontor, hvor de kunne fjerne lange statusmøder og øge effektiviteten med en 4-dages uge. Lukas er sygeplejerske på et hospital med patienter døgnet rundt, hvor man ikke kan dække nattevagterne uden ekstra hænder.}
\end{enumerate}

\textbf{Tekst B 阅读核对答案}
\begin{enumerate}
  \item \da{Forkert}（恰恰相反，芙蕾雅说公司推行后效率反而上升了：\textit{vores effektivitet er faktisk steget}）
  \item \da{Hvor meget tid de før spildte på lange statusmøder.}（发现以前在冗长例会上浪费了多少时间）
  \item \da{De kunne ikke have dækket nattevagterne.}（如果去年没有额外人手，医院根本无法覆盖夜班）
\end{enumerate}

\textbf{练习 A 答案}
\begin{enumerate}
  \item \da{havde}（过去非真实条件从句：\textit{havde indført}）
  \item \da{være}（异态动词 \textit{mislykkes} 完成时用 \textit{være}：\textit{ville være mislykkedes}）
  \item \da{havde}（否定词在从句助动词前：\textit{ikke havde ansat}）
  \item \da{Havde}（省略 \textit{hvis} 的条件倒装，置于句首大写）
  \item \da{ville}（现在/将来假设条件句对应的结果主句）
  \item \da{have}（情态虚拟：\textit{burde have testet}“当初本应该测试”）
  \item \da{var}（状态变化动词 \textit{blive} 过去完成时用 \textit{var blevet}）
\end{enumerate}

\textbf{练习 B 答案（句型改写参考答案）}
\begin{enumerate}
  \item \textit{Hvis vi \dabold{havde ændret} mødekulturen, \dabold{ville} vi \dabold{have nået} opgaverne.}（将过去未做的事情反推为虚拟条件）
  \item \textit{\dabold{Havde kommunen} planlagt vagterne bedre, \dabold{ville} personalet \dabold{have været} gladere.}（省略 \textit{hvis}，双倒装）
  \item \textit{Hvis jeg \dabold{havde vidst} det i går, \dabold{ville} jeg ikke \dabold{have taget} overarbejde.}（\textit{hvis} 从句严禁用 \textit{ville}！）
  \item \textit{Hvis forsøget ikke \dabold{var blevet} en succes, ville ledelsen aldrig have fortsat.}（\textit{blive} 必须用 \textit{var blevet}，不能用 \textit{havde blevet}）
  \item \textit{Uden ekstra personale \dabold{kunne vi ikke have dækket} alle nattevagter i fjor.}
\end{enumerate}

\textbf{练习 C 答案}
\begin{enumerate}
  \item \textit{Hvis vi havde forkortet møderne tidligere, ville medarbejderne have haft mere tid til fokus.}
  \item \textit{Uden ekstra sygeplejersker kunne hospitalet ikke have indført en 4-dages arbejdsuge.}
  \item \textit{Hvis jeg havde været teamleder dengang, ville jeg have testet dette forsøg først.}
  \item \textit{Havde virksomheden ikke ændret mødekulturen, ville planen være mislykkedes sidste år.}
  \item \textit{Havde flere støttet forslaget dengang, ville ideen være blevet mere realistisk.}
\end{enumerate}

\textbf{DU3 Modul 4 高分范文（142 词）}\\
\textit{\textbf{Fleksibilitet frem for én model til alle}}

\textit{Jeg er delvist enig med Henrik Dalgaard i, at en klassisk 4-dages arbejdsuge ikke fungerer ens i alle brancher. På hospitaler og i børnehaver kan man jo ikke bare skære åbningstiden ned uden at ansætte flere hænder. \da{Hvis} alle skoler pludselig \da{lukkede} om fredagen, \da{ville} mange forældre \da{få} svært ved at passe deres job.}

\textit{Alligevel mener jeg ikke, at ideen kun er ønsketænkning. Ganske vist passer en fast fredag fri bedst til kontorfag, men man \dabold{burde have tænkt} mere på fleksible vagtplaner i stedet for at afvise forslaget helt. \da{Havde} flere virksomheder \da{prøvet} kortere møder og mere koncentrationstid, \da{ville} mange ansatte allerede \da{have undgået} stress og sygefravær.}

\textit{Derfor bør vi ikke holde fast i 37 timer for enhver pris, men give hver arbejdsplads frihed til at finde sin egen løsning.}

\small\textbf{名师点评}：142 词结构严谨、说理克制而有层次。文章既承认了 Henrik 关于医院与托幼部门的现实限制（并运用了现在虚拟 \textit{Hvis... lukkede..., ville... få...}），又用 \textit{Ganske vist..., men...} 让步转折，自然融入情态虚拟 \textit{burde have tænkt} 与无 \textit{hvis} 过去虚拟倒装 \textit{Havde... prøvet..., ville... have undgået...}，展现出了优秀的 B2- 思辨操控力。
\end{answerbox}
\ncdlesson{9}{Skat med glæde? Debatten om topskatten}{快乐纳税？关于最高边际税的争论}
{DU3 Modul 4（B1+ $\rightarrow$ B2-）｜对标 NCE 3 \& Modultest 4 \textit{Økonomi, Velfærd \& Kausalitet}}
{\textbf{严谨因果链条副词与连词}（\textit{som følge heraf, af den grund, følgelig, hvorved}）· 句首主句 V2 倒装副词 vs 逗号后非倒装从句连词 \textit{hvorved} · \textit{som følge af}（介词）vs \textit{som følge heraf}（副词）}
{在国际媒体眼中，丹麦人以“全球最高税率之一”却依然高度支持福利国家而著称。然而，“高收入者是否应该缴纳最高边际税（\textit{topskat}）”一直是丹麦议会最激烈的经济辩论。在论证减税或加税的连锁经济后果时，光用简单的 \textit{derfor} 和 \textit{så} 远不够 B2- 水准。本课带你精准驾驭 \textbf{\textit{som følge heraf, af den grund, følgelig, hvorved}} 四大逻辑链条衔接词！}

% ---------------------------------------------------------------------
\begin{foerduleaeser}
带着下面两个问题阅读 Tekst A 和 Tekst B：
\begin{enumerate}
  \item \textit{Hvad frygter modstandere af topskatten, at skatten fører til, og hvad peger tilhængere på ifølge Tekst A?}（根据 Tekst A，最高边际税的反对者担心它会导致什么后果？而支持者又指出了什么？）
  \item \textit{Hvilket konkret valg står Thomas over for i Tekst B, og hvordan forklarer Emma topskattens rolle for velfærden?}（在 Tekst B 中，托马斯面临着什么具体抉择？埃玛如何解释最高边际税对福利的作用？）
\end{enumerate}
\end{foerduleaeser}

% ---------------------------------------------------------------------
\ncdsubsec{二、Tekster（福利国家专题与正反辩论双语篇）}

\subsection*{Tekst A：\textit{Solidaritet eller bremse på væksten?}（经济政策评论 · 约 250 词）}

\begin{tipbox}{语篇说明}
本篇以严谨的经济学因果链条剖析丹麦关于 \textit{topskat}（最高边际税）的争论。阅读时请特别区分\textbf{\color{NCDNavy}深蓝色加粗的句间因果副词}（\textit{som følge heraf, af den grund, følgelig}——触发 V2 倒装）与\textbf{\color{NCDRed}深红色加粗的从句连词 \textit{hvorved}}（后接从句不倒装）。
\end{tipbox}

Danmark hører til blandt de lande i verden med det højeste skattetryk. Til gengæld får borgerne fri adgang til hospitaler, universiteter og et bredt sikkerhedsnet. Et af de mest omdiskuterede elementer i skattesystemet er topskatten, som betales af den del af indkomsten, der ligger over en bestemt grænse.

Tilhængere ser topskatten som et nødvendigt redskab til at sikre social lighed. Når de højeste indkomster beskattes ekstra, bliver forskellen mellem rig og fattig mindre, \dabold{hvorved} sammenhængskraften i samfundet styrkes. Staten får desuden større indtægter til skoler og sygehuse. \da{Som følge heraf} kan alle borgere nyde godt af fælles velfærd, uanset hvor meget de selv tjener. \da{Af den grund} mener mange vælgere stadig, at de bredeste skuldre skal bære mest.

Modstandere peger derimod på, at høj skat på de sidste tjente kroner kan svække lysten til at arbejde ekstra. Hvis en specialist oplever, at mere end halvdelen af en lønstigning går til skat, takker vedkommende måske nej til overarbejde eller et større ansvar. Virksomhederne mister \da{følgelig} værdifuld arbejdskraft, \dabold{hvorved} både produktionen og samfundets samlede velstand kan gå glip af vækst.

Økonomer er derfor dybt uenige om, hvorvidt man bør sænke eller helt afskaffe topskatten. \da{Af den grund} foreslår et flertal ofte et kompromis med en højere beløbsgrænse. På den måde håber man at bevare tilliden til systemet, uden at drivkraften i økonomien forsvinder.

\begin{transbox}
\small
\textbf{\color{NCDTeal}[Tekst A 中文精译]}\\
丹麦位列世界上税负最高的国家之一。作为回报，公民可以免费使用医院、大学并享有广泛的社会保障网。税制中最具争议的组成部分之一便是最高边际税（topskat），即对超出某一特定门槛的那部分收入所征收的额外税款。

支持者将最高边际税视为确保社会平等的必要工具。当最高收入群体被额外征税时，贫富差距就会缩小，由此社会凝聚力得以增强。此外，国家也为学校和医院获得了更多财政收入。受此影响，所有公民无论个人挣多少钱，都能受惠于公共福利。正因如此，许多选民至今仍认为“肩膀最宽的人应当承担最重的担子”。

然而反对者指出，对多挣来的最后一笔钱征收高额税款，可能会削弱人们加班或再多工作的意愿。如果一名专业人士发现涨薪的一半以上都交了税，当事人或许就会拒绝加班或承担更大职责。企业随之就会流失宝贵的劳动力，从而导致生产以及整个社会的总体繁荣都有可能错失增长机会。

因此，经济学家们对于究竟应当下调还是彻底废除最高边际税存在严重分歧。出于这一原因，议会多数派往往会提出提高起征点门槛的折中方案。人们希望借此既能保住对福利制度的信任，又不至于让经济中的进取动力消失。
\end{transbox}

\subsection*{Tekst B：\textit{Kaffe og skattedebat mellem venner}（对话体邮箱讨论 · 约 140 词）}

\begin{emailbox}{线上讨论｜Thomas (Læge) \& Emma (Økonom)}
\textbf{Thomas:} Hej Emma! Jeg er lige blevet tilbudt flere ekstravagter på sygehuset, men så ryger hele den ekstra løn op i topskat. Jeg overvejer ærligt talt at sige nej, for der bliver jo ikke så meget tilbage efter skat. Har du som økonom ikke også den fornemmelse, at topskatten straffer flid?

\textbf{Emma:} Hej Thomas! Jeg forstår godt din frustration, og mange siger jo nej til ekstra timer \da{af den grund}. Men husk, at du kun betaler topskat af beløbet over grænsen. De ekstra penge går direkte til vores sygehuse og daginstitutioner, \dabold{hvorved} hele landet fungerer tryggere. \da{Som følge heraf} har vi et stabilt samfund, som vi begge drager nytte af hver dag.
\end{emailbox}

\textbf{Tekst B 阅读核对（Rigtigt / Forkert / Kort svar，答案见第六部分）：}
\begin{enumerate}
  \item \textit{Thomas betaler topskat af hele sin samlede løn, så snart han tager én ekstravagt.} \underline{\hspace{2.2cm}}
  \item \textit{Hvad overvejer Thomas at gøre, og hvorfor?} \underline{\hspace{4.8cm}}
  \item \textit{Hvordan gavner skatteindtægterne hele samfundet ifølge Emma?} \underline{\hspace{4.8cm}}
\end{enumerate}

\begin{tipbox}{本课发音与语调要点（Udtale \& Reduktioner）}
\begin{itemize}
  \item \textbf{三大因果连词/副词真实读音（注意静音 \textit{f} 与半元音 \textit{g}！）}：
    \begin{itemize}
      \item \textit{som følge heraf} \ipa{sʌm ˈføljə ˈhɛɐ̯ˌɑːˀ}（\textit{følge} 中的 \textit{g} 读 \ipa{j}，\textit{heraf} 的 \textit{f} 在口语常脱落读 \ipa{ɑːˀ}！）
      \item \textit{af den grund} \ipa{ɑ ˈdɛnˀ ˈɡʁɔnˀ}（\textit{af} 弱读 \ipa{ɑ}，\textit{grund} 词尾 \textit{d} 完全不发音！）
      \item \textit{følgelig} \ipa{ˈføljəli}（第一个 \textit{g} 读 \ipa{j}，词尾 \textit{-ig} 的 \textit{g} 不发音！）
      \item \textit{hvorved} \ipa{vɒˈveð}（词首 \textit{hv-} 的 \textit{h} 不发音，结尾软 \textit{d} \ipa{ð}）
    \end{itemize}
  \item \textbf{税务与经济词汇}：\textit{en topskat} \ipa{ˈtˢʌbˌsɡad}；\textit{en indkomst} \ipa{ˈenˌkʰʌmsd}（\textit{d} 不发音）；\textit{en indtægt} \ipa{ˈenˌtˢɛɡd}；\textit{en ulighed} \ipa{ˈuˌliːheðˀ}（\textit{g} 不发音）；\textit{sammenhængskraft} \ipa{ˈsɑmənˌhɛŋsˌkʰʁɑfd}。
\end{itemize}
\end{tipbox}

\begin{transbox}
\small
\textbf{\color{NCDTeal}[Tekst B 中文对照]}\\
\textbf{托马斯（医生）}：嗨，埃玛！医院刚刚给我提供了好几个额外班次，但那样一来多挣的工资全都要交最高边际税了。老实说我在考虑拒绝，因为扣完税后真剩不下多少。作为经济学家，难道你也不觉得最高边际税是在惩罚勤奋工作吗？

\textbf{埃玛（经济学家）}：嗨，托马斯！我完全理解你的沮丧，很多人也正是出于这个原因才拒绝加班。但别忘了，你只对超出起征点门槛的那部分金额缴纳最高边际税。这些额外的税收直接流向我们的医院和托幼机构，从而让整个国家运转得更有安全感。因此我们拥有了一个稳定的社会，我们俩每天都在从中受益。
\end{transbox}

% ---------------------------------------------------------------------
\ncdsubsec{三、Gloser（核心生词表）}

\begin{tipbox}{词表重点}
注意区分两个长得很像的高频词：\textit{en indtægt}（收入/财政进项，复数 \textit{indtægter}，常指国家或企业总收入）与 \textit{en indkomst}（个人所得收入）；以及名词短语中的介词搭配：\textit{at nyde godt af} / \textit{at drage nytte af}（从……中受益），\textit{at gå glip af}（错失……）。
\end{tipbox}

\begin{glosertable}
名词 & en topskat & \textit{(topskatten, kun ental)} & 最高边际税，高收入附加税 \\
名词 & en indtægt & \textit{(indtægten, indtægter, indtægterne)} & 收入，财政进项 \\
名词 & en indkomst & \textit{(indkomsten, indkomster, -erne)} & 个人收入，所得 \\
名词 & en skatteyder & \textit{(skatteyderen, skatteydere, -ne)} & 纳税人 \\
名词 & en ulighed & \textit{(uligheden, uligheder, ulighederne)} & 不平等，贫富差距 \\
名词 & en solidaritet & \textit{(solidariteten, kun ental)} & 团结互助，社会共济 \\
名词 & en drivkraft & \textit{(drivkraften, drivkræfter, -ne)} & 驱动力，进取动力 \\
名词 & en velstand & \textit{(velstanden, kun ental)} & 繁荣，富裕程度 \\
名词 & et sikkerhedsnet & \textit{(sikkerhedsnettet, -net, -nettene)} & 社会保障网 \\
名词 & et kompromis & \textit{(kompromiset, kompromiser, -erne)} & 妥协，折中方案 \\
\midrule
动词 [har] & at beskatte & \textit{(beskatter, beskattede, har beskattet)} & 向……征税 \\
动词 [har] & at sænke & \textit{(sænker, sænkede, har sænket)} & 降低，下调（税率/标准） \\
动词 [har] & at svække & \textit{(svækker, svækkede, har svækket)} & 削弱，减弱 \\
动词 [har] & at straffe & \textit{(straffer, straffede, har straffet)} & 惩罚，处罚 \\
动词 [har] & at bevare & \textit{(bevarer, bevarede, har bevaret)} & 保留，维持 \\
\midrule
连接词/副词 & følgelig / hvorved & \textit{(句间因果副词) / (从句因果连词)} & 因此，随之 / 借此，从而使 \\
\end{glosertable}

% ---------------------------------------------------------------------
\ncdsubsec{四、Grammatik \& Stilistik（核心语法、政策句式与修辞精讲）}

\begin{grammatikbox}{语法点 1：严谨因果链条四剑客——“主句副词 V2 倒装” vs “\textit{hvorved} 从句不倒装”}
写 B2- 政论文时，单纯用 \textit{derfor} 和 \textit{så} 很容易显得单调且偏口语。下面四个书面语核心衔接词必须严格分清\textbf{词性}与\textbf{语序}：

\vspace{0.3em}
\renewcommand{\arraystretch}{1.24}
{\small
\begin{tabularx}{\linewidth}{@{} >{\bfseries\raggedright\arraybackslash}p{2.7cm} >{\raggedright\arraybackslash}p{3.9cm} >{\raggedright\arraybackslash}X @{}}
\toprule
\textbf{因果衔接词} & \textbf{语法属性与语序铁律} & \textbf{标准政论例句} \\
\midrule
\da{som følge heraf}\newline （受此影响，因此） &
\textbf{句子副词}（连接两句独立主句）\newline
句首必触发 \textbf{V2 倒装}；句中放动词后 &
\textit{\da{Som følge heraf} \dabold{kan alle} nyde godt af fælles velfærd.} \\
\midrule
\da{af den grund}\newline （出于这一原因） &
\textbf{介词短语作副词}\newline
句首必触发 \textbf{V2 倒装}；也可放句尾或动词后 &
\textit{\da{Af den grund} \dabold{mener mange}, at de rigeste bør betale mest.}\newline
\textit{Mange takker nej \da{af den grund}.} \\
\midrule
\da{følgelig}\newline （随之，因而） &
\textbf{书面副词}（紧扣前因顺推后果）\newline
句首 \textbf{V2 倒装}；常紧随主句动词后 &
\textit{Virksomhederne mister \da{følgelig} arbejdskraft.}\newline
\textit{\da{Følgelig} \dabold{falder} motivationen hos nogle.} \\
\midrule
\dabold{hvorved}\newline （借此，从而使） &
\textbf{关系/从属连词}（逗号后引出自然结果）\newline
内部必须守 \textbf{从句语序（不倒装，\textit{ikke} 前置）}！ &
\textit{Forskellen bliver mindre, \dabold{hvorved} sammenhængskraften \dabold{ikke svækkes}.} \\
\bottomrule
\end{tabularx}
}
\end{grammatikbox}

\begin{grammatikbox}{语法点 2：中英对比排雷——因果链路中最危险的四大失误}
\begin{enumerate}
  \item \textbf{严禁把 \textit{som følge heraf} 当介词接名词！}：记住这一对形近词的铁打分工：
    \begin{itemize}
      \item 后面直接跟\textbf{名词}，必须用 \da{som følge af}（由于……的结果）：\ding{51}\ \textit{\da{Som følge af} den høje skat siger han nej.}
      \item 前句已经交代了原因，本句单独用来承接上文“因此/受此影响”，才用 \da{som følge heraf}（其中 \textit{heraf} 本身就等于 \textit{af dette}“由此”！）：\ding{55}\ \textit{*Som følge heraf den høje skat...}（致命重复！）。
    \end{itemize}
  \item \textbf{严禁把 \textit{hvorved} 当成主句副词而倒装！}：受 \textit{derfor} 习惯影响，同学经常写错：\ding{55}\ \textit{Skatten sættes ned, hvorved \dabold{får borgerne} flere penge.} 错！\textit{hvorved} 引导的是标准从句，主语必须留在动词前面：\ding{51}\ \textit{Skatten sættes ned, hvorved \dabold{borgerne får} flere penge.}（若是否定句：\textit{..., hvorved staten \dabold{ikke får} nok indtægter.}）。
  \item \textbf{\textit{Som følge heraf} 和 \textit{Af den grund} 放在句首绝对不要加逗号且必须倒装！}：英语写 \textit{As a result, many decline...} 会打逗号且主语在前；丹麦语必须直接动词跟上：\ding{51}\ \textit{\da{Af den grund} \dabold{siger mange} nej.}（无逗号，\textit{siger} 在第 2 位！）。
  \item \textbf{\textit{af den grund}（单数，特指上文那一个原因）vs \textit{af flere grunde}（复数）}：承接上文具体理由用单数 \textit{af den grund}；列举多个理由才用复数 \textit{af mange / flere grunde}。
\end{enumerate}
\end{grammatikbox}

% ---------------------------------------------------------------------
\ncdsubsec{五、Mønsterøvelser（句型与辩论操练）}

\textbf{练习 A：精准选填（用 \textit{som følge heraf, som følge af, af den grund, følgelig, hvorved} 补全因果链）}
\begin{enumerate}
  \item Topskatten omfordeler penge i samfundet, \underline{\hspace{2.2cm}} den økonomiske ulighed mindskes.
  \item \underline{\hspace{2.2cm}} den høje skat på overarbejde overvejer nogle læger at gå ned i tid.
  \item Mange borgere sætter pris på fri lægehjælp. \underline{\hspace{2.2cm}} accepterer et flertal stadig høje skatter.
  \item Hvis staten fjerner topskatten helt, falder skatteindtægterne. Man må \underline{\hspace{2.2cm}} spare på velfærden.
  \item Politikerne hævede grænsen sidste år, \underline{\hspace{2.2cm}} færre almindelige lønmodtagere nu betaler topskat.
  \item Han har mange udgifter til sin familie, og \underline{\hspace{2.2cm}} vil han gerne have mere i løn efter skat.
  \item Skattetrykket er højt i Danmark, men \underline{\hspace{2.2cm}} får vi til gengæld gratis uddannelse og tryghed.
\end{enumerate}

\textbf{练习 B：句型连接与纠错（合并因果两句，或修正语序与介词错误）}
\begin{enumerate}
  \item 用 \textit{Som følge heraf} 放句首合并（注意 V2 倒装）：\textit{Lønmodtagere får mindre lyst til overarbejde. Virksomhederne mangler specialister.} $\rightarrow$
  \item 用 \textit{, hvorved} 合并（注意从句语序）：\textit{Staten beskatter de højeste indkomster ekstra. Uligheden vokser ikke så hurtigt.} $\rightarrow$
  \item 介词与副词区分纠错：\textit{Som følge heraf det høje skattetryk vælger nogle at flytte til udlandet.} $\rightarrow$
  \item 从句倒装纠错：\textit{Grænsen for topskat blev flyttet op, hvorved kunne flere sygeplejersker beholde hele lønstigningen.} $\rightarrow$
  \item 用 \textit{Af den grund} 置于句首合并两句（注意删掉英语式逗号并倒装）：\textit{De ekstra penge går til hospitaler og skoler. Mange danskere betaler deres skat med glæde.} $\rightarrow$
\end{enumerate}

\textbf{练习 C：中译丹（活用严谨因果链条副词与连词）}
\begin{enumerate}
  \item 税收收入用于免费医院，从而使所有公民都能感到安全。
  \item 出于这一原因，许多高收入专家不愿接受额外的加班任务。
  \item 受此影响，如果大幅削减税收，公共财政就会面临赤字。
  \item 贫富差距随之缩小，由此社会的凝聚力变得更强。
  \item 如果税率过高，一些员工就会因此失去继续提升的动力。
\end{enumerate}

% ---------------------------------------------------------------------
\ncdsubsec{六、DU3 Modul 4 实战微任务与参考答案}

\begin{emailbox}{DU3 Modul 4 Skrivning Opgave 1 实战微任务（私人长邮件 · minimum 125 ord）}
\textbf{Situation}: Du har fået en mail fra din ven Thomas, som er læge og har fået tilbudt en lederstilling med flere vagter og højere løn, men også topskat.\\
\textbf{Opgave}: Skriv et svar til Thomas på \textbf{mindst 125 ord}. Svar på alle hans spørgsmål og brug mindst to kausale bindeord fra lektionen (\textit{som følge heraf, af den grund, følgelig, hvorved}).

\vspace{0.3em}
\textit{Hej!\\
Jeg håber, alt går godt hos dig. På mit arbejde har jeg lige fået tilbudt en stilling med mere ansvar og en højere løn. Problemet er bare, at jeg så skal arbejde mere over og betale en stor del i topskat. Nu er jeg i tvivl: Skal jeg tage imod tilbuddet for karrierens skyld, eller bør jeg hellere prioritere mere tid med min familie? Og hvad synes du egentlig om det høje danske skattesystem og topskatten generelt?\\
Kærlig hilsen, Thomas}
\end{emailbox}

\begin{answerbox}
\textbf{Før du læser 参考回答}
\begin{enumerate}
  \item \textit{Modstandere frygter, at topskatten svækker lysten til at arbejde ekstra, så virksomheder mister arbejdskraft og samfundet går glip af vækst. Tilhængere peger på, at topskatten mindsker uligheden, styrker sammenhængskraften og sikrer penge til fælles velfærd.}
  \item \textit{Thomas overvejer at sige nej til ekstravagter på sygehuset, fordi han mener, at for meget af den ekstra løn går til topskat. Emma forklarer, at man kun betaler topskat over grænsen, og at pengene går til hospitaler og daginstitutioner, som giver et trygt og stabilt samfund.}
\end{enumerate}

\textbf{Tekst B 阅读核对答案}
\begin{enumerate}
  \item \da{Forkert}（埃玛提醒他，只对超出起征点门槛的那部分收入缴纳最高边际税：\textit{kun af beløbet over grænsen}）
  \item \da{Han overvejer at sige nej til ekstravagter, fordi han synes, der bliver for lidt tilbage efter topskat.}
  \item \da{De går til sygehuse og daginstitutioner, hvilket skaber et trygt og stabilt samfund.}
\end{enumerate}

\textbf{练习 A 答案}
\begin{enumerate}
  \item \da{hvorved}（逗号后接从句，且后文动词 \textit{mindskes} 居于主语之后，呈从句语序）
  \item \da{Som følge af}（后面直接接名词短语 \textit{den høje skat}，必须用介词 \textit{af}）
  \item \da{Af den grund} / \da{Som følge heraf}（句首作副词，引发后面 \textit{accepterer et flertal} V2 倒装）
  \item \da{følgelig}（位于情态动词 \textit{må} 与实义动词 \textit{spare} 之间的句中位置）
  \item \da{hvorved}（逗号后接从句，“从而使更少的普通工薪阶层交最高税”）
  \item \da{af den grund} / \da{følgelig}（置于并列连词 \textit{og} 后面构成倒装前置项，引发 \textit{vil han} 倒装）
  \item \da{som følge heraf} / \da{af den grund}（承接上文“高税负”，引出正面的福祉结果）
\end{enumerate}

\textbf{练习 B 答案（句型改写参考答案）}
\begin{enumerate}
  \item \textit{Lønmodtagere får mindre lyst til overarbejde. \dabold{Som følge heraf mangler} virksomhederne specialister.}（句首因果副词触发 V2 倒装）
  \item \textit{Staten beskatter de højeste indkomster ekstra, \dabold{hvorved} uligheden \dabold{ikke vokser} så hurtigt.}（\textit{hvorved} 从句内部 \textit{ikke} 放在动词前，不倒装）
  \item \textit{\dabold{Som følge af} det høje skattetryk vælger nogle at flytte til udlandet.}（接名词短语必须用 \textit{Som følge af}，不能用 \textit{heraf}）
  \item \textit{Grænsen for topskat blev flyttet op, \dabold{hvorved} flere sygeplejersker \dabold{kunne beholde} hele lønstigningen.}（去掉从句里的错误倒装）
  \item \textit{De ekstra penge går til hospitaler og skoler. \dabold{Af den grund betaler} mange danskere deres skat med glæde.}（句首无逗号且必须 V2 倒装）
\end{enumerate}

\textbf{练习 C 答案}
\begin{enumerate}
  \item \textit{Skatteindtægterne bruges til gratis hospitaler, hvorved alle borgere kan føle sig trygge.}
  \item \textit{Af den grund vil mange højtlønnede specialister ikke sige ja til ekstra overarbejde.}
  \item \textit{Som følge heraf får de offentlige finanser et underskud, hvis man sænker skatten kraftigt.}
  \item \textit{Uligheden falder følgelig, hvorved sammenhængskraften i samfundet bliver endnu stærkere.}
  \item \textit{Hvis skattesatsen er for høj, mister nogle ansatte af den grund motivationen til udvikling.}
\end{enumerate}

\textbf{DU3 Modul 4 高分范文（152 词，不含署名）}\\
\textit{Hej Thomas}

\textit{Stort tillykke med det flotte tilbud! Jeg kan jo godt forstå, at du er lidt i tvivl, når mere ansvar samtidig betyder længere arbejdsdage og højere skat.}

\textit{Hvis jeg var dig, ville jeg først se på, hvor meget tid du vil have til din familie. Penge betyder jo ikke alt, og hvis arbejdet tager al din energi, mister du hurtigt glæden ved den nye stilling. \da{Af den grund} bør du måske forhandle om mere fleksible arbejdstider, før du siger ja --- så kan du få både spændende opgaver og ro derhjemme.}

\textit{Med hensyn til skattesystemet er jeg personligt ret glad for den danske model. Ganske vist føles topskatten høj på den sidste del af lønnen, men pengene finansierer jo fri uddannelse og et stærkt sundhedsvæsen, \dabold{hvorved} hele samfundet bliver mere trygt og lige. \da{Som følge heraf} synes jeg egentlig, at ordningen er rimelig, selvom grænsen måske kunne hæves lidt.}

\textit{Held og lykke med beslutningen!}

\textit{Kærlig hilsen\\
\mbox{[Dit navn]}}

\small\textbf{名师点评}：正文 152 词完整超越 125 词标准，紧扣 Thomas 的邮件三问（祝贺与共情、职业发展与家庭平衡建议、个人对高税制与 topskat 的看法）。文章不仅灵活融合了第 8 课的虚拟假设（\textit{Hvis jeg var dig, ville jeg...}）与 Unit 1 小品词（\textit{jo, egentlig}），更精准示范了本课重点 \textit{Af den grund} 句首 V2 倒装、\textit{, hvorved} 从句不倒装以及 \textit{Som følge heraf} 的因果呼应，兼顾口语亲切感与政策论证深度。
\end{answerbox}
\ncdlesson{10}{Integration og medborgerskab: Hvad kræver det?}{融入与公民身份：需要什么？}
{DU3 Modul 4（B1+ $\rightarrow$ B2-）｜对标 NCE 3 \& Modultest 4 \textit{Skrivning Opgave 2 (Debatindlæg) \& Samfundsdebat}}
{\textbf{社会议题高频动宾搭配}（\textit{stille krav om, tage stilling til, spille en afgørende rolle, bære præg af, tage ansvar for, bidrage til}）· 固定介词接口与\textbf{零冠词搭配规则} · 搭配后接从句语序}
{欢迎来到 Unit 2《福利国家的困境与辩论》的收官课！“如何实现真正的社会融入（\textit{integration}）与积极公民意识（\textit{medborgerskab}）”不仅是丹麦公共舆论场永恒的热门话题，也是 DU3 Modul 4 书面辩论与口试双人讨论的必考核心题域。在探讨这一议题时，告别泛泛的 \textit{lave regler} 或 \textit{være vigtig}，改用地道报刊动宾铁搭档：\textbf{\textit{stille krav om}}、\textbf{\textit{tage stilling til}}、\textbf{\textit{spille en afgørende rolle}} 与 \textbf{\textit{bære præg af}}，能让你的论证瞬间升格为 B2- 公民评论水准！}

% ---------------------------------------------------------------------
\begin{foerduleaeser}
带着下面两个问题阅读 Tekst A 和 Tekst B：
\begin{enumerate}
  \item \textit{Hvilke faktorer spiller ifølge Tekst A en afgørende rolle for en vellykket integration, og hvorfor må debatten ikke bære præg af fordomme?}（根据 Tekst A，哪些因素对成功融入起着决定性作用？为什么讨论不应带有偏见色彩？）
  \item \textit{Hvordan ser Ali og Sofie i Tekst B på forholdet mellem krav, rettigheder og det lokale foreningsliv?}（在 Tekst B 中，阿里和索菲如何看待要求、权利与地方社团生活之间的关系？）
\end{enumerate}
\end{foerduleaeser}

% ---------------------------------------------------------------------
\ncdsubsec{二、Tekster（福利国家专题与正反辩论双语篇）}

\subsection*{Tekst A：\textit{Integration som et fælles ansvar}（社会深度评论 · 约 220 词）}

\begin{tipbox}{语篇说明}
本篇剖析丹麦社会关于“新居民融入与积极公民身份”的政策与日常维度。阅读时请特别关注\textbf{\color{NCDNavy}深蓝色加粗的社会议题核心动宾搭配}（注意动词与固定介词之间的搭配），以及\textbf{\color{NCDRed}深红色加粗的零冠词结构与从句连接}。
\end{tipbox}

Når politikere og borgere i Danmark diskuterer fremtidens velfærdssamfund, dukker spørgsmålet om integration og aktivt medborgerskab næsten altid op. Alle partier er enige om, at et stærkt samfund bygger på tillid, men der er ofte uenighed om, hvad der kræves af både nyankomne og modtagersamfundet.

På den ene side mener et bredt flertal, at staten må \da{stille klare krav om} danskundervisning og selvforsørgelse fra dag ét. Arbejdsmarkedet og sproget \da{spiller} nemlig \da{en afgørende rolle} for, at man hurtigt lærer de uskrevne regler at kende og bliver uafhængig af offentlig hjælp. Når man har kolleger og faste opgaver, får man lettere et netværk og en følelse af tilhørsforhold.

På den anden side peger forskere på, at integration aldrig kan være en ensidig opgave. Hvis lokale virksomheder og nabolag \dabold{bærer præg af} lukkede netværk, bliver det svært for udlændinge at få foden indenfor, selvom de gerne vil \da{bidrage til fællesskabet}. Derfor bør kommuner og borgere sammen \da{tage stilling til}, hvordan man åbner dørene i både arbejdslivet og det lokale foreningsliv.

Medborgerskab handler således ikke kun om et pas eller en opholdstilladelse, men om at \da{tage ansvar for} sit eget liv og det samfund, man bor i. Hvis den offentlige debat ikke \dabold{bærer præg af} mistillid og hårde overskrifter, men i stedet af gensidig respekt, skaber vi de bedste rammer for alle.

\begin{transbox}
\small
\textbf{\color{NCDTeal}[Tekst A 中文精译]}\\
当丹麦的政治家和公民讨论未来的福利社会时，关于社会融入与积极公民身份的问题几乎总会浮现。所有政党都认同，一个稳固的社会建立在信任之上，但对于新来者和接纳社会双方各自应当承担什么要求，人们却常常存在分歧。

一方面，议会多数派认为国家必须从第一天起就对丹麦语学习与经济自立提出明确要求。因为劳动力市场和语言对于人迅速熟悉不成文的社会规则、摆脱对公共救济的依赖，起着决定性作用。当一个人有了同事和固定工作任务时，就更容易建立人脉并获得归属感。

另一方面，学者指出融入绝不可能是一项单方面的任务。如果当地企业和街区充斥着封闭圈子，那么即使外国人很愿意为社会共同体作出贡献，也很难获得踏入其中的机会。因此，各市镇与居民应当共同审视并抉择，如何在职场和当地社团生活中敞开大门。

由此可见，公民身份并不仅仅关乎一本护照或一张居留许可，更关乎为自己的生活以及自己所居住的社会承担责任。如果公共辩论不再带有不信任与尖锐标题的印记，而是体现出相互尊重，我们就能为所有人创造最好的制度环境。
\end{transbox}

\subsection*{Tekst B：\textit{Debat på lokalt borgermøde: Rettigheder og pligter}（市民大会正反对话 · 约 120 词）}

\begin{emailbox}{市民论坛对话｜Ali (Forældrerepræsentant) \& Sofie (Integrationskonsulent)}
\textbf{Ali:} Jeg har boet i Danmark i seks år, og for mig \da{spiller} den lokale fodboldklub \da{en afgørende rolle}. Her mødes vi på tværs af baggrund og sprog. Jeg mener, at vi alle skal \da{tage ansvar for} vores bydel, men desværre \dabold{bærer} mange jobopslag stadig \dabold{præg af}, at man kender nogen på forhånd.

\textbf{Sofie:} Det har du en vigtig pointe i, Ali. Som kommune skal vi selvfølgelig \da{stille tydelige krav om}, at alle lærer dansk og møder til sprogkurser, for det er nøglen til job. Men vi må samtidig \da{tage stilling til}, hvordan virksomhederne kan blive bedre til at invitere nye borgere indenfor. Først når både rettigheder og pligter følges ad, kan alle for alvor \da{bidrage til} vores fælles hverdag.
\end{emailbox}

\textbf{Tekst B 阅读核对（Rigtigt / Forkert / Kort svar，答案见第六部分）：}
\begin{enumerate}
  \item \textit{Ali mener, at bekendte aldrig betyder noget, når man søger job.} \underline{\hspace{2.2cm}}
  \item \textit{Hvilket konkret sted har spillet en afgørende rolle for Alis egen integration?} \underline{\hspace{2.2cm}}
  \item \textit{Hvad skal kommunen stille krav om, og hvad skal den tage stilling til ifølge Sofie?} \underline{\hspace{2.2cm}}
\end{enumerate}

\begin{tipbox}{本课发音与语调要点}
\begin{itemize}
  \item \textbf{公民身份与融入核心词的软 \textit{d} 与滑音}：\textit{medborgerskab} \ipa{ˈmɛðˌbɒʊ̯ˀɐˌsɡæːˀb}（\textit{med} 读软 \textit{d}，\textit{borger} 中的 \textit{g} 滑音化为 \ipa{ʊ̯ˀ}）；\textit{gensidig} \ipa{ˈɡɛnˌsiðˀi}（软 \textit{d} 加 \textit{stød}，词尾 \textit{g} 不发音）；\textit{selvforsørgelse} \ipa{ˈsɛlˀfʌˌsɶɐ̯ˀəlsə}（注意 \textit{v} 常弱读或省略）。
  \item \textbf{动宾搭配重音落点（重音全在宾语名词上！）}：丹麦语中，“动词 + 宾语”搭配读起来时，\textbf{动词轻读，主重音落在核心名词上}：
    \begin{itemize}
      \item \textit{stille krav om} \ipa{sd̥elə ˈkʰʁɑʊ̯ˀ ʌm}（重音在 \textit{krav}，带喉塞音）
      \item \textit{tage stilling til} \ipa{tˢæ ˈsd̥eleŋ tˢel}（\textit{tage} 极短缩读 \ipa{tˢæ}，重音在 \textit{stilling}）
      \item \textit{bære præg af} \ipa{bɛɐ ˈpʰʁɛːˀ ˈæːˀ}（\textit{præg} 中的 \textit{g} 几乎不发音，带长音与 \textit{stød}！）
      \item \textit{afgørende} \ipa{ˈɑʊ̯ˌɡøːˀɐnə}（前缀 \textit{af-} 读作 \ipa{ɑʊ̯}，主重音在最前面！）
    \end{itemize}
\end{itemize}
\end{tipbox}

\begin{transbox}
\small
\textbf{\color{NCDTeal}[Tekst B 中文对照]}\\
\textbf{阿里（家长代表）}：我在丹麦住了六年，对我来说当地的足球俱乐部起了决定性作用。在那里，拥有不同背景和语言的人们聚在一起。我认为我们每个人都应当为自己的街区承担责任，但遗憾的是，许多招聘启事仍然带有“事先认识熟人”的倾向。

\textbf{索菲（融入事务顾问）}：阿里，你这点说得很中肯。作为市政厅，我们当然必须对所有人学习丹麦语和参加语言课程提出明确要求，因为这是找到工作的钥匙。但与此同时，我们也必须认真思考并抉择，企业怎样才能更善于接纳新公民。只有当权利与义务并肩同行时，每个人才能真正为我们的共同日常生活作出贡献。
\end{transbox}

% ---------------------------------------------------------------------
\ncdsubsec{三、Gloser（核心生词表）}

\begin{tipbox}{词表重点}
本课动宾搭配里藏着两处极易出错的“零冠词”名词：\textit{en stilling}（立场/职位）在 \textit{tage stilling til} 里\textbf{绝对不加冠词}；\textit{et præg}（印记/特征）在 \textit{bære præg af} 里也\textbf{绝对不加冠词}！而 \textit{en rettighed}（权利）和 \textit{en pligt}（义务）通常以复数成对出现：\textit{rettigheder og pligter}。
\end{tipbox}

\begin{glosertable}
名词 & en integration & \textit{(integrationen, ingen flertal)} & 社会融入，一体化 \\
名词 & et medborgerskab & \textit{(medborgerskabet, ingen flertal)} & 公民身份，公民意识 \\
名词 & en rettighed & \textit{(rettigheden, rettigheder, -erne)} & 权利，权益 \\
名词 & en pligt & \textit{(pligten, pligter, pligterne)} & 义务，职责 \\
名词 & en selvforsørgelse & \textit{(selvforsørgelsen, ingen flertal)} & 经济自立，自力更生 \\
名词 & et tilhørsforhold & \textit{(tilhørsforholdet, -forhold, -ene)} & 归属感（单复同形） \\
名词 & et foreningsliv & \textit{(foreningslivet, ingen flertal)} & 社团生活，协会文化 \\
名词 & en fordom & \textit{(fordommen, fordomme, fordommene)} & 偏见，成见 \\
名词 & et krav & \textit{(kravet, krav, kravene)} & 要求（单复同形） \\
名词 & et præg & \textit{(præget, præg, prægene)} & 印记，色彩，特征 \\
\midrule
动词 [har] & at stille (krav) & \textit{(stiller, stillede, har stillet)} & 提出（要求、问题） \\
动词 [har] & at bære (præg af) & \textit{(bærer, bar, har båret)} & 带有，呈现出（……印记） \\
动词 [har] & at bidrage (til) & \textit{(bidrager, bidrog, har bidraget)} & 贡献，作出贡献（接 \textit{til}） \\
动词 [har] & at deltage (i) & \textit{(deltager, deltog, har deltaget)} & 参与，参加（接 \textit{i}） \\
动词 [har] & at kræve & \textit{(kræver, krævede, har krævet)} & 要求，需要 \\
\midrule
形容词 & afgørende & \textit{(不变，现在分词形容词)} & 决定性的，关键的 \\
形容词 & gensidig & \textit{(gensidigt, gensidige)} & 相互的，双向的 \\
\end{glosertable}

% ---------------------------------------------------------------------
\ncdsubsec{四、Grammatik \& Stilistik（核心语法、社会动宾搭配与修辞精讲）}

\begin{grammatikbox}{语法点 1：社会议题六大高频动宾搭配全景矩阵}
在探讨教育、就业、社会融入与福利改革时，能否用准“动词 + 名词 + 介词”黄金搭配，直接决定了你的 DU3 Modul 4 作文和口语是否具备地道报刊质感：

\vspace{0.3em}
\renewcommand{\arraystretch}{1.24}
{\small
\begin{tabularx}{\linewidth}{@{} >{\bfseries\raggedright\arraybackslash}p{3.4cm} >{\raggedright\arraybackslash}p{4.1cm} >{\raggedright\arraybackslash}X @{}}
\toprule
\textbf{核心动宾搭配} & \textbf{固定介词接口与中文语义} & \textbf{标准政论例句} \\
\midrule
\da{at stille krav om}\newline /\ \da{til} &
接 \dabold{om} 引出“要求的内容”；\newline
接 \dabold{til} 引出“被要求的對象” &
\textit{Staten \da{stiller klare krav om} dansk.}\newline
\textit{Man \da{stiller høje krav til} kommunen.} \\
\midrule
\da{at tage stilling til} &
表示“对……表明立场 / 审视抉择”\newline
（\textit{stilling} 前\textbf{零冠词}，必接 \dabold{til}） &
\textit{Borgerne må \da{tage stilling til} forslaget.}\newline
\textit{Vi skal \da{tage stilling til, om} det virker.} \\
\midrule
\da{at spille en afgørende rolle} &
表示“起决定性作用”（可替换为 \textit{central / vigtig}，后接 \dabold{for / i}） &
\textit{Sproget \da{spiller en afgørende rolle for} integrationen.} \\
\midrule
\da{at bære præg af} &
表示“带有……印记 / 体现……特点”\newline
（\textit{præg} 前\textbf{零冠词}，必接 \dabold{af}） &
\textit{Debatten \da{bærer} desværre \da{præg af} mistillid.} \\
\midrule
\da{at tage ansvar for} &
表示“对……承担责任”（零冠词 + \dabold{for}） &
\textit{Alle bør \da{tage ansvar for} deres lokalområde.} \\
\midrule
\da{at bidrage til} &
表示“为……作出贡献”（不规则过去式 \textit{bidrog}，必接 \dabold{til}） &
\textit{Frivillige \da{bidrager til} et stærkt fællesskab.} \\
\bottomrule
\end{tabularx}
}
\end{grammatikbox}

\begin{grammatikbox}{语法点 2：中英对比排雷——零冠词、介词接口与从句连接四大陷阱}
\begin{enumerate}
  \item \textbf{严禁往 \textit{tage stilling} 和 \textit{bære præg} 中间硬塞不定冠词 \textit{en / et}！}：受英语 \textit{take a stand} 或中文“采取一个立场”影响，最常见扣分错误就是多加冠词：\newline
    \ding{55}\ \textit{Vi må tage \dabold{en} stilling til spørgsmålet.} $\rightarrow$ \ding{51}\ \textit{Vi må \da{tage stilling til} spørgsmålet.}\newline
    \ding{55}\ \textit{Artiklen bærer \dabold{et} præg af frygt.} $\rightarrow$ \ding{51}\ \textit{Artiklen \da{bærer præg af} frygt.}
  \item \textbf{彻底替换“中/英式万能动词” \textit{lave / gøre / have}！}：\newline
    \ding{55}\ \textit{Politikerne laver krav om sprog.} $\rightarrow$ \ding{51}\ \textit{Politikerne \da{stiller krav om} sprog.}\newline
    \ding{55}\ \textit{Arbejdet spiller en vigtig rolle på integration.} $\rightarrow$ \ding{51}\ \textit{Arbejdet \da{spiller en vigtig rolle for} (eller \da{i}) integrationen.}（不能误用介词 \textit{på}！）
  \item \textbf{动宾搭配后面怎么接完整从句？}：当“要求的内容”或“思考的问题”是一个完整句子时，介词 \textit{om / til / af} 后面直接加逗号再接 \textit{at / om / hvordan} 从句（注意从句内不倒装）：\newline
    \ding{51}\ \textit{Kommunen stiller krav \dabold{om, at} man deltager i undervisningen.}\newline
    \ding{51}\ \textit{Vi bør tage stilling \dabold{til, hvordan} vi støtter nye kolleger.}
  \item \textbf{分辨 \textit{stille krav om} vs \textit{stille krav til}}：记住口诀——“\textbf{对人/机构提要求用 \textit{til}，要求具体事项用 \textit{om}}”。例：\textit{Samfundet stiller krav \dabold{til} alle borgere \dabold{om} aktiv deltagelse.}（同时串联两个介词接口！）。
\end{enumerate}
\end{grammatikbox}

% ---------------------------------------------------------------------
\ncdsubsec{五、Mønsterøvelser（句型与修辞操练）}

\textbf{练习 A：社会议题动宾搭配与介词精准填空（每空填入唯一准确的 1 个丹麦语词）}
\begin{enumerate}
  \item Kommunen bør stille klare \underline{\hspace{2.2cm}} om, at alle nye borgere får hurtig adgang til danskundervisning.
  \item Politikerne må nu tage \underline{\hspace{2.2cm}} til, hvordan vi kan skabe flere praktikpladser i lokalområdet.
  \item Et fast arbejde spiller en \underline{\hspace{2.2cm}} rolle for både selvforsørgelse og netværk. (决定性的)
  \item Desværre \underline{\hspace{2.2cm}} den offentlige debat ofte præg af fordomme i stedet for fakta. (bære，现在时)
  \item Det lokale foreningsliv er et godt sted, hvor alle kan bidrage \underline{\hspace{2.2cm}} fællesskabet.
  \item Samfundet stiller høje krav \underline{\hspace{2.2cm}} nye borgere, men må også tilbyde åbne døre på arbejdspladsen. (对……人/群体)
  \item Enhver voksen bør tage ansvar \underline{\hspace{2.2cm}} sit eget liv og respektere fælles spilleregler.
\end{enumerate}

\textbf{练习 B：语体升级与搭配修正（修正英语直译与冠词/介词错误，或按要求升级句子）}
\begin{enumerate}
  \item 修正万能动词与介词错误：\textit{Skolen laver høje krav på eleverne om godt dansk.} $\rightarrow$
  \item 修正多余冠词与介词错误：\textit{Byrådet skal tage en stilling om det nye integrationsprojekt i næste uge.} $\rightarrow$
  \item 修正搭配名词冠词与动词错误：\textit{Diskussionen har et præg af mistillid mellem grupperne.} $\rightarrow$
  \item 修正介词搭配错误：\textit{Sportsklubber spiller en central rolle på børnenes trivsel i kvarteret.} $\rightarrow$
  \item 动宾搭配接从句升级（用 \textit{stille krav om, at...} 合并）：\textit{Staten siger, at man skal lære dansk hurtigt.} $\rightarrow$
  \item 动宾搭配接从句升级（用 \textit{tage stilling til, hvordan...} 改写）：\textit{Virksomhederne må tænke over, hvordan de ansætter flere udlændinge.} $\rightarrow$
\end{enumerate}

\textbf{练习 C：中译丹实战（灵活运用社会议题高频动宾搭配）}
\begin{enumerate}
  \item 语言和工作对成功融入起着决定性作用。
  \item 政府应当对地方机构提出更加明确的要求。
  \item 社会共同体必须决定如何帮助新邻居建立社交网络。
  \item 遗憾的是，关于融入的辩论往往带有强烈的偏见色彩。
  \item 每位公民都可以通过志愿工作为社区作出贡献。
  \item 我们既要享有权利，也应当为自己的生活承担责任。
\end{enumerate}

% ---------------------------------------------------------------------
\ncdsubsec{六、DU3 Modul 4 实战微任务与参考答案}

\begin{emailbox}{DU3 Modul 4 Skrivning Opgave 2 实战微任务（Debatindlæg · minimum 100 ord）}
\textbf{Situation}: I en lokalavis har debattøren Martin skrevet et kort debatindlæg:\\
\textit{”Integration er udelukkende de nye borgeres eget ansvar! Staten skal bare stille strenge krav om dansk og arbejde, og ellers må folk selv finde ud af at passe ind i samfundet.”}\\[3pt]
\textbf{Opgave}: Skriv et sagligt \textbf{debatindlæg} på \textbf{mindst 100 ord} (gerne 130--155 ord), hvor du diskuterer Martins synspunkt. Skriv, om du er enig eller delvist uenig, og brug mindst \textbf{tre} faste udtryk fra denne lektion (\textit{stille krav om, tage stilling til, spille en afgørende rolle, bære præg af, bidrage til}).
\end{emailbox}

\begin{answerbox}
\textbf{Før du læser 参考回答}
\begin{enumerate}
  \item \textit{Ifølge Tekst A spiller arbejdsmarkedet og sproget en afgørende rolle for integrationen, så man lærer de uskrevne regler og får et tilhørsforhold. Debatten må ikke bære præg af fordomme eller mistillid, fordi det gør samarbejdet sværere og overser, at integration kræver gensidig respekt.}
  \item \textit{Ali peger på, at man bør tage ansvar for sin bydel og mødes i foreningslivet, men kritiserer, at jobopslag ofte bærer præg af lukkede netværk. Sofie mener, at kommunen skal stille tydelige krav om danskundervisning, samtidig med at man tager stilling til, hvordan virksomheder kan åbne sig mere, så rettigheder og pligter følges ad.}
\end{enumerate}

\textbf{Tekst B 阅读核对答案}
\begin{enumerate}
  \item \da{Forkert}（阿里明确批评许多招聘启事仍带有依赖熟人关系的倾向：\textit{bærer præg af, at man kender nogen på forhånd}）
  \item \da{Den lokale fodboldklub.}（当地足球俱乐部让他与不同背景的人交流，起到了决定性作用）
  \item \da{Krav om dansk og sprogkurser; tage stilling til, hvordan virksomheder inviterer nye borgere indenfor.}（对学丹麦语提出明确要求，并思考如何让企业更接纳新公民）
\end{enumerate}

\textbf{练习 A 答案}
\begin{enumerate}
  \item \da{krav}（固定搭配：\textit{stille klare krav om, at...} 对……提出明确要求）
  \item \da{stilling}（零冠词固定搭配：\textit{tage stilling til} 对……表态/思考抉择）
  \item \da{afgørende}（固定搭配：\textit{spille en afgørende rolle for...} 起决定性作用）
  \item \da{bærer}（固定搭配：\textit{bære præg af} 带有……色彩/印记）
  \item \da{til}（固定介词：\textit{bidrage til fællesskabet} 为共同体作贡献）
  \item \da{til}（对人或群体提要求须用介词 \textit{til}：\textit{stille krav til nye borgere}）
  \item \da{for}（固定介词：\textit{tage ansvar for} 对……承担责任）
\end{enumerate}

\textbf{练习 B 答案（参考）}
\begin{enumerate}
  \item \textit{Skolen \da{stiller} høje \da{krav til} eleverne om godt dansk.}（用动词 \textit{stille} 替代误用的 \textit{lave}，对象用介词 \textit{til}）
  \item \textit{Byrådet skal \da{tage stilling til} det nye integrationsprojekt i næste uge.}（去掉多余冠词 \textit{en}，介词改用 \textit{til}）
  \item \textit{Diskussionen \da{bærer præg af} mistillid mellem grupperne.}（使用固定搭配 \textit{bære præg af}，不加冠词 \textit{et}）
  \item \textit{Sportsklubber \da{spiller en central rolle for} børnenes trivsel i kvarteret.}（介词须用 \textit{for} 或 \textit{i}，不能误用 \textit{på}）
  \item \textit{Staten \da{stiller krav om, at} man lærer dansk hurtigt.}（利用介词 \textit{om} 连接 \textit{at} 从句）
  \item \textit{Virksomhederne må \da{tage stilling til, hvordan} de ansætter flere udlændinge.}（利用介词 \textit{til} 连接疑问词从句）
\end{enumerate}

\textbf{练习 C 答案}
\begin{enumerate}
  \item \textit{Sprog og arbejde spiller en afgørende rolle for en vellykket integration.}
  \item \textit{Regeringen bør stille tydeligere krav til de lokale institutioner.}
  \item \textit{Fællesskabet skal tage stilling til, hvordan vi hjælper nye naboer med et netværk.}
  \item \textit{Debatten om integration bærer desværre ofte præg af stærke fordomme.}
  \item \textit{Enhver borger kan bidrage til lokalsamfundet gennem frivilligt arbejde.}
  \item \textit{Vi skal både have rettigheder og tage ansvar for vores eget liv.}
\end{enumerate}

\textbf{DU3 Modul 4 高分范文（Debatindlæg · 133 词）}\\
\textit{\textbf{Vellykket integration kræver en fælles indsats}}

\textit{Jeg er kun delvist enig med Martin i hans indlæg om integration. Ganske vist er det rimeligt, at staten \da{stiller klare krav om}, at nye borgere lærer dansk og hurtigt kommer i arbejde. Alle voksne bør jo \da{tage ansvar for} deres egen selvforsørgelse og ønske at \da{bidrage til} samfundet.}

\textit{Ikke desto mindre mener jeg, at integration aldrig kan være én sides ansvar alene. Sproget og jobbet \da{spiller en afgørende rolle}, men nye borgere kan ikke øve dansk alene derhjemme. Hvis arbejdspladser og lokale foreninger \dabold{bærer præg af} lukkede kredse, bliver det svært for udlændinge at få reel kontakt til danskerne.}

\textit{Derfor må både kommuner og almindelige borgere \da{tage stilling til}, hvordan vi åbner vores netværk i hverdagen. Når faste krav kombineres med gensidig åbenhed, vinder hele fællesskabet.}

\small\textbf{名师点评}：正文含标题共 133 词，完整覆盖 Modul 4 Opgave 2 辩论评分要点。本文堪称 Unit 1--2 的综合收官秀：既复用第 4 课的高阶让步转折（\textit{Jeg er kun delvist enig... / Ganske vist... / Ikke desto mindre...}）和第 1 课小品词 \textit{jo}，又无缝串起本课五大黄金动宾搭配（\textit{stiller klare krav om, tage ansvar for, bidrage til, spiller en afgørende rolle, bærer præg af, tage stilling til}），零冠词与介词接口滴水不漏！
\end{answerbox}

% =====================================================================
% Unit 3
% =====================================================================
\ncdunit{3}{Teknologi, Medier og Etik}{科技、媒体与现代伦理}{Lektion 11--15：动词名词化（-ning/-else）· 限制与前提条件 · 价值权衡句型 · 书面长分词定语 · 批判性思维与信源核查}

\ncdlesson{11}{Kunstig intelligens i klasseværelset: Hjælp eller snyd?}{课堂里的 AI：助手还是作弊？}
{DU3 Modul 4（B1+ $\rightarrow$ B2-）｜对标 NCE 3 \& Modultest 4 \textit{Skrivning Opgave 2（Debatindlæg）\& Skriftsproglig præcision}}
{\textbf{动词名词化（Verbalsubstantiver med \textit{-ning / -else}）}：将冗长的 \textit{når man... / at man...} 从句浓缩为正式名词短语（\textit{anvendelse, indførelse, vurdering, besvarelse}）· 介词 + 动词名词化 + \textit{af} 结构与 V2 倒装}
{欢迎进入《新概念丹麦语 4》第三单元“科技、媒体与现代伦理”！从本课起，我们直面 DU3 Modul 4 最常考的现代社会辩题。在丹麦报刊评论和高分 \textit{Debatindlæg} 中，如果通篇重复口语式的 \textit{når man bruger AI...} 或 \textit{hvis man forbyder...}，文章会显得松散稚嫩。学会把动词转化为 \textbf{\textit{-ning / -else} 抽象名词}并搭配介词短语，是把书面表达从 B1 拉升至 B2- 的核心利器。}

% ---------------------------------------------------------------------
\begin{foerduleaeser}
带着下面两个问题阅读 Tekst A 和 Tekst B：
\begin{enumerate}
  \item \textit{Hvordan kan kunstig intelligens både være en hjælp og en udfordring i skolen ifølge Tekst A?}（根据 Tekst A，人工智能在学校里如何既可以是助力，又可能成为挑战？）
  \item \textit{Hvorfor mener både læreren Freja og eleven Lukas i Tekst B, at et totalt forbud mod AI ikke er den bedste løsning?}（在 Tekst B 中，为什么教师芙蕾雅和学生卢卡斯都认为全面禁止 AI 并非最佳方案？）
\end{enumerate}
\end{foerduleaeser}

% ---------------------------------------------------------------------
\ncdsubsec{二、Tekster（科技媒体专题与思辨评论双语篇）}

\subsection*{Tekst A：\textit{Kunstig intelligens i klasseværelset: Hjælp eller snyd?}（科技教育评论 · 约 250 词）}

\begin{tipbox}{语篇说明}
本篇以丹麦基础教育与高中课堂引入生成式 AI 为背景，示范“如何将冗长动作从句浓缩为报刊书面名词短语”。阅读时请重点留意\textbf{\color{NCDNavy}深蓝色加粗的 \textit{-ning / -else} 动词名词化}以及\textbf{\color{NCDRed}深红色加粗的“介词 + 名词化 + \textit{af}”紧凑结构}。
\end{tipbox}

I de danske folkeskoler og gymnasier fylder kunstig intelligens mere og mere i dagligdagen. Siden sprogmodeller som ChatGPT blev udbredte, har mange undervisere diskuteret, hvordan man bør forholde sig til den hurtige \da{udvikling} på teknologiområdet. Mens nogle ser et stort potentiale for bedre læring, frygter andre en svækkelse af elevernes egne skriftlige evner.

På den ene side viser nye erfaringer, at \dabold{ved anvendelse af} digitale værktøjer kan elever få hurtig og personlig støtte i lektielæsningen. Når en elev har svært ved et kompliceret emne, kan AI forklare svære begreber på flere niveauer og fungere som en tålmodig sparringspartner. Samtidig oplever flere lærere en lettelse i \da{forberedelsen} af undervisningen, fordi systemet kan foreslå øvelser og tekster på få sekunder.

På den anden side peger kritikere på, at let adgang til færdige tekster kan føre til \da{afhængighed} og snyd. Hvis eleverne lader en maskine skrive hele opgaven, forsvinder den vigtige \da{træning} i selvstændig tænkning. \dabold{Efter indførelsen af} kunstig intelligens er det derfor blevet sværere for skolerne at sikre en retfærdig \da{vurdering} af elevernes reelle niveau.

Af den grund foreslår flere uddannelsesforskere nye eksamensformer i stedet for et totalt forbud. \dabold{Under udarbejdelsen af} skriftlige opgaver kan man for eksempel kombinere digitale hjælpemidler med mundtlig \da{fremlæggelse}, hvor eleven selv skal forklare sine pointer. På den måde kan \dabold{en klog anvendelse af} teknologien styrke læringen uden at erstatte elevens egen \da{forståelse}.

\begin{transbox}
\small
\textbf{\color{NCDTeal}[Tekst A 中文精译]}\\
在丹麦的公立中小学和高中里，人工智能在日常教学中所占的分量越来越重。自从像 ChatGPT 这样的语言模型普及以来，许多教育工作者一直在讨论，应当如何应对科技领域的这一快速发展。一些人看到了提升学习效果的巨大潜力，而另一些人则担忧学生自身的书面写作能力会因此削弱。

一方面，最新实践表明，通过运用数字工具，学生在做功课时能获得迅速且个性化的支持。当学生对某个复杂课题感到吃力时，AI 可以从多个层面解释晦涩概念，充当一名耐心的切磋伙伴。与此同时，不少教师在备课方面也感到压力有所减轻，因为系统只需几秒钟就能提供练习与阅读材料建议。

另一方面，批评者指出，唾手可得的现成文本可能导致依赖心理与作弊行为。如果学生放任机器包办整篇作业，独立思考这一至关重要的训练便会消失。因此，在人工智能引入校园之后，学校要确保对学生真实水平作出公正评估变得比以往更加困难。

基于这一原因，多位教育研究人员建议采用新的考试形式，而非全面禁止。\textbf{在撰写}书面作业期间，学校例如可以把数字辅助工具与口头陈述结合起来，要求学生亲口阐述自己的核心论点。如此一来，明智地运用技术便能够在不取代学生自身理解的前提下，强化学习成效。
\end{transbox}

\subsection*{Tekst B：\textit{Debat på lærerværelset: Skal AI forbydes til eksamen?}（教育论坛对话 · 约 140 词）}

\begin{emailbox}{校园辩论节选｜Freja（高中丹麦语教师）\& Lukas（高中三年级学生）}
\textbf{Freja:} Vi bliver nødt til at drøfte \da{anvendelsen} af AI før sommerens eksamen. Uden klare regler risikerer vi, at \da{besvarelsen} af stile mister sin værdi, fordi vi ikke ved, hvem der har skrevet teksten.\\
\textbf{Lukas:} Det forstår jeg godt, men et totalt forbud er urealistisk. I vores kommende arbejdsliv bliver \da{betjeningen} af nye værktøjer jo en naturlig del af hverdagen.\\
\textbf{Freja:} Du har ret i, at forbud sjældent virker på lang sigt. Måske skal vi i stedet lægge større vægt på den mundtlige \da{forklaring} efter afleveringen, så man skal forsvare sine egne argumenter.\\
\textbf{Lukas:} Ja, præcis! \dabold{Ved kombinationen af} skriftlige opgaver og en kort samtale kan skolen både teste vores viden og lære os en ansvarlig brug af teknologien.
\end{emailbox}

\textbf{Tekst B 阅读核对（Rigtigt / Forkert / Kort svar，答案见第六部分）：}
\begin{enumerate}
  \item \textit{Freja foreslår et totalt forbud mod alle digitale værktøjer til sommerens eksamen.} \underline{\hspace{2.2cm}}
  \item \textit{Hvorfor mener Lukas, at et totalt forbud mod AI er urealistisk?} \underline{\hspace{2.2cm}}
  \item \textit{Hvilken løsning bliver Freja og Lukas enige om til sidst?} \underline{\hspace{2.2cm}}
\end{enumerate}

\begin{tipbox}{本课发音与语调要点（Udtale \& Reduktioner）}
\begin{itemize}
  \item \textbf{\textit{-ning} 与 \textit{-else} 后缀的固定韵律}：丹麦语中，动词名词化后缀永远\textbf{不重读}，主重音留在词干，\textit{-ning} 读作轻短鼻音 \ipa{-neŋ}，\textit{-else} 读作弱元音 \ipa{-əlsə}：
    \begin{itemize}
      \item \textit{en anvendelse} \ipa{ˈanˌvɛnˀəlsə}（\textit{-nd-} 中的 \textit{d} 不发音，带 \textit{stød}！）
      \item \textit{en indførelse} \ipa{ˈenˌføːˀɐlsə}（重音在前缀 \textit{ind-}，\textit{d} 不发音）
      \item \textit{en vurdering} \ipa{vuˈdeːˀɐeŋ}（重音在第二音节 \textit{-de-}，带长音和 \textit{stød}）
      \item \textit{en besvarelse} \ipa{beˈsvɑːˀɐlsə}（前缀 \textit{be-} 轻读）
    \end{itemize}
  \item \textbf{科技伦理高频词与静音字母}：\textit{kunstig intelligens} \ipa{ˈkʰonsdi entˢeliˈɡɛnˀs}（\textit{kunstig} 词尾 \textit{g} 不发音）；\textit{et hjælpemiddel} \ipa{ˈjɛlbəˌmiðˀəl}（词首 \textit{hj-} 的 \textit{h} 不发音，\textit{-dd-} 为软 d）；\textit{et værktøj} \ipa{ˈvɛɐ̯ɡˌtˢʌɪ̯}。
\end{itemize}
\end{tipbox}

\begin{transbox}
\small
\textbf{\color{NCDTeal}[Tekst B 中文对照]}\\
\textbf{芙蕾雅}：我们必须在夏季考试前讨论一下人工智能的使用问题。如果没有明确规则，论文答卷就有可能失去价值，因为我们不知道究竟是谁写了文章。\\
\textbf{卢卡斯}：我完全理解，但全面禁止并不现实。在我们未来的职业生涯中，操作新工具本来就会成为日常工作的自然组成部分。\\
\textbf{芙蕾雅}：你说得对，禁令从长远来看很少奏效。或许我们反而应该更重视提交作业后的口头阐释，让每个人必须亲自为自己的论点辩护。\\
\textbf{卢卡斯}：没错！通过把书面作业和简短面谈结合起来，学校既能考查我们的知识，又能教会我们负责任地使用技术。
\end{transbox}

% ---------------------------------------------------------------------
\ncdsubsec{三、Gloser（核心生词表）}

\begin{tipbox}{词表重点}
请成组记忆“动词 $\leftrightarrow$ 名词化”的配对关系（如 \textit{at anvende} $\leftrightarrow$ \textit{en anvendelse}、\textit{at vurdere} $\leftrightarrow$ \textit{en vurdering}）。请牢记一条黄金词性铁律：几乎所有以 \textbf{\textit{-ning}} 和 \textbf{\textit{-else}} 结尾的动词派生动词名词，全都是共性 \textbf{\textit{en}} 名词！
\end{tipbox}

\begin{glosertable}
名词 & en anvendelse & \textit{(anvendelsen, anvendelser, anvendelserne)} & 应用，使用 \\
名词 & en udvikling & \textit{(udviklingen, udviklinger, udviklingerne)} & 发展，演变 \\
名词 & en vurdering & \textit{(vurderingen, vurderinger, vurderingerne)} & 评估，评价 \\
名词 & en indførelse & \textit{(indførelsen, ingen flertal)} & 引入，推行 \\
名词 & en forberedelse & \textit{(forberedelsen, forberedelser, -erne)} & 准备，备课 \\
名词 & en besvarelse & \textit{(besvarelsen, besvarelser, besvarelserne)} & 作答，答卷 \\
名词 & en udarbejdelse & \textit{(udarbejdelsen, ingen flertal)} & 拟定，撰写，编制 \\
名词 & en forståelse & \textit{(forståelsen, forståelser, forståelserne)} & 理解，领悟 \\
名词 & en fremlæggelse & \textit{(fremlæggelsen, fremlæggelser, -erne)} & 展示，口头陈述 \\
名词 & et hjælpemiddel & \textit{(hjælpemidlet, hjælpemidler, -midlerne)} & 辅助工具 \\
名词 & et værktøj & \textit{(værktøjet, værktøjer, værktøjerne)} & 工具，手段 \\
名词 & et forbud & \textit{(forbuddet, forbud, forbuddene)} & 禁令，禁止 \\
\midrule
动词 [har] & at anvende & \textit{(anvender, anvendte, har anvendt)} & 使用，运用（书面） \\
动词 [har] & at vurdere & \textit{(vurderer, vurderede, har vurderet)} & 评估，评判 \\
动词 [har] & at indføre & \textit{(indfører, indførte, har indført)} & 引进，实行 \\
动词 [har] & at erstatte & \textit{(erstatter, erstattede, har erstattet)} & 取代，替代 \\
动词 [har] & at snyde & \textit{(snyder, snød, har snydt)} & 作弊，欺骗 \\
\midrule
形容词 & kunstig & \textit{(kunstigt, kunstige)} & 人工的，人造的 \\
\end{glosertable}

% ---------------------------------------------------------------------
\ncdsubsec{四、Grammatik \& Stilistik（核心语法、伦理表达与修辞精讲）}

\begin{grammatikbox}{语法点 1：动词名词化（Verbalsubstantiver med \textit{-ning / -else}）的构词规律}
将动词转为抽象名词，是丹麦语报刊专栏与政策讨论中最标志性的书面语体特征。最核心的两大后缀是 \da{-ning}（或词干以 \textit{r} 结尾直接加 \textit{-ing}）与 \da{-else}：

\vspace{0.3em}
\renewcommand{\arraystretch}{1.24}
{\small
\begin{tabularx}{\linewidth}{@{} >{\bfseries\raggedright\arraybackslash}p{2.3cm} >{\raggedright\arraybackslash}p{4.2cm} >{\raggedright\arraybackslash}X @{}}
\toprule
\textbf{派生后缀} & \textbf{常见动词来源特征} & \textbf{高频动词 $\rightarrow$ 名词化转换} \\
\midrule
\da{-ning / -ing}\newline （一律 \textit{en} 类） &
多数规则动词及以 \textit{-ere}、\textit{-r} 结尾的动词（去词尾 \textit{-e} 加后缀） &
\textit{at vurdere} $\rightarrow$ \textit{en vurdering}（评估）\newline
\textit{at udvikle} $\rightarrow$ \textit{en udvikling}（发展）\newline
\textit{at løse} $\rightarrow$ \textit{en løsning}（解决）\newline
\textit{at træne} $\rightarrow$ \textit{en træning}（训练） \\
\midrule
\da{-else}\newline （一律 \textit{en} 类） &
带前缀 \textit{an-, be-, for-, ind-, ud-, op-} 的大量双音节/多音节动词 &
\textit{at anvende} $\rightarrow$ \textit{en anvendelse}（应用）\newline
\textit{at indføre} $\rightarrow$ \textit{en indførelse}（引入）\newline
\textit{at forberede} $\rightarrow$ \textit{en forberedelse}（准备）\newline
\textit{at besvare} $\rightarrow$ \textit{en besvarelse}（作答）\newline
\textit{at udarbejde} $\rightarrow$ \textit{en udarbejdelse}（撰写） \\
\bottomrule
\end{tabularx}
}
\end{grammatikbox}

\begin{grammatikbox}{语法点 2：把“冗长 \textit{når man... / at man...} 从句”浓缩为正式名词短语}
在 B1 阶段，我们习惯用连词引导完整从句；迈入 B1+ $\rightarrow$ B2- 后，可以用\textbf{“介词 + 动词名词化 + \textit{af} + 宾语”}三步浓缩法，让表达更紧凑有力：
\begin{enumerate}
  \item \textbf{替代时间/方式从句（\textit{når man... / ved at man...}）}：\\
    公式：\textbf{介词（\textit{ved / under / efter / før / gennem}）+ 名词化 + \textit{af} + 原动词宾语}。\\
    \ding{172} 口语展开：\textit{Når man anvender AI, kan elever få hurtig hjælp.}\\
    \ding{173} 书面浓缩：\textit{\dabold{Ved anvendelse af} AI kan elever få hurtig hjælp.}（注意：介词短语放句首，后面主句仍须 \textbf{V2 倒装} \textit{kan elever}！）\\
    \ding{174} 表示“在……之后/期间”：\textit{Efter at skolen har indført nye regler...} $\rightarrow$ \textit{\dabold{Efter indførelsen af} nye regler...}
  \item \textbf{替代作主语或宾语的 \textit{at} 从句}：\\
    \ding{172} 从句作主语：\textit{At man forbereder undervisningen grundigt, tager tid.}\\
    \ding{173} 名词短语作主语：\textit{\dabold{En grundig forberedelse af} undervisningen tager tid.}
  \item \textbf{如果原动词没有宾语怎么办？}直接省略 \textit{af}，或用所有格 \textit{-s} 引出动作发出者：\\
    \textit{Når teknologien udvikler sig hurtigt...} $\rightarrow$ \textit{Med \dabold{teknologiens} hurtige \dabold{udvikling}...} / \textit{Ved \dabold{elevernes brug af} værktøjet...}
\end{enumerate}
\end{grammatikbox}

\begin{grammatikbox}{修辞避坑指南 \& AI 伦理核心搭配（Stilistik \& Faste vendinger）}
\begin{enumerate}
  \item \textbf{千万不要“词性硬凑”，并记清必须接 \textit{af} 的原则}：把 \textit{anvende AI} 名词化后，\textit{anvendelse} 不能直接跟名词，中间必须用介词 \dabold{af} 搭桥：\ding{51}\ \textit{anvendelsen \dabold{af} AI}（绝不能写成 \ding{55}\ \textit{anvendelsen AI}！）。
  \item \textbf{别忘了 \textit{en} 类定指后缀 \textit{-n / -en}}：以 \textit{-else} 结尾加 \textit{-n}（\textit{anvendelse $\rightarrow$ anvendelsen}）；以 \textit{-ning} 结尾加 \textit{-en}（\textit{vurdering $\rightarrow$ vurderingen}）。泛指整体概念时常用定指形式（如 \textit{efter indførelsen af loven}），泛指方式时常用不定指（如 \textit{ved anvendelse af}）。
  \item \textbf{警惕“英语动名词 \textit{-ing}”直译病}：英语写 \textit{By using AI...}；丹麦语绝对不能用 \textit{*ved brugende AI}，而应使用动词名词化 \ding{51}\ \textit{ved anvendelse af AI}（或不定式 \textit{ved at bruge AI}）。
\end{enumerate}

\textbf{科技与教育辩论高频搭配（Faste vendinger）：}
\begin{itemize}
  \item \da{at stille nye krav til (skolerne)}（对学校提出新要求）
  \item \da{at forholde sig til (den teknologiske udvikling)}（面对、审视科技发展）
  \item \da{at sikre en retfærdig vurdering}（确保公正评估）
  \item \da{et tveægget sværd}（一把双刃剑） / \da{en ansvarlig brug af teknologi}（负责任地使用技术）
\end{itemize}
\end{grammatikbox}

% ---------------------------------------------------------------------
\ncdsubsec{五、Mønsterøvelser（句型与修辞操练）}

\textbf{练习 A：词形转化与介词填空（将括号内动词变为正确的 \textit{-ning / -else} 名词化形式，或填入准确介词）}
\begin{enumerate}
  \item Ved \underline{\hspace{2.2cm}} (anvende) af kunstig intelligens kan lærere spare tid på rutineopgaver.
  \item Efter \underline{\hspace{2.2cm}} (indføre, 定指) af nye eksamensregler skal alle kilder oplyses tydeligt.
  \item En grundig \underline{\hspace{2.2cm}} (forberede) af timen giver eleverne mere udbytte.
  \item Mange forskere frygter, at let adgang til AI kan svække elevernes \underline{\hspace{2.2cm}} (forstå) af svære tekster.
  \item En retfærdig \underline{\hspace{2.2cm}} (vurdere) af skriftlige stile kræver klare retningslinjer.
  \item Under \underline{\hspace{2.2cm}} (udarbejde, 定指) af opgaven må kursisterne kun bruge godkendte ordbøger.
  \item Den hurtige \underline{\hspace{2.2cm}} (udvikle) inden for sprogmodeller stiller store krav til uddannelsessystemet.
  \item Skolen ønsker at kombinere digitale værktøjer med en mundtlig \underline{\hspace{2.2cm}} (besvare).
\end{enumerate}

\textbf{练习 B：句型浓缩手术（将下列口语化的 \textit{når / hvis / efter at / at} 从句改写为正式名词短语，注意 V2 倒装）}
\begin{enumerate}
  \item \textit{Når man anvender digitale hjælpemidler, bliver lektierne mere overskuelige.} $\rightarrow$ \textit{Ved...}
  \item \textit{Efter at skolerne har indført nye regler, er der mindre snyd til prøverne.} $\rightarrow$ \textit{Efter...}
  \item \textit{Når læreren vurderer opgaven, lægger hun vægt på elevens egne ideer.} $\rightarrow$ \textit{Ved lærerens...}
  \item \textit{Mens eleverne udarbejder projektet, må de gerne søge inspiration på nettet.} $\rightarrow$ \textit{Under...}
  \item \textit{At man forbyder AI fuldstændigt, løser ikke problemet på lang sigt.} $\rightarrow$ \textit{Et totalt forbud mod...}
  \item \textit{Hvis man kombinerer skriftligt arbejde med en mundtlig samtale, kan man teste elevens viden bedre.} $\rightarrow$ \textit{Ved en kombination af...}
\end{enumerate}

\textbf{练习 C：中译丹（灵活运用动词名词化与科技教育词汇）}
\begin{enumerate}
  \item 通过使用人工智能力量，学生能获得快速的帮助。（用 \textit{Ved anvendelse af...}）
  \item 在新规则引入之后，考试变得更加公平了。（用 \textit{Efter indførelsen af...}）
  \item 公正地评估书面作业对老师来说非常重要。（用 \textit{En retfærdig vurdering af...}）
  \item 人工智能的快速发展对丹麦学校提出了新的要求。
  \item 全面禁止新工具无法取代对批判性思考的训练。（用 \textit{Et forbud...}）
  \item 在准备演讲期间，学生应当理解所有重要概念。（用 \textit{Under forberedelsen af...}）
\end{enumerate}

% ---------------------------------------------------------------------
\ncdsubsec{六、DU3 Modul 4 实战微任务与参考答案}

\begin{emailbox}{DU3 Modul 4 微型写作任务（对标 Skrivning Opgave 2 · Debatindlæg）}
\textbf{Situation}: I en onlineavis skriver bloggeren Tobias:\\
\textit{”Kunstig intelligens ødelægger unges evne til at tænke selv. Alle skoler og gymnasier bør derfor indføre et totalt forbud mod AI i undervisningen og til eksamen.”}\\[4pt]
\textbf{Opgave}: Skriv et debatindlæg på \textbf{mindst 100 ord} (anbefalet 125--150 ord), hvor du diskuterer Tobias' synspunkt. Fortæl, om du er enig eller uenig, og begrund din holdning. Brug mindst \textbf{tre} verbalsubstantiver på \textit{-ning / -else} fra denne lektion.
\end{emailbox}

\begin{answerbox}
\textbf{Før du læser 参考回答}
\begin{enumerate}
  \item \textit{AI kan hjælpe elever med lektielæsningen som en tålmodig sparringspartner og lette lærernes forberedelse, men kan også føre til afhængighed, snyd og svække elevernes selvstændige tænkning samt gøre vurderingen sværere.}
  \item \textit{Lukas mener, at AI-værktøjer vil være en naturlig del af fremtidens arbejdsliv, og Freja erkender, at forbud sjældent virker på lang sigt. De foreslår i stedet at kombinere skriftlige opgaver med mundtlig forklaring.}
\end{enumerate}

\textbf{Tekst B 阅读核对答案}
\begin{enumerate}
  \item \da{Forkert}（她认为禁令从长远看很少奏效：\textit{forbud sjældent virker på lang sigt}，转而提议加强口头阐释）
  \item \da{Fordi nye værktøjer bliver en naturlig del af arbejdslivet.}（未来职场中操作新工具是日常的一部分）
  \item \da{At kombinere skriftlige opgaver med en mundtlig forklaring eller samtale.}（将书面作业与口头阐述/面谈结合）
\end{enumerate}

\textbf{练习 A 答案}
\begin{enumerate}
  \item \da{anvendelse}（固定搭配：\textit{ved anvendelse af}，泛指方式用不定指）
  \item \da{indførelsen}（明确特指“新规则的引入”，定指形式加 \textit{-n}）
  \item \da{forberedelse}（修饰语 \textit{en grundig} 表明此处须用不定指单数名词）
  \item \da{forståelse}（所有格 \textit{elevernes} 后接不带冠词的派生名词）
  \item \da{vurdering}（\textit{en retfærdig vurdering af...} 对……的公正评估）
  \item \da{udarbejdelsen}（\textit{Under udarbejdelsen af opgaven} 在撰写作业期间）
  \item \da{udvikling}（\textit{den hurtige udvikling} 定指前加形容词，但名词自身用基础形）
  \item \da{besvarelse}（\textit{en mundtlig besvarelse} 口头作答）
\end{enumerate}

\textbf{练习 B 答案（句型浓缩参考答案）}
\begin{enumerate}
  \item \textit{\da{Ved anvendelse af} digitale hjælpemidler bliver lektierne mere overskuelige.}（介词短语放句首，主句维持 V2 倒装 \textit{bliver lektierne}）
  \item \textit{\da{Efter indførelsen af} nye regler er der mindre snyd til prøverne.}（将整个时间从句浓缩为 \textit{Efter + 名词化 + af}）
  \item \textit{\da{Ved lærerens vurdering af} opgaven lægger hun vægt på elevens egne ideer.}（保留原主语时使用所有格 \textit{lærerens}）
  \item \textit{\da{Under udarbejdelsen af} projektet må eleverne gerne søge inspiration på nettet.}（\textit{mens man udarbejder} 浓缩为 \textit{Under udarbejdelsen af}）
  \item \textit{\da{Et totalt forbud mod} AI løser ikke problemet på lang sigt.}（把 \textit{At man forbyder...} 转为干净利落的名词短语作主语）
  \item \textit{\da{Ved en kombination af} skriftligt arbejde og en mundtlig samtale kan man teste elevens viden bedre.}（双名词并用）
\end{enumerate}

\textbf{练习 C 答案}
\begin{enumerate}
  \item \textit{Ved anvendelse af kunstig intelligens kan elever få hurtig hjælp.}
  \item \textit{Efter indførelsen af de nye regler er eksamen blevet mere retfærdig.}
  \item \textit{En retfærdig vurdering af skriftlige opgaver er meget vigtig for læreren.}
  \item \textit{Den hurtige udvikling af kunstig intelligens stiller nye krav til danske skoler.}
  \item \textit{Et totalt forbud mod nye værktøjer kan ikke erstatte træning i kritisk tænkning.}
  \item \textit{Under forberedelsen af fremlæggelsen bør eleverne forstå alle vigtige begreber.}
\end{enumerate}

\textbf{DU3 Modul 4 高分范文（Debatindlæg · 139 词）}\\
\textit{\textbf{Forbud mod AI er den forkerte vej}}

\textit{Tobias har ret i, at ukritisk brug af kunstig intelligens kan skabe problemer i skolen. Hvis elever lader en maskine skrive deres opgaver, mister de nemlig træningen i at formulere egne tanker.}

\textit{Ikke desto mindre er jeg uenig i, at et totalt forbud er løsningen. Den teknologiske \da{udvikling} stopper jo ikke ved skolens dør, og unge har brug for at lære en ansvarlig \da{anvendelse} af moderne værktøjer før deres arbejdsliv. \dabold{Ved forberedelsen af} svære emner kan AI faktisk fungere som en nyttig sparringspartner.}

\textit{Derfor bør skolerne hellere ændre eksamensformen end at forbyde teknologien. En klog \da{kombination} af skriftlige afleveringer og mundtlige \da{besvarelser} sikrer stadig en retfærdig \da{vurdering} af elevernes niveau. På den måde styrker vi både digital dannelse og selvstændig tænkning.}

\small\textbf{考官点评}：全文共 139 词（远超 Modul 4 Opgave 2 的 100 词底线），结构紧扣“礼貌承认对方合理担忧 $\rightarrow$ 用 \textit{Ikke desto mindre} 反驳总禁令 $\rightarrow$ 提出替代方案”。全篇自然嵌入了 \textit{udvikling, anvendelse, ved forberedelsen af, kombination, besvarelser, vurdering} 六处名词化表达与情态小品词 \textit{jo}，紧凑度与思辨力均达到 B2- 优异水准。
\end{answerbox}

\ncdlesson{12}{Børn, skærmtid og sociale medier}{儿童、屏幕时间与社交媒体}
{DU3 Modul 4（B1+ $\rightarrow$ B2-）｜对标 NCE 3 \& Modultest 4 \textit{Skrivning Opgave 1（125+ 词回信）\& 严密前提论证}}
{\textbf{限制与前提条件连词}（\textit{medmindre, forudsat at, såfremt, i det omfang}）· 从句否定词前置语序与句首条件从句触发的 V2 倒装 · 与 \textit{hvis} / \textit{selvom} 的边界辨析}
{儿童的屏幕时间、夜间睡眠以及社交媒体低龄化，是当今丹麦家庭与政策辩论中最受关注的话题之一。在 Modul 4 写作与口语辩论中，把观点推向 B2- 水准的关键，在于不再只会用笼统的 \textit{hvis}，而是懂得用 \textbf{\textit{medmindre}}（除非）、\textbf{\textit{forudsat at}}（前提是）、\textbf{\textit{såfremt}}（倘若/只要，书面政策语体）和 \textbf{\textit{i det omfang}}（在……限度内）精确划定立场的边界与前提条件。}

% ---------------------------------------------------------------------
\begin{foerduleaeser}
带着下面两个问题阅读 Tekst A 和 Tekst B：
\begin{enumerate}
  \item \textit{Hvorfor advarer eksperter mod for meget skærmtid hos børn ifølge Tekst A, og på hvilke betingelser kan skærmen stadig være positiv?}（根据 Tekst A，为什么专家对儿童过多屏幕时间提出警告？在什么前提下，屏幕仍然可以是积极的？）
  \item \textit{Hvilke aftaler foreslår Maja og Henrik for børnenes mobilbrug i klassen i Tekst B?}（在 Tekst B 中，玛雅和亨里克针对班里孩子的手机使用提出了哪些约定方案？）
\end{enumerate}
\end{foerduleaeser}

% ---------------------------------------------------------------------
\ncdsubsec{二、Tekster（科技媒体专题与思辨评论双语篇）}

\subsection*{Tekst A：\textit{Børn, skærmtid og sociale medier}（社会政策述评 · 约 250 词）}

\begin{tipbox}{语篇说明}
本篇围绕丹麦关于儿童手机禁令与社交媒体年龄限制的公共辩论展开。阅读时请重点标出\textbf{\color{NCDNavy}深蓝色加粗的四大限制/前提条件连词}，并留意\textbf{\color{NCDRed}深红色加粗的从句否定词位置与主句 V2 倒装}。
\end{tipbox}

I de seneste år har debatten om børns skærmtid og adgang til sociale medier fyldt meget i Danmark. Flere pædagoger og børnelæger peger på, at konstant brug af telefoner og tablets kan gå ud over både nattesøvnen, koncentrationen og det fysiske samvær i hverdagen. Mange mener derfor, at problemet vil vokse, \da{medmindre} voksne sætter langt tydeligere rammer omkring børns digitale vaner.

Alligevel advarer forskere imod at se alle skærme som entydigt skadelige. Digitale platforme kan nemlig støtte kreativitet, sprogudvikling og venskaber, \da{forudsat at} brugen bliver styret med omtanke. For en genert teenager kan et trygt onlinefællesskab endda give mod og selvtillid, \da{i det omfang} det supplerer --- og \dabold{ikke erstatter} --- det virkelige liv.

Spørgsmålet handler altså sjældent om et totalt nej til teknologi, men om klare betingelser. I mange danske folkeskoler har man indført mobilfri skoledage, så telefonerne bliver lagt i en kasse om morgenen. En sådan regel skaber mere ro i frikvartererne, \da{såfremt} både lærere og forældre bakker konsekvent op om ordningen.

På politisk plan diskuteres der nu strammere aldersgrænser for sociale medier, fordi unge ofte bliver udsat for algoritmer, der er designet til at fastholde deres opmærksomhed. Kritikere understreger dog, at lovgivning alene ikke hjælper, \da{medmindre} techgiganterne samtidig tvinges til at tage ansvar. \dabold{Forudsat at} skoler, forældre og platforme samarbejder, \dabold{kan børn} lære at bruge nettet på en sund måde.

\begin{transbox}
\small
\textbf{\color{NCDTeal}[Tekst A 中文精译]}\\
近年来，关于儿童屏幕时间与使用社交媒体的争论在丹麦占据了很大篇幅。多位教育员和儿科医生指出，持续使用手机和平板电脑不仅会损害夜间睡眠，还会影响注意力和日常的线下相处。因此许多人认为，除非成年人为孩子的数字习惯设立清晰得多的边界，否则这一问题还将继续恶化。

尽管如此，研究人员仍提醒不要把所有屏幕都一概视为有害。只要使用方式经过深思熟虑的引导，数字平台其实能够促进创造力、语言发展和友谊维系。对一名内向的青少年而言，只要安全的线上社群是对现实生活的补充——而非取代现实生活——它甚至还能赋予勇气和自信。

因此，关键很少在于全盘否定技术，而在于明确前提条件。在许多丹麦公立学校，人们已经推行“无手机上学日”，早晨就把手机统一放进收纳盒。倘若老师和家长双方都始终如一地支持这项制度，这样的规定就能为课间休息带来更多安宁。

在政策层面，丹麦目前正讨论针对社交媒体设立更严格的年龄门槛，因为年轻人常常暴露在旨在牢牢抓住注意力的算法之下。不过批评人士强调，除非同时迫使科技巨头承担责任，否则单靠立法也无济于事。前提只要学校、家长和平台三方协作，孩子们依然能够学会以健康的方式使用网络。
\end{transbox}

\subsection*{Tekst B：\textit{Forældresnak efter klassemødet}（家长群讨论节选 · 约 140 词）}

\begin{emailbox}{5 年级家长讨论｜Maja \& Henrik}
\textbf{Maja:} Jeg er faktisk glad for skolens nye forslag om at vente med sociale medier til 13-årsalderen. Det kan skabe mere tryghed i klassen, \da{forudsat at} de fleste forældre følger aftalen.\\
\textbf{Henrik:} Enig! Ingen børn tør jo sige nej alene, \da{medmindre} vennerne også holder pause fra de samme apps. Men hvad med klassens chattråd om fodbold og lektier?\\
\textbf{Maja:} Børnene kan sagtens have en fælles gruppe på SMS, \da{i det omfang} en voksen lige holder øje med tonen. Så slipper vi også for reklamer og skadelige videoer.\\
\textbf{Henrik:} God idé! Og \da{såfremt} alle skriver under på en fælles forældreaftale i næste uge, bliver det meget nemmere at håndhæve reglerne derhjemme.
\end{emailbox}

\textbf{Tekst B 阅读核对（Rigtigt / Forkert / Kort svar，答案见第六部分）：}
\begin{enumerate}
  \item \textit{Henrik mener, at børn sagtens kan droppe sociale medier alene uden støtte fra kammeraterne.} \underline{\hspace{1.6cm}}
  \item \textit{Hvilken betingelse stiller Maja for, at børnene kan have en fælles chatgruppe på SMS?} \underline{\hspace{2.2cm}}
  \item \textit{Hvad vil ifølge Henrik gøre det lettere at overholde skærmreglerne hjemme hos familierne?} \underline{\hspace{2.2cm}}
\end{enumerate}

\begin{tipbox}{本课发音与语调要点（Udtale \& Reduktioner）}
\begin{itemize}
  \item \textbf{四大条件连词的地道读音与停顿}：
    \begin{itemize}
      \item \textit{medmindre} \ipa{mɛˈmenˀdʁə}（重音在第二音节 \textit{-min-}，第一个 \textit{d} 完全不发音，\textit{med} 弱读为 \ipa{mɛ}！）
      \item \textit{forudsat at} \ipa{ˈfɒːuðˌsad ad}（重音在最前面 \textit{for-}，\textit{ud} 读软 d，带前重音）
      \item \textit{såfremt} \ipa{sʌˈfʁɛmˀd}（重音在后音节 \textit{-fremt} 且带 \textit{stød}）
      \item \textit{i det omfang} \ipa{i de ˈʌmˌfɑŋˀ}（\textit{det} 弱读不带 \textit{t}，重音落在 \textit{om-}，词尾 \textit{ng} 带 \textit{stød}）
    \end{itemize}
  \item \textbf{媒体与政策高频词}：\textit{en skærmtid} \ipa{ˈsgɛɐ̯mˌtˢiðˀ}；\textit{en aldersgrænse} \ipa{ˈalˀɐsˌɡʁɛnsə}（\textit{alder} 中的 \textit{d} 不发音！）；\textit{en koncentration} \ipa{kʰʌnsɛntˢʁɑˈɕoːˀn}；\textit{at håndhæve} \ipa{ˈhʌnˌhɛːˀvə}（\textit{hånd} 中的 \textit{d} 不发音）。
\end{itemize}
\end{tipbox}

\begin{transbox}
\small
\textbf{\color{NCDTeal}[Tekst B 中文对照]}\\
\textbf{玛雅}：我其实很支持学校提出的新倡议，等到 13 岁再让孩子使用社交媒体。只要大多数家长都遵守这项约定，这就能为班级营造更多的安全感。\\
\textbf{亨里克}：同意！毕竟除非朋友们也一起暂停使用那些应用，否则没有哪个孩子敢一个人说“不”。不过，班里讨论足球和作业的聊天群怎么办呢？\\
\textbf{玛雅}：只要有位成年人顺便留意一下发言语气，孩子们完全可以在短信上建一个公共群聊。这样一来，我们也能避开广告和有害视频了。\\
\textbf{亨里克}：好主意！倘若大家下周都能在一份共同的家长公约上签字，我们在家里执行起这些规则就会容易得多。
\end{transbox}

% ---------------------------------------------------------------------
\ncdsubsec{三、Gloser（核心生词表）}

\begin{tipbox}{词表重点}
注意区分两个极易混淆的高频词：\textit{et samvær}（线下相处/共聚时光，不可数）与 \textit{et fællesskab}（社群/集体归属感）。另外，\textit{at udsætte} 搭配介词 \textit{for} 表示“使暴露于/遭受”（如 \textit{udsætte børn for reklamer}），被动语态常用 \textit{blive udsat for}。
\end{tipbox}

\begin{glosertable}
名词 & en skærmtid & \textit{(skærmtiden, ingen flertal)} & 屏幕时间，用屏时长 \\
名词 & en aldersgrænse & \textit{(aldersgrænsen, -grænser, -erne)} & 年龄限制，年龄门槛 \\
名词 & en koncentration & \textit{(koncentrationen, ingen flertal)} & 注意力，专注力 \\
名词 & et samvær & \textit{(samværet, ingen flertal)} & 相处，共聚时光 \\
名词 & et fællesskab & \textit{(fællesskabet, fællesskaber, -erne)} & 社群，集体，共同体 \\
名词 & en nattesøvn & \textit{(nattesøvnen, ingen flertal)} & 夜间睡眠 \\
名词 & en opmærksomhed & \textit{(opmærksomheden, ingen flertal)} & 注意力，关注度 \\
名词 & en ramme & \textit{(rammen, rammer, rammerne)} & 框架，边界，规范 \\
名词 & en aftale & \textit{(aftalen, aftaler, aftalerne)} & 约定，协议 \\
\midrule
动词 [har] & at begrænse & \textit{(begrænser, begrænsede, har begrænset)} & 限制，限定 \\
动词 [har] & at udsætte (for) & \textit{(udsætter, udsatte, har udsat)} & 使暴露于，使遭受 \\
动词 [har] & at håndhæve & \textit{(håndhæver, håndhævede, har håndhævet)} & 执行，维护（规则） \\
动词 [har] & at bakke op (om) & \textit{(bakker, bakkede, har bakket op om)} & 支持，声援 \\
动词 [har] & at fastholde & \textit{(fastholder, fastholdt, har fastholdt)} & 留住，坚持，维持 \\
\midrule
连词 & medmindre & \textit{(不变)} & 除非（后接从句，含否定排除意） \\
连词短语 & forudsat at & \textit{(不变)} & 前提是，只要（设定必要条件） \\
连词 & såfremt & \textit{(不变)} & 倘若，如果（正式书面/制度条件） \\
连词短语 & i det omfang & \textit{(不变)} & 在……程度上，在……限度内 \\
\end{glosertable}

% ---------------------------------------------------------------------
\ncdsubsec{四、Grammatik \& Stilistik（核心语法、伦理表达与修辞精讲）}

\begin{grammatikbox}{语法点 1：四大“限制与前提条件连词”精准分工表}
丹麦语中的条件句远不止 \textit{hvis}。在 Modul 4 辩论中，精准选用下面四个从属连词，可以立刻展示你论证的严密性与分寸感：

\vspace{0.3em}
\renewcommand{\arraystretch}{1.25}
{\small
\begin{tabularx}{\linewidth}{@{} >{\bfseries\raggedright\arraybackslash}p{2.6cm} >{\raggedright\arraybackslash}p{3.9cm} >{\raggedright\arraybackslash}X @{}}
\toprule
\textbf{条件连词} & \textbf{逻辑边界与语域} & \textbf{典型 B1+ / B2- 论证例句} \\
\midrule
\da{medmindre}\newline （除非） &
等价于 \textit{hvis ikke}；设定唯一例外情形（口语/书面通用） &
\textit{Reglen virker ikke, \da{medmindre} alle forældre støtter den.}\newline
（除非所有家长都支持，否则规定无效。） \\
\midrule
\da{forudsat at}\newline （前提是，只要） &
设定明确的前提条件，强调“只有满足此条件才成立” &
\textit{Børn må gerne spille, \da{forudsat at} de har lavet lektier.}\newline
（只要孩子们完成了作业，玩游戏也无妨。） \\
\midrule
\da{såfremt}\newline （倘若，如果） &
比 \textit{hvis} 更正式，常用于政策讨论、规章制度与报刊评论 &
\textit{\da{Såfremt} skolen indfører mobilfri dage, falder uroen.}\newline
（倘若学校实行无手机日，课堂躁动便会减少。） \\
\midrule
\da{i det omfang}\newline （在……限度内） &
限定程度、范围或比例，避免观点绝对化（极具思辨感！） &
\textit{Teknologi er gavnlig, \da{i det omfang} den støtter læringen.}\newline
（技术只在促进学习的限度内才是有益的。） \\
\bottomrule
\end{tabularx}
}
\end{grammatikbox}

\begin{grammatikbox}{语法点 2：从句语序铁律 \& 句首条件从句触发的主句 V2 倒装}
这四个词全部是\textbf{从属连词（ledsætningskonjunktioner）}，因此无论它们放在句中还是句首，都必须严格遵守两条丹麦语核心铁律：
\begin{enumerate}
  \item \textbf{从句内部语序（严禁倒装，中心副词前置）}：\\
    公式：\textbf{连词 + 主语 +（\textit{ikke / altid / kun / ofte}）+ 变位动词}。\\
    例：\textit{Forudsat at platformene \dabold{ikke indsamler} børns data, er appen tryg.}（注意 \textit{ikke} 放在动词 \textit{indsamler} 前面！）
  \item \textbf{条件从句放句首 $\rightarrow$ 主句立刻 V2 倒装！}：\\
    公式：\textbf{[\textit{Forudsat at / Såfremt / Medmindre / I det omfang} + 从句], + \dabold{主句动词} + \dabold{主句主语} + ...}\\
    \ding{172} \textit{Såfremt alle \dabold{samarbejder}, \dabold{bliver} hverdagen nemmere.}（不是 \textit{*hverdagen bliver}！）\\
    \ding{173} \textit{I det omfang unge \dabold{bruger} skærme kreativt, \dabold{kan} det være en fordel.}
\end{enumerate}
\end{grammatikbox}

\begin{grammatikbox}{英丹对比避坑指南 \& 辩论地道表达（Stilistik \& Faste vendinger）}
\begin{enumerate}
  \item \textbf{\textit{medmindre} 的“双重否定”致命陷阱}：因为 \textit{medmindre} 本身已包含“除非/如果不（= \textit{unless}）”的否定条件，里面\textbf{千万不能再多加 \textit{ikke}}！\\
    \ding{55}\ \textit{Han må ikke få en telefon, medmindre han \dabold{ikke} er fyldt 13 år.}（语义彻底颠倒！）\\
    \ding{51}\ \textit{Han må ikke få en telefon, \da{medmindre} han er fyldt 13 år.}
  \item \textbf{不要画蛇添足写 \textit{*forudsat hvis} 或 \textit{*såfremt at}}：\textit{forudsat at} 后面直接接主语（不能接 \textit{hvis}）；\textit{såfremt} 本身就是完整连词，现代标准书面语中后面\textbf{不加 \textit{at}}（\ding{51}\ \textit{Såfremt loven vedtages...}）。
  \item \textbf{严格区分 \textit{medmindre}（除非）与 \textit{selvom}（尽管）}：\textit{medmindre} 改变的是“前提成立与否”；\textit{selvom} 表示“即便发生某事结果依然如此”。
\end{enumerate}

\textbf{儿童与数字媒体高频搭配（Faste vendinger）：}
\begin{itemize}
  \item \da{at sætte klare rammer (omkring)}（为……设定清晰边界）
  \item \da{at gå ud over (søvnen / koncentrationen)}（损害、影响到……）
  \item \da{at bakke op om en aftale}（全力支持某项协议）
  \item \da{at indføre en aldersgrænse på 13 år}（设立 13 岁年龄门槛）
\end{itemize}
\end{grammatikbox}

% ---------------------------------------------------------------------
\ncdsubsec{五、Mønsterøvelser（句型与修辞操练）}

\textbf{练习 A：语境逻辑选填（从 \textit{medmindre, forudsat at, såfremt, i det omfang, ikke} 中选择最恰当的词填空）}
\begin{enumerate}
  \item Skærme i skolen kan være et godt værktøj, \underline{\hspace{2.2cm}} de bruges til faglige opgaver og ikke til spil.
  \item Børn under 13 år bør ikke oprette profiler på sociale medier, \underline{\hspace{2.2cm}} forældrene giver en udtrykkelig tilladelse.
  \item \underline{\hspace{2.2cm}} kommunen ønsker bedre trivsel i klasserne, må den støtte fælles fritidstilbud efter skole.
  \item Mobiler kan være tilladt i pausen, \underline{\hspace{2.2cm}} eleverne \underline{\hspace{2.2cm}} filmer hinanden uden samtykke.
  \item Ingen aftale om skærmtid virker i praksis, \underline{\hspace{2.2cm}} de voksne selv lægger telefonen væk ved middagsbordet.
  \item \underline{\hspace{2.2cm}} sociale medier styrker eksisterende venskaber, kan de spille en positiv rolle for unge.
  \item Eleverne må bruge tablets i gruppearbejdet, \underline{\hspace{2.2cm}} læreren har godkendt appen på forhånd.
  \item \underline{\hspace{2.2cm}} nye regler træder i kraft efter sommerferien, skal alle familier informeres i god tid.
\end{enumerate}

\textbf{练习 B：句型升级与纠错手术（按括号要求改写条件句或修正语法错误，注意从句语序与主句 V2 倒装）}
\begin{enumerate}
  \item 用 \textit{medmindre} 改写：\textit{Et forbud hjælper ikke, hvis techfirmaerne ikke ændrer deres algoritmer.} $\rightarrow$
  \item 用 \textit{forudsat at} 置句首合并（注意倒装）：\textit{Børn sover nok om natten. Så påvirker skærmtid ikke deres energi i skolen.} $\rightarrow$
  \item 升级为正式书面语 \textit{såfremt}（放句首）：\textit{Hvis alle forældre i klassen bliver enige, kan vi vente med smartphones til 6. klasse.} $\rightarrow$
  \item 用 \textit{i det omfang} 精准限定范围：\textit{Digitale medier er nyttige, men kun så længe de støtter børns kreativitet.} $\rightarrow$
  \item 纠错改写（双重否定错误）：\textit{Familien ser ikke film sammen, medmindre alle ikke har tid om aftenen.} $\rightarrow$
  \item 纠错改写（句首条件从句后缺 V2 倒装及 \textit{ikke} 位置错误）：\textit{Forudsat at forældrene snakker ikke i telefon ved bordet, børnene lærer hurtigt gode vaner.} $\rightarrow$
\end{enumerate}

\textbf{练习 C：中译丹（重点演练四大条件连词与 V2 倒装）}
\begin{enumerate}
  \item 除非成年人设定清晰的边界，否则孩子会花太多时间在屏幕上。
  \item 只要大家按时熄灯，手机就可以放在客厅里。
  \item 倘若学校推行无手机政策，课间气氛就会变得更安静。
  \item 社交媒体只有在补充现实友谊的限度内才是有益的。
  \item 前提是孩子没有忘记睡觉，适度玩游戏并非问题。（用 \textit{Forudsat at...}）
  \item 除非所有家长都支持这项规定，否则很难执行它。
\end{enumerate}

% ---------------------------------------------------------------------
\ncdsubsec{六、DU3 Modul 4 实战微任务与参考答案}

\begin{emailbox}{DU3 Modul 4 微型写作任务（对标 Skrivning Opgave 1 · 125+ 词私人邮件）}
\textbf{Situation}: Din ven Jonas har skrevet til dig, fordi hans 10-årige datter gerne vil have sin første smartphone og en profil på sociale medier ligesom nogle af sine veninder. Jonas er i tvivl om, hvad han skal gøre, og spørger også, hvordan debatten om børns skærmtid er i Danmark.\\[4pt]
\textbf{Opgave}: Skriv en mail til Jonas på \textbf{mindst 125 ord} (anbefalet 145--170 ord), hvor du:
\begin{itemize}
  \item reagerer på Jonas' dilemma og fortæller, hvad du selv synes, han bør gøre
  \item giver konkrete betingelser for telefon- og skærmbrug i familien
  \item fortæller kort om den danske debat om børns skærmtid og sociale medier.
\end{itemize}
\textit{（要求：自然融入至少 \textbf{三个} 本课条件连词 \textit{medmindre, forudsat at, såfremt, i det omfang} 以及情态小品词。）}
\end{emailbox}

\begin{answerbox}
\textbf{Før du læser 参考回答}
\begin{enumerate}
  \item \textit{Eksperter advarer om, at for meget skærmtid kan skade børns nattesøvn, koncentration og fysiske samvær. Skærmen kan dog være positiv, forudsat at brugen styres med omtanke, og i det omfang den supplerer det virkelige liv.}
  \item \textit{Maja og Henrik foreslår at vente med sociale medier til 13-årsalderen, at oprette en fælles SMS-gruppe under voksenopsyn og at underskrive en fælles forældreaftale i klassen.}
\end{enumerate}

\textbf{Tekst B 阅读核对答案}
\begin{enumerate}
  \item \da{Forkert}（他明确指出没有孩子敢一个人说“不”，除非朋友也一起暂停使用这些应用）
  \item \da{At en voksen holder øje med tonen i gruppen.}（只要有成年人留意群里的语气即可）
  \item \da{At alle skriver under på en fælles forældreaftale.}（大家共同签署一份家长公约）
\end{enumerate}

\textbf{练习 A 答案}
\begin{enumerate}
  \item \da{forudsat at}（设定积极前提：只要用于学科任务而非游戏）
  \item \da{medmindre}（除非家长明确许可）
  \item \da{Såfremt}（市政政策层面的正式条件“倘若”）
  \item \da{forudsat at}；\da{ikke}（条件从句内否定词 \textit{ikke} 必须置于动词 \textit{filmer} 之前）
  \item \da{medmindre}（除非成年人自己也放下手机）
  \item \da{I det omfang}（从程度与限度上界定：在能增进既有友谊的限度内）
  \item \da{forudsat at}（前提是老师已提前批准该应用）
  \item \da{Såfremt}（正式书面/规章场景：倘若新规在暑假后生效）
\end{enumerate}

\textbf{练习 B 答案（句型改写参考答案）}
\begin{enumerate}
  \item \textit{Et forbud hjælper ikke, \da{medmindre} techfirmaerne ændrer deres algoritmer.}（去掉冗余的 \textit{ikke}，用 \textit{medmindre} 表达“除非”）
  \item \textit{\da{Forudsat at} børn sover nok om natten, \dabold{påvirker skærmtid ikke} deres energi i skolen.}（从句置首，主句动词 \textit{påvirker} 紧跟逗号之后形成 V2 倒装）
  \item \textit{\da{Såfremt} alle forældre i klassen bliver enige, \dabold{kan vi} vente med smartphones til 6. klasse.}（正式连词替代 \textit{hvis}）
  \item \textit{Digitale medier er nyttige, \da{i det omfang} de støtter børns kreativitet.}（精准体现“在……限度内”）
  \item \textit{Familien ser ikke film sammen, \da{medmindre} alle har tid om aftenen.}（删去从句中多余的 \textit{ikke}，消除双重否定错误）
  \item \textit{\da{Forudsat at} forældrene \dabold{ikke snakker} i telefon ved bordet, \dabold{lærer børnene} hurtigt gode vaner.}（从句中 \textit{ikke} 放在动词前，主句执行 V2 倒装 \textit{lærer børnene}）
\end{enumerate}

\textbf{练习 C 答案}
\begin{enumerate}
  \item \textit{Børn bruger for meget tid på skærme, medmindre voksne sætter tydelige rammer.}
  \item \textit{Telefonerne må ligge i stuen, forudsat at alle slukker lyset til tiden.}
  \item \textit{Såfremt skolen indfører en mobilfri politik, bliver stemningen i frikvartererne mere rolig.}
  \item \textit{Sociale medier er kun gavnlige, i det omfang de supplerer virkelige venskaber.}
  \item \textit{Forudsat at barnet ikke glemmer at sove, er lidt spil ikke noget problem.}
  \item \textit{Det er svært at håndhæve reglen, medmindre alle forældre bakker op om den.}
\end{enumerate}

\textbf{DU3 Modul 4 高分范文（Skrivning Opgave 1 · 154 词，不含署名）}\\
\textit{Til: Jonas \quad|\quad Emne: Sv: Smartphone og sociale medier til din datter}

\textit{Hej Jonas}

\textit{Tak for din mail! Jeg kan jo sagtens forstå, at du er i tvivl, for når veninderne allerede har en mobil, vil din datter naturligvis ikke føle sig udenfor.}

\textit{Personligt synes jeg, at ti år stadig er tidligt til sociale medier. På den anden side kan en almindelig telefon til opkald og beskeder være fin, \da{forudsat at} I laver nogle klare aftaler derhjemme. Skærmen bør for eksempel blive ude af soveværelset om aftenen. Børn mister let nattesøvnen, \da{medmindre} der er faste rammer for slukketid.}

\textit{Her i Danmark fylder debatten om børns skærmtid rigtig meget for tiden. Mange skoler har indført mobilfri skoledage, og flere eksperter foreslår strammere aldersgrænser for sociale platforme. \da{Såfremt} forældrene i klassen kan blive enige om en fælles linje, bliver presset på det enkelte barn meget mindre. Digitale medier er nemlig kun en hjælp, \da{i det omfang} de styrker fællesskabet uden at stjæle børnenes ro.}

\textit{Held og lykke med snakken derhjemme!}

\textit{Mange hilsner\\
\mbox{[Dit navn]}}

\small\textbf{考官点评}：全信共 154 词（不含署名与标题），完全达到 Modul 4 \textit{Skrivning Opgave 1} 对 125+ 词私人深度回复的篇幅与语域要求：开头用情态小品词 \textit{jo} 表达共情，随后用 \textit{forudsat at}、\textit{medmindre}、\textit{Såfremt}（句首触发主句 V2 倒装 \textit{bliver presset}）和 \textit{i det omfang} 层层设定前提边界，个人建议与丹麦社会大讨论结合紧密自然。
\end{answerbox}

\ncdlesson{13}{Overvågning i det offentlige rum: Tryghed kontra privatliv}{公共监控：安全感 vs 隐私}
{DU3 Modul 4（B1+ $\rightarrow$ B2-）｜对标 NCE 3 \& Modultest 4 \textit{Skrivning Opgave 2（Debatindlæg）\& Mundtlig Opgave 2（利弊权衡）}}
{\textbf{价值权衡核心句型}（\textit{afveje A over for B, gå på kompromis med, ske på bekostning af}）· 衍生权衡表达（\textit{veje tungere end, opveje, i bytte for}）· \textit{tryghed} 与 \textit{sikkerhed} 文化辨析}
{当车站、广场和商业街安装越来越多监控摄像头时，丹麦社会最常展开一场围绕\textbf{安全感（\textit{tryghed}）与隐私权（\textit{privatliv}）}的伦理大辩论。在 DU3 Modul 4 写作与双人辩论中，考官看重的不是情绪化的“支持或反对”，而是你能否像成熟评论员那样，运用 \textbf{\textit{afveje A over for B}}、\textbf{\textit{gå på kompromis med}} 和 \textbf{\textit{ske på bekostning af}} 精准表达两难取舍！}

% ---------------------------------------------------------------------
\begin{foerduleaeser}
带着下面两个问题阅读 Tekst A 和 Tekst B：
\begin{enumerate}
  \item \textit{Hvorfor ønsker politiet og nogle politikere mere videoovervågning i byerne ifølge Tekst A, og hvad advarer juristerne imod?}（根据 Tekst A，为什么警方和一些政界人士希望在城市增加视频监控？法学专家又对此提出了什么警告？）
  \item \textit{Hvilket konkret krav stiller Kasper i sit debatindlæg i Tekst B, før man sætter flere kameraer op?}（在 Tekst B 的读者来信中，卡斯珀在安装更多摄像头前提出了什么具体前提要求？）
\end{enumerate}
\end{foerduleaeser}

% ---------------------------------------------------------------------
\ncdsubsec{二、Tekster（科技媒体专题与思辨评论双语篇）}

\subsection*{Tekst A：\textit{Overvågning i det offentlige rum: Tryghed kontra privatliv}（社会伦理专题 · 约 255 词）}

\begin{tipbox}{语篇说明}
本篇剖析丹麦公共场所扩大视频监控的政策争议。阅读时请重点观察\textbf{\color{NCDNavy}深蓝色加粗的价值权衡句型}（注意介词 \textit{over for / med / af} 的固定搭配），以及\textbf{\color{NCDRed}深红色加粗的立场让步与补充权重表达}。
\end{tipbox}

På danske togstationer, gågader og i nattelivet hænger der i dag tusindvis af kameraer. Spørgsmålet om, hvor meget vi skal overvåges i det offentlige rum, skaber jævnligt en hed politisk debat. Her står et vigtigt dilemma i centrum: Hvordan kan samfundet \da{afveje} hensynet til borgernes tryghed \da{over for} retten til et privatliv?

Tilhængere af udvidet videoovervågning --- herunder flere kommuner og politifolk --- fremhæver, at kameraer gør det lettere at opklare alvorlig kriminalitet og skabe hurtigere indsats ved vold eller tyveri. For mange borgere giver synlige kameraer på mørke stationer desuden en følelse af sikkerhed om aftenen. I deres øjne \dabold{vejer} bekæmpelsen af kriminalitet \dabold{tungere end} ulempen ved at blive filmet på gaden.

Modstandere og retssikkerhedseksperter ser anderledes på sagen. De påpeger, at konstant overvågning kan ændre vores adfærd i det fri, hvis vi altid føler os observeret af myndighederne. En øget tryghed for nogle må derfor ikke \da{ske på bekostning af} alle borgeres grundlæggende frihed. Samtidig viser internationale undersøgelser, at kameraer sjældent forebygger impulsiv vold, selvom de hjælper bagefter.

Derfor advarer flere forskere imod at \da{gå på kompromis med} databeskyttelsen uden meget klare grænser. \dabold{En reel afvejning} af sikkerhed \da{over for} privatliv kræver nemlig gennemsigtighed: Borgerne skal vide, hvem der gemmer optagelserne, hvor længe de opbevares, og hvornår hensynet til efterforskning virkelig \dabold{opvejer} indgrebet i den enkeltes hverdag.

\begin{transbox}
\small
\textbf{\color{NCDTeal}[Tekst A 中文精译]}\\
如今在丹麦的火车站、步行街和夜生活区，悬挂着成千上万个摄像头。在公共空间里我们究竟应该被监控多少，经常引发激烈的政治争论。争论的核心摆着一道重要难题：社会应当如何在安全感需求与隐私权之间作出权衡？

支持扩大视频监控的人——包括多个市政厅和警界人士——强调，摄像头让侦破严重犯罪变得更容易，并在发生暴力或盗窃时能更迅速出警。此外，对许多市民来说，昏暗车站里清晰可见的摄像头在夜间也带来了一定的安全感。在他们看来，打击犯罪比在街上被拍摄的不便分量更重。

反对者和法治保障专家对此事看法不同。他们指出，如果我们始终感觉自己在被当局注视，持续的监控可能会改变我们在公共场合的行为方式。因此，为一部分人增加安全感，绝不能以牺牲全体公民的基本自由为代价。同时，国际研究表明，尽管摄像头有助于事后侦破，但很少能真正预防冲动型暴力。

正因如此，多位研究人员警告不要在缺乏极明确边界的情况下对数据保护作出妥协。对安全与隐私进行真正的权衡，需要透明度作支撑：公民必须知道是谁保存录像、录像会留存多久，以及侦查破案的需要究竟何时才真正抵得过对个人日常生活的干预。
\end{transbox}

\subsection*{Tekst B：\textit{Debatindlæg: Tryghed må ikke koste vores frihed}（报刊读者来信 · 约 140 词）}

\begin{emailbox}{读者来信投稿｜Skribent: Kasper Holm}
Flere politikere ønsker nu endnu flere kameraer i vores bymidte. Jeg forstår godt ønsket om at skabe større tryghed efter mørkets frembrud. Alligevel mener jeg, at vi aldrig må \da{gå på kompromis med} borgernes retssikkerhed.

Vi skal naturligvis hjælpe politiet med at opklare forbrydelser, men en hurtig løsning må ikke \da{ske på bekostning af} privatlivet for tusindvis af helt almindelige mennesker på vej hjem fra arbejde. Desuden kan kameraer alene ikke skabe ægte tillid mellem naboer.

Derfor bør kommunen først \da{afveje} den reelle effekt \da{over for} de demokratiske omkostninger. Hvis man sætter nye kameraer op, skal optagelserne slettes automatisk efter få dage, og pengene bør også gå til bedre gadebelysning og flere sociale indsatser.
\end{emailbox}

\textbf{Tekst B 阅读核对（Rigtigt / Forkert / Kort svar，答案见第六部分）：}
\begin{enumerate}
  \item \textit{Kasper mener, at øget overvågning automatisk skaber større tillid mellem naboer.} \underline{\hspace{2.2cm}}
  \item \textit{Hvad kræver Kasper, der skal ske med optagelserne, hvis man sætter nye kameraer op?} \underline{\hspace{2.2cm}}
  \item \textit{Hvilke alternative tiltag foreslår Kasper i stedet for kun at bruge penge på kameraer?} \underline{\hspace{2.2cm}}
\end{enumerate}

\begin{tipbox}{本课发音与语调要点（Udtale \& Reduktioner）}
\begin{itemize}
  \item \textbf{权衡核心词与不发音辅音}：
    \begin{itemize}
      \item \textit{en overvågning} \ipa{ˈʌʊ̯ɐˌvɔʊ̯ˀneŋ}（两个 \textit{v} 与 \textit{åg} 均滑音化，词干带 \textit{stød}！）
      \item \textit{en tryghed} \ipa{ˈtˢʁøɡˌheðˀ}（与 \textit{sikkerhed} 区分，词尾为软 \textit{d} 加喉塞音）
      \item \textit{et privatliv} \ipa{pʰʁiˈvæːˀdˌliʊ̯ˀ}（重音在第二音节 \textit{-vat-}，末尾 \textit{-v} 读 \ipa{ʊ̯ˀ}）
      \item \textit{at afveje} \ipa{ˈɑʊ̯ˌvɑi̯ə}（\textit{af-} 读 \ipa{ɑʊ̯}，\textit{ej} 读滑音 \ipa{ɑi̯}）
      \item \textit{et kompromis} \ipa{kʰʌmpʰʁoˈmi}（注意法语借词！\textbf{词尾 \textit{s} 在丹麦语单数中完全不发音}！）
      \item \textit{på bekostning af} \ipa{pʰʌ beˈkʰʌsdneŋ ˈæːˀ}（口语常弱读为 \ipa{ɑ}）
    \end{itemize}
  \item \textbf{其他辩论高频词}：\textit{retssikkerhed} \ipa{ˈʁɛdsˌseɡɐˌheðˀ}；\textit{kriminalitet} \ipa{kʰʁiminɑliˈtˢeːˀd}（重音在最后 \textit{-tet}）；\textit{tillid} \ipa{ˈtˢeˌliðˀ}。
\end{itemize}
\end{tipbox}

\begin{transbox}
\small
\textbf{\color{NCDTeal}[Tekst B 中文对照]}\\
\textbf{读者投稿（卡斯珀·霍尔姆）}：多位政界人士如今希望在我们市中心安装更多摄像头。我完全理解大家想在天黑后营造更大安全感的愿望。然而我认为，我们绝不能对公民的法治安全保障作出妥协。

我们自然应当协助警方侦破犯罪，但一项权宜之计绝不能以牺牲成千上万下班回家的普通民众的隐私为代价。此外，单靠摄像头本身也无法在邻里之间建立真正的信任。

因此，市政厅应当首先把实际效果放在民主代价面前仔细权衡。如果确实安装了新摄像头，录像必须在几天后自动删除，而且资金也应该投入到更好的街道照明和更多社会帮扶措施上。
\end{transbox}

% ---------------------------------------------------------------------
\ncdsubsec{三、Gloser（核心生词表）}

\begin{tipbox}{词表重点}
\textit{et kompromis} 的发音与复数最容易出错：单数时词尾 \textit{s} 不发音，但定指（\textit{kompromiset}）和复数（\textit{kompromiser}）中 \textit{s} 要恢复发音！另外请注意丹麦语中 \textit{tryghed}（主观心理的安心与托底感）与 \textit{sikkerhed}（客观物理层面的防护安全）的细微区分。
\end{tipbox}

\begin{glosertable}
名词 & en overvågning & \textit{(overvågningen, ingen flertal)} & 监控，监视 \\
名词 & en tryghed & \textit{(trygheden, ingen flertal)} & 安全感，安心感 \\
名词 & et privatliv & \textit{(privatlivet, ingen flertal)} & 隐私，私人生活 \\
名词 & et kompromis & \textit{(kompromiset, kompromiser, -erne)} & 妥协，折中方案 \\
名词 & en bekostning & \textit{(bekostningen, ingen flertal)} & 代价，费用（\textit{på bekostning af}） \\
名词 & en afvejning & \textit{(afvejningen, afvejninger, -erne)} & 权衡，斟酌 \\
名词 & en retssikkerhed & \textit{(retssikkerheden, ingen flertal)} & 法治安全，法律程序保障 \\
名词 & en kriminalitet & \textit{(kriminaliteten, ingen flertal)} & 犯罪，犯罪率 \\
名词 & en tillid & \textit{(tilliden, ingen flertal)} & 信任，信赖 \\
名词 & en gennemsigtighed & \textit{(gennemsigtigheden, ingen flertal)} & 透明度，公开性 \\
名词 & en optagelse & \textit{(optagelsen, optagelser, optagelserne)} & 录像，录音 \\
\midrule
动词 [har] & at afveje & \textit{(afvejer, afvejede, har afvejet)} & 权衡，衡量（\textit{over for}） \\
动词 [har] & at opveje & \textit{(opvejer, opvejede, har opvejet)} & 抵消，弥补，胜过 \\
动词 [har] & at forebygge & \textit{(forebygger, forebyggede, har forebygget)} & 预防，防范 \\
动词 [har] & at opklare & \textit{(opklarer, opklarede, har opklaret)} & 侦破，查明 \\
动词 [har] & at beskytte & \textit{(beskytter, beskyttede, har beskyttet)} & 保护，维护 \\
\midrule
形容词 & offentlig & \textit{(offentligt, offentlige)} & 公共的，公开的 \\
\end{glosertable}

% ---------------------------------------------------------------------
\ncdsubsec{四、Grammatik \& Stilistik（核心语法、伦理表达与修辞精讲）}

\begin{grammatikbox}{语法点 1：价值权衡三大黄金句型与精准介词骨架}
在探讨科技、法律与社会政策时，极少有“非黑即白”的答案。下面这三组句型是 DU3 Modul 4 辩论中展现思辨分寸的王牌武器：

\vspace{0.3em}
\renewcommand{\arraystretch}{1.25}
{\small
\begin{tabularx}{\linewidth}{@{} >{\bfseries\raggedright\arraybackslash}p{3.0cm} >{\raggedright\arraybackslash}p{4.0cm} >{\raggedright\arraybackslash}X @{}}
\toprule
\textbf{核心句型} & \textbf{固定结构与介词铁律} & \textbf{标准 B1+ / B2- 示范句} \\
\midrule
\da{at afveje A over for B}\newline （在 B 面前权衡 A） &
必须用双词介词 \textbf{\textit{over for}}（切勿拆开或误作 \textit{til / med}） &
\textit{Politikerne skal \da{afveje} tryghed \da{over for} privatliv.}\newline
\textit{（可名词化：\textbf{en afvejning af A over for B}）} \\
\midrule
\da{at gå på kompromis med A}\newline （对 A 作出妥协） &
固定搭配 \textbf{\textit{gå på kompromis med}}（常用否定表达底线） &
\textit{Vi må aldrig \da{gå på kompromis med} retssikkerheden.}\newline
（我们绝不能对法治保障作出妥协。） \\
\midrule
\da{at ske på bekostning af B}\newline （以牺牲 B 为代价） &
主语常为举措或成果；介词固定为 \textbf{\textit{på ... af}} &
\textit{Øget sikkerhed må ikke \da{ske på bekostning af} friheden.}\newline
（也可以直接作状语：\textit{... \da{på bekostning af} privatlivet.}） \\
\bottomrule
\end{tabularx}
}
\end{grammatikbox}

\begin{grammatikbox}{语法点 2：进阶权重表达与让步融合术（B2- 提分组合）}
光说“A 和 B 都很重要”停留在 B1 水平；结合权重对比动词与让步句，就能立刻写出高水准辩证结论：
\begin{enumerate}
  \item \textbf{用 \da{veje tungere end}（比……分量更重）表态倾斜}：\\
    \textit{For mange borgere \da{vejer} hensynet til tryghed \da{tungere end} frygten for kameraer.}
  \item \textbf{用 \da{opveje}（抵消 / 胜过）评估利弊净值}：\\
    \textit{Spørgsmålet er, om fordelene ved overvågning virkelig \da{opvejer} ulemperne.}
  \item \textbf{“肯定一点 + 守住底线”的让步黄金公式}：\\
    公式：\textbf{\textit{Selvom} + [正面价值] + , må det ikke \dabold{ske på bekostning af} / \dabold{gå på kompromis med} + [核心底线]。}\\
    例：\textit{Selvom kameraer hjælper politiet, må overvågning aldrig \dabold{ske på bekostning af} borgernes privatliv.}
\end{enumerate}
\end{grammatikbox}

\begin{grammatikbox}{英丹/中丹对比排雷 \& 社会伦理地道搭配（Stilistik \& Faste vendinger）}
\begin{enumerate}
  \item \textbf{介词搭配零容忍错误}：
    \begin{itemize}
      \item \ding{55}\ \textit{afveje A \dabold{med / mod} B} $\rightarrow$ \ding{51}\ \textit{afveje A \da{over for} B}（\textit{over for} 分开写两个词！）
      \item \ding{55}\ \textit{gå på kompromis \dabold{på / om} vores frihed} $\rightarrow$ \ding{51}\ \textit{gå på kompromis \da{med} vores frihed}
      \item \ding{55}\ \textit{ske på bekostning \dabold{til} privatlivet} $\rightarrow$ \ding{51}\ \textit{ske på bekostning \da{af} privatlivet}
    \end{itemize}
  \item \textbf{\textit{tryghed}（安全感）vs \textit{sikkerhed}（安全性）}：丹麦社会极其看重 \textit{tryghed}，它指人内心感到“被保护、放心、生活无忧”（如在车站感到安心）；而 \textit{sikkerhed} 侧重客观物理、技术或数据层面的安全防护（如 \textit{datasikkerhed, trafiksikkerhed}）。
  \item \textbf{分清 \textit{overvåge}（动词）与 \textit{overvågning}（名词）}：\textit{overvågning} 通常不可数（直接说 \textit{mere overvågning}，不加 \textit{-er}）。
\end{enumerate}

\textbf{高频伦理辩论搭配（Faste vendinger）：}
\begin{itemize}
  \item \da{hensynet til (tryghed / privatlivet)}（对……的考量/顾及）
  \item \da{i det offentlige rum}（在公共空间中）
  \item \da{i bytte for (større sikkerhed)}（以换取更高的安全性）
  \item \da{at sætte klare grænser for}（为……划定明确边界）
\end{itemize}
\end{grammatikbox}

% ---------------------------------------------------------------------
\ncdsubsec{五、Mønsterøvelser（句型与修辞操练）}

\textbf{练习 A：价值权衡精准填空（填入恰当的介词或核心词汇：\textit{over for, med, af, på, tungere, end, kompromis, opveje}）}
\begin{enumerate}
  \item Myndighederne er nødt til at afveje politiets behov \underline{\hspace{2.2cm}} borgernes ret til privatliv.
  \item En hurtig løsning på gadekriminalitet må ikke ske \underline{\hspace{2.2cm}} bekostning \underline{\hspace{2.2cm}} vores retssikkerhed.
  \item Vi må aldrig gå på \underline{\hspace{2.2cm}} med grundlæggende demokratiske rettigheder.
  \item For mange pendlere vejer følelsen af tryghed \underline{\hspace{2.2cm}} \underline{\hspace{2.2cm}} ulempen ved et kamera på stationen.
  \item Det er tvivlsomt, om flere kameraer i nattelivet kan \underline{\hspace{2.2cm}} indgrebet i den personlige frihed.
  \item En ansvarlig lovgivning kræver en grundig afvejning \underline{\hspace{2.2cm}} sikkerhed over for gennemsigtighed.
  \item Nogle borgere er villige til at acceptere mere overvågning i bytte \underline{\hspace{2.2cm}} hurtigere politihjælp.
  \item Politiet må ikke gå på kompromis \underline{\hspace{2.2cm}} databeskyttelsen, når optagelser gemmes.
\end{enumerate}

\textbf{练习 B：句型重构手术（用本课指定的价值权衡句型，将下列简单对立句升级为成熟思辨句）}
\begin{enumerate}
  \item 用 \textit{afveje A over for B} 改写：\textit{Samfundet skal både tænke på tryghed og på privatliv.} $\rightarrow$
  \item 用 \textit{ske på bekostning af} 改写：\textit{Vi vil have mere sikkerhed, men det må ikke ødelægge vores frihed.} $\rightarrow$
  \item 用 \textit{gå på kompromis med} 改写：\textit{Staten må ikke acceptere lavere standarder for borgernes retssikkerhed.} $\rightarrow$
  \item 用 \textit{veje tungere end} 改写：\textit{For kritikere er privatlivet vigtigere end muligheden for at fange flere tyve.} $\rightarrow$
  \item 合并让步与底线（\textit{Selvom..., må... ikke ske på bekostning af...}）：\textit{Kameraer hjælper med opklaringen. De må ikke fjerne tilliden i samfundet.} $\rightarrow$
  \item 纠错改写（两处介词错误）：\textit{Vi bør afveje tryghed med privatliv og ikke gå på kompromis om friheden.} $\rightarrow$
\end{enumerate}

\textbf{练习 C：中译丹（活用本课三大价值权衡句型与主题词汇）}
\begin{enumerate}
  \item 社会必须在安全感与个人隐私之间进行权衡。（用 \textit{afveje ... over for}）
  \item 更多的监控绝不能以牺牲基本自由为代价。（用 \textit{ske på bekostning af}）
  \item 我们绝不能在公民的数据保护上作出妥协。（用 \textit{gå på kompromis med}）
  \item 对一些人来说，夜间车站上的安全感分量更重。（用 \textit{veje tungere}）
  \item 问题是，这些优势是否真的抵得过对隐私的干预。（用 \textit{opveje}）
  \item 在公共空间里，清晰的数据保留边界是必要的。
\end{enumerate}

% ---------------------------------------------------------------------
\ncdsubsec{六、DU3 Modul 4 实战微任务与参考答案}

\begin{emailbox}{DU3 Modul 4 微型写作任务（对标 Skrivning Opgave 2 · Debatindlæg）}
\textbf{Situation}: På et lokalt debatsite skriver borgeren Mads:\\
\textit{”Hvis man ikke har noget at skjule, gør det jo ikke noget, at der hænger kameraer på alle gader og pladser. Jo mere overvågning vi har, jo tryggere bliver vores by.”}\\[4pt]
\textbf{Opgave}: Skriv et debatindlæg på \textbf{mindst 100 ord} (anbefalet 130--155 ord), hvor du forholder dig til Mads' synspunkt. Gør rede for både fordele og ulemper ved overvågning i det offentlige rum, og brug mindst \textbf{tre} værdiafvejningsudtryk fra denne lektion (\textit{afveje ... over for, gå på kompromis med, ske på bekostning af, veje tungere end}).
\end{emailbox}

\begin{answerbox}
\textbf{Før du læser 参考回答}
\begin{enumerate}
  \item \textit{Politiet og nogle politikere ønsker mere overvågning for at opklare alvorlig kriminalitet og øge trygheden, mens jurister advarer om, at konstant overvågning kan ændre vores adfærd og true den grundlæggende frihed og databeskyttelsen.}
  \item \textit{Kasper kræver, at man først afvejer effekten over for de demokratiske omkostninger, at optagelser slettes automatisk efter få dage, og at kommunen også investerer i gadebelysning og sociale indsatser.}
\end{enumerate}

\textbf{Tekst B 阅读核对答案}
\begin{enumerate}
  \item \da{Forkert}（他指出单靠摄像头本身无法在邻里之间建立真正的信任）
  \item \da{De skal slettes automatisk efter få dage.}（录像必须在几天后自动删除）
  \item \da{Bedre gadebelysning og flere sociale indsatser.}（更好的街道照明和更多的社会帮扶措施）
\end{enumerate}

\textbf{练习 A 答案}
\begin{enumerate}
  \item \da{over for}（固定搭配：\textit{afveje A over for B}，分开写两个词）
  \item \da{på}；\da{af}（固定结构：\textit{ske på bekostning af}）
  \item \da{kompromis}（\textit{gå på kompromis med} 对……妥协）
  \item \da{tungere}；\da{end}（\textit{veje tungere end} 比……更重要）
  \item \da{opveje}（情态动词 \textit{kan} 后接原形：抵消、胜过）
  \item \da{af}（名词化结构：\textit{en afvejning af A over for B}）
  \item \da{for}（固定搭配：\textit{i bytte for} 以换取）
  \item \da{med}（固定介词：\textit{kompromis med}）
\end{enumerate}

\textbf{练习 B 答案（价值权衡改写参考答案）}
\begin{enumerate}
  \item \textit{Samfundet er nødt til at \da{afveje} hensynet til tryghed \da{over for} hensynet til privatliv.}（用 \textit{afveje ... over for} 替换简单并列）
  \item \textit{Ønsket om mere sikkerhed må ikke \da{ske på bekostning af} vores frihed.}（以“以牺牲……为代价”整合对比两句）
  \item \textit{Staten må aldrig \da{gå på kompromis med} borgernes retssikkerhed.}
  \item \textit{For kritikere \da{vejer} privatlivet \da{tungere end} muligheden for at fange flere tyve.}
  \item \textit{Selvom kameraer hjælper med opklaringen, må overvågning ikke \da{ske på bekostning af} tilliden i samfundet.}（从句置首后主句维持 V2 倒装 \textit{må overvågning}）
  \item \textit{Vi bør afveje tryghed \dabold{over for} privatliv og ikke gå på kompromis \dabold{med} friheden.}（修正 \textit{*med} 为 \textit{over for}，修正 \textit{*om} 为 \textit{med}）
\end{enumerate}

\textbf{练习 C 答案}
\begin{enumerate}
  \item \textit{Samfundet skal afveje tryghed over for det personlige privatliv.}
  \item \textit{Mere overvågning må aldrig ske på bekostning af den grundlæggende frihed.}
  \item \textit{Vi må aldrig gå på kompromis med databeskyttelsen for borgerne.}
  \item \textit{For nogle mennesker vejer trygheden på mørke stationer tungere om aftenen.}
  \item \textit{Spørgsmålet er, om fordelene virkelig opvejer indgrebet i privatlivet.}
  \item \textit{I det offentlige rum er klare grænser for opbevaring af data nødvendige.}
\end{enumerate}

\textbf{DU3 Modul 4 高分范文（Debatindlæg · 144 词）}\\
\textit{\textbf{Tryghed kræver mere end kameraer}}

\textit{Mads har ret i, at flere kameraer kan hjælpe politiet med at opklare kriminalitet og skabe større tryghed på mørke stationer. Ingen ønsker jo at føle sig utryg på vej hjem om aftenen.}

\textit{Alligevel mener jeg, at vi skal passe på med at sætte kameraer op overalt. Vi må aldrig \da{gå på kompromis med} borgernes grundlæggende privatliv, og kampen mod kriminalitet bør ikke \da{ske på bekostning af} vores frihed til at færdes i det offentlige rum uden konstant kontrol. Det handler altså om at \da{afveje} hensynet til sikkerhed \da{over for} retten til at være i fred.}

\textit{Hvis politikerne vil have en tryggere by, er bedre gadebelysning, synligt politi og sociale indsatser ofte en bedre løsning. Her \da{vejer} forebyggelse og tillid efter min mening \da{tungere end} endeløs overvågning.}

\small\textbf{考官点评}：全篇 144 词，精准反驳了论坛上“问心无愧就不怕监控”的常见片面论调。文章首先温和肯定监控的侦破与夜行安全感作用（用 \textit{jo} 自然拉近读者），随后依次调度 \textit{gå på kompromis med}、\textit{ske på bekostning af}、\textit{afveje ... over for} 与 \textit{vejer ... tungere end} 提出兼具原则性与替代对策的均衡结论，完全达到 B2- 优异水准。
\end{answerbox}

\ncdlesson{14}{Køb og smid væk-kulturen under pres}{受冲击的“买了就扔”消费文化}
{DU3 Modul 4（B1+ $\rightarrow$ B2-）｜对标 NCE 3 \& Modultest 4 \textit{Læsning（书面长定语拆解）\& Skrivning Opgave 1（125+ 词回信）}}
{\textbf{正式书面语前置长分词定语}（如 \textit{den af regeringen vedtagne lov}）的识别与自然改写 · 不规则过去分词的定指词尾变化（\textit{vedtagne, fremlagte, oversete, trufne}）· 书面语体与日常沟通转换}
{阅读丹麦报刊时政评论、环保法条以及 DU3 Modul 4 阅读真题时，你会频繁遇见“把整个被动从句塞进定冠词和名词之间”的长定语结构，例如 \textit{den af regeringen vedtagne lov}（由政府通过的那部法律）。掌握这套结构“先拆中心词、再还原施动者”的解码法，并学会在自己写邮件或辩论文章时将它改写为自然流畅的关系从句，你就攻克了 Modul 4 阅读最硬核的拦路虎！}

% ---------------------------------------------------------------------
\begin{foerduleaeser}
带着下面两个问题阅读 Tekst A 和 Tekst B：
\begin{enumerate}
  \item \textit{Hvorfor er brug-og-smid-væk-kulturen kommet under pres i dag ifølge Tekst A, og hvilke politiske tiltag nævnes?}（根据 Tekst A，为什么“用了就扔”的消费文化如今正面临压力？文中提到了哪些政策举措？）
  \item \textit{Hvorfor beslutter Sofie i Tekst B sig for at få repareret sin kaffemaskine i stedet for at købe en ny?}（在 Tekst B 中，索菲为什么决定去修好自己的咖啡机，而不是直接买一台新的？）
\end{enumerate}
\end{foerduleaeser}

% ---------------------------------------------------------------------
\ncdsubsec{二、Tekster（科技媒体专题与思辨评论双语篇）}

\subsection*{Tekst A：\textit{Køb og smid væk-kulturen under pres}（消费与环保评论 · 约 255 词）}

\begin{tipbox}{语篇说明}
本篇探讨欧洲与丹麦正兴起的“维修权”和反过度消费浪潮。阅读时请特别捕捉\textbf{\color{NCDNavy}深蓝色加粗的前置长分词定语}（观察“定冠词 $\rightarrow$ 施动者/状语 $\rightarrow$ 分词 \textit{-e} $\rightarrow$ 核心名词”的语序链），并与\textbf{\color{NCDRed}深红色加粗的自然 \textit{som} 关系从句}作语体对比。
\end{tipbox}

I årtier har det moderne forbrugersamfund været præget af en ”køb og smid væk”-kultur. Når en brødrister, en støvsuger eller et par høretelefoner går i stykker efter et par år, vælger mange forbrugere automatisk at købe nyt, fordi reparation ofte virker dyrere end udskiftning. Denne vane har skabt \da{den hurtigt voksende mængde} af elektronikaffald, som i dag udgør en alvorlig klimabelastning.

Nu møder denne kultur dog stigende modstand fra både politikere og borgere. \da{De af EU vedtagne regler} om ret til reparation skal fremover gøre det lettere og billigere at få hverdagsprodukter fikset. Formålet med \da{den nyligt fremlagte lovgivning} er, at producenterne forpligtes til at levere reservedele i op til ti år og designe apparater, \dabold{som nemt kan skilles ad}.

Samtidig ser man et markant holdningsskift ude i samfundet. I mange danske byer åbner frivillige reparationscaféer, hvor lokale borgere mødes for at give nyt liv til \da{de af ejerne ellers kasserede ting}. Her handler det ikke kun om penge, men også om at bryde \da{det af reklamer skabte behov} for hele tiden at eje den nyeste model.

Alligevel peger økonomer på \da{et ofte overset problem}: Hvis reservedele og arbejdsløn fortsat er høje, vælger travle familier stadig den hurtigste løsning i butikken. Derfor kræver en reel grøn omstilling ikke blot gode intentioner, men også økonomiske incitamenter, \dabold{der gør det bæredygtige valg billigere i hverdagen}.

\begin{transbox}
\small
\textbf{\color{NCDTeal}[Tekst A 中文精译]}\\
几十年来，现代消费社会一直深受“买了就扔”文化的烙印影响。当烤面包机、吸尘器或一副耳机用了几年出现故障时，许多消费者会自动选择买新的，因为维修往往显得比更换新机更贵。这种习惯制造了如今构成严重气候负担的、迅速增长的电子垃圾总量。

不过如今，这种文化正遭到来自政界和民众双方越来越大的抵制。欧盟新近通过的关于“维修权”的法规，未来应当会让日常用品的维修变得更容易且更便宜。这项最新提交的立法的目的，在于迫使生产商在长达十年内必须提供备件，并设计出易于拆解的电器设备。

与此同时，人们在社会中也看到了明显的态度转变。在许多丹麦城市，志愿维修咖啡馆纷纷开张，当地居民聚在一起，为那些原本会被主人丢弃的物品赋予新生命。在这里，这不仅关乎省钱，更关乎打破由商业广告营造出来的、让人不断想要拥有最新款式的消费需求。

尽管如此，经济学家也指出了一个常被忽视的问题：如果备件和人工费依然高昂，忙碌的家庭在商店里还是会选最快捷的方案。因此，真正的绿色转型不仅需要良好的意愿，还需要能够让可持续选择在日常生活中变得更实惠的经济激励措施。
\end{transbox}

\subsection*{Tekst B：\textit{Kaffemaskinen i reparationscaféen}（日常交流对话 · 约 140 词）}

\begin{emailbox}{邻居闲聊｜Sofie \& Mikkel}
\textbf{Mikkel:} Hej Sofie! Hvad skal du med den gamle kaffemaskine under armen? Skal den ikke bare på genbrugspladsen?\\
\textbf{Sofie:} Nej, slet ikke! Jeg er på vej hen til den lokale reparationscafé. Jeg vil prøve at redde den i stedet for at smide den ud.\\
\textbf{Mikkel:} Kan det virkelig betale sig? Jeg troede, at den slags maskiner var umulige at åbne.\\
\textbf{Sofie:} Jo, jeg har faktisk læst \da{den af kommunen udgivne guide} om reparation. Det er ofte kun en lille pakning, der skal skiftes, og den reservedel koster næsten ingenting.\\
\textbf{Mikkel:} Smart! Jeg har selv en gammel lampe i kælderen, \dabold{som jeg aldrig fik fikset}. Måske skulle jeg tage med dig næste lørdag i stedet for at købe en ny.
\end{emailbox}

\textbf{Tekst B 阅读核对（Rigtigt / Forkert / Kort svar，答案见第六部分）：}
\begin{enumerate}
  \item \textit{Sofie vil aflevere sin gamle kaffemaskine på genbrugspladsen.} \underline{\hspace{2.2cm}}
  \item \textit{Hvad er ifølge Sofie årsagen til, at det godt kan betale sig at reparere kaffemaskinen?} \underline{\hspace{2.2cm}}
  \item \textit{Hvilken ting overvejer Mikkel nu selv at tage med til reparationscaféen næste lørdag?} \underline{\hspace{2.2cm}}
\end{enumerate}

\begin{tipbox}{本课发音与语调要点（Udtale \& Reduktioner）}
\begin{itemize}
  \item \textbf{不规则分词定语重音与弱音}：
    \begin{itemize}
      \item \textit{vedtagne} \ipa{ˈveðˌtˢɑːˀnə}（原形 \textit{vedtage} 变化时 \textit{g} 消失，\textit{ved} 读软 d！）
      \item \textit{fremlagte} \ipa{ˈfʁɛmˌlɑɡdə}（由 \textit{fremlægge} 变来，元音由 \textit{æ} 转 \textit{a}）
      \item \textit{oversete} \ipa{ˈʌʊ̯ɐˌseːˀdə}（重音在词首 \textit{over-}，\textit{v} 滑音化为 \ipa{ʊ̯}）
      \item \textit{kasserede} \ipa{kʰɑˈseːˀɐəðə}（外来动词后缀 \textit{-ere} 重音在 \textit{-se-}）
    \end{itemize}
  \item \textbf{绿色消费高频词}：\textit{bæredygtig} \ipa{ˈbɛːɐˌdøɡdi}（词尾 \textit{g} 不发音）；\textit{en reparation} \ipa{ʁepʰɑʁɑˈɕoːˀn}；\textit{et genbrug} \ipa{ˈɡɛnˌbʁuːˀ}（\textit{brug} 末尾 \textit{g} 在此作无声长音伴随 \textit{stød}）；\textit{en levetid} \ipa{ˈleːvəˌtˢiðˀ}。
\end{itemize}
\end{tipbox}

\begin{transbox}
\small
\textbf{\color{NCDTeal}[Tekst B 中文对照]}\\
\textbf{米克尔}：嗨，索菲！你胳膊底下夹着那台旧咖啡机准备去哪儿呀？它不是该直接送去回收站吗？\\
\textbf{索菲}：才不是呢！我正打算去社区维修咖啡馆。我想试着抢救一下它，而不是直接扔掉。\\
\textbf{米克尔}：这真的划算吗？我一直以为那种机器根本打不开呢。\\
\textbf{索菲}：划算呀，我其实刚读了市政厅编发的那份维修指南。通常只是里面一个小密封圈需要更换，而那个备件几乎不花什么钱。\\
\textbf{米克尔}：真聪明！我地下室里自己也有一盏一直没修好的旧台灯。下周六或许我该跟你一起去，而不是重新买一盏新的。
\end{transbox}

% ---------------------------------------------------------------------
\ncdsubsec{三、Gloser（核心生词表）}

\begin{tipbox}{词表重点}
本课重点掌握不规则强变化动词的过去分词定指形式：\textit{at vedtage $\rightarrow$ vedtaget $\rightarrow$ \dabold{vedtagne}}（而不是 \textit{*vedtagede}！）、\textit{at fremlægge $\rightarrow$ fremlagt $\rightarrow$ \dabold{fremlagte}}、\textit{at overse $\rightarrow$ overset $\rightarrow$ \dabold{oversete}}。记住 \textit{at skille ad}（拆开）和 \textit{at smide ud}（扔掉）均为高频可分动词短语。
\end{tipbox}

\begin{glosertable}
名词 & en forbrugskultur & \textit{(forbrugskulturen, -kulturer, -erne)} & 消费文化 \\
名词 & en reparation & \textit{(reparationen, reparationer, -erne)} & 维修，修理 \\
名词 & et affald & \textit{(affaldet, ingen flertal)} & 垃圾，废弃物 \\
名词 & en levetid & \textit{(levetiden, levetider, levetiderne)} & 使用寿命，生命期 \\
名词 & en lovgivning & \textit{(lovgivningen, ingen flertal)} & 立法，法规 \\
名词 & en garanti & \textit{(garantien, garantier, garantierne)} & 保修期，担保 \\
名词 & en reservedel & \textit{(reservedelen, reservedele, -ene)} & 备件，零配件 \\
名词 & en producent & \textit{(producenten, producenter, -erne)} & 生产商，制造商 \\
名词 & et genbrug & \textit{(genbruget, ingen flertal)} & 再利用，二手循环 \\
名词 & en ressource & \textit{(ressourcen, ressourcer, ressourcerne)} & 资源 \\
\midrule
动词 [har] & at vedtage & \textit{(vedtager, vedtog, har vedtaget)} & 通过，表决通过（法案） \\
动词 [har] & at fremlægge & \textit{(fremlægger, fremlagde, har fremlagt)} & 提交，公布（方案/法案） \\
动词 [har] & at kassere & \textit{(kasserer, kasserede, har kasseret)} & 报废，丢弃，淘汰 \\
动词 [har] & at reparere & \textit{(reparerer, reparerede, har repareret)} & 修理，维修 \\
动词 [har] & at forlænge & \textit{(forlænger, forlængede, har forlænget)} & 延长，延展 \\
\midrule
形容词 & bæredygtig & \textit{(bæredygtigt, bæredygtige)} & 可持续的，环保的 \\
形容词 & holdbar & \textit{(holdbart, holdbare)} & 耐用的，持久的 \\
\end{glosertable}

% ---------------------------------------------------------------------
\ncdsubsec{四、Grammatik \& Stilistik（核心语法、伦理表达与修辞精讲）}

\begin{grammatikbox}{语法点 1：拆解正式书面语“前置长分词定语”（Foranstillede participier）}
丹麦语正式新闻、法律文书和 B2 阅读题常常把一个包含施动者或状语的被动关系从句，直接压缩塞到\textbf{定冠词与中心名词之间}。看懂它的秘密只有一句话：\textbf{最左边是定冠词，最右边才是真正的主角名词！}

\vspace{0.3em}
\renewcommand{\arraystretch}{1.25}
{\small
\begin{tabularx}{\linewidth}{@{} >{\bfseries\raggedright\arraybackslash}p{2.3cm} >{\raggedright\arraybackslash}p{4.8cm} >{\raggedright\arraybackslash}X @{}}
\toprule
\textbf{结构组件} & \textbf{前置长分词定语（书面压缩体）} & \textbf{自然关系从句（日常/通用改写）} \\
\midrule
定冠词 + 施动者 + 过去分词 + 名词 &
\textit{\da{den} \dabold{af regeringen} \da{vedtagne} lov}\newline
（直译：那-由政府-通过的-法律） &
\textit{den lov, \dabold{som regeringen har vedtaget}}\newline
（或：\textit{den lov, der er vedtaget af regeringen}） \\
\midrule
复数定冠词 + 状语/施动者 + 分词 + 名词 &
\textit{\da{de} \dabold{af ejerne} kasserede ting}\newline
（那些被主人丢弃的物品） &
\textit{de ting, \dabold{som ejerne har kasseret}} \\
\midrule
定冠词 + 副词 + 现在分词（\textit{-ende}）+ 名词 &
\textit{\da{den} \dabold{hurtigt} \da{voksende} mængde affald}\newline
（那迅速增长的垃圾量） &
\textit{den mængde affald, \dabold{der vokser hurtigt}} \\
\bottomrule
\end{tabularx}
}

\vspace{0.3em}
\textbf{阅读三步“拆弹法”}：\ding{172} 跳过 \textit{den / det / de} 后面的长串修饰语，直接锁定紧挨后面的\textbf{中心名词}（如 \textit{lov / ting / problem}）；\ding{173} 把 \textit{af + 名词} 识别为“动作的发出者（原主语）”；\ding{174} 用“\textit{som + [施动者] + har/er + [分词]}”在脑中自然还原！
\end{grammatikbox}

\begin{grammatikbox}{语法点 2：分词定语的词尾变形规则 \& 语体适切性原则}
\begin{enumerate}
  \item \textbf{定指或复数时，过去分词必须变位（通常加 \textit{-e}）}：
    \begin{itemize}
      \item \textbf{规则弱变化动词（\textit{-et $\rightarrow$ -ede}）}：\textit{et kasseret produkt} $\rightarrow$ \textit{de \da{kasserede} produkter}；\textit{repareret $\rightarrow$ \da{reparerede}}。
      \item \textbf{不规则及 \textit{-t} 结尾分词（特殊变体，必背！）}：\textit{vedtaget $\rightarrow$ \da{vedtagne}}（去 \textit{-et} 加 \textit{-ne}）；\textit{skrevet $\rightarrow$ \da{skrevne}}；\textit{fundet $\rightarrow$ \da{fundne}}；\textit{fremlagt $\rightarrow$ \da{fremlagte}}；\textit{overset $\rightarrow$ \da{oversete}}（只加 \textit{-e}）。
      \item \textbf{不定指单数保持原形}：\textit{et ofte \da{overset} problem}（一个常被忽视的问题）。
    \end{itemize}
  \item \textbf{现在分词（\textit{-ende}）永远不变词尾}：无论单复数还是定指，一律保持 \textit{-ende}（如 \textit{den voksende mængde, de stigende priser}）。
  \item \textbf{什么时候用“长定语”，什么时候改回“关系从句”？}：\\
    在自己写 DU3 Modul 4 私人邮件（Opgave 1）或面向大众的论坛帖子（Opgave 2）时，\textbf{切忌生硬堆砌超过 4 个词的官僚长定语}（丹麦语称为 \textit{kancellisprog}）。阅读能秒懂长定语，写作主动改用轻巧的 \textit{som/der} 从句或简短的“副词 + 分词”（如 \textit{den nye lov, som Folketinget vedtog i går}），反而最地道！
\end{enumerate}
\end{grammatikbox}

\begin{grammatikbox}{英丹对比避坑指南 \& 可持续消费表达（Stilistik \& Faste vendinger）}
\begin{enumerate}
  \item \textbf{英语靠“后置分词短语”，丹麦语要么“前置全包”要么用 \textit{som} 从句}：英语常说 \textit{the law passed by the government}（名词 + 后置分词）。在丹麦语正式书面语里往往必须前置：\ding{51}\ \textit{den af regeringen vedtagne lov}，或者后接关系从句：\ding{51}\ \textit{den lov, som regeringen har vedtaget}（或 \textit{den lov, der blev vedtaget af regeringen}）。
  \item \textbf{漏掉定指分词词尾 \textit{-e}}：前面一旦出现定冠词 \textit{den / det / de}，后面的分词定语绝不能保留原形：\ding{55}\ \textit{den af EU vedtaget regel} $\rightarrow$ \ding{51}\ \textit{den af EU \dabold{vedtagne} regel}！
\end{enumerate}

\textbf{可持续消费高频搭配（Faste vendinger）：}
\begin{itemize}
  \item \da{at købe og smide væk}（买了就扔）/ \da{en brug-og-smid-væk-kultur}（用完即弃文化）
  \item \da{det kan ikke betale sig at reparere}（修理不划算）
  \item \da{at give gamle ting nyt liv}（赋予旧物新生命）/ \da{at forlænge levetiden}（延长使用寿命）
\end{itemize}
\end{grammatikbox}

% ---------------------------------------------------------------------
\ncdsubsec{五、Mønsterøvelser（句型与修辞操练）}

\textbf{练习 A：分词定语词尾变位填空（用括号内动词或分词的正确前置定语形式填空）}
\begin{enumerate}
  \item Den af parlamentet \underline{\hspace{2.2cm}} (vedtage) lov skal mindske mængden af elektronikaffald.
  \item De nyligt \underline{\hspace{2.2cm}} (fremlægge) forslag giver forbrugerne bedre ret til reparation.
  \item Mange frivillige reparerer de af ejerne \underline{\hspace{2.2cm}} (kassere) cykler og lamper.
  \item Det hurtigt \underline{\hspace{2.2cm}} (vokse) bjerg af affald bekymrer både forskere og borgere.
  \item Prisen på reservedele er et ofte \underline{\hspace{2.2cm}} (overse) problem i klimadebatten.
  \item De \underline{\hspace{2.2cm}} (reparere) kaffemaskiner kan ofte holde i flere år endnu.
  \item Den af kommunen \underline{\hspace{2.2cm}} (udgive) vejledning forklarer, hvordan man finder et værksted.
  \item Vi bør også læse de af forbrugerorganisationen \underline{\hspace{2.2cm}} (skrive) rapporter om kvalitet.
\end{enumerate}

\textbf{练习 B：语体“拆弹”与互转手术（1--4 题将正式前置长定语改写为自然的 \textit{som/der} 关系从句；5--6 题修正病句或逆向压缩）}
\begin{enumerate}
  \item 拆解为自然从句：\textit{Den af EU vedtagne lovgivning skal beskytte miljøet.} $\rightarrow$
  \item 拆解为自然从句：\textit{De af forbrugerne kasserede telefoner indeholder værdifulde materialer.} $\rightarrow$
  \item 拆解为自然从句：\textit{Det af reklamerne skabte forbrugspres påvirker især unge.} $\rightarrow$
  \item 拆解为自然从句：\textit{Den hurtigt voksende genbrugsbevægelse ændrer vores indkøbsvaner.} $\rightarrow$
  \item 纠错改写（分词未变位与语序错误）：\textit{Den af regeringen fremlagt plan skal skabe flere grønne arbejdspladser.} $\rightarrow$
  \item 压缩为书面长定语：\textit{De regler, som kommunen har vedtaget, træder i kraft til januar.} $\rightarrow$
\end{enumerate}

\textbf{练习 C：中译丹（活用本课消费环保词汇与分词定语/关系从句）}
\begin{enumerate}
  \item 人们丢弃了太多仍然可以修理的电子产品。
  \item 这项由政府通过的新法律要求生产商提供备件。
  \item 在许多丹麦城市，志愿者们正在赋予旧物品新的生命。
  \item 有时候维修旧设备甚至比买一台新机器更便宜。
  \item 延长产品的寿命对地球的资源非常重要。
  \item 迅速增长的电子垃圾量是一个常被忽视的问题。
\end{enumerate}

% ---------------------------------------------------------------------
\ncdsubsec{六、DU3 Modul 4 实战微任务与参考答案}

\begin{emailbox}{DU3 Modul 4 微型写作任务（对标 Skrivning Opgave 1 · 125+ 词私人邮件）}
\textbf{Situation}: Din ven Emma har skrevet til dig, fordi hendes to år gamle kaffemaskine pludselig larmer og ikke varmer ordentligt. Hun overvejer bare at smide den ud og købe en helt ny model, men hun ved også, at du går meget op i bæredygtighed.\\[4pt]
\textbf{Opgave}: Skriv en mail til Emma på \textbf{mindst 125 ord} (anbefalet 145--170 ord), hvor du:
\begin{itemize}
  \item reagerer på Emmas situation og råder hende til, hvad hun bør prøve først
  \item fortæller om dine egne erfaringer med at reparere ting eller købe genbrug
  \item kort beskriver, hvordan køb og smid væk-kulturen diskuteres i Danmark i dag.
\end{itemize}
\textit{（提示：写私人邮件时请把长定语改写为流畅自然的 \textit{som/der} 从句，并搭配情态小品词！）}
\end{emailbox}

\begin{answerbox}
\textbf{Før du læser 参考回答}
\begin{enumerate}
  \item \textit{Kulturen er kommet under pres på grund af den voksende mængde elektronikaffald og klimabelastningen. Af politiske tiltag nævnes de af EU vedtagne regler om ret til reparation samt lovgivning, der forpligter producenter til at levere reservedele.}
  \item \textit{Sofie vil redde kaffemaskinen i den lokale reparationscafé, fordi hun har læst kommunens guide og ved, at det ofte kun er en billig pakning, der skal skiftes.}
\end{enumerate}

\textbf{Tekst B 阅读核对答案}
\begin{enumerate}
  \item \da{Forkert}（她根本不想扔进回收站，而是带去维修咖啡馆抢救它）
  \item \da{Det er ofte kun en lille pakning, der skal skiftes, og den koster næsten ingenting.}（通常只需更换一个几乎不花钱的小密封圈）
  \item \da{En gammel lampe fra kælderen.}（地下室里的一盏旧台灯）
\end{enumerate}

\textbf{练习 A 答案}
\begin{enumerate}
  \item \da{vedtagne}（不规则动词 \textit{vedtage} 的定指过去分词：\textit{vedtaget $\rightarrow$ vedtagne}）
  \item \da{fremlagte}（定指复数变化：\textit{fremlagt $\rightarrow$ fremlagte}）
  \item \da{kasserede}（规则第一类弱变化动词的复数定指：\textit{kasseret $\rightarrow$ kasserede}）
  \item \da{voksende}（现在分词永远不变形，保持 \textit{-ende}）
  \item \da{overset}（不定指中性单数名词 \textit{et ... problem}，保留分词原形 \textit{overset}）
  \item \da{reparerede}（定指复数变化：\textit{repareret $\rightarrow$ reparerede}）
  \item \da{udgivne}（不规则变化 \textit{udgivet $\rightarrow$ udgivne}，由 \textit{den ... vejledning} 决定定指变位）
  \item \da{skrevne}（不规则变化 \textit{skrevet $\rightarrow$ skrevne}）
\end{enumerate}

\textbf{练习 B 答案（拆解与改写参考答案）}
\begin{enumerate}
  \item \textit{\da{Den lovgivning, som EU har vedtaget}, skal beskytte miljøet.}（将前置被动短语还原为日常自然的 \textit{som} 关系从句）
  \item \textit{\da{De telefoner, som forbrugerne har kasseret}, indeholder værdifulde materialer.}（\textit{af forbrugerne} 还原为关系从句主语）
  \item \textit{\da{Det forbrugspres, som reklamerne har skabt}, påvirker især unge.}
  \item \textit{\da{Den genbrugsbevægelse, der vokser hurtigt}, ændrer vores indkøbsvaner.}（现在分词拆解为主动进行式关系从句）
  \item \textit{Den af regeringen \dabold{fremlagte} plan skal skabe flere grønne arbejdspladser.}（定指结构 \textit{Den ... plan} 后，分词必须加 \textit{-e} 变为 \textit{fremlagte}）
  \item \textit{\da{De af kommunen vedtagne regler} træder i kraft til januar.}（将关系从句压缩为正式书面前置长定语）
\end{enumerate}

\textbf{练习 C 答案}
\begin{enumerate}
  \item \textit{Man kasserer alt for mange elektroniske varer, som stadig kan repareres.}
  \item \textit{Den nye lov, som regeringen har vedtaget, kræver reservedele fra producenterne.}
  \item \textit{I mange danske byer giver frivillige nu gamle ting nyt liv.}
  \item \textit{Nogle gange er det endda billigere at reparere et gammelt apparat end at købe nyt.}
  \item \textit{At forlænge produkters levetid er meget vigtigt for jordens ressourcer.}
  \item \textit{Den hurtigt voksende mængde elektronikaffald er et ofte overset problem.}
\end{enumerate}

\textbf{DU3 Modul 4 高分范文（Skrivning Opgave 1 · 157 词，不含署名）}\\
\textit{Til: Emma \quad|\quad Emne: Sv: Vent med at smide kaffemaskinen ud!}

\textit{Hej Emma}

\textit{Tak for din besked! Jeg kan da godt forstå, at det er irriterende med kaffemaskinen, men du bør faktisk ikke smide den ud endnu. Som regel er det jo bare en lille kalkprop eller en billig reservedel, der driller, når en maskine kun er to år gammel.}

\textit{Hvis jeg var dig, ville jeg først tjekke garantien eller tage den med i en lokal reparationscafé. Jeg var selv afsted sidste måned med en støvsuger, \da{som jeg næsten havde opgivet}. De frivillige fandt fejlen på ti minutter, og nu kører den perfekt igen!}

\textit{Her i Danmark er den gamle køb og smid væk-kultur virkelig kommet under pres i de senere år. Flere nye regler, \da{som både EU og Folketinget har vedtaget}, skal gøre det nemmere at få produkter fikset i stedet for at købe nyt hele tiden. Derudover oplever mange, at det både sparer ressourcer og penge at forlænge tingenes levetid.}

\textit{Skriv endelig, om du vil have adressen på caféen!}

\textit{Kærlig hilsen\\
\mbox{[Dit navn]}}

\small\textbf{考官点评}：全信共 157 词（不含署名与标题），语域极其贴合对好友艾玛的回信：主动使用了情态小品词 \textit{da / jo} 以及非真实条件句（\textit{Hvis jeg var dig, ville jeg...}）提出真诚建议；同时在私人邮件体裁中刻意将“欧盟与议会通过的法规”用自然的 \textit{som} 关系从句写出，而非生硬堆砌公文前置长定语，展现了精准的 B2- 语体切换自觉。
\end{answerbox}

\ncdlesson{15}{Nyheder i algoritmernes tidsalder: Ekkokamre og kildekritik}{算法时代的新闻：回音室与信源批判}
{DU3 Modul 4（B1+ $\rightarrow$ B2-）｜对标 NCE 3 \& Modultest 4 \textit{Læsning \& Skrivning Opgave 2（批判性思辨辩论）}}
{\textbf{批判性思维核心表达}（\textit{sætte spørgsmålstegn ved, forholde sig kritisk til, nuancere debatten}）· 媒体素养惯用语（\textit{tage for gode varer, skille fakta fra holdninger}）· Unit 3 收官综合思辨}
{欢迎来到 Unit 3“科技、媒体与现代伦理”的收官课！在算法主宰信息流的今天，人们极易被困在只放大同温层声音的“回音室（\textit{ekkokammer}）”中。在 DU3 Modul 4 考试中，如何用丹麦语冷静地提出质疑、区分事实与立场、并使讨论摆脱偏激的非黑即白？掌握本课的 \textbf{\textit{sætte spørgsmålstegn ved}}、\textbf{\textit{forholde sig kritisk til}} 与 \textbf{\textit{nuancere debatten}}，你就真正具备了丹麦语 B2- 公民思辨力！}

% ---------------------------------------------------------------------
\begin{foerduleaeser}
带着下面两个问题阅读 Tekst A 和 Tekst B：
\begin{enumerate}
  \item \textit{Hvordan opstår ekkokamre på sociale medier ifølge Tekst A, og hvorfor kan de skade den demokratiske samtale?}（根据 Tekst A，社交媒体上的“回音室”是如何产生的？为什么它们可能损害民主对话？）
  \item \textit{Hvilke tre konkrete vaner foreslår journalisten Sara i Tekst B, når man møder en opsigtsvækkende nyhed på nettet?}（在 Tekst B 中，记者萨拉建议大家在网上看到一条耸人听闻的新闻时，养成哪三个具体习惯？）
\end{enumerate}
\end{foerduleaeser}

% ---------------------------------------------------------------------
\ncdsubsec{二、Tekster（科技媒体专题与思辨评论双语篇）}

\subsection*{Tekst A：\textit{Nyheder i algoritmernes tidsalder: Ekkokamre og kildekritik}（媒体思辨专栏 · 约 255 词）}

\begin{tipbox}{语篇说明}
本篇聚焦算法推荐机制下的信息茧房与信源批判意识。阅读时请重点标出\textbf{\color{NCDNavy}深蓝色加粗的批判性思维核心动词短语}（留意介词 \textit{ved} 和反身代词 \textit{sig}），以及\textbf{\color{NCDRed}深红色加粗的习惯搭配与从句衔接}。
\end{tipbox}

For tyve år siden fik de fleste danskere deres daglige nyheder fra trykte aviser og faste tv-udsendelser. I dag orienterer især mange unge sig hurtigt gennem personlige feeds på sociale medier. Det giver nem og øjeblikkelig adgang til verdens begivenheder, men eksperter opfordrer borgerne til at \da{forholde sig kritisk til} den måde, nyhederne bliver udvalgt på.

Bag ethvert nyhedsfeed arbejder nemlig usynlige algoritmer. De registrerer, hvilke opslag vi klikker på, og viser os derefter endnu mere indhold, som ligner det samme. På den måde risikerer brugerne at havne i såkaldte ekkokamre, hvor deres egne fordomme hele tiden bliver bekræftet, uden at de møder modargumenter.

Når stærke følelser og hurtige delinger belønnes med flest klik, bliver det sværere at \dabold{skille fakta fra holdninger}. Derfor bør enhver læser \da{sætte spørgsmålstegn ved}, hvem der står bag en påstand, og hvad afsenderens formål egentlig er. Man må simpelthen aldrig \dabold{tage} en dramatisk overskrift \dabold{for gode varer}, før man har tjekket oplysningerne hos uafhængige medier.

Heldigvis kan man aktivt bryde ud af sin egen boble. Redaktionelle medier og bred public service-journalistik spiller fortsat en vigtig rolle, fordi uddannede journalister arbejder med at efterprøve kilder og \da{nuancere debatten}. Jo bedre vi træner vores kildekritik i skolen og i hverdagen, desto stærkere står det oplyste demokrati.

\begin{transbox}
\small
\textbf{\color{NCDTeal}[Tekst A 中文精译]}\\
二十年前，大多数丹麦人通过纸质报纸和固定时段的电视节目获取每日新闻。而如今，尤其是许多年轻人习惯通过社交媒体上的个性化信息流快速了解时事。这虽然提供了便捷即时的世界资讯入口，但专家呼吁公民对新闻被筛选的方式保持批判态度。

因为在每一条信息流的背后，都有看不见的算法在运转。算法会记录我们点击了哪些帖子，随后向我们推送更多相似的内容。如此一来，用户就可能落入所谓的“回音室”——在其中，他们自己的偏见不断得到印证，却很难接触到相反的论点。

当强烈情绪与快速转发能换来最多点击量时，将事实与观点区分开来就变得愈发困难。因此，每位读者都应当对某个论断背后是谁、发布者真正的意图是什么提出质疑。人们绝不应在尚未向独立媒体核查信息之前，就轻信一个耸动的标题。

幸好，人们可以主动打破自己的信息气泡。有专业采编团队的媒体以及覆盖面广的公共服务新闻依然扮演着重要角色，因为受过专业训练的记者致力于核实信源，并让公共讨论呈现出多维层次。我们在学校和日常生活中把信源批判力训练得越好，建立在充分知情基础上的民主制度就越稳固。
\end{transbox}

\subsection*{Tekst B：\textit{Sådan træner du din kildekritik}（媒体素养访谈节选 · 约 140 词）}

\begin{emailbox}{校园广播访谈｜Jonas（大学生）\& Sara（新闻记者）}
\textbf{Jonas:} Sara, mange unge scroller bare hurtigt forbi ti overskrifter i minuttet. Hvordan undgår man egentlig at blive snydt af falske nyheder?\\
\textbf{Sara:} Først og fremmest skal man lære at stoppe op og \da{sætte spørgsmålstegn ved}, om en historie lyder for utrolig til at være sand. Vær ekstra skeptisk, hvis en tekst gør dig vred med det samme.\\
\textbf{Jonas:} Hvad gør jeg så, hvis jeg gerne vil tjekke en påstand, før jeg deler den med vennerne?\\
\textbf{Sara:} Så bør du altid \da{forholde dig kritisk til} kilden og tjekke mindst to uafhængige nyhedsmedier. Derudover hjælper det at læse analyser, der forsøger at \da{nuancere debatten}, i stedet for kun at følge korte videoklip, der giver dig ret.
\end{emailbox}

\textbf{Tekst B 阅读核对（Rigtigt / Forkert / Kort svar，答案见第六部分）：}
\begin{enumerate}
  \item \textit{Sara anbefaler, at man straks deler nyheder, hvis de vækker stærke følelser eller vrede.} \underline{\hspace{2.2cm}}
  \item \textit{Hvor mange uafhængige nyhedsmedier bør man mindst tjekke ifølge Sara?} \underline{\hspace{2.2cm}}
  \item \textit{Hvorfor bør man vælge analyser frem for kun korte videoklip?} \underline{\hspace{2.2cm}}
\end{enumerate}

\begin{tipbox}{本课发音与语调要点（Udtale \& Reduktioner）}
\begin{itemize}
  \item \textbf{批判性思维核心词与软 \textit{d} / 不发音辅音}：
    \begin{itemize}
      \item \textit{et spørgsmålstegn} \ipa{ˈsbɶɐ̯smɔːlsˌtˢɑi̯ˀn}（\textit{spørgs-} 中的 \textit{g} 不发音；\textit{tegn} 读滑音 \ipa{tˢɑi̯ˀn}，\textit{g} 消失化作长音与喉塞音！）
      \item \textit{at nuancere} \ipa{nyɑŋˈseːˀɐ}（法语借词！\textit{nu-} 读轻柔的 \ipa{ny-}，\textit{-an-} 读鼻音 \ipa{-ɑŋ-}，重音在 \textit{-ce-}）
      \item \textit{en fordom} \ipa{ˈfɒːˌdʌmˀ}（重音在前面 \textit{for-}，不同于英文的后缀重音！）
      \item \textit{en påstand} \ipa{ˈpʰʌˌsdanˀ}（词尾 \textit{d} 不发音，带 \textit{stød}！）
    \end{itemize}
  \item \textbf{媒体与信息高频词}：\textit{en algoritme} \ipa{ɑlɡoˈʁedmə}；\textit{et ekkokammer} \ipa{ˈɛɡoˌkʰɑmˀɐ}；\textit{kildekritik} \ipa{ˈkʰeləˌkʰʁitˢiɡ}（第一个 \textit{d} 不发音！）；\textit{troværdig} \ipa{tˢʁoˈvɛɐ̯ˀdi}（词尾 \textit{g} 不发音）。
\end{itemize}
\end{tipbox}

\begin{transbox}
\small
\textbf{\color{NCDTeal}[Tekst B 中文对照]}\\
\textbf{乔纳斯}：萨拉，很多年轻人一分钟内就快速划过十来条新闻标题。到底怎样才能避免被假新闻欺骗呢？\\
\textbf{萨拉}：首先，人必须学会停下来并提出质疑：一个故事听起来是不是好得或离奇得不像真的。如果一段文字让你立刻感到愤怒，就要格外保持怀疑。\\
\textbf{乔纳斯}：那如果我想在转发给朋友之前核实某个说法，该怎么做呢？\\
\textbf{萨拉}：那你就应当始终对信息来源保持批判态度，并且至少核对两家独立的新闻媒体。此外，阅读那些试图让讨论更全面深入的分析文章会很有帮助，而不要只关注一味顺着你心意的短视频片段。
\end{transbox}

% ---------------------------------------------------------------------
\ncdsubsec{三、Gloser（核心生词表）}

\begin{tipbox}{词表重点}
注意词尾双写与变音规则：\textit{en fordom}（偏见）定指和复数要双写 \textit{m}（\textit{fordommen, fordomme, fordommene}）；\textit{at forholde sig til} 是反身动词短语，绝对不能漏掉反身代词 \textit{sig}；而 \textit{en kildekritik}（信源批判能力）是不可数共性名词。
\end{tipbox}

\begin{glosertable}
名词 & en algoritme & \textit{(algoritmen, algoritmer, algoritmerne)} & 算法 \\
名词 & et ekkokammer & \textit{(ekkokammeret, ekkokamre, -kamrene)} & 回音室，信息茧房 \\
名词 & en kildekritik & \textit{(kildekritikken, ingen flertal)} & 信源批判，信息甄别能力 \\
名词 & en kilde & \textit{(kilden, kilder, kilderne)} & 来源，信源，消息源 \\
名词 & et spørgsmålstegn & \textit{(spørgsmålstegnet, spørgsmålstegn, -ene)} & 问号，质疑（单复同形） \\
名词 & en fordom & \textit{(fordommen, fordomme, fordommene)} & 偏见，成见 \\
名词 & en påstand & \textit{(påstanden, påstande, påstandene)} & 论断，主张，说法 \\
名词 & en dækning & \textit{(dækningen, dækninger, dækningerne)} & 新闻报道，覆盖面 \\
名词 & en troværdighed & \textit{(troværdigheden, ingen flertal)} & 可信度，公信力 \\
名词 & en offentlighed & \textit{(offentligheden, ingen flertal)} & 公众舆论界，社会公众 \\
\midrule
动词 [har] & at nuancere & \textit{(nuancerer, nuancerede, har nuanceret)} & 使具层次感，多维度阐明 \\
动词 [har] & at forholde sig (til) & \textit{(forholder, forholdt, har forholdt sig til)} & 对……持……态度，应对 \\
动词 [har] & at bekræfte & \textit{(bekræfter, bekræftede, har bekræftet)} & 证实，印证，强化 \\
动词 [har] & at efterprøve & \textit{(efterprøver, efterprøvede, har efterprøvet)} & 核查，检验，复核 \\
动词 [har] & at skille & \textit{(skiller, skilte, har skilt)} & 区分，分离（\textit{skille A fra B}） \\
\midrule
形容词 & troværdig & \textit{(troværdigt, troværdige)} & 可信的，可靠的 \\
形容词 & ensidig & \textit{(ensidigt, ensidige)} & 片面的，单方面的 \\
\end{glosertable}

% ---------------------------------------------------------------------
\ncdsubsec{四、Grammatik \& Stilistik（核心语法、伦理表达与修辞精讲）}

\begin{grammatikbox}{语法点 1：批判性思维三大王牌句型与介词骨架}
进入 DU3 Modul 4 辩论与阅读深水区，你必须精准驾驭下面三组核心批判性思维动词短语。请把它们的\textbf{固定介词、反身成分与从句接口}背成肌肉记忆：

\vspace{0.3em}
\renewcommand{\arraystretch}{1.25}
{\small
\begin{tabularx}{\linewidth}{@{} >{\bfseries\raggedright\arraybackslash}p{3.3cm} >{\raggedright\arraybackslash}p{3.9cm} >{\raggedright\arraybackslash}X @{}}
\toprule
\textbf{核心表达} & \textbf{句法结构与接续规则} & \textbf{标准 B1+ / B2- 示范句} \\
\midrule
\da{at sætte spørgsmålstegn ved}\newline （对……提出质疑） &
后接名词，或直接接 \textbf{\textit{, om / hvorvidt}} 从句；介词固定为 \textbf{\textit{ved}} &
\textit{Man bør \da{sætte spørgsmålstegn ved} kildens formål.}\newline
\textit{Jeg \da{sætter spørgsmålstegn ved, om} tallene passer.} \\
\midrule
\da{at forholde sig kritisk til}\newline （对……保持批判态度） &
必须带反身代词（\textit{mig / dig / sig / os}）+ 介词 \textbf{\textit{til}} &
\textit{Unge skal lære at \da{forholde sig kritisk til} nyheder på nettet.}\newline
\textit{（可替换副词：\textit{forholde sig skeptisk / neutralt til}）} \\
\midrule
\da{at nuancere debatten}\newline （使讨论更具层次/避免片面） &
直接及物动词（不加介词！），常见宾语：\textit{debatten, billedet, diskussionen} &
\textit{Artiklen er med til at \da{nuancere debatten} om algoritmer.}\newline
（也可以用分词定语：\textit{en \da{nuanceret} dækning}） \\
\bottomrule
\end{tabularx}
}
\end{grammatikbox}

\begin{grammatikbox}{语法点 2：媒体批判常用习语与 Unit 3 思辨句法联动}
把本课批判思维表达与前面几课的武器结合起来，就能写出真正 B2- 水平的 \textit{Debatindlæg}：
\begin{enumerate}
  \item \textbf{\da{at tage noget for gode varer}（不加思考地轻信某事）}：极地道的丹麦语惯用语（多用于否定或警示句）：\\
    \textit{Vi må ikke bare \da{tage} opslag på sociale medier \da{for gode varer}.}
  \item \textbf{\da{at skille fakta fra holdninger}（把事实与观点剥离开来）}：\\
    \textit{En god kildekritik hjælper os med at \da{skille fakta fra holdninger}.}
  \item \textbf{联动 Unit 3 前几课句法升级论证}：\\
    \ding{172} 结合 L11 名词化：\textit{\dabold{Ved efterprøvelse af} kilderne undgår man at sprede usande påstande.}\\
    \ding{173} 结合 L12 前提条件：\textit{Sociale medier informerer hurtigt, \dabold{forudsat at} vi forholder os kritisk til indholdet.}
\end{enumerate}
\end{grammatikbox}

\begin{grammatikbox}{英丹/中丹对比避坑指南（Stilistik \& Pitfalls）}
\begin{enumerate}
  \item \textbf{\textit{sætte spørgsmålstegn} 后面只能接 \dabold{ved}}：英语说 \textit{put a question mark on / question something}，中文说“对……画问号”，丹麦语介词固定且唯一是 \dabold{ved}：\ding{55}\ \textit{*sætte spørgsmålstegn på / til påstanden} $\rightarrow$ \ding{51}\ \textit{sætte spørgsmålstegn \da{ved} påstanden}。后面接“是否……”时，必须用 \da{om} 或 \da{hvorvidt}，严禁写成 \textit{*ved hvis}！
  \item \textbf{千万别漏掉反身代词 \dabold{sig}（或 \textit{mig / os}）}：丹麦语中“持某种态度面对某事”是反身结构 \textit{at forholde \dabold{sig} [副词] til}：\ding{55}\ \textit{*Jeg forholder kritisk til medierne} $\rightarrow$ \ding{51}\ \textit{Jeg forholder \dabold{mig} kritisk til medierne}。注意此处的 \textit{kritisk} 作副词修饰动作，因此遇到复数主语时（如 \textit{Vi forholder os \dabold{kritisk} til...}）通常保持基本形 \textit{kritisk}（不加 \textit{-e}）！
  \item \textbf{\textit{nuancere} 不是“中立/和稀泥”}：它表示“展现事物的复杂性与多面性，避免单向偏激”，相当于英文的 \textit{add nuance to}，而非“削弱立场”。
\end{enumerate}
\end{grammatikbox}

% ---------------------------------------------------------------------
\ncdsubsec{五、Mønsterøvelser（句型与修辞操练）}

\textbf{练习 A：介词与批判性思维核心词选填（从 \textit{ved, til, sig, fra, for, nuancere, fordomme, efterprøve} 中选出最恰当的词）}
\begin{enumerate}
  \item Læserne bør altid sætte spørgsmålstegn \underline{\hspace{2.2cm}} anonyme påstande på sociale medier.
  \item Som borger i et demokrati skal man lære at forholde \underline{\hspace{2.2cm}} kritisk \underline{\hspace{2.2cm}} nyhedsdækningen.
  \item En grundig baggrundsartikel kan være med til at \underline{\hspace{2.2cm}} debatten om kunstig intelligens.
  \item Man må aldrig tage dramatiske overskrifter \underline{\hspace{2.2cm}} gode varer uden at tjekke kilden.
  \item Algoritmer viser os ofte historier, der blot bekræfter vores egne \underline{\hspace{2.2cm}}.
  \item Kildekritik handler blandt andet om at skille dokumenterede fakta \underline{\hspace{2.2cm}} personlige holdninger.
  \item Dygtige journalister bruger tid på at \underline{\hspace{2.2cm}} oplysninger hos flere uafhængige kilder.
  \item Jeg sætter spørgsmålstegn \underline{\hspace{2.2cm}}, hvorvidt korte videoklip kan erstatte rigtige aviser.
\end{enumerate}

\textbf{练习 B：思辨表达升级手术（用本课批判性思维表达改写下列口语化简单句，或修正语法错误）}
\begin{enumerate}
  \item 用 \textit{sætte spørgsmålstegn ved, om...} 升级：\textit{Jeg tror ikke på, at denne video viser hele sandheden.} $\rightarrow$
  \item 用 \textit{forholde sig kritisk til} 升级：\textit{Unge bør ikke bare tro på alt, hvad de ser på deres feed.} $\rightarrow$
  \item 用 \textit{nuancere debatten} 升级：\textit{Vi har brug for artikler, som viser både fordele og ulemper i diskussionen.} $\rightarrow$
  \item 用 \textit{tage noget for gode varer} 改写：\textit{Du må ikke acceptere en overskrift uden at tjekke den først.} $\rightarrow$
  \item 纠错改写（漏了反身代词与介词错误）：\textit{Mange borgere forholder kritisk om de hurtige nyheder i ekkokamre.} $\rightarrow$
  \item 纠错改写（介词错误与误用 \textit{hvis}）：\textit{Eksperten sætter spørgsmålstegn på, hvis algoritmer kan være helt neutrale.} $\rightarrow$
\end{enumerate}

\textbf{练习 C：中译丹（综合调用批判性思维表达与新闻词汇）}
\begin{enumerate}
  \item 我们应当始终对网上来历不明的说法提出质疑。（用 \textit{sætte spørgsmålstegn ved}）
  \item 年轻人需要学会对社交媒体上的新闻保持批判态度。（用 \textit{forholde sig kritisk til}）
  \item 独立的报道有助于让公共讨论变得更有层次。（用 \textit{nuancere debatten}）
  \item 在核查信源之前，绝不能轻信耸人听闻的标题。（用 \textit{tage ... for gode varer}）
  \item 算法往往会强化我们自己的偏见，并制造信息茧房。
  \item 学会区分事实与个人观点是现代公民的重要能力。
\end{enumerate}

% ---------------------------------------------------------------------
\ncdsubsec{六、DU3 Modul 4 实战微任务与参考答案}

\begin{emailbox}{DU3 Modul 4 微型写作任务（对标 Skrivning Opgave 2 · Debatindlæg）}
\textbf{Situation}: På et ungdomsmagasins hjemmeside skriver bloggeren Freja:\\
\textit{”Jeg får alle mine nyheder på fem minutter gennem korte videoer på sociale medier. Traditionelle aviser og tv-nyheder er langsomme, gammeldags og helt unødvendige i dag.”}\\[4pt]
\textbf{Opgave}: Skriv et debatindlæg på \textbf{mindst 100 ord} (anbefalet 130--160 ord), hvor du tager stilling til Frejas påstand. Forklar både muligheder og farer ved nyheder på sociale medier, og brug mindst \textbf{tre} udtryk for kritisk tænkning fra denne lektion (\textit{sætte spørgsmålstegn ved, forholde sig kritisk til, nuancere debatten, skille fakta fra holdninger, tage noget for gode varer}).
\end{emailbox}

\begin{answerbox}
\textbf{Før du læser 参考回答}
\begin{enumerate}
  \item \textit{Ekkokamre opstår, fordi algoritmer registrerer vores klik og derefter viser mere indhold, som bekræfter vores egne fordomme. Det skader demokratiet, fordi borgerne ikke længere møder modargumenter og har sværere ved at skille fakta fra holdninger.}
  \item \textit{Sara foreslår, at man stopper op og sætter spørgsmålstegn ved utrolige historier, tjekker mindst to uafhængige nyhedsmedier samt læser analyser, der nuancerer debatten.}
\end{enumerate}

\textbf{Tekst B 阅读核对答案}
\begin{enumerate}
  \item \da{Forkert}（相反，当一条文字让人立刻感到愤怒时，她建议人要格外保持怀疑）
  \item \da{Mindst to uafhængige nyhedsmedier.}（至少核查两家独立的新闻媒体）
  \item \da{Fordi analyser forsøger at nuancere debatten.}（因为分析文章试图让讨论更具层次感，而不是只一味附和你的观点）
\end{enumerate}

\textbf{练习 A 答案}
\begin{enumerate}
  \item \da{ved}（固定介词：\textit{sætte spørgsmålstegn ved}）
  \item \da{sig}；\da{til}（反身代词 \textit{sig} 呼应泛指主语 \textit{man}，介词固定为 \textit{til}）
  \item \da{nuancere}（\textit{at nuancere debatten} 使讨论更全面深入）
  \item \da{for}（习语：\textit{tage noget for gode varer} 轻信）
  \item \da{fordomme}（复数形式双写 \textit{m}：\textit{egne fordomme}）
  \item \da{fra}（固定搭配：\textit{skille A fra B} 把 A 从 B 中区分开）
  \item \da{efterprøve}（\textit{at efterprøve oplysninger} 核查核实信息）
  \item \da{ved}（接 \textit{hvorvidt} 从句前仍必须保留介词 \textit{ved}！）
\end{enumerate}

\textbf{练习 B 答案（句型升级参考答案）}
\begin{enumerate}
  \item \textit{Jeg \da{sætter spørgsmålstegn ved, om} denne video viser hele sandheden.}（用 \textit{ved, om...} 升级客观质疑）
  \item \textit{Unge bør \da{forholde sig kritisk til} alt, hvad de ser på deres feed.}（准确使用反身代词 \textit{sig} 与介词 \textit{til}）
  \item \textit{Vi har brug for artikler, som \da{nuancerer debatten}.}（用凝练动词 \textit{nuancere} 替代松散描述）
  \item \textit{Du må ikke \da{tage} en overskrift \da{for gode varer} uden at tjekke den først.}（地道习语替换）
  \item \textit{Mange borgere \dabold{forholder sig} kritisk \dabold{til} de hurtige nyheder i ekkokamre.}（补齐反身代词 \textit{sig}，并将错介词 \textit{om} 改为 \textit{til}）
  \item \textit{Eksperten sætter spørgsmålstegn \dabold{ved, hvorvidt} algoritmer kan være helt neutrale.}（介词 \textit{ved} + 表达“是否”必须用 \textit{hvorvidt} 或 \textit{om}，不可用 \textit{hvis}）
\end{enumerate}

\textbf{练习 C 答案}
\begin{enumerate}
  \item \textit{Vi bør altid sætte spørgsmålstegn ved anonyme påstande på nettet.}
  \item \textit{Unge skal lære at forholde sig kritisk til nyheder på sociale medier.}
  \item \textit{Uafhængig journalistik hjælper med at nuancere den offentlige debat.}
  \item \textit{Man må aldrig tage dramatiske overskrifter for gode varer før kildekritik.}
  \item \textit{Algoritmer bekræfter ofte vores egne fordomme og skaber lukkede ekkokamre.}
  \item \textit{At skille fakta fra personlige holdninger er en vigtig evne for borgere.}
\end{enumerate}

\textbf{DU3 Modul 4 高分范文（Debatindlæg · 150 词）}\\
\textit{\textbf{Hurtige nyheder gør os ikke klogere}}

\textit{Freja har jo ret i, at sociale medier giver et hurtigt overblik over dagens begivenheder, og mange unge foretrækker korte videoer frem for lange artikler i en travl hverdag.}

\textit{Alligevel mener jeg, at traditionelle nyhedsmedier stadig er helt uundværlige. På sociale medier styrer algoritmer nemlig, hvad vi ser, og vi risikerer let at havne i ekkokamre, der kun bekræfter vores egne fordomme. Derfor bør man altid \da{forholde sig kritisk til} korte klip og \da{sætte spørgsmålstegn ved}, om de overhovedet fortæller hele sandheden.}

\textit{Uddannede journalister arbejder med at efterprøve kilder og \da{skille fakta fra holdninger}. Hvis vi stopper med at læse dybdegående journalistik, mister vi de analyser, der kan \da{nuancere debatten}. Sociale medier kan altså supplere, men aldrig erstatte troværdige aviser.}

\small\textbf{考官点评}：全篇 150 词（不含标题），完美展示 Unit 3 收官思辨水准！开篇先用情态小品词 \textit{jo} 认可社交媒体“快速概览”的便利，继而指出算法与回音室隐患，连续精准调度 \textit{forholde sig kritisk til}、\textit{sætte spørgsmålstegn ved, om...}、\textit{skille fakta fra holdninger} 与 \textit{nuancere debatten} 四大武器，论证逻辑层层递进、收束利落有力。
\end{answerbox}


% =====================================================================
% Unit 4
% =====================================================================
\ncdunit{4}{Uddannelse, Klima og Verden}{教育、气候与全球视野}{Lektion 16--20：多重归因与相关性 · 政策举措与妥协 · 双刃剑议论文模板 · 设问推进与强调 · 第四册总汇演（Debatindlæg）}

\ncdlesson{16}{Karakterræs og trivselskrisen blandt unge}{分数竞争与青少年心理危机}
{DU3 Modul 4（B1+ $\rightarrow$ B2-）｜对标 NCE 3 \& Modultest 4 \textit{Uddannelse, Unge og Trivsel}}
{\textbf{多重归因与相关性表达}（\textit{hænge nøje sammen med, en kombination af flere faktorer, kan i høj grad tilskrives, spille sammen med}）· “相关性 $\neq$ 单一因果”的思辨论证句式}
{在探讨教育压力、青少年焦虑等复杂的社会现象时，B1+ 与 B2- 最大的区别在于：不再用“因为 A 所以 B”做简单武断的单一归因，而是学会用严谨的书面表达揭示\textbf{多重因素的交互作用}与\textbf{相关性}。掌握本课的归因工具箱，你在 Modul 4 辩论与书信写作中就能写出真正具有专家视角的深度分析！}

% ---------------------------------------------------------------------
\begin{foerduleaeser}
带着下面两个问题阅读 Tekst A 和 Tekst B：
\begin{enumerate}
  \item \textit{Hvorfor kan trivselskrisen blandt unge ifølge Tekst A ikke forklares ud fra én enkelt årsag, og hvilke faktorer spiller sammen?}（根据 Tekst A，为什么青少年心理健康危机不能用单一原因来解释？有哪些因素在相互作用？）
  \item \textit{Hvad er Sofies hovedbekymring i læserbrevet i Tekst B, og hvilket konkret forslag stiller hun til politikerne?}（在 Tekst B 的读者来信中，索菲最担忧的是什么？她向政治家提出了什么具体建议？）
\end{enumerate}
\end{foerduleaeser}

% ---------------------------------------------------------------------
\ncdsubsec{二、Tekster（社会议题专栏与辩论来信双语篇）}

\subsection*{Tekst A：\textit{Karakterræs og trivselskrisen blandt unge}（社会深度专栏 · 约 255 词）}

\begin{tipbox}{语篇说明}
本篇剖析丹麦高中生与大学生的“分数内卷”（\textit{karakterræs}）与心理困扰议题。阅读时请重点留意\textbf{\color{NCDNavy}深蓝色标注的多重归因与相关性骨架}，以及\textbf{\color{NCDRed}深红色标注的限定条件与否定审慎表达}。
\end{tipbox}

I de senere år har trivselskrisen blandt danske unge fyldt stadig mere i den offentlige debat. Flere nationale undersøgelser dokumenterer, at en stigende andel af eleverne i gymnasiet og på ungdomsuddannelserne oplever vedvarende stress, søvnmangel og eksamensangst. Men når forskere og pædagoger analyserer problemet, understreger de næsten samstemmende, at udviklingen \dabold{ikke kan forklares ud fra én enkelt årsag}.

For det første viser mange studier, at den voksende mistrivsel \da{hænger nøje sammen med} et intensiveret karakterræs. Når adgangskravene til populære videregående uddannelser konstant stiger, opfatter mange unge hver eneste aflevering og prøve som afgørende for deres fremtid. \da{Sammenhængen mellem} høje forventninger og psykisk pres bliver særlig tydelig hos de såkaldte ”12-talspiger” og andre pligtopfyldende elever, der sjældent tillader sig selv at begå fejl.

Alligevel er det for simpelt kun at give karaktersystemet skylden. Problemet \da{skyldes snarere en kombination af flere faktorer}, der forstærker hinanden i hverdagen. Det øgede præstationspres i skolen \da{spiller nemlig sammen med} en konstant sammenligning på sociale medier, hvor unge døgnet rundt præsenteres for polerede billeder af succes, krop og fritidsliv. Derudover \da{kan} unges usikkerhed \da{i høj grad tilskrives} en bredere perfektionskultur i hele samfundet.

At to fænomener optræder samtidig, \dabold{betyder imidlertid ikke automatisk, at} det ene direkte forårsager det andet. Eksperter advarer derfor mod hurtige standardløsninger og peger i stedet på, at øget trivsel kræver både færre unødvendige test, mere åbenhed om sårbare følelser og en kultur, hvor der er plads til at fejle.

\begin{transbox}
\small
\textbf{\color{NCDTeal}[Tekst A 中文精译]}\\
近年来，丹麦青少年的心理福祉危机在公共讨论中占据了越来越大的比重。多项全国性调查证实，越来越高比例的高中及中等职业教育学生正经历着持续的压力、睡眠不足与考试焦虑。然而，当学者与教育工作者分析这一问题时，他们几乎一致强调，这一趋势无法仅凭单一原因来解释。

首先，许多研究表明，不断加剧的心理不适与愈演愈烈的分数竞争密切相关。当热门高等教育的录取分数线持续攀升时，许多年轻人会把每一次作业提交和随堂测验都视为决定自己未来的关键节点。高期待与心理压力之间的关联在所谓的“满分（12分）女生”以及其他极度自律守责的学生身上显得尤为明显，他们很少允许自己犯错。

尽管如此，仅仅把责任归咎于评分体系是过于简单化的。这个问题更确切地说是多种因素共同作用的结果，这些因素在日常生活中相互强化。学校里日益增长的绩效压力，实际上正与社交媒体上的持续攀比相互作用——在社交平台上，年轻人全天候面对着经过精心修饰的成功、身材与休闲生活图景。此外，年轻人的不安全感在很大程度上也可以归因于整个社会中更广泛的完美主义文化。

不过，两个现象同时出现并不自动意味着其中一个直接导致了另一个。因此，专家们警告不要采取仓促的一刀切方案，而是指出，提升心理健康需要减少不必要的测试、对脆弱情绪保持更多坦诚，以及一种允许犯错的校园与社会文化。
\end{transbox}

\subsection*{Tekst B：\textit{Læserbrev: ”Vi må tale om mere end bare karakterer”}（报刊读者来信 · 约 135 词）}

\textbf{阅读任务（理解与核对）：}\textit{Læs Sofies læserbrev i Politiken om karakterpres i gymnasiet, og besvar spørgsmålene nedenfor.}

\begin{emailbox}{Læserbrev｜Fra: Sofie Holm (gymnasielærer i Aarhus) \quad Emne: Mindre pres, mere læring}
Som lærer gennem femten år ser jeg hver uge, hvordan dygtige unge brænder ud længe før studentereksamen. Når eleverne mister arbejdsglæden, \da{hænger det} efter min mening \da{direkte sammen med}, at næsten alle samtaler i dag handler om gennemsnit.

\da{En kombination af} høje adgangskrav, bekymrede forældre og manglende hvile gør, at mange unge konstant føler sig utilstrækkelige. Det \da{skyldes sjældent én ting alene}, men et system, der belønner fejlfrihed frem for nysgerrighed. Selvom sociale medier også spiller ind, \dabold{kan man ikke} blot give skærmene hele ansvaret.

Jeg vil derfor opfordre politikerne til at skabe flere alternative optagelsesveje uden et tårnhøjt karakterkrav. Først når læring bliver vigtigere end tal på et papir, kan vi for alvor vende den negative udvikling.
\end{emailbox}

\textbf{Tekst B 阅读核对（Rigtigt / Forkert / Kort svar，答案见第六部分）：}
\begin{enumerate}
  \item \textit{Sofie mener, at sociale medier alene har hele ansvaret for de unges mistrivsel.} \underline{\hspace{2.2cm}}
  \item \textit{Hvilke tre faktorer nævner Sofie som en kombination, der får unge til at føle sig utilstrækkelige?} \underline{\hspace{2.2cm}}
  \item \textit{Hvad opfordrer Sofie politikerne til at gøre i slutningen af sit læserbrev?} \underline{\hspace{2.2cm}}
\end{enumerate}

\begin{tipbox}{本课发音与语调要点（Udtale \& Intonation）}
\begin{itemize}
  \item \textbf{归因与核心抽象词读音}：\textit{trivselskrise} \ipa{ˈtˢʁiʊ̯ˀsəlsˌkʰʁiːsə}（\textit{trivsel} 中的 \textit{v} 滑音化为 \ipa{ʊ̯ˀ}）；\textit{karakterræs} \ipa{kʰɑʁɑɡˈtˢeɐ̯ˀˌʁɛːˀs}（注意主重音在 \textit{-ter-} 上！）；\textit{kombination} \ipa{kʰʌmbinaˈɕoːˀn}（\textit{-tion} 读作 \ipa{ɕoːˀn}）；\textit{faktor} \ipa{ˈfɑɡtˢʌ} / 复数 \textit{faktorer} \ipa{fɑɡˈtˢoːˀɐ}（重音在复数后移！）。
  \item \textbf{静音辅音与连读弱化}：\textit{sammenhæng} \ipa{ˈsɑmənˌhɛŋˀ}（无软 d，尾音为鼻音 \ipa{ŋˀ}）；\textit{tilskrive} \ipa{ˈtˢelˌsgʁiːʊ̯ə}（\textit{v} 在口语中常滑音化）；\textit{årsag} \ipa{ˈɒːˌsɑːˀ}（结尾 \textit{g} 不发音）；\textit{skyldes} \ipa{ˈsgyləs}（\textit{-d-} 不发音！）。
  \item \textbf{否定强调的语调重音}：读 \textit{ikke kan forklares ud fra én enkelt årsag} 时，重音刻意落在数词 \textit{\textbf{én}} \ipa{ˈeːˀn} 和 \textit{\textbf{en}kelt} 上，突出“并非仅仅单一原因”的对比语气。
\end{itemize}
\end{tipbox}

\begin{transbox}
\small
\textbf{\color{NCDTeal}[Tekst B 中文对照]}\\
\textbf{读者来信（致《政治报》｜少一点分数压力，多一点真实学习）：}作为一名执教十五年的高中老师，我每周都能看到优秀的年轻人在高中毕业考之前很久就已精疲力竭。在我看来，当学生们失去学习乐趣时，这与如今几乎所有对话都围绕平均分展开有着直接关联。

高企的录取门槛、焦虑的父母以及缺乏休息，这几者的结合让许多年轻人时刻觉得自己不够好。这很少仅仅源于某一件事情，而是源于一个奖励“零失误”而非“好奇心”的体系。尽管社交媒体也起了一定作用，但人们绝不能把全部责任都推给屏幕。

因此我呼吁政治家们开辟更多不依赖天价分数线的替代性升学路径。只有当学习本身变得比纸面上的数字更重要时，我们才能真正扭转这一负面趋势。
\end{transbox}

% ---------------------------------------------------------------------
\ncdsubsec{三、Gloser（核心生词表）}

\begin{tipbox}{词表重点}
注意区分动词 \textit{at mistrives}（感到不适、陷入抑郁/低迷，与 \textit{at trives} 相反）——它们都是典型的\textbf{异态动词（Deponent verber）}，词尾永远带 \textit{-s}，但表达主动含义。另外，\textit{en faktor} 的复数重音后移且带 \textit{stød}：\textit{faktorer}。
\end{tipbox}

\begin{glosertable}
名词 & et karakterræs & \textit{(karakterræset, ingen flertal)} & 分数竞争，唯分数内卷 \\
名词 & en trivsel & \textit{(trivselen, ingen flertal)} & 身心健康，幸福感，适应状态 \\
名词 & en mistrivsel & \textit{(mistrivselen, ingen flertal)} & 心理不适，低迷抑郁状态 \\
名词 & en faktor & \textit{(faktoren, faktorer, faktorerne)} & 因素，要素 \\
名词 & en årsag & \textit{(årsagen, årsager, årsagerne)} & 原因，起因 \\
名词 & en sammenhæng & \textit{(sammenhængen, sammenhænge, -ene)} & 关联，相关性；上下文 \\
名词 & et samspil & \textit{(samspillet, samspil, samspillene)} & 相互作用，联动配合 \\
名词 & et adgangskrav & \textit{(adgangskravet, adgangskrav, -ene)} & 录取门槛，准入要求 \\
名词 & et præstationspres & \textit{(præstationspresset, ingen flertal)} & 绩效压力，表现压力 \\
名词 & en søvnmangel & \textit{(søvnmanglen, ingen flertal)} & 睡眠不足 \\
名词 & et gennemsnit & \textit{(gennemsnittet, gennemsnit, -tene)} & 平均分，平均值 \\
\midrule
动词 [har] & at hænge sammen & \textit{(hænger, hang, har hængt sammen)} & 相关联，有内在联系（接 \textit{med}） \\
动词 [har] & at tilskrive & \textit{(tilskriver, tilskrev, har tilskrevet)} & 把……归因于 \\
动词 [har] & at forårsage & \textit{(forårsager, forårsagede, har forårsaget)} & 导致，引发（直接因果） \\
动词 [har] & at præstere & \textit{(præsterer, præsterede, har præsteret)} & 表现，做出成绩 \\
异态动词 & at mistrives & \textit{(mistrives, mistrivedes, har mistrivedes)} & 身心不适，不适应 \\
\midrule
形容词 & afgørende & \textit{(afgørende, 不变)} & 决定性的，至关重要的 \\
\end{glosertable}

% ---------------------------------------------------------------------
\ncdsubsec{四、Grammatik \& Stilistik（核心语法、论证结构与修辞精讲）}

\begin{grammatikbox}{语法点 1：拒绝“单一归因”——多重归因四大黄金句型骨架}
在 DU3 Modul 4 思辨写作与口试辩论中，遇到复杂社会现象时，最忌讳用单一的 \textit{Det er kun fordi...} 下死结论。以下四组骨架能瞬间提升论证的严谨度与客观性：

\vspace{0.3em}
\renewcommand{\arraystretch}{1.24}
{\small
\begin{tabularx}{\linewidth}{@{} >{\bfseries\raggedright\arraybackslash}p{3.0cm} >{\raggedright\arraybackslash}p{4.3cm} >{\raggedright\arraybackslash}X @{}}
\toprule
\textbf{论证功能} & \textbf{核心句型公式} & \textbf{标准书面示范} \\
\midrule
\textbf{1. 否定单一原因} & \textit{kan ikke forklares ud fra én enkelt årsag} / \textit{skyldes sjældent én ting alene} & \textit{Udviklingen \dabold{kan ikke forklares ud fra én enkelt årsag}.} \\
\midrule
\textbf{2. 多因素组合} & \textit{skyldes en kombination af A, B og C} & \textit{Krisen skyldes \da{en kombination af flere faktorer}, herunder søvnmangel og pres.} \\
\midrule
\textbf{3. 相互联动作用} & \textit{A spiller sammen med B} / \textit{i samspil med} & \textit{Skolepresset \da{spiller sammen med} en konstant sammenligning på nettet.} \\
\midrule
\textbf{4. 高阶被动归因} & \textit{A kan i høj grad tilskrives B}（A 在很大程度上可归因于 B） & \textit{Unges usikkerhed \da{kan i høj grad tilskrives} en bredere perfektionskultur.} \\
\bottomrule
\end{tabularx}
}
\end{grammatikbox}

\begin{grammatikbox}{语法点 2：相关性（\textit{Sammenhæng}）vs 直接因果（\textit{Årsag}）的精密区分}
英语和中文母语者写议论文时，常把“相关”直接写成“导致”。Modul 4 考官特别看重你能否区分两者的逻辑强度：
\begin{enumerate}
  \item \textbf{表达强相关性：\da{at hænge nøje sammen med} + 名词/从句}：
    \begin{itemize}
      \item 动词短语拆合规则：\textit{hænge sammen} 是可分动词短语！主句现在时中，主语 + \dabold{hænger} + 程度副词（\textit{nøje / tæt / direkte}）+ \dabold{sammen med} + 对象。例：\textit{Mistrivsel \dabold{hænger nøje sammen med} et højt pres.}
      \item 若在从句中带否定词 \textit{ikke}，否定词依然置于变位动词之前：\textit{... at det \dabold{ikke hænger} direkte sammen med skærmtid.}
      \item 名词化进阶（复现第 11 课名词化技巧）：\da{sammenhængen mellem A og B}（A 与 B 之间的关联）。例：\textit{Forskere undersøger \da{sammenhængen mellem} søvnmangel og stress.}
    \end{itemize}
  \item \textbf{审慎界定“相关 $\neq$ 因果”}：使用固定转折框架：\textit{At A og B optræder samtidig, \dabold{betyder imidlertid ikke automatisk, at} A forårsager B.}（A 与 B 同时出现，然而并不自动意味着 A 导致了 B。）其中 \textit{imidlertid}（然而）置于变位动词 \textit{betyder} 之后，是极为地道的 B2- 书面语序！
\end{enumerate}
\end{grammatikbox}

\begin{grammatikbox}{英丹对比排雷 \& 辩论修辞搭配（Stilistik）}
\begin{enumerate}
  \item \textbf{千万别受英语 \textit{reason / cause} 混用误导介词}：丹麦语中，“……的原因”固定用 \da{årsagen til}（而不是 \textit{*årsagen af / for}！）：\ding{51}\ \textit{Hvad er årsagen \dabold{til} krisen?}；“与……相关”必须用 \da{sammenhæng med / mellem}，绝不能写 \ding{55}\ \textit{*hænger sammen til}。
  \item \textbf{\textit{skyldes} 本身就是被动含义，绝不能再用 \textit{blive + skyldes}！}：很多同学模仿英语 \textit{is caused by} 误写成 \ding{55}\ \textit{*Problemet bliver skyldes...}。记住 \da{skyldes}（源于、归因于）后面直接接名词原因：\ding{51}\ \textit{Problemet \da{skyldes} flere faktorer.}
  \item \textbf{记住 \textit{tilskrive} 的两种语态方向}：主动为“专家把 A 归因于 B”（\textit{Eksperter tilskriver presset karakterræset}），而议论文中更常用被动“A 可归因于 B”（\textit{Presset kan \da{tilskrives} karakterræset}，施动方向切勿写反！）。
\end{enumerate}

\textbf{本课辩论高频修辞搭配（Faste vendinger）：}
\begin{itemize}
  \item \da{at give noget skylden (for...)}（把……的责任归咎于某事物）
  \item \da{at sætte fokus på trivsel}（聚焦心理健康）/ \da{at vende den negative udvikling}（扭转负面趋势）
  \item \da{at brænde ud}（精疲力竭、职业/学业倦怠）/ \da{at føle sig utilstrækkelig}（觉得自己不够格/不够好）
\end{itemize}
\end{grammatikbox}

% ---------------------------------------------------------------------
\ncdsubsec{五、Mønsterøvelser（句型与辩论操练）}

\textbf{练习 A：多重归因与相关性核心结构填空（每空填入最准确的 1--2 个丹麦语词）}
\begin{enumerate}
  \item Mange psykologer mener, at unges stress \underline{\hspace{2.2cm}} nøje sammen \underline{\hspace{2.2cm}} et stigende karakterpres.
  \item Den negative udvikling kan ikke forklares ud fra én enkelt \underline{\hspace{2.2cm}}.
  \item Trivselskrisen skyldes snarere en \underline{\hspace{2.2cm}} af flere forskellige faktorer.
  \item En stor del af de unges søvnmangel kan i høj grad \underline{\hspace{2.2cm}} sene timer foran skærmen.
  \item En ny rapport undersøger \underline{\hspace{2.2cm}} mellem høje adgangskrav og eksamensangst i gymnasiet.
  \item Sociale medier står ikke alene, men \underline{\hspace{2.2cm}} ofte sammen med høje forventninger fra både skole og hjem.
  \item At to tendenser ses samtidig, betyder \underline{\hspace{2.2cm}} ikke automatisk, at den ene forårsager den anden.
\end{enumerate}

\textbf{练习 B：语体升级与病句改错（将单一武断句升级为严谨的多重归因句，或修正介词/语态错误）}
\begin{enumerate}
  \item 【纠错改写】\textit{Problemet bliver skyldes for meget pres i skolen.} $\rightarrow$
  \item 【纠错改写】\textit{Hvad er den største årsag for den stigende stress blandt unge?} $\rightarrow$
  \item 【纠错改写】\textit{Forskerne siger, at søvnmangel hænger tæt sammen til sociale medier.} $\rightarrow$
  \item 【升级改写】用 \textit{kan ikke forklares ud fra én enkelt årsag} 升级：\textit{Det er ikke kun én ting, der skaber krisen.} $\rightarrow$
  \item 【升级改写】用 \textit{en kombination af} 合并两句：\textit{Stressen kommer fra karakterpres. Den kommer også fra søvnmangel.} $\rightarrow$ \textit{Stressen skyldes...}
  \item 【升级改写】用 \textit{kan i høj grad tilskrives} 改写被动归因：\textit{Perfektionskulturen forklarer en stor del af presset.} $\rightarrow$ \textit{En stor del af presset...}
\end{enumerate}

\textbf{练习 C：中译丹（活用多重归因、相关性表达与否定审慎句式）}
\begin{enumerate}
  \item 青少年心理不适与不断提高的录取要求密切相关。
  \item 这一复杂的问题不能仅仅从单一原因来解释。
  \item 许多学生的压力是由多种因素共同造成的。（用 \textit{en kombination af}）
  \item 年轻人的不安全感在很大程度上可以归因于完美主义文化。
  \item 两个现象同时出现并不意味着其中一个直接导致了另一个。
\end{enumerate}

% ---------------------------------------------------------------------
\ncdsubsec{六、DU3 Modul 4 实战微任务与参考答案}

\begin{emailbox}{DU3 Modul 4 Skrivning Opgave 1 实战任务（私人长邮件 · mindst 125 ord）}
\textbf{Situation}: Du har fået en mail fra din ven, Kasper, hvis lillesøster går i 2. g og er meget stresset over sine karakterer. Kasper spørger dig, hvorfor så mange unge i dag føler et stort pres, og hvad man kan gøre for at hjælpe dem.\\
\textbf{Opgave}: Skriv et svar til Kasper på \textbf{135--170 ord} (mindst 125 ord). Du skal:
\begin{itemize}
  \item vise forståelse for Kaspers bekymring (brug gerne et modalpartikel som \textit{jo} eller \textit{egentlig})
  \item forklare, at mistrivsel skyldes et samspil mellem flere faktorer (brug mindst to udtryk fra denne lektion)
  \item give et konkret råd til, hvordan familie og skole kan mindske presset.
\end{itemize}
\end{emailbox}

\begin{answerbox}
\textbf{Før du læser 参考回答}
\begin{enumerate}
  \item \textit{Ifølge Tekst A kan trivselskrisen ikke forklares ud fra én enkelt årsag, fordi den skyldes en kombination af flere faktorer: et intensiveret karakterræs med høje adgangskrav, konstant sammenligning på sociale medier og en bredere perfektionskultur i samfundet.}
  \item \textit{Sofie er bekymret over, at næsten alle samtaler i gymnasiet handler om gennemsnit, hvilket får eleverne til at miste arbejdsglæden. Hun opfordrer politikerne til at skabe flere alternative optagelsesveje uden tårnhøje karakterkrav.}
\end{enumerate}

\textbf{Tekst B 阅读核对答案}
\begin{enumerate}
  \item \da{Forkert}（她明确指出：\textit{Selvom sociale medier også spiller ind, kan man ikke blot give skærmene hele ansvaret.}）
  \item \da{Høje adgangskrav, bekymrede forældre og manglende hvile.}
  \item \da{At skabe flere alternative optagelsesveje uden et tårnhøjt karakterkrav.}
\end{enumerate}

\textbf{练习 A 答案}
\begin{enumerate}
  \item \da{hænger}；\da{med}（固定动词短语 \textit{hænge nøje sammen med}，主句现在时动词变位为 \textit{hænger}）
  \item \da{årsag}（\textit{én enkelt årsag} 单一原因）
  \item \da{kombination}（\textit{en kombination af flere faktorer} 多种因素的组合）
  \item \da{tilskrives}（情态动词 \textit{kan} 后接 \textit{-s} 被动原形：\textit{kan tilskrives}）
  \item \da{sammenhængen}（特指两者之间的关联，用定指形式 \textit{sammenhængen mellem A og B}）
  \item \da{spiller}（\textit{spille sammen med} 与……相互作用，主语为复数仍用现在时 \textit{spiller}）
  \item \da{imidlertid}（书面转折副词，置于变位动词 \textit{betyder} 之后）
\end{enumerate}

\textbf{练习 B 答案}
\begin{enumerate}
  \item \textit{Problemet skyldes for meget pres i skolen.}（\textit{skyldes} 本身已是被动含义，绝不可与 \textit{bliver} 连用）
  \item \textit{Hvad er den største årsag til den stigende stress blandt unge?}（名词 \textit{årsag} 后固定搭配介词 \textit{til}）
  \item \textit{Forskerne siger, at søvnmangel hænger tæt sammen med sociale medier.}（固定搭配是 \textit{hænge sammen med}，不能用 \textit{til}）
  \item \textit{Krisen kan ikke forklares ud fra én enkelt årsag.}（用更紧凑凝练的被动归因短语取代松散分裂句）
  \item \textit{Stressen skyldes en kombination af karakterpres og søvnmangel.}（用 \textit{skyldes en kombination af A og B} 合并多重原因）
  \item \textit{En stor del af presset kan i høj grad tilskrives perfektionskulturen.}（利用被动归因凸显主体对象）
\end{enumerate}

\textbf{练习 C 答案}
\begin{enumerate}
  \item \textit{Unge menneskers mistrivsel hænger nøje sammen med stadig højere adgangskrav.}
  \item \textit{Dette komplekse problem kan ikke forklares ud fra én enkelt årsag.}
  \item \textit{Mange elevers stress skyldes en kombination af flere forskellige faktorer.}
  \item \textit{Unges usikkerhed kan i høj grad tilskrives en udbredt perfektionskultur.}
  \item \textit{At to fænomener optræder samtidig, betyder ikke, at det ene forårsager det andet.}
\end{enumerate}

\textbf{DU3 Modul 4 Skrivning Opgave 1 高分范文（150 词，不含署名）}\\
\textit{Til: Kasper \quad|\quad Emne: Sv: Din lillesøster og karakterpresset}

\textit{Hej Kasper}

\textit{Tak for din mail. Jeg kan virkelig godt forstå, at du er bekymret for din lillesøster. Det er jo desværre blevet ret almindeligt, at unge i gymnasiet føler sig alt for stressede.}

\textit{Når man ser på trivselskrisen blandt unge, kan udviklingen efter min mening ikke forklares ud fra én enkelt årsag. Den øgede mistrivsel \da{hænger} ganske vist \da{nøje sammen med} det hårde karakterræs og de høje adgangskrav til universiteterne. Men problemet \da{skyldes} snarere \da{en kombination af flere faktorer}. Skolepresset \da{spiller nemlig sammen med} sociale medier, hvor man hele tiden sammenligner sig med andre. Derudover \da{kan} presset \da{i høj grad tilskrives} en kultur, hvor der sjældent er plads til at begå fejl.}

\textit{Jeg synes egentlig, at I som familie kan hjælpe hende ved at tale åbent om, at karakterer ikke definerer hendes værdi. Samtidig bør skolerne mindske antallet af test og skabe mere ro i hverdagen.}

\textit{Skriv endelig, hvis du vil vende det over en kop kaffe!}

\textit{De bedste hilsner\\
\mbox{[Dit navn]}}

\small\textbf{考官点评}：全信 150 词（不含署名），完美兼顾私人邮件的温暖语气（巧用情态小品词 \textit{jo} 与 \textit{egentlig}）和 B1+/B2- 的客观多维归因（自然融入 \textit{kan ikke forklares ud fra én enkelt årsag}、\textit{hænger nøje sammen med}、\textit{en kombination af flere faktorer}、\textit{spiller sammen med} 及 \textit{kan i høj grad tilskrives}），让步与建议层次分明，稳居 Modul 4 高分档。
\end{answerbox}

\ncdlesson{17}{Klimamål og landbrugets grønne omstilling}{气候目标与农业绿色转型}
{DU3 Modul 4（B1+ $\rightarrow$ B2-）｜对标 NCE 3 \& Modultest 4 \textit{Klima, Politik og Landbrug}}
{\textbf{政策举措与妥协词汇}（\textit{iværksætte tiltag, stramme lovgivningen, indgå et bredt forlig, pålægge en afgift, yde økonomisk støtte, opnå et kompromis}）· 丹麦“共识政治”与政论条件让步句式}
{气候危机与农业转型是丹麦近年来最激烈的公共辩论核心，也是 DU3 Modul 4 阅读与辩论的高频政治议题。在丹麦政治文化中，重大改革极少靠一方单边强推，而是靠各党派与利益方\textbf{“达成广泛协议（\textit{indgå et bredt forlig}）”}并\textbf{“达成妥协（\textit{opnå et kompromis}）”}。掌握本课的政策动宾强搭配，能让你的政论与读者来信立刻脱离口语化，跃升至标准报刊社论水准！}

% ---------------------------------------------------------------------
\begin{foerduleaeser}
带着下面两个问题阅读 Tekst A 和 Tekst B：
\begin{enumerate}
  \item \textit{Hvorfor er landbrugets grønne omstilling et politisk dilemma i Danmark ifølge Tekst A, og hvilke virkemidler indgår i et bredt forlig?}（根据 Tekst A，为什么农业绿色转型在丹麦是一个政治两难议题？广泛协议中通常包含哪些政策手段？）
  \item \textit{Hvilke krav stiller landmanden Henrik i Tekst B, før han mener, at de nye klimatiltag kan fungere i praksis?}（在 Tekst B 中，农民亨里克对新的气候举措提出了哪些前提条件，认为满足后才能在实践中行得通？）
\end{enumerate}
\end{foerduleaeser}

% ---------------------------------------------------------------------
\ncdsubsec{二、Tekster（社会议题专栏与辩论来信双语篇）}

\subsection*{Tekst A：\textit{Klimamål og landbrugets grønne omstilling}（时政深度专栏 · 约 265 词）}

\begin{tipbox}{语篇说明}
本篇聚焦丹麦减排目标与传统农业大省之间的利益博弈。阅读时请重点留意\textbf{\color{NCDNavy}深蓝色标注的政策举措与政治妥协动宾搭配}，以及\textbf{\color{NCDRed}深红色标注的限制前提与被动政策句式}。
\end{tipbox}

Danmark har sat et ambitiøst mål om at reducere sin samlede udledning af drivhusgasser markant frem mod 2030 og blive klimaneutral i 2045. Hvis de mål skal nås, er alle eksperter enige om, at landbruget spiller en helt central rolle, da erhvervet i dag står for en stor del af landets udledninger. Netop derfor er debatten om, hvordan man skal \da{iværksætte tiltag} på området, blevet et af de sværeste politiske dilemmaer.

På den ene side kræver miljøorganisationer, at Folketinget hurtigt skal \da{stramme lovgivningen} og \da{pålægge} landbruget \da{en CO2-afgift}. De peger på, at frivillige aftaler tidligere ikke har sænket udledningen hurtigt nok, og at klare regler er nødvendige for at beskytte både klimaet og vandmiljøet i de danske fjorde.

På den anden side advarer landmænd og erhvervsorganisationer om, at for hårde krav kan true arbejdspladser i landdistrikterne og tvinge mindre gårde til at lukke. Hvis produktionen blot flytter til udlandet, hvor miljøkravene er lavere, gavner det jo slet ikke det globale klima. Derfor har regeringen, landbruget og grønne organisationer gennem langvarige forhandlinger arbejdet på at \da{indgå et bredt forlig}, så alle parter kan \da{opnå et kompromis}, der holder i mange år.

En langtidsholdbar løsning kræver nemlig, at man samtidig \da{yder økonomisk støtte} til nye teknologier og etablerer en fair overgangsordning. \dabold{Forudsat at} omstillingen ledsages af rådgivning og investeringer i grøn forskning, \dabold{kan} både klimaet og et konkurrencedygtigt landbrug \dabold{sikres} i fremtiden.

\begin{transbox}
\small
\textbf{\color{NCDTeal}[Tekst A 中文精译]}\\
丹麦制定了一项雄心勃勃的目标，即到 2030 年大幅削减温室气体总排放量，并在 2045 年实现气候中和。如果这些目标要得以实现，所有专家都一致认为农业扮演着绝对核心的角色，因为该行业目前占了全国相当大一部分排放量。正因如此，关于应当如何在该领域启动政策举措的争论，已成为最棘手的政治两难议题之一。

一方面，环保组织要求丹麦议会迅速收紧立法，并对农业征收二氧化碳排放税。他们指出，过去的自愿协议未能足够快地降低排放，因此清晰的法规对于保护气候以及丹麦峡湾的水域生态都是必不可少的。

另一方面，农民和行业协会警告称，过于严苛的要求可能会威胁乡村地区的就业岗位，并迫使规模较小的农场倒闭。如果生产只是转移到环保标准更低的国外去，那对全球气候岂不是毫无益处。因此，政府、农业界与环保组织通过漫长的谈判努力达成广泛政治协议，好让各方能够取得一个长期有效的妥协方案。

一个能够经得起时间考验的解决方案，实际上要求政府同时为新技术提供财政支持，并设立公平的过渡期安排。只要转型过程伴随着专业指导以及对绿色科研的投入，气候目标与具有竞争力的农业就能在未来同时得到保障。
\end{transbox}

\subsection*{Tekst B：\textit{Debatindlæg: ”Vi vil gerne omstille, men på fair vilkår”}（农业从业者投稿 · 约 135 词）}

\textbf{阅读任务（理解与核对）：}\textit{Læs landmanden Henriks indlæg om den grønne trepart og lovgivning, og besvar spørgsmålene nedenfor.}

\begin{emailbox}{Debatindlæg｜Fra: Henrik Dalgaard (landmand i Midtjylland) \quad Emne: Et retfærdigt klimaforlig}
Mange i byerne tror, at vi landmænd er imod den grønne omstilling. Det passer ikke. Vi mærker jo selv tørke og ekstremt regnvejr på markerne hvert år. Derfor er jeg sådan set glad for, at man nu vil \da{indgå et bredt forlig} i stedet for at lave kortsigtede symbollove.

Men politikerne kan ikke bare \da{stramme lovgivningen} fra den ene dag til den anden uden en plan for økonomien. Hvis staten vil \da{pålægge} os \da{en ny afgift}, må den også \da{yde økonomisk støtte} til dyre investeringer i ny teknik og skovrejsning. Ellers vil mange unge familier ikke turde overtage en gård.

Vi kan sagtens \da{opnå et kompromis}, \dabold{medmindre} man glemmer virkeligheden på landet. Giv os klare rammer og en reel overgangsordning, så skal vi nok levere resultater.
\end{emailbox}

\textbf{Tekst B 阅读核对（Rigtigt / Forkert / Kort svar，答案见第六部分）：}
\begin{enumerate}
  \item \textit{Henrik er imod, at man indgår et bredt forlig om landbrugets grønne omstilling.} \underline{\hspace{2.2cm}}
  \item \textit{Hvad skal staten gøre, hvis den vil pålægge landmændene en ny afgift ifølge Henrik?} \underline{\hspace{2.2cm}}
  \item \textit{Hvad frygter Henrik vil ske for unge familier på landet, hvis der ikke kommer støtte?} \underline{\hspace{2.2cm}}
\end{enumerate}

\begin{tipbox}{本课发音与语调要点（Udtale \& Intonation）}
\begin{itemize}
  \item \textbf{政策高频动词的喉塞音与静音字母}：\textit{iværksætte} \ipa{iˈvɛɐ̯ɡˌsɛdə}（重音在第二音节 \textit{-værk-}！）；\textit{indgå} \ipa{ˈenˌɡɔːˀ}（\textit{ind} 的 \textit{d} 不发音！）；\textit{forlig} \ipa{fʌˈliːˀ}（\textit{-ig} 中的 \textit{g} 完全不发音，长元音带 \textit{stød}）；\textit{pålægge} \ipa{ˈpʰʌˌlɛɡə}；\textit{yde} \ipa{ˈyːðə}（软 d）。
  \item \textbf{政治与妥协词汇读音}：\textit{lovgivning} \ipa{ˈlʌʊ̯ˌɡiʊ̯ˀneŋ}（两个 \textit{v} 分别滑音化为 \ipa{ʌʊ̯} 与 \ipa{iʊ̯ˀ}！）；\textit{kompromis} \ipa{kʰʌmpʰʁoˈmi}（注意词尾 \textit{s} 在标准丹麦语中\textbf{通常不发音}！或仅轻微吐气）；\textit{afgift} \ipa{ˈɑʊ̯ˌɡifd}（\textit{af} 读 \ipa{ɑʊ̯}）；\textit{overgangsordning} \ipa{ˈʌʊ̯ɐɡɑŋsˌɒːˀdneŋ}。
  \item \textbf{分清 \textit{landmand} 的读音减音}：单数 \textit{en landmand} \ipa{ˈlanˌmanˀ}（两个 \textit{d} 均不发音！）；复数 \textit{landmænd} \ipa{ˈlanˌmɛnˀ}。
\end{itemize}
\end{tipbox}

\begin{transbox}
\small
\textbf{\color{NCDTeal}[Tekst B 中文对照]}\\
\textbf{读者投稿（中日德兰农民亨里克｜一项公正的气候协议）：}城里很多人以为我们农民反对绿色转型，其实并非如此。我们自己每年在田间地头都能切身感受到干旱与极端暴雨。因此，我其实很高兴现在大家愿意达成一项广泛政治协议，而不是制定短视的作秀式法律。

但政治家们不能脱离经济现实，一夜之间就收紧法规。如果国家要向我们征收新的税费，就也必须为昂贵的新技术投资和造林项目提供财政补贴。否则，许多年轻家庭将不敢接手农场。

我们完全可以达成妥协，前提是别忘记乡村的实际情况。只要给我们清晰的制度框架和切实的过渡期安排，我们一定会交出成果。
\end{transbox}

% ---------------------------------------------------------------------
\ncdsubsec{三、Gloser（核心生词表）}

\begin{tipbox}{词表重点}
丹麦语政治与政策词汇的\textbf{动词—名词搭配极其固定}：动词 \textit{at indgå}（签订、达成）不规则变化为 \textit{indgik, har indgået}；\textit{at pålægge}（课以、征收）不规则变化为 \textit{pålægger, pålagde, har pålagt}；而表达“提供经济支持”时，高分书面语常用 \textit{at yde økonomisk støtte}，比口语化的 \textit{give penge} 正式得多！
\end{tipbox}

\begin{glosertable}
名词 & et tiltag & \textit{(tiltaget, tiltag, tiltagene)} & 举措，措施（单复同形） \\
名词 & en lovgivning & \textit{(lovgivningen, lovgivninger, -erne)} & 立法，法规 \\
名词 & et forlig & \textit{(forliget, forlig, forligene)} & 政治协议，跨党派和解协议 \\
名词 & et kompromis & \textit{(kompromiset, kompromiser, -ene)} & 妥协，折中方案 \\
名词 & en afgift & \textit{(afgiften, afgifter, afgifterne)} & 税费，专项附加税 \\
名词 & en udledning & \textit{(udledningen, udledninger, -erne)} & 排放（量） \\
名词 & et landbrug & \textit{(landbruget, landbrug, landbrugene)} & 农业；农场 \\
名词 & en landmand & \textit{(landmanden, landmænd, landmændene)} & 农民，农场主 \\
名词 & en overgangsordning & \textit{(overgangsordningen, -er, -erne)} & 过渡期安排，过渡制度 \\
名词 & en støtte & \textit{(støtten, ingen flertal)} & 支持，补贴，资助 \\
名词 & en enighed & \textit{(enigheden, ingen flertal)} & 一致意见，共识 \\
\midrule
动词 [har] & at iværksætte & \textit{(iværksætter, iværksatte, har iværksat)} & 启动，实施（举措/计划） \\
动词 [har] & at stramme & \textit{(strammer, strammede, har strammet)} & 收紧，从严规定 \\
动词 [har] & at indgå & \textit{(indgår, indgik, har indgået)} & 达成，缔结（协议/妥协） \\
动词 [har] & at pålægge & \textit{(pålægger, pålagde, har pålagt)} & 征收（税费），强加（义务） \\
动词 [har] & at yde & \textit{(yder, ydede, har ydet)} & 提供，给予（支持/帮助） \\
\midrule
形容词 & langtidsholdbar & \textit{(langtidsholdbart, -bare)} & 长期可行的，持久耐用的 \\
\end{glosertable}

% ---------------------------------------------------------------------
\ncdsubsec{四、Grammatik \& Stilistik（核心语法、论证结构与修辞精讲）}

\begin{grammatikbox}{语法点 1：丹麦时政六大核心动宾搭配（Politisk kollokationsmatrix）}
写丹麦气候、教育、医疗等社会政策类文章时，决定文章像不像母语者社论的关键，在于能不能准确调动下面六组“动词 + 政策名词”的黄金搭档：

\vspace{0.3em}
\renewcommand{\arraystretch}{1.24}
{\small
\begin{tabularx}{\linewidth}{@{} >{\bfseries\raggedright\arraybackslash}p{3.4cm} >{\raggedright\arraybackslash}p{4.1cm} >{\raggedright\arraybackslash}X @{}}
\toprule
\textbf{核心动宾搭配} & \textbf{准确中文语义与介词} & \textbf{地道政论例句（注意语序与被动）} \\
\midrule
\da{at iværksætte tiltag} & 启动、落实具体政策举措 & \textit{Regeringen må \da{iværksætte konkrete tiltag} nu.} \\
\midrule
\da{at stramme lovgivningen} & 收紧法律/从严监管（也可用 \textit{reglerne}） & \textit{Flere partier ønsker at \da{stramme lovgivningen} på området.} \\
\midrule
\da{at indgå et bredt forlig} & 达成广泛跨党派协议（接 \textit{om}） & \textit{Partierne \da{indgik et bredt forlig om} den grønne omstilling.} \\
\midrule
\da{at opnå et kompromis} & 经谈判各方退让后取得妥协 & \textit{Efter lange forhandlinger \da{opnåede parterne et kompromis}.} \\
\midrule
\da{at pålægge X en afgift} & 对某群体/商品征收专项税（双宾语） & \textit{Staten vil \da{pålægge} landbruget \da{en CO2-afgift}.} \\
\midrule
\da{at yde økonomisk støtte} & 提供财政补贴/经费支持（接 \textit{til}） & \textit{Det er nødvendigt at \da{yde økonomisk støtte til} ny teknologi.} \\
\bottomrule
\end{tabularx}
}
\end{grammatikbox}

\begin{grammatikbox}{语法点 2：政策妥协句式——条件限制连词与情态被动策论}
在论证政策方案为什么需要妥协与配套机制时，请把本册第 12 课的条件连词（\textit{forudsat at, medmindre}）与被动语态结合起来使用：
\begin{enumerate}
  \item \textbf{\da{forudsat at}（前提条件：“只要/前提是……”）+ 完整从句}：用来设定改革成功的必备配套条件。若从句置于句首，后面主句发生\textbf{V2 倒装}！\newline
  例：\textit{\da{Forudsat at} staten yder støtte, \dabold{kan} de fleste landmænd \dabold{støtte} planen.}（注意主句动词 \textit{kan} 紧跟逗号倒装！）
  \item \textbf{\da{medmindre}（反面限制：“除非……”）+ 从句语序}：相当于 \textit{hvis ikke}。切记：\textit{medmindre} 本身已含否定意味，从句内不要再重复加 \textit{ikke}，且保持“主语 + 副词 + 动词”不倒装！\newline
  例：\textit{Mange gårde må lukke, \da{medmindre} man \dabold{hurtigt finder} en fair overgangsordning.}
  \item \textbf{策论建议被动：\textit{bør / må} + \textit{-s} 被动}：在提建议时隐藏具体个人，使语气更客观权威：\textit{Der \da{bør iværksættes} nye tiltag, og lovgivningen \da{må strammes} gradvist.}
\end{enumerate}
\end{grammatikbox}

\begin{grammatikbox}{英丹对比排雷 \& 词汇精细辨析（Stilistik）}
\begin{enumerate}
  \item \textbf{严禁中/英式“万能动词”直译（\textit{*lave / *putte / *gøre}）}：
    \begin{itemize}
      \item \ding{55}\ \textit{Partierne \dabold{lavede} et kompromis.} $\rightarrow$ \ding{51}\ \textit{Partierne \da{opnåede} (eller \da{indgik}) et kompromis.}
      \item \ding{55}\ \textit{Regeringen \dabold{putter} en skat på kød.} $\rightarrow$ \ding{51}\ \textit{Regeringen \da{pålægger} kød en afgift.}（对特定行业/排放课税用 \textit{afgift}，而非泛指所得税 \textit{skat}！）
      \item \ding{55}\ \textit{De vil \dabold{gøre} loven hårdere.} $\rightarrow$ \ding{51}\ \textit{De vil \da{stramme lovgivningen}.}
    \end{itemize}
  \item \textbf{分清 \textit{et forlig} vs \textit{en aftale} vs \textit{et kompromis}}：\textit{et forlig} 特指议会多个政党（或政府与行业代表）达成的、具有长期政治约束力的“跨党派大协议”（常见 \textit{forligspartierne}）；\textit{en aftale} 泛指日常或商业协议；\textit{et kompromis} 强调双方互让一步后的“折中结果”。
  \item \textbf{双宾语结构 \textit{pålægge [人/行业] [税费]} 的被动改写}：主动句 \textit{Staten pålægger landbruget en afgift} 变为被动时，既可把人/行业作主语（\textit{Landbruget bliver pålagt en afgift}），也可用无人称（\textit{Der pålægges en afgift på udledningen}）。
\end{enumerate}

\textbf{本课政论高频搭配（Faste vendinger）：}
\begin{itemize}
  \item \da{fra den ene dag til den anden}（一夜之间、仓促地）
  \item \da{at nå til enighed om noget}（就某事达成一致意见）/ \da{at stille krav til nogen}（对某人提出要求）
  \item \da{at skabe klare rammer for...}（为……制定清晰框架）/ \da{at sikre arbejdspladser i landdistrikterne}（保障乡村就业）
\end{itemize}
\end{grammatikbox}

% ---------------------------------------------------------------------
\ncdsubsec{五、Mønsterøvelser（句型与辩论操练）}

\textbf{练习 A：政策举措与妥协搭配精准填空（每空填入最准确的 1 个丹麦语词）}
\begin{enumerate}
  \item Hvis Danmark skal nå sine klimamål, skal regeringen hurtigt \underline{\hspace{2.2cm}} effektive tiltag i landbruget.
  \item Miljøorganisationer kræver, at Folketinget vil \underline{\hspace{2.2cm}} lovgivningen for at sænke udledningen.
  \item Efter lange forhandlinger lykkedes det endelig partierne at \underline{\hspace{2.2cm}} et bredt forlig.
  \item Det bliver svært at \underline{\hspace{2.2cm}} et kompromis, medmindre alle parter er villige til at give sig lidt.
  \item Staten overvejer at \underline{\hspace{2.2cm}} landbruget en ny afgift på udledning af drivhusgasser.
  \item For at hjælpe mindre gårde med omstillingen må staten \underline{\hspace{2.2cm}} økonomisk støtte til grøn teknologi.
  \item Landmændene ønsker en fair \underline{\hspace{2.2cm}}, så de har tid til at tilpasse produktionen gradvist.
\end{enumerate}

\textbf{练习 B：语体升级与病句改错（修正英语直译动词、介词或语序错误，或按要求改写）}
\begin{enumerate}
  \item 【纠错改写】\textit{De politiske partier lavede endelig et kompromis i går aftes.} $\rightarrow$
  \item 【纠错改写】\textit{Regeringen vil putte en ny afgift til alle virksomheder, der udleder meget CO2.} $\rightarrow$
  \item 【纠错改写】\textit{Forudsat at staten yder støtte, landmændene kan investere i nye maskiner.} $\rightarrow$
  \item 【纠错改写】\textit{Klimamålene bliver ikke nået, medmindre strammer politikerne lovgivningen.} $\rightarrow$
  \item 【升级改写】用 \textit{bør} + \textit{-s} 被动改写为客观策论句：\textit{Man skal stramme reglerne nu.} $\rightarrow$ \textit{Lovgivningen...}
  \item 【升级改写】用 \textit{nå til enighed om} 升级句子：\textit{Alle partier sagde ja til den samme klimaplan.} $\rightarrow$
\end{enumerate}

\textbf{练习 C：中译丹（灵活调用政策动宾搭配与条件让步句型）}
\begin{enumerate}
  \item 政府必须在不久的将来启动针对农业的具体举措。
  \item 环保组织要求议会进一步收紧现行的气候立法。
  \item 各政党经过漫长谈判后终于达成了一项广泛协议。
  \item 只要国家提供经济支持，许多农民都愿意推进绿色转型。
  \item 如果不设立公平的过渡期安排，我们很难达成妥协。
  \item 对碳排放征税可能会影响部分乡村地区的就业岗位。（用 \textit{afgift}）
\end{enumerate}

% ---------------------------------------------------------------------
\ncdsubsec{六、DU3 Modul 4 实战微任务与参考答案}

\begin{emailbox}{DU3 Modul 4 Skrivning Opgave 2 实战任务（Debatindlæg · mindst 100 ord）}
\textbf{Opgave}: På et dansk debatforum skriver en bruger følgende kommentar:\newline
\textit{”Landbruget udleder alt for meget CO2. Staten bør straks pålægge landmændene høje klimaafgifter uden støtteordninger, så vi hurtigt kan nå vores klimamål!”}\\[3pt]
Skriv et \textbf{debatindlæg på 120--155 ord} (mindst 100 ord), hvor du forholder dig til kommentaren. Skriv, om du er enig eller uenig, begrund din holdning med argumenter for både klima og økonomi, og brug mindst tre politiske udtryk fra denne lektion.
\end{emailbox}

\begin{answerbox}
\textbf{Før du læser 参考回答}
\begin{enumerate}
  \item \textit{Dilemmaet består i, at landbruget står for en stor del af udledningen, samtidig med at for hårde krav kan true arbejdspladser i landdistrikterne og flytte produktionen til udlandet. Et bredt forlig omfatter derfor både nye tiltag eller en afgift, økonomisk støtte til ny teknologi og en fair overgangsordning.}
  \item \textit{Henrik kræver, at hvis staten pålægger landbruget en ny afgift, skal den samtidig yde økonomisk støtte til ny teknik og skovrejsning samt give klare rammer og en reel overgangsordning.}
\end{enumerate}

\textbf{Tekst B 阅读核对答案}
\begin{enumerate}
  \item \da{Forkert}（他明确表示：\textit{Derfor er jeg sådan set glad for, at man nu vil indgå et bredt forlig...}）
  \item \da{Den må yde økonomisk støtte til nye investeringer og skovrejsning.}
  \item \da{At mange unge familier ikke vil turde overtage en gård.}
\end{enumerate}

\textbf{练习 A 答案}
\begin{enumerate}
  \item \da{iværksætte}（固定搭配：\textit{iværksætte tiltag} 启动举措）
  \item \da{stramme}（固定搭配：\textit{stramme lovgivningen} 收紧法规）
  \item \da{indgå}（固定搭配：\textit{indgå et bredt forlig} 达成广泛协议）
  \item \da{opnå} / \da{indgå}（固定搭配：\textit{opnå / indgå et kompromis} 达成妥协）
  \item \da{pålægge}（双宾动词：\textit{pålægge landbruget en ny afgift} 向农业征收新税）
  \item \da{yde}（正式动宾搭配：\textit{yde økonomisk støtte} 提供财政支持）
  \item \da{overgangsordning}（复合名词：过渡期安排）
\end{enumerate}

\textbf{练习 B 答案}
\begin{enumerate}
  \item \textit{De politiske partier opnåede endelig et kompromis i går aftes.}（表示“达成妥协”须用 \textit{opnå} 或 \textit{indgå}，忌用 \textit{*lave}）
  \item \textit{Regeringen vil pålægge alle virksomheder, der udleder meget CO2, en ny afgift.}（或用 \textit{en ny CO2-afgift}，固定动词为 \textit{pålægge [对象] [税费]}，不用 \textit{*putte... til}）
  \item \textit{Forudsat at staten yder støtte, kan landmændene investere i nye maskiner.}（前提从句置于句首时，主句必须发生 V2 倒装：\textit{kan landmændene}）
  \item \textit{Klimamålene bliver ikke nået, medmindre politikerne strammer lovgivningen.}（\textit{medmindre} 引导从句，内部维持“主语 + 动词”不倒装语序）
  \item \textit{Lovgivningen bør strammes nu.}（用情态动词 + \textit{-s} 被动使策论表达更客观正式）
  \item \textit{Alle partier nåede til enighed om den samme klimaplan.}（使用政治高频搭配 \textit{nå til enighed om}）
\end{enumerate}

\textbf{练习 C 答案}
\begin{enumerate}
  \item \textit{Regeringen skal i den nærmeste fremtid iværksætte konkrete tiltag inden for landbruget.}
  \item \textit{Miljøorganisationer kræver, at Folketinget strammer den nuværende klimalovgivning yderligere.}
  \item \textit{Partierne indgik endelig et bredt forlig efter lange forhandlinger.}
  \item \textit{Forudsat at staten yder økonomisk støtte, vil mange landmænd gerne omstille.}
  \item \textit{Vi kan næppe opnå et kompromis, medmindre vi skaber en fair overgangsordning.}
  \item \textit{En afgift på udledningen kan true visse arbejdspladser i landdistrikterne.}
\end{enumerate}

\textbf{DU3 Modul 4 Skrivning Opgave 2 高分范文（138 词）}\\
\textit{Jeg er kun delvist enig i kommentaren. Ganske vist er det nødvendigt at \da{stramme lovgivningen} og \da{iværksætte} effektive \da{tiltag}, hvis Danmark skal nå sine klimamål. Alle bør jo bidrage til den grønne omstilling, og en målrettet afgift kan motivere til hurtigere forandring.}

\textit{Ikke desto mindre mener jeg, at det ville være kortsigtet at \da{pålægge} landbruget høje afgifter fra den ene dag til den anden uden støtteordninger. Hvis for mange danske gårde går konkurs, flytter produktionen blot til andre lande med lavere krav, hvorved klimaet slet ikke vinder noget.}

\textit{Derfor bør politikerne arbejde for at \da{indgå et bredt forlig} og \da{opnå et kompromis}, der kombinerer klare klimakrav med en fair overgangsordning. \dabold{Forudsat at} staten samtidig \da{yder økonomisk støtte} til grøn teknologi, kan vi både beskytte klimaet og bevare stærke lokalsamfund på landet.}

\small\textbf{考官点评}：这篇 138 词的辩论贴立场鲜明且充满丹麦式的理性权衡（部分赞同 + 顾及现实后果）。自然嵌入了 \textit{stramme lovgivningen, iværksætte tiltag, pålægge afgifter, indgå et bredt forlig, opnå et kompromis, yde økonomisk støtte} 六大核心搭配，还顺畅呼应了第四册的 \textit{Ganske vist..., Ikke desto mindre...}、\textit{hvorved} 与 \textit{Forudsat at...} V2 倒装，堪称 Modul 4 Opgave 2 满分示范。
\end{answerbox}

\ncdlesson{18}{Turisme i København: Velsignelse eller byrde?}{哥本哈根旅游业：福音还是负担？}
{DU3 Modul 4（B1+ $\rightarrow$ B2-）｜对标 NCE 3 \& Modultest 4 \textit{Byudvikling, Turisme \& Argumenterende Essay}}
{\textbf{标准双刃剑议论文四段式模板}（背景引入 $\rightarrow$ 积极面 $\rightarrow$ 消极面 $\rightarrow$ 权衡结论与建议）· 篇章路标连接词（\textit{På den ene side, På den anden side, Alt i alt, Sammenfattende}）· 城市经济与承压词汇}
{在 DU3 Modul 4 书面写作（\textit{Opgave 2}）和双人口试辩论（\textit{Mundtlig Opgave 2}）中，绝大多数社会议题都是“利弊并存”的\textbf{双刃剑话题}——例如大城市旅游业、远程办公或无车市中心。考官最忌讳情绪化的一边倒吐槽，最青睐结构严谨、正反均衡且有具体可行建议的论述。掌握本课的\textbf{四段式双刃剑模板}，你在任何辩论题上都能做到逻辑滴水不漏！}

% ---------------------------------------------------------------------
\begin{foerduleaeser}
带着下面两个问题阅读 Tekst A 和 Tekst B：
\begin{enumerate}
  \item \textit{Hvilke positive og negative konsekvenser af den stigende turisme i København fremhæves i Tekst A, og hvordan bør kommunen reagere ifølge artiklen?}（根据 Tekst A，文章强调了哥本哈根旅游业增长带来了哪些积极与消极后果？根据文章，市政府应当如何应对？）
  \item \textit{Hvorfor er hverdagen i Indre By blevet en byrde for Mette i Tekst B, og hvilke to konkrete tiltag foreslår hun?}（在 Tekst B 中，为什么内城区的日常生活对梅特来说变成了一种负担？她提出了哪两项具体举措？）
\end{enumerate}
\end{foerduleaeser}

% ---------------------------------------------------------------------
\ncdsubsec{二、Tekster（双刃剑思辨专栏与市民来信双语篇）}

\subsection*{Tekst A：\textit{Turisme i København: Velsignelse eller byrde?}（四段式议论范文 · 约 245 词）}

\begin{tipbox}{语篇说明}
本篇严格按照\textbf{“背景引入 $\rightarrow$ 积极面 $\rightarrow$ 消极面 $\rightarrow$ 权衡结论与建议”}四段式双刃剑模板写就。阅读时请特别关注\textbf{\color{NCDNavy}深蓝色标注的篇章路标词与城市经济词汇}，以及\textbf{\color{NCDRed}深红色标注的悬念引入与策论建议框架}。
\end{tipbox}

I de senere år er København blevet et af Nordeuropas mest populære rejsemål. Hver sommer ankommer millioner af overnattende turister og krydstogtgæster til hovedstaden for at opleve Nyhavn, cykelkulturen og de historiske gader. \dabold{Spørgsmålet er, om} den kraftigt voksende turisme først og fremmest er en økonomisk \da{velsignelse} eller en tung \da{byrde} for byen.

\da{På den ene side} skaber turismen store værdier for hele samfundet. De mange internationale besøgende \da{øger omsætningen} markant i restauranter, hoteller, museer og lokale specialbutikker. \da{Samtidig} skaber erhvervet tusindvis af arbejdspladser, ikke mindst for unge og studerende, og giver kommunen højere skatteindtægter til kultur og offentlig transport. Et rigt gadeliv med gæster udefra gør desuden hovedstaden mere åben og levende.

\da{På den anden side} oplever mange københavnere nu også bagsiden af medaljen. Når enorme menneskemængder samles i de smalle gader i Indre By, \da{slider} det hårdt \da{på byrummet} og skaber daglig \da{trængsel} på fortove og cykelstier. \da{Derudover presser} den omfattende \da{korttidsudlejning} via digitale platforme \da{boligmarkedet}, så almindelige familier får sværere ved at finde en betalelig lejlighed. For flere beboere føles det, som om deres eget kvarter langsomt forvandles til en turistkulisse.

\da{Alt i alt} rummer turismen altså både store muligheder og alvorlige udfordringer. \dabold{Derfor bør} Københavns Kommune \dabold{ikke} forsøge at stoppe gæsterne, \dabold{men i stedet regulere} udviklingen klogt. Ved at sætte klare grænser for korttidsudlejning, sprede besøgende ud til nyere bydele og fremme bæredygtig turisme hele året kan hovedstaden \da{afbalancere} økonomisk vækst og høj livskvalitet for de lokale borgere.

\begin{transbox}
\small
\textbf{\color{NCDTeal}[Tekst A 中文精译]}\\
近年来，哥本哈根已成为北欧最受欢迎的旅游目的地之一。每年夏天，数以百万计的过夜游客和邮轮旅客抵达这座首都，体验新港、骑行文化与历史街区。问题在于，迅猛增长的旅游业对这座城市而言究竟主要是一种经济福音，还是一种沉重负担。

一方面，旅游业为整个社会创造了巨大价值。众多国际访客显著提高了餐厅、酒店、博物馆和本地特色商店的营业额。与此同时，该行业创造了成千上万个就业岗位——尤其惠及年轻人和在校大学生——并为市政厅带来了更多税收，用于支持文化事业和公共交通。此外，有来自世界各地的游客参与的丰富街头生活，也让首都显得更加开放与充满活力。

另一方面，许多哥本哈根市民如今也切身感受到了硬币的反面。当庞大的人潮聚集在内城区狭窄的街道上时，这会严重损耗城市公共空间，并在人行道和自行车道上造成日常拥堵。此外，通过数字平台进行的大规模短租挤占了住房市场，使得普通家庭越来越难以找到负担得起的公寓。对不少居民而言，他们感觉自己的街区正慢慢变成一座旅游布景板。

总而言之，旅游业既蕴含着巨大机遇，也伴随着严峻挑战。因此，哥本哈根市政府不应当试图阻止游客到来，而应当明智地调控这一发展进程。通过为短租设定明确上限、将游客分流到新兴城区并推动全年可持续旅游，首都就能在经济增长与本地居民的高生活质量之间取得平衡。
\end{transbox}

\subsection*{Tekst B：\textit{Læserbrev: ”København skal stadig være en by for os, der bor her”}（市民来信 · 约 125 词）}

\textbf{阅读任务（理解与核对）：}\textit{Læs Mettes læserbrev om hverdagen i Indre By, og besvar spørgsmålene nedenfor.}

\begin{emailbox}{Læserbrev｜Fra: Mette Lindholm (beboer i Indre By) \quad Emne: Grænser for masseturismen}
Som københavner gennem tyve år er jeg stolt af, at gæster fra hele verden sætter pris på vores by. \da{På den ene side} giver turismen liv i gaderne og holder mange små caféer og lokale butikker i gang hele sommeren.

\da{På den anden side} er hverdagen omkring Nyhavn efterhånden blevet en reel \da{byrde}. Hvert år er fortovene fyldt med store turistgrupper, og flere lejligheder i min opgang fungerer nærmest som hoteller via \da{korttidsudlejning}. Det \da{presser boligpriserne} op og svækker det trygge naboskab.

\da{Sammenfattende} mener jeg, at byen mangler balance. Kommunen bør \dabold{ikke} lukke af for gæster, \dabold{men i stedet} \da{regulere} antallet af udlejningsdage strengere og dirigere store turistbusser uden om centrum, så København forbliver et rart sted at leve --- ikke kun at besøge.
\end{emailbox}

\textbf{Tekst B 阅读核对（Rigtigt / Forkert / Kort svar，答案见第六部分）：}
\begin{enumerate}
  \item \textit{Mette mener, at turismen udelukkende skader små caféer og lokale butikker.} \underline{\hspace{2.2cm}}
  \item \textit{Hvilke to konkrete gener i hverdagen beskriver Mette i sit kvarter og i sin opgang?} \underline{\hspace{2.2cm}}
  \item \textit{Hvilke to tiltag foreslår Mette, at kommunen gennemfører for at skabe bedre balance?} \underline{\hspace{2.2cm}}
\end{enumerate}

\begin{tipbox}{本课发音与语调要点（Udtale \& Intonation）}
\begin{itemize}
  \item \textbf{本课核心抽象名词与静音字母}：\textit{en velsignelse} \ipa{vɛlˈsiˀnəlsə}（中间的 \textit{g} 完全不发音！带长元音与 \textit{stød}）；\textit{en byrde} \ipa{ˈbyɐ̯ðə}（软 \textit{d}）；\textit{en trængsel} \ipa{ˈtˢʁɛŋˀsəl}（鼻音带 \textit{stød}）；\textit{et krydstogt} \ipa{ˈkʰʁysˌtˢʌɡd}（\textit{kryds} 中的 \textit{d} 不发音！）。
  \item \textbf{长复合词重音与滑音}：\textit{korttidsudlejning} \ipa{ˈkʰɒːdˌtˢiðˀsˌuðlɑi̯ˀneŋ}（主重音落在最前面的 \textit{kort-}，\textit{tid} 与 \textit{ud-} 均含软 \textit{d}）；\textit{en omsætning} \ipa{ˈʌmˌsɛdneŋ}；\textit{et boligmarked} \ipa{ˈboːliˌmɑːɡəð}（\textit{bolig} 词尾 \textit{g} 不发音！）。
  \item \textbf{双刃剑过渡短语的连读与语调}：\textit{bagsiden af medaljen} \ipa{ˈbɑʊ̯ˌsiːðən ɑ meˈdɑljən}（\textit{bag-} 读 \ipa{bɑʊ̯}，\textit{af} 弱读为 \ipa{ɑ}）；朗读 \textit{På den ene side} 与 \textit{På den anden side} 时，重音分别落在数词 \textit{\textbf{e}ne} 和 \textit{\textbf{an}den} 上，形成清晰的对比节奏。
\end{itemize}
\end{tipbox}

\begin{transbox}
\small
\textbf{\color{NCDTeal}[Tekst B 中文对照]}\\
\textbf{读者来信（梅特·林霍尔姆，内城区居民｜为大众旅游设定边界）：}作为住了二十年的哥本哈根人，我为世界各地的游客如此喜爱我们的城市感到自豪。一方面，旅游业让街头充满生机，也让许多小咖啡馆和本地商店整个夏天都能持续运转。

另一方面，新港一带的日常生活如今渐渐成了一种实实在在的负担。每年夏天人行道都被大型旅游团挤满，我所住楼道里的好几套公寓也通过短租几乎变成了旅馆。这推高了住房价格，并削弱了令人安心的邻里关系。

综上所述，我认为这座城市缺少平衡。市政府不应将游客拒之门外，而应当更严格地调控允许出租的天数，并将大型旅游巴士疏导绕开市中心，好让哥本哈根继续是一个宜居之地——而不仅仅是一个观光去处。
\end{transbox}

% ---------------------------------------------------------------------
\ncdsubsec{三、Gloser（核心生词表）}

\begin{tipbox}{词表重点}
注意区分两个易混的“压力/损耗”动词：\textit{at presse} 是直接及物动词（\textit{presse boligmarkedet} 挤占/施压住房市场，后面不加介词）；而 \textit{at slide} 表示“磨损、消耗”，必须搭配介词 \dabold{på}（\textit{slide på byrummet} 损耗城市公共空间）。此外，\textit{en trængsel}（拥堵）与 \textit{en turisme}（旅游业）均为不可数共性名词（\textit{kun ental}）。
\end{tipbox}

\begin{glosertable}
名词 & en velsignelse & \textit{(velsignelsen, velsignelser, velsignelserne)} & 福音，莫大的恩惠 \\
名词 & en byrde & \textit{(byrden, byrder, byrderne)} & 负担，重负 \\
名词 & en turisme & \textit{(turismen, kun ental)} & 旅游业 \\
名词 & et rejsemål & \textit{(rejsemålet, rejsemål, rejsemålene)} & 旅游目的地（单复同形） \\
名词 & en omsætning & \textit{(omsætningen, omsætninger, -erne)} & 营业额，销售额 \\
名词 & en trængsel & \textit{(trængslen, kun ental)} & 拥挤，人潮拥堵 \\
名词 & en korttidsudlejning & \textit{(korttidsudlejningen, -er, -erne)} & 短期出租，短租 \\
名词 & et boligmarked & \textit{(boligmarkedet, -markeder, -ederne)} & 住房市场 \\
名词 & et byrum & \textit{(byrummet, byrum, byrummene)} & 城市公共空间 \\
名词 & en kulisse & \textit{(kulissen, kulisser, kulisserne)} & 舞台布景，空壳景观 \\
名词 & en livskvalitet & \textit{(livskvaliteten, kun ental)} & 生活质量 \\
\midrule
动词 [har] & at øge & \textit{(øger, øgede, har øget)} & 增加，提高 \\
动词 [har] & at presse & \textit{(presser, pressede, har presset)} & 挤压，向……施压 \\
动词 [har] & at slide (på) & \textit{(slider, sled, har slidt på)} & 磨损，过度消耗 \\
动词 [har] & at regulere & \textit{(regulerer, regulerede, har reguleret)} & 调控，规范，管制 \\
动词 [har] & at afbalancere & \textit{(afbalancerer, -ede, har afbalanceret)} & 使平衡，兼顾权衡 \\
\midrule
形容词 & betalelig & \textit{(betaleligt, betalelige)} & 负担得起的（房价/租金） \\
形容词 & omfattende & \textit{(omfattende, 不变)} & 大规模的，广泛的 \\
\end{glosertable}

% ---------------------------------------------------------------------
\ncdsubsec{四、Grammatik \& Stilistik（核心语法、双刃剑模板与修辞精讲）}

\begin{grammatikbox}{语法点 1：标准“双刃剑”议论文四段式模板（Firetrinsmodellen）}
在 DU3 Modul 4 面对“旅游业、人工智能、居家办公、无车市中心”等经典利弊议题时，只要套用下面这套四段式模板，就能确保结构清晰、立场稳健、字数达标：

\vspace{0.3em}
\renewcommand{\arraystretch}{1.24}
{\small
\begin{tabularx}{\linewidth}{@{} >{\bfseries\raggedright\arraybackslash}p{2.9cm} >{\raggedright\arraybackslash}p{4.8cm} >{\raggedright\arraybackslash}X @{}}
\toprule
\textbf{段落环节与目标词数} & \textbf{核心句型公式（路标词）} & \textbf{典型书面示范句} \\
\midrule
\da{1. 背景引入}\newline （约 45--60 词） &
\textit{I de senere år er... blevet...}\newline
\textit{\dabold{Spørgsmålet er, om} X er en velsignelse eller en byrde.} &
\textit{Spørgsmålet er, om den stigende turisme er en velsignelse eller en byrde.} \\
\midrule
\da{2. 积极面 / 机遇}\newline （约 65--75 词） &
\textit{\da{På den ene side} + 动词 + 主语...}\newline
\textit{\da{Samtidig / Desuden} + 动词 + 主语...} &
\textit{På den ene side \dabold{øger turismen} omsætningen og skaber arbejdspladser.} \\
\midrule
\da{3. 消极面 / 挑战}\newline （约 65--75 词） &
\textit{\da{På den anden side} + 动词 + 主语...}\newline
\textit{\da{Derudover} + 动词 + 主语...} &
\textit{På den anden side \dabold{presser} korttidsudlejning boligmarkedet og skaber trængsel.} \\
\midrule
\da{4. 权衡结论与建议}\newline （约 55--65 词） &
\textit{\da{Alt i alt / Sammenfattende} + 动词 + 主语...}\newline
\textit{Man bør \dabold{ikke}..., \dabold{men i stedet}...} &
\textit{Alt i alt rummer turismen både fordele og ulemper. Man bør ikke forbyde gæster, men i stedet regulere udlejningen.} \\
\bottomrule
\end{tabularx}
}
\end{grammatikbox}

\begin{grammatikbox}{语法点 2：篇章路标词的 V2 倒装铁律 \& 平衡策论句式}
写四段式议论文时，语法扣分最集中的地方恰恰是段首的路标短语与建议句：
\begin{enumerate}
  \item \textbf{所有段首路标短语放在句首时，必须立刻触发 V2 倒装，且丹麦语不打逗号！}：\\
    \ding{172} \textit{\da{På den ene side} \dabold{skaber} turismen job.}（不能受英语 \textit{On the one hand, tourism creates...} 影响而在后面打逗号或不倒装！）\\
    \ding{173} \textit{\da{På den anden side} \dabold{oplever} mange borgere trængsel.} \\
    \ding{174} \textit{\da{Alt i alt} \dabold{rummer} udviklingen både muligheder og udfordringer.} / \textit{\da{Sammenfattende} \dabold{mener} jeg, at...}
  \item \textbf{建设性平衡建议结构：\da{ikke..., men i stedet...} + \da{ved at} + 不定式}：\\
    避免生硬地说“必须彻底禁止”，改用“不应一刀切禁止 A，而应通过 B 来调控”：\\
    例：\textit{Kommunen bør \dabold{ikke} forbyde turister, \dabold{men i stedet} regulere udlejningen \dabold{ved at} sætte klare grænser.}（注意 \textit{ved at} 后面接动词原形 \textit{sætte}！）
  \item \textbf{背景疑问句 \textit{Spørgsmålet er, om...} 的语序}：\textit{om}（是否）引导宾语/表语从句，从句内保持“主语 + 动词”正常语序，千万别倒装！
\end{enumerate}
\end{grammatikbox}

\begin{grammatikbox}{英丹对比避坑指南 \& 双刃剑写作高频搭配（Stilistik）}
\begin{enumerate}
  \item \textbf{杜绝“英语式逗号 + SVO”双连错}：\ding{55}\ \textit{På den anden side, korttidsudlejning presser priserne.} $\rightarrow$ \ding{51}\ \textit{På den anden side \dabold{presser korttidsudlejning} priserne.}
  \item \textbf{分清 \textit{presse}（无介词）与 \textit{slide på}（必加 \textit{på}）}：\ding{55}\ \textit{Turismen presser på boligmarkedet.}（书面语中 \textit{presse} 直接接宾语：\ding{51}\ \textit{Turismen \da{presser} boligmarkedet}；只有使用名词结构 \textit{lægge pres} 时才接 \textit{på}：\textit{lægge pres på boligmarkedet}）。而“损耗设施/空间”必须用 \da{slide på}：\ding{51}\ \textit{Det \da{slider på} byrummet.}
  \item \textbf{禁止把“权衡结论”写成只重复一边的立场}：Modul 4 口试与作文均要求兼顾双面。结论句中先用 \textit{Alt i alt} 承认正反双方均真实存在，再用 \textit{Derfor bør...} 给出折中方案，最易获考官认可。
\end{enumerate}

\textbf{本课城市与辩论高频搭配（Faste vendinger）：}
\begin{itemize}
  \item \da{bagsiden af medaljen}（硬币的反面，伴随的弊端）
  \item \da{at sætte klare grænser for noget}（为某事设定清晰上限）
  \item \da{at skabe arbejdspladser}（创造就业岗位）/ \da{at øge omsætningen}（提高营业额）
  \item \da{at finde en bedre balance mellem A og B}（在 A 与 B 之间找到更好的平衡）
\end{itemize}
\end{grammatikbox}

% ---------------------------------------------------------------------
\ncdsubsec{五、Mønsterøvelser（句型与辩论操练）}

\textbf{练习 A：双刃剑议论文路标词与核心搭配精准填空（每空填入最准确的 1 个丹麦语词）}
\begin{enumerate}
  \item I de \underline{\hspace{2.2cm}} år er København blevet et af Nordeuropas mest populære rejsemål.
  \item På den \underline{\hspace{2.2cm}} side øger turismen omsætningen i både restauranter og museer.
  \item På den \underline{\hspace{2.2cm}} side skaber de mange gæster trængsel i de smalle gader.
  \item En stor strøm af besøgende \underline{\hspace{2.2cm}} hårdt på byrummet i Indre By.
  \item Mange lokale beboere frygter, at deres kvarter langsomt forvandles til en \underline{\hspace{2.2cm}}.
  \item Alt i \underline{\hspace{2.2cm}} rummer den voksende turisme både store muligheder og tydelige udfordringer.
  \item Kommunen bør derfor \underline{\hspace{2.2cm}} korttidsudlejningen i stedet for at indføre et totalt forbud.
\end{enumerate}

\textbf{练习 B：语序纠错与四段式句型升级（修正路标语序或介词错误，或按提示改写）}
\begin{enumerate}
  \item 【纠错改写】\textit{På den ene side, turismen skaber mange arbejdspladser for unge i hovedstaden.} $\rightarrow$
  \item 【纠错改写】\textit{Den omfattende korttidsudlejning presser på boligmarkedet i centrum.} $\rightarrow$
  \item 【纠错改写】\textit{Alt i alt, byen har brug for en mere bæredygtig turisme i fremtiden.} $\rightarrow$
  \item 【纠错改写】\textit{De mange krydstogtgæster slider meget byrummet i sommermånederne.} $\rightarrow$
  \item 【升级改写】用 \textit{ikke..., men i stedet...} 合并为平衡建议句：\textit{Politikerne skal ikke stoppe turismen. De skal kontrollere korttidsudlejning.} $\rightarrow$
  \item 【升级改写】用 \textit{Spørgsmålet er, om ... en velsignelse eller en byrde} 引出背景悬念：\textit{Mange diskuterer, om turismen er god eller dårlig for København.} $\rightarrow$
\end{enumerate}

\textbf{练习 C：中译丹（灵活调用四段式路标词与城市旅游词汇）}
\begin{enumerate}
  \item 近年来，哥本哈根接待了越来越多的国际游客。
  \item 一方面，旅游业显著提高了餐厅和博物馆的营业额。
  \item 另一方面，大规模的短租挤压了市中心的住房市场。
  \item 问题在于，这一发展究竟是一种福音还是一种负担。
  \item 市政府应当调控增长，而不是全面禁止游客。
  \item 总而言之，我们必须兼顾经济增长与居民的生活质量。
\end{enumerate}

% ---------------------------------------------------------------------
\ncdsubsec{六、DU3 Modul 4 实战微任务与参考答案}

\begin{emailbox}{DU3 Modul 4 Skrivning Opgave 2 实战任务（Debatindlæg · mindst 100 ord）}
\textbf{Situation}: På et lokalt debatforum i København skriver borgeren Jonas følgende indlæg:\newline
\textit{”Vi bør sætte en fast grænse for antallet af turister i København og helt forbyde korttidsudlejning af private lejligheder via nettet. Byen er jo ved at drukne i gæster!”}\\[3pt]
\textbf{Opgave}: Skriv et \textbf{debatindlæg på 140--170 ord} (mindst 100 ord), hvor du forholder dig til Jonas' forslag. Brug firetrinsmodellen (baggrund $\rightarrow$ fordele $\rightarrow$ ulemper $\rightarrow$ afbalanceret konklusion med forslag), og inddrag lektionens centrale udtryk.
\end{emailbox}

\begin{answerbox}
\textbf{Før du læser 参考回答}
\begin{enumerate}
  \item \textit{Ifølge Tekst A skaber turismen omsætning, arbejdspladser og skatteindtægter, men samtidig slider den på byrummet, skaber trængsel og presser boligmarkedet via korttidsudlejning. Kommunen bør ikke stoppe gæsterne, men regulere udlejningen, sprede besøgende til nye bydele og fremme bæredygtig turisme.}
  \item \textit{Mette mener, at hverdagen er blevet en byrde, fordi fortovene blokeres af turistgrupper, og flere lejligheder i hendes opgang fungerer som hoteller, hvilket presser priserne op. Hun foreslår at regulere antallet af udlejningsdage strengere og dirigere store turistbusser uden om centrum.}
\end{enumerate}

\textbf{Tekst B 阅读核对答案}
\begin{enumerate}
  \item \da{Forkert}（她指出旅游业带来了街头活力，让许多小咖啡馆和本地商店整个夏天都能持续运转）
  \item \da{Blokerede fortove med turistgrupper og korttidsudlejning af lejligheder i opgangen.}
  \item \da{At regulere antallet af udlejningsdage strengere og dirigere turistbusser uden om centrum.}
\end{enumerate}

\textbf{练习 A 答案}
\begin{enumerate}
  \item \da{senere}（背景时间短语：\textit{I de senere år} 近年来）
  \item \da{ene}（对偶路标词：\textit{På den ene side} 一方面）
  \item \da{anden}（对偶路标词：\textit{På den anden side} 另一方面）
  \item \da{slider}（固定搭配：\textit{slide på byrummet} 损耗城市公共空间）
  \item \da{kulisse}（比喻名词：旅游布景板，注意前文不定冠词 \textit{en}）
  \item \da{alt}（总结路标词：\textit{Alt i alt} 总而言之）
  \item \da{regulere}（策论动词：\textit{regulere korttidsudlejningen} 调控短租）
\end{enumerate}

\textbf{练习 B 答案}
\begin{enumerate}
  \item \textit{På den ene side skaber turismen mange arbejdspladser for unge i hovedstaden.}（去掉英语式逗号，且句首路标触发 V2 倒装：\textit{skaber turismen}）
  \item \textit{Den omfattende korttidsudlejning presser boligmarkedet i centrum.}（动词 \textit{presse} 直接接宾语，不用加介词 \textit{på}）
  \item \textit{Alt i alt har byen brug for en mere bæredygtig turisme i fremtiden.}（去掉逗号并实行 V2 倒装：\textit{har byen}）
  \item \textit{De mange krydstogtgæster slider meget på byrummet i sommermånederne.}（动词 \textit{slide} 表达“损耗某物”必须保留介词 \textit{på}）
  \item \textit{Politikerne bør ikke stoppe turismen, men i stedet regulere korttidsudlejningen.}（用 \textit{ikke..., men i stedet...} 构建有分寸的替代方案）
  \item \textit{Spørgsmålet er, om turismen er en velsignelse eller en byrde for København.}（用双刃剑标准疑问句优雅开启论题）
\end{enumerate}

\textbf{练习 C 答案}
\begin{enumerate}
  \item \textit{I de senere år har København modtaget stadig flere internationale turister.}
  \item \textit{På den ene side øger turismen omsætningen markant i restauranter og museer.}
  \item \textit{På den anden side presser den omfattende korttidsudlejning boligmarkedet i centrum.}
  \item \textit{Spørgsmålet er, om denne udvikling er en velsignelse eller en byrde.}
  \item \textit{Kommunen bør regulere væksten i stedet for at forbyde turister helt.}
  \item \textit{Alt i alt skal vi afbalancere økonomisk vækst og beboernes livskvalitet.}
\end{enumerate}

\textbf{DU3 Modul 4 Skrivning Opgave 2 高分范文（150 词）}\\
\textit{I de senere år er turismen i København vokset hurtigt, og flere borgere diskuterer nu, om udviklingen er blevet en byrde. Jeg forstår godt Jonas' frustration, men et totalt forbud mod korttidsudlejning går efter min mening for langt.}

\textit{\da{På den ene side} er turismen en vigtig økonomisk \da{velsignelse} for hovedstaden. Gæsterne \da{øger omsætningen} i hoteller, restauranter og museer, samtidig med at branchen skaber mange arbejdspladser og giver byen et internationalt præg.}

\textit{\da{På den anden side} kan jeg godt genkende bagsiden af medaljen. Når alt for mange lejligheder lejes ud til turister, \da{presser} det \da{boligmarkedet} for almindelige familier. Desuden \da{slider} de store menneskemængder \da{på byrummet} og skaber trængsel i Indre By.}

\textit{\da{Alt i alt} handler debatten om at finde den rette balance. Kommunen bør \dabold{ikke} afvise turisterne, \dabold{men i stedet regulere} korttidsudlejningen med et årligt loft og sprede besøgende ud til flere bydele. Sådan sikrer vi både livskvalitet og grøn vækst.}

\small\textbf{考官点评}：这篇 150 词的范文严格兑现四段式双刃剑蓝图：首段引入并温和回应发帖人 Jonas；第二段用 \textit{På den ene side} V2 倒装切入经济收益；第三段用 \textit{På den anden side} 与 \textit{Desuden} 倒装剖析住房挤压与拥堵；末段以 \textit{Alt i alt} 收束并给出“限天调控 + 分流”具体建议，结构工整、用词地道精准，稳获 Modul 4 高分评价。
\end{answerbox}

\ncdlesson{19}{Dansk sprog i en engelsktalende verden}{英语世界中的丹麦语}
{DU3 Modul 4（B1+ $\rightarrow$ B2-）｜对标 NCE 3 \& Modultest 4 \textit{Sprog, Globalisering \& Retorisk Argumentation}}
{\textbf{设问推进与强调句}（\textit{Spørgsmålet er imidlertid, hvorvidt... / Det rejser spørgsmålet om, om... / Det står uden for enhver tvivl, at... / Det er netop..., der/at...}）· 间接疑问从句 \textit{hvorvidt} vs \textit{om} vs \textit{hvis} 辨析}
{当议论文讨论到“丹麦大学和跨国企业越来越普遍使用英语，这是否会造成丹麦语‘语域流失（\textit{domænetab}）’”时，真正打动 DU3 Modul 4 考官的，不是单纯罗列现象，而是用\textbf{“设问推进”}引出深层矛盾，再用\textbf{“权威强调与分裂句”}锚定核心立场。掌握本课的 \textbf{\textit{Spørgsmålet er imidlertid, hvorvidt...}} 与 \textbf{\textit{Det står uden for enhver tvivl, at...}}，你的议论文章法将充满报刊评论员的思辨张力！}

% ---------------------------------------------------------------------
\begin{foerduleaeser}
带着下面两个问题阅读 Tekst A 和 Tekst B：
\begin{enumerate}
  \item \textit{Hvad forstår man ved ”domænetab” ifølge Tekst A, og hvilke konsekvenser kan det få for både videnskab og demokrati?}（根据 Tekst A，人们如何理解“语域流失（domænetab）”？它可能给科研和民主带来哪些后果？）
  \item \textit{Hvorfor sætter Emma spørgsmålstegn ved, at engelsk alene bruges som arbejdssprog på Kaspers arbejdsplads i Tekst B?}（在 Tekst B 中，为什么埃玛对卡斯珀的公司只把英语当作工作语言提出质疑？）
\end{enumerate}
\end{foerduleaeser}

% ---------------------------------------------------------------------
\ncdsubsec{二、Tekster（语言政策专栏与职场对话双语篇）}

\subsection*{Tekst A：\textit{Dansk sprog i en engelsktalende verden}（文化与语言思辨专栏 · 约 225 词）}

\begin{tipbox}{语篇说明}
本篇探讨全球化背景下丹麦语面临的英语冲击与“语域流失”之争。阅读时请特别留意\textbf{\color{NCDNavy}深蓝色标记的设问推进句型}（\textit{Spørgsmålet er imidlertid, hvorvidt... / Det rejser spørgsmålet om...}），以及\textbf{\color{NCDRed}深红色标记的强调与分裂句型}（\textit{Det står uden for enhver tvivl, at... / Det er netop..., at...}）。
\end{tipbox}

Overalt i det danske samfund vinder engelsk hurtigt frem. På landets universiteter udbydes mange kandidatuddannelser nu udelukkende på engelsk, og i internationale virksomheder i København og Aarhus fungerer engelsk ofte som fast koncernsprog fra dag ét. Samtidig blander især unge danskere et væld af engelske låneord ind i både daglig tale og opslag på sociale medier.

\dabold{Det står uden for enhver tvivl, at} stærke engelskkundskaber er en afgørende fordel for et lille land som Danmark. Uden et fælles internationalt arbejdssprog ville det være langt sværere at tiltrække dygtige forskere, ingeniører og udenlandske investeringer. \dabold{Det er netop derfor, at} mange betragter sprogets globalisering som en både naturlig og nødvendig udvikling.

\da{Spørgsmålet er imidlertid, hvorvidt} denne udvikling på længere sigt kan svække det danske sprog gennem såkaldt domænetab. Hvis forskere kun skriver og underviser på engelsk, mister dansk gradvist de præcise fagtermer, som er nødvendige for at beskrive ny viden og teknologi. \da{Det rejser også spørgsmålet om, hvorvidt} alle borgere fortsat kan følge med i den demokratiske samtale, når centrale rapporter og beslutninger ikke længere formuleres på dansk.

\dabold{Det kan ikke understreges nok, at} et sprog hænger uløseligt sammen med kultur og sammenhængskraft. Løsningen er næppe at lukke døren for engelsk, men derimod at føre en bevidst sprogpolitik, hvor man både udvikler et levende dansk fagsprog og giver internationale medarbejdere reelle muligheder for at lære dansk i hverdagen.

\begin{transbox}
\small
\textbf{\color{NCDTeal}[Tekst A 中文精译]}\\
在丹麦社会的各个角落，英语正在迅速普及。在全国多所大学里，许多硕士专业如今完全用英语授课；而在哥本哈根和奥胡斯的跨国企业中，英语从入职第一天起就常常作为固定的集团工作语言。与此同时，尤其是丹麦年轻人在日常口语和社交媒体发帖中，混用了大量英语借词。

毫无疑问的是，扎实的英语能力对丹麦这样一个小国来说是一项决定性优势。如果没有一门共通的国际工作语言，吸引优秀的学者、工程师和海外投资就会困难得多。正因如此，许多人将语言的全球化视为一种既自然又必要的发展。

然而问题在于，从长远来看，这一发展是否会通过所谓的“语域流失”而削弱丹麦语。如果学者们只用英语撰写论文和授课，丹麦语就会逐渐丧失描述新知识与新技术所必需的精确专业术语。这也引发了一个问题：当重要报告与政策决策不再用丹麦语表述时，所有公民是否依然能够跟得上公共民主讨论。

再怎么强调也不过分的是，一门语言与文化和凝聚力密不可分。解决之道恐怕不是对英语关上大门，而是推行一种自觉的语言政策——既培育充满活力的丹麦语专业术语体系，又为外籍员工提供在日常生活中学习丹麦语的切实机会。
\end{transbox}

\subsection*{Tekst B：\textit{Skal arbejdssproget på danske kontorer altid være engelsk?}（职场双人对谈 · 约 130 词）}

\textbf{阅读任务（理解与核对）：}\textit{Læs samtalen mellem Kasper og Emma om sproget på arbejdspladsen, og besvar spørgsmålene nedenfor.}

\begin{emailbox}{职场播客对谈｜Kasper（IT-udvikler）\& Emma（HR-chef）}
\textbf{Kasper:} Hos os taler vi næsten kun engelsk hele dagen, fordi vi har kolleger fra tolv forskellige lande. \dabold{Det står uden for enhver tvivl, at} et fælles arbejdssprog skaber et åbent miljø, hvor nye internationale medarbejdere føler sig velkomne fra starten.\\
\textbf{Emma:} Det er jeg helt enig i, Kasper. \da{Spørgsmålet er imidlertid, hvorvidt} engelsk alene er nok, hvis man virkelig vil føle sig hjemme i Danmark. Hvis internationale ansatte aldrig hører dansk i kantinen, går de let glip af mange kulturelle nuancer.\\
\textbf{Kasper:} Det har du måske ret i. \da{Hvad betyder det egentlig for} vores fællesskab, hvis vi aldrig øver dansk sammen?\\
\textbf{Emma:} \dabold{Det er netop gennem sproget, at} man opbygger dybere tillid. Derfor bør virksomheder både bruge engelsk i faglige projekter og give plads til dansk i pauserne.
\end{emailbox}

\textbf{Tekst B 阅读核对（Rigtigt / Forkert / Kort svar，答案见第六部分）：}
\begin{enumerate}
  \item \textit{Kasper arbejder i en virksomhed, hvor alle medarbejdere kommer fra Danmark.} \underline{\hspace{2.2cm}}
  \item \textit{Hvad risikerer internationale ansatte at gå glip af ifølge Emma, hvis de aldrig hører dansk på arbejdet?} \underline{\hspace{2.2cm}}
  \item \textit{Hvilken løsning foreslår Emma til sidst for danske virksomheder?} \underline{\hspace{2.2cm}}
\end{enumerate}

\begin{tipbox}{本课发音与语调要点（Udtale \& Intonation）}
\begin{itemize}
  \item \textbf{设问与强调高频词的静音字母与 \textit{stød}}：\textit{hvorvidt} \ipa{vɒˈvid}（词首 \textit{hv-} 中的 \textit{h} 与词尾 \textit{d} 均不发音，重音落在后音节 \textit{-vidt}！）；\textit{en tvivl} \ipa{ˈtˢviʊ̯ˀl}（\textit{v} 滑音化为 \ipa{ʊ̯ˀ}）；短语 \textit{uden for enhver tvivl} \ipa{ˈuːðən fʌ enˈvɛɐ̯ˀ ˈtˢviʊ̯ˀl}（\textit{enhver} 中的 \textit{h} 不发音）。
  \item \textbf{学术与语言讨论词汇}：\textit{et domænetab} \ipa{doˈmɛːnəˌtˢæːˀb}（重音在 \textit{-mæ-}，\textit{tab} 带 \textit{stød}）；\textit{imidlertid} \ipa{iˈmiðˀlɐˌtˢiðˀ}（两个软 \textit{d} 均带 \textit{stød}）；\textit{en videnskab} \ipa{ˈviːðənˌsɡæːˀb}；\textit{uløseligt} \ipa{uˈløːˀsəli}（\textit{g} 与 \textit{t} 不发音！）。
  \item \textbf{设问推进句的语调落点}：朗读 \textit{Spørgsmålet er imidlertid, hvorvidt...} 时，重音刻意落在 \textit{\textbf{Spørgs}målet} 与 \textit{imidler\textbf{tid}} 上，逗号处略作停顿后再平稳引出间接疑问从句，营造理性思辨的说服力。
\end{itemize}
\end{tipbox}

\begin{transbox}
\small
\textbf{\color{NCDTeal}[Tekst B 中文对照]}\\
\textbf{卡斯珀（IT 开发工程师）：}在我们公司，大家整天几乎只说英语，因为我们有来自十二个不同国家的同事。毫无疑问的是，一种共同的工作语言能够营造开放的氛围，让新来的外籍员工从一开始就感到受欢迎。\\
\textbf{埃玛（人力资源总监）：}这一点我完全同意，卡斯珀。然而问题在于，如果一个人真的想在丹麦产生归属感，仅仅靠英语是否足够。如果外籍员工在食堂里从听不到丹麦语，他们就很容易错过许多文化上的细微之处。\\
\textbf{卡斯珀：}你也许说得对。如果我们从来不一起练习丹麦语，这对我们的团队凝聚力究竟意味着什么呢？\\
\textbf{埃玛：}正是通过语言，人们才能建立起更深层次的信任。因此，企业应当既在专业项目中使用英语，又在休息时间给丹麦语留出空间。
\end{transbox}

% ---------------------------------------------------------------------
\ncdsubsec{三、Gloser（核心生词表）}

\begin{tipbox}{词表重点}
注意书面语中表示“提出（问题/质疑）”的动词搭配是 \textit{at rejse et spørgsmål}（而非英语直译的 \textit{lave/hæve et spørgsmål}）；表示“提出问题……在于是否……”常说 \textit{rejse spørgsmålet om, hvorvidt...}。另外，\textit{hvorvidt} 是正式书面语中的“是否”，常接在 \textit{spørgsmålet er}、\textit{i tvivl om} 或 \textit{afhænge af} 之后。
\end{tipbox}

\begin{glosertable}
名词 & et domænetab & \textit{(domænetabet, domænetab, -ene)} & 语域流失（语言退出某领域） \\
名词 & et koncernsprog & \textit{(koncernsproget, koncernsprog, -ene)} & 集团通用工作语言 \\
名词 & et arbejdssprog & \textit{(arbejdssproget, arbejdssprog, -ene)} & 工作语言 \\
名词 & et låneord & \textit{(låneordet, låneord, låneordene)} & 外来借词（单复同形） \\
名词 & en fagterm & \textit{(fagtermen, fagtermer, fagtermerne)} & 专业术语 \\
名词 & et fagsprog & \textit{(fagsproget, fagsprog, fagsprogene)} & 专业用语，行业语言 \\
名词 & en sprogpolitik & \textit{(sprogpolitikken, -politikker, -erne)} & 语言政策 \\
名词 & en videnskab & \textit{(videnskaben, videnskaber, -erne)} & 科学，学术研究 \\
名词 & en globalisering & \textit{(globaliseringen, kun ental)} & 全球化 \\
名词 & en tvivl & \textit{(tvivlen, kun ental)} & 怀疑，疑问 \\
\midrule
动词 [har] & at tiltrække & \textit{(tiltrækker, tiltrak, har tiltrukket)} & 吸引（人才/投资） \\
动词 [har] & at understrege & \textit{(understreger, -ede, har understreget)} & 强调，着重指出 \\
动词 [har] & at rejse (et spørgsmål) & \textit{(rejser, rejste, har rejst)} & 提出（问题/质疑） \\
动词 [har] & at opbygge & \textit{(opbygger, opbyggede, har opbygget)} & 建立，逐步构建 \\
\midrule
形容词 & uløselig & \textit{(uløseligt, uløselige)} & 密不可分的，难解的 \\
形容词 & bevidst & \textit{(bevidst, bevidste)} & 有意识的，自觉的 \\
连词 & hvorvidt & \textit{(从属连词，不变)} & 是否，究竟是否（正式书面） \\
\end{glosertable}

% ---------------------------------------------------------------------
\ncdsubsec{四、Grammatik \& Stilistik（核心语法、设问推进与强调句精讲）}

\begin{grammatikbox}{语法点 1：设问推进与立场强调核心句型矩阵}
写高水平思辨议论文时，学会用“先抑后扬的设问”转入核心矛盾，再用“不容置疑的强调框架”定调，能瞬间把文章从 B1 平铺直叙推上 B2- 报刊水准：

\vspace{0.3em}
\renewcommand{\arraystretch}{1.24}
{\small
\begin{tabularx}{\linewidth}{@{} >{\bfseries\raggedright\arraybackslash}p{2.9cm} >{\raggedright\arraybackslash}p{5.2cm} >{\raggedright\arraybackslash}X @{}}
\toprule
\textbf{修辞功能} & \textbf{核心句型公式} & \textbf{标准 B1+/B2- 示范句} \\
\midrule
\da{1. 转折设问推进}\newline （将论题推向深层） &
\textit{\da{Spørgsmålet er imidlertid, hvorvidt}...}\newline
\textit{\da{Det rejser spørgsmålet om, hvorvidt/om}...} &
\textit{Spørgsmålet er imidlertid, hvorvidt dansk kan bevare sit fagsprog.} \\
\midrule
\da{2. 追问核心影响}\newline （聚焦实质利害） &
\textit{\da{Hvad betyder det egentlig for}... ?} &
\textit{Hvad betyder det egentlig for demokratiet, hvis rapporter kun skrives på engelsk?} \\
\midrule
\dabold{3. 权威立场强调}\newline （斩钉截铁定调） &
\textit{\dabold{Det står uden for enhver tvivl, at}...}\newline
\textit{\dabold{Det kan ikke understreges nok, at}...} &
\textit{Det står uden for enhver tvivl, at engelsk er en styrke for Danmark.} \\
\midrule
\dabold{4. 分裂强调聚焦}\newline （锁定关键要素） &
\textit{\dabold{Det er netop} [状语], \dabold{at}...}\newline
\textit{\dabold{Det er først og fremmest} [主语], \dabold{der/som}...} &
\textit{Det er netop gennem sproget, at man opbygger dybere tillid.} \\
\bottomrule
\end{tabularx}
}
\end{grammatikbox}

\begin{grammatikbox}{语法点 2：三大关键句法规则（从句语序 + 分裂句连接词 + \textit{imidlertid} 位置）}
\begin{enumerate}
  \item \textbf{\da{hvorvidt} 引导间接疑问从句（= 正式体 \textit{om}）：内部绝不倒装！}：\\
    \textit{hvorvidt} 是从属连词，专门表达“究竟是否……”。从句内部必须严格遵守\textbf{主语 +（\textit{ikke/aldrig}）+ 变位动词}：\\
    \ding{51}\ \textit{Spørgsmålet er imidlertid, hvorvidt dansk \dabold{ikke mister} vigtige fagtermer.}（\ding{55}\ \textit{...hvorvidt mister dansk...}）
  \item \textbf{分裂强调句（Kløvning）何时用 \dabold{der/som}，何时用 \dabold{at}？}：\\
    \ding{172} 当被强调部分是\textbf{人或事物名词（作从句主语或宾语）}时，用关系代词 \da{der}（仅作主语）或 \da{som}：\\
    例：\textit{Det er først og fremmest unge, \dabold{der/som} bruger mange låneord.}\\
    \ding{173} 当被强调部分是\textbf{时间、地点、原因或方式状语}（如 \textit{derfor, i kantinen, gennem sproget}）时，正式书面语使用连接词 \dabold{at}（绝不能用 \textit{der}！）：\\
    例：\textit{Det er netop derfor, \dabold{at} mange firmaer vælger engelsk.} / \textit{Det er netop gennem sproget, \dabold{at} man forstår kulturen.}
  \item \textbf{副词 \da{imidlertid}（然而）的准确落点}：在 \textit{Spørgsmålet er imidlertid, hvorvidt...} 中，\textit{imidlertid} 必须紧跟主句系动词 \textit{er} 之后，绝对不能插到 \textit{hvorvidt} 后面。
\end{enumerate}
\end{grammatikbox}

\begin{grammatikbox}{英丹对比避坑指南 \& 语言政策高频搭配（Stilistik）}
\begin{enumerate}
  \item \textbf{严禁受英语影响把“是否”写成 \textit{hvis}！}：英语日常口语常用 \textit{if} 代替 \textit{whether}，但丹麦语的 \textit{hvis} 永远只表示条件“如果”，绝对不能接在 \textit{Spørgsmålet er} 后面：\ding{55}\ \textit{Spørgsmålet er, hvis dansk overlever.} $\rightarrow$ \ding{51}\ \textit{Spørgsmålet er, \da{hvorvidt / om} dansk overlever.}
  \item \textbf{\textit{uden for} 必须分写，且搭配 \textit{enhver tvivl}}：表示“毫无疑问”，固定书写为 \da{uden for enhver tvivl}（注意 \textit{uden for} 介词后带宾语时分写，且 \textit{tvivl} 为共性单数名词，配 \textit{enhver} 而非 \textit{alle} 或 \textit{nogen}！）。
  \item \textbf{区分 \textit{et fagsprog} 与 \textit{en fagterm}}：\textit{en fagterm}（可数）指单个专业术语；\textit{et fagsprog} 指某个学科或行业的整套专业语言体系。
\end{enumerate}

\textbf{本课全球化与语言政策高频搭配（Faste vendinger）：}
\begin{itemize}
  \item \da{at vinde frem}（日益普及、逐渐占据上风）
  \item \da{at gå glip af noget}（错过、错失某事物）
  \item \da{at følge med i debatten}（跟进/参与公共讨论）
  \item \da{at føre en bevidst sprogpolitik}（推行自觉清晰的语言政策）
\end{itemize}
\end{grammatikbox}

% ---------------------------------------------------------------------
\ncdsubsec{五、Mønsterøvelser（句型与辩论操练）}

\textbf{练习 A：设问推进与强调句关键结构填空（每空填入最准确的 1 个丹麦语词）}
\begin{enumerate}
  \item Det står \underline{\hspace{2.2cm}} for enhver tvivl, at stærke engelskkundskaber gavner dansk erhvervsliv.
  \item Spørgsmålet er \underline{\hspace{2.2cm}}, hvorvidt dansk mister vigtige fagtermer på universiteterne.
  \item Det rejser spørgsmålet \underline{\hspace{2.2cm}}, om alle borgere fortsat kan følge med i den offentlige debat.
  \item Det er netop gennem det fælles sprog, \underline{\hspace{2.2cm}} man opbygger dybere tillid mellem kolleger.
  \item Det er først og fremmest unge på sociale medier, \underline{\hspace{2.2cm}} bruger mange engelske låneord.
  \item Det kan ikke \underline{\hspace{2.2cm}} nok, at sprog hænger uløseligt sammen med kultur.
  \item Mange forskere er i tvivl om, \underline{\hspace{2.2cm}} engelsksprogede uddannelser på sigt skader dansk videnskab.
\end{enumerate}

\textbf{练习 B：设问改写与强调句纠错手术（修正英语直译、分裂句连接词或语序错误，或按要求升级）}
\begin{enumerate}
  \item 【纠错改写】\textit{Spørgsmålet er imidlertid, hvis dansk kan overleve som videnskabssprog.} $\rightarrow$
  \item 【纠错改写】\textit{Det står udenfor enhver tvivl, at internationale studerende beriger universitetet.} $\rightarrow$
  \item 【纠错改写】\textit{Det er netop på arbejdspladsen, der man lærer de uskrevne kulturelle regler.} $\rightarrow$
  \item 【纠错改写】\textit{Spørgsmålet er, hvorvidt kan vi bevare både dansk og engelsk i undervisningen.} $\rightarrow$
  \item 【升级改写】用 \textit{Det står uden for enhver tvivl, at...} 升级平淡句：\textit{Det er helt sikkert, at engelsk er et vigtigt arbejdssprog.} $\rightarrow$
  \item 【升级改写】用 \textit{Spørgsmålet er imidlertid, hvorvidt...} 引出转折追问：\textit{Men vi ved ikke, om koncernsproget skal være engelsk i alle møder.} $\rightarrow$
\end{enumerate}

\textbf{练习 C：中译丹（灵活运用设问推进、强调句与语言政策词汇）}
\begin{enumerate}
  \item 毫无疑问的是，良好的英语能力对丹麦来说是一项优势。
  \item 然而问题在于，这是否会在长期内削弱丹麦语专业术语。
  \item 这也引发了关于所有公民是否仍能参与讨论的问题。
  \item 正是通过在工作中使用丹麦语，新同事才会感到归属感。
  \item 再怎么强调语域流失对科研产生的影响也不为过。
  \item 在丹麦大学里，越来越多专业完全使用英语授课。
\end{enumerate}

% ---------------------------------------------------------------------
\ncdsubsec{六、DU3 Modul 4 实战微任务与参考答案}

\begin{emailbox}{DU3 Modul 4 Skrivning Opgave 2 实战任务（Debatindlæg · mindst 100 ord）}
\textbf{Situation}: På et universitets debatside skriver ingeniørstuderende Lukas følgende indlæg:\newline
\textit{”Danske universiteter og større virksomheder bør skifte helt til engelsk. Dansk er simpelthen et alt for lille sprog til moderne forskning og global konkurrence!”}\\[3pt]
\textbf{Opgave}: Skriv et \textbf{debatindlæg på 135--170 ord} (mindst 100 ord), hvor du tager stilling til Lukas' synspunkt. Brug mindst én \textbf{设问推进句}（如 \textit{Spørgsmålet er imidlertid, hvorvidt... / Det rejser spørgsmålet om...}）og mindst én \textbf{强调句}（如 \textit{Det står uden for enhver tvivl, at... / Det er netop..., at...}）。
\end{emailbox}

\begin{answerbox}
\textbf{Før du læser 参考回答}
\begin{enumerate}
  \item \textit{Ved ”domænetab” forstås, at dansk gradvist mister præcise fagtermer, når forskning og undervisning kun foregår på engelsk. Sprogforskere frygter, at dansk ikke længere kan bruges til at beskrive ny viden, og at ikke alle borgere kan følge med i den demokratiske samtale.}
  \item \textit{Emma mener ikke, at engelsk alene er nok, fordi internationale ansatte let går glip af kulturelle nuancer og dybere tillid, hvis de aldrig hører dansk i hverdagen og i kantinen.}
\end{enumerate}

\textbf{Tekst B 阅读核对答案}
\begin{enumerate}
  \item \da{Forkert}（他所在的公司有来自十二个不同国家的同事）
  \item \da{Mange kulturelle nuancer.}（会错过许多文化上的细微之处）
  \item \da{At virksomheder både bruger engelsk i faglige projekter og giver plads til dansk i pauserne.}
\end{enumerate}

\textbf{练习 A 答案}
\begin{enumerate}
  \item \da{uden}（固定强调表达：\textit{Det står uden for enhver tvivl, at...} 毫无疑问……）
  \item \da{imidlertid}（设问推进句中的核心转折副词，位于系动词 \textit{er} 之后）
  \item \da{om}（固定搭配：\textit{rejse spørgsmålet om, om / hvorvidt...}）
  \item \da{at}（强调方式状语 \textit{gennem det fælles sprog} 时，分裂句连接词须用 \textit{at}）
  \item \da{der} / \da{som}（强调主语 \textit{unge} 时，关系代词可用 \textit{der} 或 \textit{som}）
  \item \da{understreges}（强调句型：\textit{Det kan ikke understreges nok, at...} 再怎么强调也不为过）
  \item \da{hvorvidt} / \da{om}（表示“是否”，位于介词 \textit{om} 之后，书面语首选 \textit{hvorvidt}）
\end{enumerate}

\textbf{练习 B 答案}
\begin{enumerate}
  \item \textit{Spørgsmålet er imidlertid, hvorvidt dansk kan overleve som videnskabssprog.}（表示“是否”必须用 \textit{hvorvidt} 或 \textit{om}，绝不能用条件连词 \textit{hvis}）
  \item \textit{Det står uden for enhver tvivl, at internationale studerende beriger universitetet.}（此处 \textit{uden for} 后面接名词短语，必须拆开写成两个词）
  \item \textit{Det er netop på arbejdspladsen, at man lærer de uskrevne kulturelle regler.}（被强调的 \textit{på arbejdspladsen} 是地点状语，连接词应用 \textit{at} 而非 \textit{der}）
  \item \textit{Spørgsmålet er, hvorvidt vi kan bevare både dansk og engelsk i undervisningen.}（\textit{hvorvidt} 引导间接疑问从句，内部必须保持“主语 + 动词”不倒装语序）
  \item \textit{Det står uden for enhver tvivl, at engelsk er et vigtigt arbejdssprog.}（用权威强调框架提升书面说服力）
  \item \textit{Spørgsmålet er imidlertid, hvorvidt koncernsproget bør være engelsk i alle møder.}（用设问推进句优雅引导下文辩论）
\end{enumerate}

\textbf{练习 C 答案}
\begin{enumerate}
  \item \textit{Det står uden for enhver tvivl, at gode engelskkundskaber er en styrke for Danmark.}
  \item \textit{Spørgsmålet er imidlertid, hvorvidt dette på længere sigt kan svække danske fagtermer.}
  \item \textit{Det rejser også spørgsmålet om, hvorvidt alle borgere fortsat kan deltage i debatten.}
  \item \textit{Det er netop ved at tale dansk på arbejdet, at nye kolleger føler sig hjemme.}
  \item \textit{Det kan ikke understreges nok, hvordan domænetab påvirker dansk forskning.}
  \item \textit{På danske universiteter undervises stadig flere uddannelser udelukkende på engelsk.}
\end{enumerate}

\textbf{DU3 Modul 4 Skrivning Opgave 2 高分范文（135 词，不含标题）}\\
\textit{\textbf{Engelsk er vigtigt --- men dansk må ikke forsvinde}}

\textit{Jeg er kun delvist enig med Lukas i hans debatindlæg. \dabold{Det står uden for enhver tvivl, at} stærke engelskkundskaber er en afgørende fordel for Danmark. Et fælles internationalt arbejdssprog gør det lettere både at forske på tværs af grænser og at tiltrække dygtige eksperter fra udlandet.}

\textit{\da{Spørgsmålet er imidlertid, hvorvidt} vi vinder noget ved at afskaffe dansk på universiteter og arbejdspladser. Hvis nye fagtermer kun udvikles på engelsk, risikerer vi et alvorligt domænetab. \da{Det rejser spørgsmålet om, hvorvidt} almindelige borgere fortsat kan følge med i den demokratiske debat om ny teknologi og forskning.}

\textit{Samtidig \dabold{er det netop gennem det danske sprog, at} internationale kolleger bliver en reel del af fællesskabet i hverdagen. Derfor bør vi føre en bevidst sprogpolitik, der forener et stærkt engelsk arbejdssprog med et levende dansk fagsprog og god danskundervisning til udlændinge.}

\small\textbf{考官点评}：全篇 135 词（不含标题），在紧凑篇幅内完美展示了设问推进与强调句的修辞节奏：先用 \textit{Det står uden for enhver tvivl, at...} 肯定 Lukas 的合理之处，再用 \textit{Spørgsmålet er imidlertid, hvorvidt...} 与 \textit{Det rejser spørgsmålet om, hvorvidt...} 层层追问语域流失与民主参与隐患，最后用分裂强调句 \textit{Det er netop gennem det danske sprog, at...} 锁定双重语言政策建议，逻辑严密、语体成熟。
\end{answerbox}

\ncdlesson{20}{Kunsten at skrive et overbevisende debatindlæg}{如何写一篇有说服力的辩论投稿 [综合]}
{DU3 Modul 4（B1+ $\rightarrow$ B2-）结业总汇演｜对标 NCE 3 结业 \& Modultest 4 \textit{Skrivning Opgave 2（Debatindlæg）}}
{\textbf{第四册语法总汇演}（情态小品词 \textit{jo/da/nok/vist/egentlig} · 让步反驳 \textit{Ganske vist..., men ikke desto mindre...} · 复杂从句 \textit{medmindre/forudsat at/hvorved} · 虚拟推演、名词化与设问强调）· \textbf{350--400 词报刊读者来信架构}}
{祝贺你抵达《新概念丹麦语 4 · 思辨篇》的收官总汇演！如果说前三册让你能够流利叙事并阐释社会制度，第四册则赋予你用丹麦语参与公共辩论、精准拿捏语气分寸与逻辑链条的 B2- 公民表达力。一篇真正有说服力的报刊读者来信（\textit{Debatindlæg}）从不靠提高嗓门取胜，而是巧妙融合\textbf{共情小品词、礼貌让步、严谨条件与因果链条}。精读本课完整的 360 词范文与反驳回信，并完成综合汇演，你就完全具备了冲刺 DU3 Modul 4 乃至更高阶 PD3 写作的实力！}

% ---------------------------------------------------------------------
\begin{foerduleaeser}
带着下面两个宏观问题阅读 Tekst A 和 Tekst B：
\begin{enumerate}
  \item \textit{Hvordan kombinerer Sara i Tekst A anerkendelse af modargumenter med konkrete betingelser og forslag i sit debatindlæg om gratis offentlig transport?}（在 Tekst A 中，萨拉在关于免费公共交通的辩论投稿里，是如何将“承认反方论点”与“具体前提条件及建议”结合起来的？）
  \item \textit{Hvorfor mener trafikøkonomen Henrik i Tekst B, at gratis billetter alene ikke vil få pendlere uden for centrum til at droppe bilen?}（在 Tekst B 中，交通经济学家亨里克为什么认为仅仅免除车票并不能让市中心以外的通勤者放弃私家车？）
\end{enumerate}
\end{foerduleaeser}

% ---------------------------------------------------------------------
\ncdsubsec{二、Tekster（报刊辩论投稿典范与反驳回信双语篇）}

\subsection*{Tekst A：\textit{Gratis bybusser er en klog investering --- hvis vi gør det rigtigt}（完整报刊 Debatindlæg · 约 360 词）}

\begin{tipbox}{语篇说明}
本篇是一份完整的 360 词丹麦报刊读者来信（\textit{Debatindlæg}），系统汇演 NCD 4 全册六大核心语法引擎：请重点观察\textbf{\color{NCDNavy}深蓝色标记的让步、设问、强调与政策搭配}，以及\textbf{\color{NCDRed}深红色标记的情态小品词、条件限制、因果链与名词化倒装}。
\end{tipbox}

Hver morgen sidder tusindvis af danske pendlere fast i lange bilkøer på vej ind mod de største byer, mens luftforureningen og støjen stiger. Det er \dabold{jo} en velkendt udfordring, som vi alle mærker i hverdagen. Derfor dukker forslaget om at gøre bybusser og metro helt gratis op igen og igen i den offentlige debat. Efter min mening bør vi faktisk tage det skridt --- men kun som en del af en gennemtænkt plan.

\da{Ganske vist} indvender modstandere med rette, at gratis offentlig transport vil koste kommunerne mange hundrede millioner kroner om året i tabte billetindtægter. Andre peger desuden på risikoen for overfyldte busser, hvis der ikke samtidig indsættes flere afgange. \da{Ikke desto mindre} \dabold{står det uden for enhver tvivl, at} den nuværende bilisme på længere sigt koster samfundet langt dyrere i både klimaskader, slidte veje og tabte arbejdstimer i trafikpropper.

\dabold{Ved indførelsen af} gratis kollektiv trafik i byzonerne får almindelige borgere et tydeligt incitament til at lade bilen stå derhjemme. Når færre private biler kører gennem centrum, falder det samlede udslip af skadelige partikler, \dabold{hvorved} både bymiljøet og folkesundheden styrkes markant. \dabold{Som følge heraf} får børnefamilier, ældre og cyklister langt tryggere og renere gader at færdes i. Samtidig skaber billigere mobilitet større social lighed for lavindkomstgrupper, der ellers bruger en stor del af deres budget på daglige billetter.

\da{Spørgsmålet er imidlertid, hvorvidt} gratis billetter alene kan overbevise bilisterne om at ændre vaner. Svaret afhænger af en kombination af pris og kvalitet. Ordningen vil næppe fungere efter hensigten, \dabold{medmindre} byerne samtidig udvider antallet af afgange og skaber bedre forbindelser mellem forstæder og arbejdspladser. \dabold{Hvis kommunerne havde prioriteret} hurtigere letbaner og busbaner for tyve år siden, \dabold{ville langt færre pendlere i dag være afhængige} af bilen. Offentlig transport bliver altså kun et reelt alternativ, \dabold{forudsat at} man ikke går på kompromis med pålideligheden.

\da{Sammenfattende} mener jeg, at vi bør \da{afveje} den høje udgift \da{over for} de langsigtede gevinster for klimaet, folkesundheden og bylivet. Det er netop gennem et bredt politisk forlig om både hyppigere afgange og trinvis gratis kollektiv transport i de største byer, at vi kan afhjælpe trængslen og skabe grønnere byer for alle borgere.

\begin{transbox}
\small
\textbf{\color{NCDTeal}[Tekst A 中文精译]}\\
每天早晨，成千上万名丹麦通勤者在驶向大城市的路上被困在漫长的汽车长龙中，与此同时空气污染与噪音不断加剧。这毕竟是我们所有人在日常生活中都能切身感受到的熟悉难题。正因如此，将城市公交和地铁完全免费的提议在公共讨论中屡屡被提起。在我看来，我们其实应当迈出这一步——但前提是将其作为周密规划的一部分。

诚然，反对者有理有据地反驳称，免费公共交通每年会让各市镇损失数亿克朗的车票收入。此外，还有人指出，如果不同步增加班次，公交车可能会面临过度拥挤的风险。尽管如此，毫无疑问的是，长远来看当前的私家车出行模式在气候损害、道路损耗以及堵车损失的工时方面，让社会付出的代价高昂得多。

通过在城区推行免费公共交通，普通市民将获得明确的激励，从而把私家车留在家里。当穿过市中心的私家车减少时，有害颗粒物的总排放量就会下降，由此城市环境与公共健康都能得到显著增强。受此影响，有孩子的家庭、老年人以及骑行者将拥有更加安全、洁净的街道可以通行。与此同时，更低廉的出行成本也为低收入群体创造了更大的社会平等——否则他们不得不把相当大一部分预算花在日常车票上。

然而问题在于，仅仅免费车票本身是否足以说服驾车族改变习惯。答案取决于价格与服务质量的结合。除非城市同时增加发车班次，并在郊区与工作地点之间建立更顺畅的接驳线路，否则这项制度在实践中几乎难以如愿运转。倘若各市镇二十年前就优先发展更快捷的轻轨与公交专用道，今天对汽车产生依赖的通勤者原本就会少得多。因此，公共交通只有在不向准点可靠性妥协的前提下，才能成为真正的替代选择。

综上所述，我认为我们应当将高额财政支出与气候、公共健康及城市生活的长期收益放在一起权衡。正是通过就“更密集的发车班次”与“在大城市分阶段实施免费公共交通”达成广泛的政治协议，我们才能缓解交通拥堵，并为所有市民打造更绿色的城市。
\end{transbox}

\subsection*{Tekst B：\textit{Modsvar: ”Gratis billetter skaber ikke flere busser”}（报刊反驳短评 · 约 90 词）}

\textbf{阅读任务（理解与核对）：}\textit{Læs trafikøkonomen Henriks korte modsvar til Saras debatindlæg, og besvar spørgsmålene nedenfor.}

\begin{emailbox}{Modsvar i avisen｜Fra: Henrik Madsen (trafikøkonom) \quad Emne: Sv: Gratis billetter skaber ikke flere busser}
Jeg forstår \dabold{jo} godt Saras ønske om grønnere byer. Det kan \dabold{godt være, at} gratis bybusser lyder attraktivt på papiret, men vores undersøgelser viser, at pris kun er en mindre del af forklaringen.

Mange pendlere uden for centrum vælger nemlig bilen, fordi busserne kører for sjældent om aftenen og i weekenden. Gratis billetter hjælper ikke disse borgere, \dabold{medmindre} kommunen først udvider rutenettet. Hvis man bruger hele driftsbudgettet på nul kroner i billetpris, risikerer man blot dårligere service. Billigere kollektiv trafik giver kun mening, \dabold{i det omfang} kvaliteten følger med.
\end{emailbox}

\textbf{Tekst B 阅读核对（Rigtigt / Forkert / Kort svar，答案见第六部分）：}
\begin{enumerate}
  \item \textit{Henrik er fuldstændig uenig i, at det er ønskeligt at skabe grønnere byer.} \underline{\hspace{2.2cm}}
  \item \textit{Hvorfor vælger mange pendlere uden for centrum stadig bilen ifølge Henrik?} \underline{\hspace{2.2cm}}
  \item \textit{Hvad risikerer man ifølge Henrik, hvis hele driftsbudgettet bruges på gratis billetter?} \underline{\hspace{2.2cm}}
\end{enumerate}

\begin{tipbox}{本课发音与语调要点（Udtale \& Prosodi）}
\begin{itemize}
  \item \textbf{政论与交通高频词减音规律}：\textit{offentlig} \ipa{ˈʌfənli}（\textit{t} 与结尾 \textit{g} 均不发音！）；\textit{kollektiv} \ipa{kʰʌlɛɡˈtˢiʊ̯ˀ}（词尾 \textit{v} 滑音化为 \ipa{ʊ̯ˀ}）；\textit{en pendler} \ipa{ˈpʰɛnlɐ}（\textit{d} 不发音！）；\textit{afhængig} \ipa{ɑʊ̯ˈhɛŋˀi}（\textit{af} 读 \ipa{ɑʊ̯}，词尾 \textit{g} 不发音）。
  \item \textbf{第四册核心衔接语读音复盘}：\textit{ganske vist} \ipa{ˈɡansgə ˈvesd}；\textit{ikke desto mindre} \ipa{ˈeɡə dɛsd̥o ˈmenˀdʁə}（\textit{d} 在 \textit{mindre} 中常弱化消失，带 \textit{stød}）；\textit{at overbevise} \ipa{ˈʌʊ̯ɐbeˌviːˀsə}；\textit{en investering} \ipa{enveˈsd̥eːˀɐeŋ}。
  \item \textbf{长篇辩论投稿朗读节奏}：句首路标短语（\textit{Ganske vist, Ikke desto mindre, Som følge heraf, Sammenfattende}）须读出清晰重音；而情态小品词（\textit{jo, egentlig, vel}）在句中必须轻快一带而过，才能听出丹麦母语者的自然语流。
\end{itemize}
\end{tipbox}

\begin{transbox}
\small
\textbf{\color{NCDTeal}[Tekst B 中文对照]}\\
\textbf{报刊反驳来信（亨里克·马森，交通经济学家｜免费车票变不出更多公交车）：}我当然很理解萨拉对更绿色城市的向往。免费城市公交在纸面上听起来确实挺诱人，但我们的研究表明，票价只是人们出行选择的很小一部分原因。

因为许多市中心以外的通勤者之所以选择开车，是因为公交车在晚上和周末的发车班次太少。除非市政府先扩展公交线路网络，否则免费车票帮不到这些市民。如果把全部运营预算都花在“零票价”上，结果只会面临服务质量下滑的风险。更便宜的公共交通只有在服务质量同步跟上的限度内才有意义。
\end{transbox}

% ---------------------------------------------------------------------
\ncdsubsec{三、Gloser（核心生词表）}

\begin{tipbox}{词表重点}
注意动词 \textit{at afhjælpe}（缓解、改善困难）属于强变化不规则动词（\textit{afhjælper, afhjalp, har afhjulpet}），与 \textit{at hjælpe} 变化完全一致！而形容词 \textit{afhængig} 后面固定搭配介词 \dabold{af}（依赖……的）；中性名词 \textit{et udslip}（排放）与 \textit{et modsvar}（反驳回应）均为单复同形。
\end{tipbox}

\begin{glosertable}
名词 & et debatindlæg & \textit{(debatindlægget, debatindlæg, -ene)} & 辩论投稿，报刊评论来信 \\
名词 & et modsvar & \textit{(modsvaret, modsvar, modsvarene)} & 反驳回复，回应文章 \\
名词 & en pendler & \textit{(pendleren, pendlere, pendlerne)} & 通勤者，往返上班族 \\
名词 & en trængsel & \textit{(trængslen, kun ental)} & 交通拥堵，人潮拥挤 \\
名词 & et udslip & \textit{(udslippet, udslip, udslippene)} & 排放（废气/颗粒物） \\
名词 & en bilisme & \textit{(bilismen, kun ental)} & 汽车出行，私家车交通 \\
名词 & en forbindelse & \textit{(forbindelsen, forbindelser, -erne)} & 交通接驳线路，联系 \\
名词 & en afgang & \textit{(afgangen, afgange, afgangene)} & 班次，发车时刻 \\
名词 & et driftsbudget & \textit{(driftsbudgettet, -budgetter, -terne)} & 运营预算 \\
名词 & en lighed & \textit{(ligheden, ligheder, lighederne)} & 平等，均等 \\
\midrule
动词 [har] & at overbevise & \textit{(overbeviser, overbeviste, har overbevist)} & 说服，使信服 \\
动词 [har] & at afhjælpe & \textit{(afhjælper, afhjalp, har afhjulpet)} & 缓解，改善（问题/拥堵） \\
动词 [har] & at udvide & \textit{(udvider, udvidede, har udvidet)} & 扩展，增加（线路/规模） \\
动词 [har] & at prioritere & \textit{(prioriterer, prioriterede, har prioriteret)} & 优先考虑，排定优先次序 \\
\midrule
形容词 & overbevisende & \textit{(overbevisende, 不变)} & 有说服力的，令人信服的 \\
形容词 & afhængig (af) & \textit{(afhængigt, afhængige)} & 依赖……的，取决于……的 \\
副词短语 & ikke desto mindre & \textit{(不变，句间转折副词)} & 尽管如此，即便这样 \\
\end{glosertable}

% ---------------------------------------------------------------------
\ncdsubsec{四、Grammatik \& Stilistik（第四册语法总汇演与 350--400 词 Debatindlæg 写作指南）}

\begin{grammatikbox}{全景复盘：NCD 4（DU3 Modul 4）六大高阶思辨语法支柱总表}
在结束《新概念丹麦语 4》全部 20 课的学习之前，请借助下面这张总览表，完成全册核心句法与修辞武器的系统串联：

\vspace{0.3em}
\renewcommand{\arraystretch}{1.24}
{\small
\begin{tabularx}{\linewidth}{@{} >{\bfseries\raggedright\arraybackslash}p{3.0cm} >{\raggedright\arraybackslash}p{5.2cm} >{\raggedright\arraybackslash}X @{}}
\toprule
\textbf{NCD 4 语法支柱} & \textbf{核心规则与语序铁律} & \textbf{Tekst A / 考场高分示范} \\
\midrule
\textbf{1. 情态小品词}\newline （L1--L2） & \textit{jo, da, nok, vist, lige, vel, bare, egentlig}；主句放动词（及轻代词）后，从句放动词前。 & \textit{Det er \dabold{jo} en velkendt udfordring.} \\
\midrule
\textbf{2. 让步、对立与双刃剑}\newline （L4, L7, L18） & \textit{Ganske vist..., men ikke desto mindre...} 触发主句 V2 倒装；\textit{På den ene/anden side} 同样倒装。 & \textit{\da{Ganske vist} indvender modstandere... \da{Ikke desto mindre} står det...} \\
\midrule
\textbf{3. 限制条件、因果与虚拟}\newline （L8, L9, L12） & \textit{medmindre}（禁加 \textit{ikke}）/\textit{forudsat at}；\textit{Som følge heraf}（V2）vs \textit{, hvorved}（不倒装）；\textit{havde... ville have...}。 & \textit{... \dabold{hvorved} bymiljøet styrkes.}\newline \textit{... \dabold{medmindre} byerne udvider...} \\
\midrule
\textbf{4. 名词化与长定语拆解}\newline （L3, L11, L14） & 消除口语左移外置与填充词；用 \textit{ved/efter + 名词化 + af} 浓缩从句；拆解前置分词定语。 & \textit{\dabold{Ved indførelsen af} gratis kollektiv trafik får borgere...} \\
\midrule
\textbf{5. 设问推进、强调与归因}\newline （L6, L16, L19） & \textit{Spørgsmålet er imidlertid, hvorvidt...}；\textit{Det står uden for enhver tvivl, at...}；多重因素归因。 & \textit{\da{Spørgsmålet er imidlertid, hvorvidt} gratis billetter kan overbevise...} \\
\midrule
\textbf{6. 权衡、批判与政策妥协}\newline （L10, L13, L15, L17） & \textit{afveje A over for B}；\textit{gå på kompromis med}；\textit{sætte spørgsmålstegn ved}；\textit{indgå et bredt forlig}。 & \textit{... \da{afveje} den høje udgift \da{over for} de langsigtede gevinster.} \\
\bottomrule
\end{tabularx}
}
\end{grammatikbox}

\begin{grammatikbox}{高阶写作蓝图：350--400 词完整报刊 Debatindlæg 五步架构（vs Modul 4 浓缩版）}
Tekst A 全篇 360 词，完整示范了如何将单课所学的零件组装成一篇有温度、有深度、无懈可击的报刊辩论文章；在 DU3 Modul 4 \textit{Opgave 2}（100 词以上，推荐 150--180 词）考场实战中，只需按比例浓缩各段即可：

\vspace{0.3em}
\renewcommand{\arraystretch}{1.24}
{\small
\begin{tabularx}{\linewidth}{@{} >{\bfseries\raggedright\arraybackslash}p{2.9cm} >{\raggedright\arraybackslash}p{5.3cm} >{\raggedright\arraybackslash}X @{}}
\toprule
\textbf{写作阶段与篇幅} & \textbf{论证任务} & \textbf{必选 NCD 4 修辞武器} \\
\midrule
\textbf{1. 引言立论}\newline （约 60 词 / 考场 35 词） & 从生活痛点切入 + 运用 \textit{jo} 拉近读者共情 + 亮明分寸立场（\textit{Efter min mening...}）。 & 情态小品词 \textit{jo} + 清晰论点句 \\
\midrule
\textbf{2. 礼貌让步反驳}\newline （约 70 词 / 考场 45 词） & 承认对手最有力的一项顾虑，随即用强转折和权威强调拉回己方立论。 & \textit{Ganske vist... men ikke desto mindre...} + \textit{Det står uden for enhver tvivl, at...} \\
\midrule
\textbf{3. 正面深度论证}\newline （约 85 词 / 考场 45 词） & 展开核心收益：用名词化凝练主干，用严谨因果链推导连锁积极影响。 & \textit{Ved indførelsen af...} + \textit{, hvorved...} + \textit{Som følge heraf...} \\
\midrule
\textbf{4. 设问与条件边界}\newline （约 85 词 / 考场 30 词） & 主动设问检验前提条件与历史教训，展现“不走极端”的成熟思辨。 & \textit{Spørgsmålet er imidlertid, hvorvidt...} + \textit{medmindre / forudsat at} + 虚拟句 \\
\midrule
\textbf{5. 权衡收束呼吁}\newline （约 60 词 / 考场 25 词） & 总结双面价值权衡，以分裂强调句与共识政策搭配收尾呼吁。 & \textit{Sammenfattende...} + \textit{afveje A over for B} + \textit{Det er netop..., at...} \\
\bottomrule
\end{tabularx}
}
\end{grammatikbox}

\begin{grammatikbox}{第四册终极考场排雷清单（Top 5 Pitfalls）}
\begin{enumerate}
  \item \textbf{\ding{172} 句首副词与前置从句的 V2 倒装铁律}：\textit{Ganske vist, Ikke desto mindre, Som følge heraf, Af den grund, På den anden side, Sammenfattende} 以及前置条件从句（\textit{Forudsat at..., Hvis man...}）位于句首时，后面主句的第一顺位必须是\textbf{变位动词}（且丹麦语绝不在这些副词后打逗号！）。
  \item \textbf{\ding{173} 从句内部语序不可逆转}：在 \textit{at, fordi, hvorvidt, hvorved, medmindre, forudsat at} 引导的从句中，\textbf{严禁倒装}，而且否定词 \textit{ikke} 与情态小品词（\textit{jo, vist, nok}）必须站在变位动词\textbf{前面}！
  \item \textbf{\ding{174} 表达“是否”永远禁用 \textit{hvis}}：在 \textit{Spørgsmålet er, ...} 或 \textit{sætte spørgsmålstegn ved, ...} 后只能接 \da{hvorvidt} 或 \da{om}，绝不能写 \ding{55}\ \textit{Spørgsmålet er, hvis...}。
  \item \textbf{\ding{175} \textit{medmindre} 禁止再加 \textit{ikke}}：\textit{medmindre} 本身即“如果不 / 除非”，再加 \textit{ikke} 会导致逻辑完全颠倒。
  \item \textbf{\ding{176} 语体高度统一}：在报刊 \textit{Debatindlæg} 中杜绝 \textit{altså, mega, sådan set} 等松散口头禅与左移外置重复代词，用客观动词搭配（\textit{afveje, afhjælpe, prioritere, gå på kompromis med}）保持冷静专业的评论员腔调。
\end{enumerate}
\end{grammatikbox}

% ---------------------------------------------------------------------
\ncdsubsec{五、Mønsterøvelser（第四册全册语法综合演练）}

\textbf{练习 A：NCD 4 综合语法与思辨连接词精准填空（每空填入最准确的 1 个丹麦语词）}
\begin{enumerate}
  \item Du ved \underline{\hspace{2.2cm}} godt, at vores bybusser ofte er overfyldte i myldretiden.
  \item \underline{\hspace{2.2cm}} vist koster gratis offentlig transport mange penge, men det mindsker også bilkøerne.
  \item Ikke desto \underline{\hspace{2.2cm}} bør kommunen investere langt mere i grøn transport i byerne.
  \item Færre biler i centrum giver renere luft, \underline{\hspace{2.2cm}} både bymiljøet og folkesundheden bliver bedre.
  \item Ordningen virker næppe efter hensigten, \underline{\hspace{2.2cm}} kommunen samtidig udvider antallet af afgange.
  \item Spørgsmålet er \underline{\hspace{2.2cm}}, hvorvidt skatteyderne er villige til at finansiere hele forslaget.
  \item Det står uden \underline{\hspace{2.2cm}} enhver tvivl, at kollektiv trafik spiller en nøglerolle i den grønne omstilling.
  \item Politikere og borgere må afveje klimahensyn \underline{\hspace{2.2cm}} for kommunens stramme budget.
\end{enumerate}

\textbf{练习 B：第四册六大语法支柱综合纠错与升级（修正典型考场病句，或按提示改写）}
\begin{enumerate}
  \item 【小品词与从句语序纠错】\textit{Han sagde, at ordningen virker jo ikke uden flere afgange om aftenen.} $\rightarrow$
  \item 【让步倒装纠错】\textit{Ganske vist projektet er dyrt, men ikke desto mindre det kan betale sig på sigt.} $\rightarrow$
  \item 【因果连词与倒装纠错】\textit{Billetpriserne falder, hvorved vælger flere pendlere bussen i hverdagen.} $\rightarrow$
  \item 【条件与双重否定纠错】\textit{Pendlere skifter ikke til toget, medmindre afgangene ikke bliver mere pålidelige.} $\rightarrow$
  \item 【名词化与 V2 升级】用 \textit{Ved indførelsen af} 浓缩句首从句：\textit{Når man indfører gratis bybusser, falder antallet af biler i centrum.} $\rightarrow$
  \item 【虚拟推演升级】改写为过去非真实虚拟句：\textit{Kommunen investerede ikke i metroen i tide, så trængslen voksede.}（用 \textit{Hvis kommunen havde..., ville trængslen ikke være...}）$\rightarrow$
\end{enumerate}

\textbf{练习 C：中译丹（综合调用 NCD 4 核心句型撰写高分辩论句）}
\begin{enumerate}
  \item 尽管该项目花费巨大，但从长远来看仍然值得投资。
  \item 然而问题在于，免费的公交车票是否能真正吸引驾车者。
  \item 毫无疑问的是，我们必须在环境保护和财政预算之间做出权衡。
  \item 除非同时增加发车班次，否则许多通勤者依然会选择开车。
  \item 城市中心车流减少，从而使空气质量和安全感都得到提升。
  \item 只要各党派能达成广泛协议，这座城市就能实现更绿色的未来。
\end{enumerate}

% ---------------------------------------------------------------------
\ncdsubsec{六、DU3 Modul 4 实战微任务与参考答案}

\begin{emailbox}{DU3 Modul 4 结业综合写作任务（Debatindlæg · mindst 100 ord，建议 165--195 词 / 350 词报刊挑战）}
\textbf{Situation}: På et debatforum om fremtidens byer skriver Marco følgende kommentar:\newline
\textit{”Gratis offentlig transport lyder sødt, men det koster milliarder, og folk i forstæderne bor alligevel for langt væk fra stoppestederne. Brug hellere pengene på flere motorveje og lavere skat!”}\\[3pt]
\textbf{Opgave}: Skriv et overbevisende \textbf{debatindlæg på mindst 100 ord} (anbefalet eksamensniveau \textbf{165--195 ord}, eller udvid efter Tekst A's femtrinsmodel til et fuldt avisindlæg på 350--400 ord), hvor du svarer Marco. Brug mindst \textbf{fem} virkemidler fra NCD 4 (fx modalpartikel, høflig indrømmelse, retorisk spørgsmål, emfatisk sætning, kausalforbindelse, betingelsesled eller nominalisering).
\end{emailbox}

\begin{answerbox}
\textbf{Før du læser 参考回答}
\begin{enumerate}
  \item \textit{Sara anerkender modstandernes argument om, at gratis kollektiv trafik koster kommunerne mange penge og kan skabe overfyldte busser, men hun fremhæver derefter klimagevinster, renere luft og social lighed. Til sidst stiller hun konkrete betingelser om flere afgange og bedre forbindelser samt foreslår et bredt politisk forlig.}
  \item \textit{Henrik mener, at pris kun er en mindre del af forklaringen, fordi mange pendlere uden for centrum vælger bilen på grund af for få busser om aftenen og i weekenden. Hvis man bruger hele driftsbudgettet på gratis billetter uden at udvide rutenettet, risikerer man dårligere service.}
\end{enumerate}

\textbf{Tekst B 阅读核对答案}
\begin{enumerate}
  \item \da{Forkert}（他一开头就表示理解萨拉对更绿色城市的向往：\textit{Jeg forstår jo godt Saras ønske om grønnere byer.}）
  \item \da{Fordi busserne kører for sjældent om aftenen og i weekenden.}
  \item \da{Dårligere service.}（如果把全部运营预算都用来免票，只会面临服务质量变差的风险）
\end{enumerate}

\textbf{练习 A 答案}
\begin{enumerate}
  \item \da{jo}（诉诸双方共识的情态小品词，置于主句动词之后）
  \item \da{Ganske}（让步固定开头：\textit{Ganske vist} 诚然、固然）
  \item \da{mindre}（让步转折副词短语：\textit{Ikke desto mindre} 尽管如此）
  \item \da{hvorved}（逗号后接非倒装从句引出自然结果：“从而使……”）
  \item \da{medmindre}（限制条件连词：“除非……”，后接肯定从句）
  \item \da{imidlertid}（设问推进核心副词：\textit{Spørgsmålet er imidlertid, hvorvidt...}）
  \item \da{for}（强调固定搭配：\textit{uden for enhver tvivl} 毫无疑问）
  \item \da{over}（价值权衡搭配：\textit{afveje A over for B} 权衡 A 与 B）
\end{enumerate}

\textbf{练习 B 答案}
\begin{enumerate}
  \item \textit{Han sagde, at ordningen jo ikke virker uden flere afgange om aftenen.}（在 \textit{at} 从句中，情态小品词 \textit{jo} 与否定词 \textit{ikke} 均须位于变位动词 \textit{virker} 之前）
  \item \textit{Ganske vist er projektet dyrt, men ikke desto mindre kan det betale sig på sigt.}（句首 \textit{Ganske vist} 与转折后的 \textit{ikke desto mindre} 都必须触发 V2 倒装）
  \item \textit{Billetpriserne falder, hvorved flere pendlere vælger bussen i hverdagen.}（\textit{hvorved} 引导从句，内部维持“主语 + 动词”不倒装语序）
  \item \textit{Pendlere skifter ikke til toget, medmindre afgangene bliver mere pålidelige.}（\textit{medmindre} 本身已含否定意味，从句内须删去多余的 \textit{ikke}）
  \item \textit{Ved indførelsen af gratis bybusser falder antallet af biler i centrum.}（介词名词化短语放句首，后面主句发生 V2 倒装 \textit{falder antallet}）
  \item \textit{Hvis kommunen havde investeret i metroen i tide, ville trængslen ikke være vokset.}（注意变化动词 \textit{vokse} 在结果句中须用 \textit{ville ... være vokset}）
\end{enumerate}

\textbf{练习 C 答案}
\begin{enumerate}
  \item \textit{Ganske vist koster projektet meget, men ikke desto mindre er det værd at investere.}
  \item \textit{Spørgsmålet er imidlertid, hvorvidt gratis busbilletter virkelig kan tiltrække bilister.}
  \item \textit{Det står uden for enhver tvivl, at vi skal afveje miljø og økonomi.}
  \item \textit{Mange pendlere vælger stadig bilen, medmindre antallet af afgange samtidig øges.}
  \item \textit{Trafikken i centrum falder, hvorved både luftkvaliteten og trygheden bliver bedre.}
  \item \textit{Forudsat at partierne indgår et bredt forlig, kan byen sikre en grønnere fremtid.}
\end{enumerate}

\textbf{DU3 Modul 4 结业满分 Debatindlæg 范文（165 词，不含标题）}\\
\textit{\textbf{Grøn mobilitet kræver både bedre busser og kloge investeringer}}

\textit{Marco peger \da{jo} på et vigtigt problem, når han nævner de høje udgifter og de lange afstande i forstæderne. \da{Ganske vist} kan gratis offentlig transport ikke stå alene, og det ville være uansvarligt at love alt uden en holdbar finansiering.}

\textit{\da{Ikke desto mindre} \dabold{står det uden for enhver tvivl, at} flere motorveje på længere sigt kun skaber endnu mere bilisme i byerne. \dabold{Ved indførelsen af} billigere eller gradvist gratis kollektiv trafik bliver færre borgere afhængige af bilen, \dabold{hvorved} både luftkvaliteten, bymiljøet og klimaet vinder.}

\textit{\da{Spørgsmålet er imidlertid, hvorvidt} pendlere uden for centrum faktisk vil skifte til bussen i hverdagen. Svaret afhænger af pålideligheden: Ordningen lykkes næppe i praksis, \dabold{medmindre} kommunerne samtidig udvider rutenettet og sikrer langt hyppigere afgange. \dabold{Hvis byerne havde prioriteret} hurtige busbaner tidligere, \dabold{ville} mange forstadsfamilier allerede \dabold{have valgt} toget.}

\textit{\da{Sammenfattende} bør vi \da{afveje} prisen \da{over for} fremtidens trængsel og sundhed. Gennem et bredt politisk forlig kan vi styrke de lokale forbindelser først og gøre grøn transport til det letteste valg i hverdagen.}

\small\textbf{结业总评}：这篇 165 词的结业佳作（不含标题）浓缩了整部《新概念丹麦语 4》的精华！全篇既符合 DU3 Modul 4 Opgave 2 的考试节奏（也可按 Tekst A 五步法轻松扩写为 350--400 词报刊投稿），又天衣无缝地调度了全册十大武器：情态小品词 \textit{jo}、让步反驳 \textit{Ganske vist... Ikke desto mindre...}、强调句 \textit{står det uden for enhver tvivl, at...}、动词名词化 \textit{Ved indførelsen af...}、非倒装因果 \textit{, hvorved...}、设问推进 \textit{Spørgsmålet er imidlertid, hvorvidt...}、限制条件 \textit{medmindre}、非真实虚拟 \textit{Hvis byerne havde..., ville... have valgt} 以及权衡妥协搭配 \textit{afveje ... over for ...} 与 \textit{et bredt forlig}。恭喜你圆满毕业，祝在 DU3 Modul 4 考场稳操胜券！
\end{answerbox}


\end{document}
